% A_robustness.tex -- online appendix: robustness of the main design (2026-10-26, integration pass for JFEc; APPENDIX
% agent). \input inside \section{Robustness} (label app:rh_robust_stat) of rehaul/10_appendix.tex; \subsection's only.
% Bullets only (author, 2026-10-26). House rules: every number from a generated macro; replaced text is parked in dated
% comments where it was removed.
% Contents and sources:
%   - Estimation window (app:rh_robust_window): decision D8 and audit correctness-2 (W was fixed after out-of-sample
%     scoring of six windows, results/window_sweep_2014start_2026-10-22/part2, commits d8050a1 and a334d0f) and
%     correctness-3 (late-start columns in the tuning fold). The item "Estimation window" of the former section
%     "Robustness of the statistical results" (10_appendix.tex) moved here, with Table tab:rh_robust_w3500 (\rhrb...,
%     writeup/make_rehaul_robust_window.py). The every-window table of the linear models is not generated yet: it is
%     \input when generated/rehaul_window_sweep_table.tex exists (label tab:rh_window_sweep, requested from NUMBERS).
%   - Training sample (app:rh_robust_train): decision D5, drafts/2026-10-26_train_arms.tex (\rhta...,
%     writeup/make_rehaul_train_arms.py), less its CLOSE results, which are in the final section (rehaul/A_closing.tex,
%     sec:rh_econ_last30; D3). The table tab:rh_train_arms (Panels A-F) is \input here.
%   - NYSE-holiday sessions (app:rh_robust_holiday): decision D15, results/fomcfix_2026-10-26 (RULE.txt addendum A, the
%     linear 2 x 2 of FOMC fixes and holiday sessions); generated/rehaul_holiday_table.tex from NUMBERS (label
%     tab:rh_holiday requested), \input when it exists.
%   - The tuning disclosure (D14) is rehaul/A_methods.tex, app:rh_tuning, with Table tab:rh_tuning (METHODS agent);
%     the section Robustness points there, and this file does not \input the tuning table again.
% FOMC fixes (D13): the window sweep, the robustness window and the training-sample arms ran on the design before the
% FOMC fixes; the text marks them so (DATA relay: app:rh_dp_fomc says the pre-patch arms are marked where reported).
\ifdefined\rhtaNdays\else% AUTO-GENERATED by writeup/make_rehaul_train_arms.py -- do not edit; re-run the script.
% \rhta... macros: the pre-registered comparison of training samples (tab:rh_train_arms; draft
% writeup/drafts/2026-10-26_train_arms.tex).  Names: <Stat><Model><Arm><Bar>; Arm Rth (RTH13, the main design),
% All (ALL48), Cl (CLOSE), ClNat (CLOSE_NAT), AllCt / RthCt (penalties chosen at 16:00), AllThreeK / RthThreeK /
% AllThirtyNineK (the last 3,000 / 39,000 rows); Bar Close (16:00) or Bar (every bar); X... = first arm vs second
% (full set); Sh... = Sharpe ratio (D19: the paper's pooled ratio; p by Ledoit-Wolf); Ws... = the close-only
% window sweep.
% Sources (git show, pinned): insurance-selling-2026-09-19-e34-train3way 0106ccb; insurance-selling-2026-09-19-e34-train3way-trees 3501ab7; insurance-selling-2026-09-19-e34-closeonly-wsweep f4abb69.
\providecommand{\rhtaNdays}{}\renewcommand{\rhtaNdays}{2{,}579}  % trading days scored, 2 January 2014 .. 30 April 2024 (early closes dropped)
\providecommand{\rhtaNbars}{}\renewcommand{\rhtaNbars}{33{,}527}  % regular-hours bars scored (every bar; RV > 0)
\providecommand{\rhtaStart}{}\renewcommand{\rhtaStart}{2 January 2014}  % first evaluation day
\providecommand{\rhtaEnd}{}\renewcommand{\rhtaEnd}{30 April 2024}  % last evaluation day
\providecommand{\rhtaHacLag}{}\renewcommand{\rhtaHacLag}{8}  % Newey-West lag of the daily Diebold-Mariano tests
\providecommand{\rhtaNIntervals}{}\renewcommand{\rhtaNIntervals}{48}  % half-hour intervals of the futures' day (ALL48)
\providecommand{\rhtaNBarsDay}{}\renewcommand{\rhtaNBarsDay}{13}  % regular-hours bars of a day, bar ends 10:00 .. 16:00 ET (RTH13)
\providecommand{\rhtaWinDays}{}\renewcommand{\rhtaWinDays}{3{,}000}  % rolling window of the equal-calendar arms, trading days
\providecommand{\rhtaRowsAll}{}\renewcommand{\rhtaRowsAll}{144{,}000}  % training rows of ALL48 (3,000 days x 48 intervals)
\providecommand{\rhtaRowsRth}{}\renewcommand{\rhtaRowsRth}{39{,}000}  % training rows of RTH13 (3,000 days x 13 bars)
\providecommand{\rhtaRowsCl}{}\renewcommand{\rhtaRowsCl}{3{,}000}  % training rows of CLOSE (3,000 closes)
\providecommand{\rhtaRowsSmall}{}\renewcommand{\rhtaRowsSmall}{3{,}000}  % training rows of the small equal-row arms (ALL48_N3k, RTH13_N3k; CLOSE)
\providecommand{\rhtaRowsLarge}{}\renewcommand{\rhtaRowsLarge}{39{,}000}  % training rows of the large equal-row arm ALL48_N39k (= RTH13's)
\providecommand{\rhtaSpanAllSmall}{}\renewcommand{\rhtaSpanAllSmall}{62.5}  % days of 48 intervals spanned by 3,000 rows (about)
\providecommand{\rhtaSpanRthSmall}{}\renewcommand{\rhtaSpanRthSmall}{231}  % sessions of 13 bars spanned by 3,000 rows (about)
\providecommand{\rhtaSpanClSmall}{}\renewcommand{\rhtaSpanClSmall}{12}  % years of closes spanned by 3,000 closes (about)
\providecommand{\rhtaSpanAllLarge}{}\renewcommand{\rhtaSpanAllLarge}{812.5}  % days of 48 intervals spanned by 39,000 rows (about)
\providecommand{\rhtaValYear}{}\renewcommand{\rhtaValYear}{2013}  % validation year of the pre-sample fold on which every arm's linear penalties are chosen
\providecommand{\rhtaOptDays}{}\renewcommand{\rhtaOptDays}{614}  % post-warmup option days of the Sharpe ratio (Section 7; D19 pooled ratio)
\providecommand{\rhtaNClasses}{}\renewcommand{\rhtaNClasses}{6}  % model classes (ridge, lasso, elastic net, LightGBM, XGBoost, random forest)
\providecommand{\rhtaWsMin}{}\renewcommand{\rhtaWsMin}{1{,}000}  % smallest window of the close-only sweep, closes
\providecommand{\rhtaWsMax}{}\renewcommand{\rhtaWsMax}{3{,}500}  % largest window of the close-only sweep, closes
\providecommand{\rhtaPctRidgeRthClose}{}\renewcommand{\rhtaPctRidgeRthClose}{0.8}  % QLIKE % change, full set vs HAR benchmark: Ridge, RTH13 (main design: the 13 regular-hours bars, 3,000 days), 16:00 bar
\providecommand{\rhtaDmRidgeRthClose}{}\renewcommand{\rhtaDmRidgeRthClose}{0.31}  % ... its Diebold-Mariano statistic (negative: full more accurate)
\providecommand{\rhtaPRidgeRthClose}{}\renewcommand{\rhtaPRidgeRthClose}{\ensuremath{p = 0.76}}  % ... its p statement
\providecommand{\rhtaQHarRidgeRthClose}{}\renewcommand{\rhtaQHarRidgeRthClose}{0.1269}  % ... mean QLIKE, HAR benchmark
\providecommand{\rhtaQFullRidgeRthClose}{}\renewcommand{\rhtaQFullRidgeRthClose}{0.1279}  % ... mean QLIKE, full set
\providecommand{\rhtaPctLassoRthClose}{}\renewcommand{\rhtaPctLassoRthClose}{3.3}  % QLIKE % change, full set vs HAR benchmark: Lasso, RTH13 (main design: the 13 regular-hours bars, 3,000 days), 16:00 bar
\providecommand{\rhtaDmLassoRthClose}{}\renewcommand{\rhtaDmLassoRthClose}{1.04}  % ... its Diebold-Mariano statistic (negative: full more accurate)
\providecommand{\rhtaPLassoRthClose}{}\renewcommand{\rhtaPLassoRthClose}{\ensuremath{p = 0.30}}  % ... its p statement
\providecommand{\rhtaQHarLassoRthClose}{}\renewcommand{\rhtaQHarLassoRthClose}{0.1269}  % ... mean QLIKE, HAR benchmark
\providecommand{\rhtaQFullLassoRthClose}{}\renewcommand{\rhtaQFullLassoRthClose}{0.1310}  % ... mean QLIKE, full set
\providecommand{\rhtaPctEnetRthClose}{}\renewcommand{\rhtaPctEnetRthClose}{\ensuremath{-}0.5}  % QLIKE % change, full set vs HAR benchmark: Elastic net, RTH13 (main design: the 13 regular-hours bars, 3,000 days), 16:00 bar
\providecommand{\rhtaDmEnetRthClose}{}\renewcommand{\rhtaDmEnetRthClose}{\ensuremath{-}0.21}  % ... its Diebold-Mariano statistic (negative: full more accurate)
\providecommand{\rhtaPEnetRthClose}{}\renewcommand{\rhtaPEnetRthClose}{\ensuremath{p = 0.83}}  % ... its p statement
\providecommand{\rhtaQHarEnetRthClose}{}\renewcommand{\rhtaQHarEnetRthClose}{0.1269}  % ... mean QLIKE, HAR benchmark
\providecommand{\rhtaQFullEnetRthClose}{}\renewcommand{\rhtaQFullEnetRthClose}{0.1262}  % ... mean QLIKE, full set
\providecommand{\rhtaPctLgbmRthClose}{}\renewcommand{\rhtaPctLgbmRthClose}{\ensuremath{-}2.7}  % QLIKE % change, full set vs HAR benchmark: LightGBM, RTH13 (main design: the 13 regular-hours bars, 3,000 days), 16:00 bar
\providecommand{\rhtaDmLgbmRthClose}{}\renewcommand{\rhtaDmLgbmRthClose}{\ensuremath{-}1.81}  % ... its Diebold-Mariano statistic (negative: full more accurate)
\providecommand{\rhtaPLgbmRthClose}{}\renewcommand{\rhtaPLgbmRthClose}{\ensuremath{p = 0.071}}  % ... its p statement
\providecommand{\rhtaQHarLgbmRthClose}{}\renewcommand{\rhtaQHarLgbmRthClose}{0.1235}  % ... mean QLIKE, HAR benchmark
\providecommand{\rhtaQFullLgbmRthClose}{}\renewcommand{\rhtaQFullLgbmRthClose}{0.1201}  % ... mean QLIKE, full set
\providecommand{\rhtaPctXgbRthClose}{}\renewcommand{\rhtaPctXgbRthClose}{\ensuremath{-}3.4}  % QLIKE % change, full set vs HAR benchmark: XGBoost, RTH13 (main design: the 13 regular-hours bars, 3,000 days), 16:00 bar
\providecommand{\rhtaDmXgbRthClose}{}\renewcommand{\rhtaDmXgbRthClose}{\ensuremath{-}2.18}  % ... its Diebold-Mariano statistic (negative: full more accurate)
\providecommand{\rhtaPXgbRthClose}{}\renewcommand{\rhtaPXgbRthClose}{\ensuremath{p = 0.030}}  % ... its p statement
\providecommand{\rhtaQHarXgbRthClose}{}\renewcommand{\rhtaQHarXgbRthClose}{0.1234}  % ... mean QLIKE, HAR benchmark
\providecommand{\rhtaQFullXgbRthClose}{}\renewcommand{\rhtaQFullXgbRthClose}{0.1192}  % ... mean QLIKE, full set
\providecommand{\rhtaPctRfRthClose}{}\renewcommand{\rhtaPctRfRthClose}{\ensuremath{-}1.4}  % QLIKE % change, full set vs HAR benchmark: Random forest, RTH13 (main design: the 13 regular-hours bars, 3,000 days), 16:00 bar
\providecommand{\rhtaDmRfRthClose}{}\renewcommand{\rhtaDmRfRthClose}{\ensuremath{-}0.92}  % ... its Diebold-Mariano statistic (negative: full more accurate)
\providecommand{\rhtaPRfRthClose}{}\renewcommand{\rhtaPRfRthClose}{\ensuremath{p = 0.36}}  % ... its p statement
\providecommand{\rhtaQHarRfRthClose}{}\renewcommand{\rhtaQHarRfRthClose}{0.1266}  % ... mean QLIKE, HAR benchmark
\providecommand{\rhtaQFullRfRthClose}{}\renewcommand{\rhtaQFullRfRthClose}{0.1248}  % ... mean QLIKE, full set
\providecommand{\rhtaPctRidgeRthBar}{}\renewcommand{\rhtaPctRidgeRthBar}{\ensuremath{-}5.4}  % QLIKE % change, full set vs HAR benchmark: Ridge, RTH13 (main design: the 13 regular-hours bars, 3,000 days), every bar
\providecommand{\rhtaDmRidgeRthBar}{}\renewcommand{\rhtaDmRidgeRthBar}{\ensuremath{-}5.63}  % ... its Diebold-Mariano statistic (negative: full more accurate)
\providecommand{\rhtaPRidgeRthBar}{}\renewcommand{\rhtaPRidgeRthBar}{\ensuremath{p < 0.001}}  % ... its p statement
\providecommand{\rhtaQHarRidgeRthBar}{}\renewcommand{\rhtaQHarRidgeRthBar}{0.1893}  % ... mean QLIKE, HAR benchmark
\providecommand{\rhtaQFullRidgeRthBar}{}\renewcommand{\rhtaQFullRidgeRthBar}{0.1791}  % ... mean QLIKE, full set
\providecommand{\rhtaPctLassoRthBar}{}\renewcommand{\rhtaPctLassoRthBar}{\ensuremath{-}6.2}  % QLIKE % change, full set vs HAR benchmark: Lasso, RTH13 (main design: the 13 regular-hours bars, 3,000 days), every bar
\providecommand{\rhtaDmLassoRthBar}{}\renewcommand{\rhtaDmLassoRthBar}{\ensuremath{-}4.72}  % ... its Diebold-Mariano statistic (negative: full more accurate)
\providecommand{\rhtaPLassoRthBar}{}\renewcommand{\rhtaPLassoRthBar}{\ensuremath{p < 0.001}}  % ... its p statement
\providecommand{\rhtaQHarLassoRthBar}{}\renewcommand{\rhtaQHarLassoRthBar}{0.1893}  % ... mean QLIKE, HAR benchmark
\providecommand{\rhtaQFullLassoRthBar}{}\renewcommand{\rhtaQFullLassoRthBar}{0.1775}  % ... mean QLIKE, full set
\providecommand{\rhtaPctEnetRthBar}{}\renewcommand{\rhtaPctEnetRthBar}{\ensuremath{-}6.0}  % QLIKE % change, full set vs HAR benchmark: Elastic net, RTH13 (main design: the 13 regular-hours bars, 3,000 days), every bar
\providecommand{\rhtaDmEnetRthBar}{}\renewcommand{\rhtaDmEnetRthBar}{\ensuremath{-}6.98}  % ... its Diebold-Mariano statistic (negative: full more accurate)
\providecommand{\rhtaPEnetRthBar}{}\renewcommand{\rhtaPEnetRthBar}{\ensuremath{p < 0.001}}  % ... its p statement
\providecommand{\rhtaQHarEnetRthBar}{}\renewcommand{\rhtaQHarEnetRthBar}{0.1893}  % ... mean QLIKE, HAR benchmark
\providecommand{\rhtaQFullEnetRthBar}{}\renewcommand{\rhtaQFullEnetRthBar}{0.1778}  % ... mean QLIKE, full set
\providecommand{\rhtaPctLgbmRthBar}{}\renewcommand{\rhtaPctLgbmRthBar}{\ensuremath{-}14.0}  % QLIKE % change, full set vs HAR benchmark: LightGBM, RTH13 (main design: the 13 regular-hours bars, 3,000 days), every bar
\providecommand{\rhtaDmLgbmRthBar}{}\renewcommand{\rhtaDmLgbmRthBar}{\ensuremath{-}10.32}  % ... its Diebold-Mariano statistic (negative: full more accurate)
\providecommand{\rhtaPLgbmRthBar}{}\renewcommand{\rhtaPLgbmRthBar}{\ensuremath{p < 0.001}}  % ... its p statement
\providecommand{\rhtaQHarLgbmRthBar}{}\renewcommand{\rhtaQHarLgbmRthBar}{0.1875}  % ... mean QLIKE, HAR benchmark
\providecommand{\rhtaQFullLgbmRthBar}{}\renewcommand{\rhtaQFullLgbmRthBar}{0.1613}  % ... mean QLIKE, full set
\providecommand{\rhtaPctXgbRthBar}{}\renewcommand{\rhtaPctXgbRthBar}{\ensuremath{-}13.8}  % QLIKE % change, full set vs HAR benchmark: XGBoost, RTH13 (main design: the 13 regular-hours bars, 3,000 days), every bar
\providecommand{\rhtaDmXgbRthBar}{}\renewcommand{\rhtaDmXgbRthBar}{\ensuremath{-}9.94}  % ... its Diebold-Mariano statistic (negative: full more accurate)
\providecommand{\rhtaPXgbRthBar}{}\renewcommand{\rhtaPXgbRthBar}{\ensuremath{p < 0.001}}  % ... its p statement
\providecommand{\rhtaQHarXgbRthBar}{}\renewcommand{\rhtaQHarXgbRthBar}{0.1888}  % ... mean QLIKE, HAR benchmark
\providecommand{\rhtaQFullXgbRthBar}{}\renewcommand{\rhtaQFullXgbRthBar}{0.1627}  % ... mean QLIKE, full set
\providecommand{\rhtaPctRfRthBar}{}\renewcommand{\rhtaPctRfRthBar}{\ensuremath{-}11.1}  % QLIKE % change, full set vs HAR benchmark: Random forest, RTH13 (main design: the 13 regular-hours bars, 3,000 days), every bar
\providecommand{\rhtaDmRfRthBar}{}\renewcommand{\rhtaDmRfRthBar}{\ensuremath{-}9.85}  % ... its Diebold-Mariano statistic (negative: full more accurate)
\providecommand{\rhtaPRfRthBar}{}\renewcommand{\rhtaPRfRthBar}{\ensuremath{p < 0.001}}  % ... its p statement
\providecommand{\rhtaQHarRfRthBar}{}\renewcommand{\rhtaQHarRfRthBar}{0.1895}  % ... mean QLIKE, HAR benchmark
\providecommand{\rhtaQFullRfRthBar}{}\renewcommand{\rhtaQFullRfRthBar}{0.1685}  % ... mean QLIKE, full set
\providecommand{\rhtaPctRidgeAllClose}{}\renewcommand{\rhtaPctRidgeAllClose}{\ensuremath{-}2.2}  % QLIKE % change, full set vs HAR benchmark: Ridge, ALL48 (all 48 intervals, 3,000 days), 16:00 bar
\providecommand{\rhtaDmRidgeAllClose}{}\renewcommand{\rhtaDmRidgeAllClose}{\ensuremath{-}1.18}  % ... its Diebold-Mariano statistic (negative: full more accurate)
\providecommand{\rhtaPRidgeAllClose}{}\renewcommand{\rhtaPRidgeAllClose}{\ensuremath{p = 0.24}}  % ... its p statement
\providecommand{\rhtaQHarRidgeAllClose}{}\renewcommand{\rhtaQHarRidgeAllClose}{0.1262}  % ... mean QLIKE, HAR benchmark
\providecommand{\rhtaQFullRidgeAllClose}{}\renewcommand{\rhtaQFullRidgeAllClose}{0.1234}  % ... mean QLIKE, full set
\providecommand{\rhtaPctLassoAllClose}{}\renewcommand{\rhtaPctLassoAllClose}{\ensuremath{-}2.6}  % QLIKE % change, full set vs HAR benchmark: Lasso, ALL48 (all 48 intervals, 3,000 days), 16:00 bar
\providecommand{\rhtaDmLassoAllClose}{}\renewcommand{\rhtaDmLassoAllClose}{\ensuremath{-}1.45}  % ... its Diebold-Mariano statistic (negative: full more accurate)
\providecommand{\rhtaPLassoAllClose}{}\renewcommand{\rhtaPLassoAllClose}{\ensuremath{p = 0.15}}  % ... its p statement
\providecommand{\rhtaQHarLassoAllClose}{}\renewcommand{\rhtaQHarLassoAllClose}{0.1264}  % ... mean QLIKE, HAR benchmark
\providecommand{\rhtaQFullLassoAllClose}{}\renewcommand{\rhtaQFullLassoAllClose}{0.1232}  % ... mean QLIKE, full set
\providecommand{\rhtaPctEnetAllClose}{}\renewcommand{\rhtaPctEnetAllClose}{\ensuremath{-}1.8}  % QLIKE % change, full set vs HAR benchmark: Elastic net, ALL48 (all 48 intervals, 3,000 days), 16:00 bar
\providecommand{\rhtaDmEnetAllClose}{}\renewcommand{\rhtaDmEnetAllClose}{\ensuremath{-}0.96}  % ... its Diebold-Mariano statistic (negative: full more accurate)
\providecommand{\rhtaPEnetAllClose}{}\renewcommand{\rhtaPEnetAllClose}{\ensuremath{p = 0.34}}  % ... its p statement
\providecommand{\rhtaQHarEnetAllClose}{}\renewcommand{\rhtaQHarEnetAllClose}{0.1259}  % ... mean QLIKE, HAR benchmark
\providecommand{\rhtaQFullEnetAllClose}{}\renewcommand{\rhtaQFullEnetAllClose}{0.1236}  % ... mean QLIKE, full set
\providecommand{\rhtaPctLgbmAllClose}{}\renewcommand{\rhtaPctLgbmAllClose}{\ensuremath{-}6.1}  % QLIKE % change, full set vs HAR benchmark: LightGBM, ALL48 (all 48 intervals, 3,000 days), 16:00 bar
\providecommand{\rhtaDmLgbmAllClose}{}\renewcommand{\rhtaDmLgbmAllClose}{\ensuremath{-}4.45}  % ... its Diebold-Mariano statistic (negative: full more accurate)
\providecommand{\rhtaPLgbmAllClose}{}\renewcommand{\rhtaPLgbmAllClose}{\ensuremath{p < 0.001}}  % ... its p statement
\providecommand{\rhtaQHarLgbmAllClose}{}\renewcommand{\rhtaQHarLgbmAllClose}{0.1250}  % ... mean QLIKE, HAR benchmark
\providecommand{\rhtaQFullLgbmAllClose}{}\renewcommand{\rhtaQFullLgbmAllClose}{0.1174}  % ... mean QLIKE, full set
\providecommand{\rhtaPctXgbAllClose}{}\renewcommand{\rhtaPctXgbAllClose}{\ensuremath{-}6.5}  % QLIKE % change, full set vs HAR benchmark: XGBoost, ALL48 (all 48 intervals, 3,000 days), 16:00 bar
\providecommand{\rhtaDmXgbAllClose}{}\renewcommand{\rhtaDmXgbAllClose}{\ensuremath{-}5.03}  % ... its Diebold-Mariano statistic (negative: full more accurate)
\providecommand{\rhtaPXgbAllClose}{}\renewcommand{\rhtaPXgbAllClose}{\ensuremath{p < 0.001}}  % ... its p statement
\providecommand{\rhtaQHarXgbAllClose}{}\renewcommand{\rhtaQHarXgbAllClose}{0.1251}  % ... mean QLIKE, HAR benchmark
\providecommand{\rhtaQFullXgbAllClose}{}\renewcommand{\rhtaQFullXgbAllClose}{0.1170}  % ... mean QLIKE, full set
\providecommand{\rhtaPctRfAllClose}{}\renewcommand{\rhtaPctRfAllClose}{\ensuremath{-}5.6}  % QLIKE % change, full set vs HAR benchmark: Random forest, ALL48 (all 48 intervals, 3,000 days), 16:00 bar
\providecommand{\rhtaDmRfAllClose}{}\renewcommand{\rhtaDmRfAllClose}{\ensuremath{-}4.11}  % ... its Diebold-Mariano statistic (negative: full more accurate)
\providecommand{\rhtaPRfAllClose}{}\renewcommand{\rhtaPRfAllClose}{\ensuremath{p < 0.001}}  % ... its p statement
\providecommand{\rhtaQHarRfAllClose}{}\renewcommand{\rhtaQHarRfAllClose}{0.1265}  % ... mean QLIKE, HAR benchmark
\providecommand{\rhtaQFullRfAllClose}{}\renewcommand{\rhtaQFullRfAllClose}{0.1194}  % ... mean QLIKE, full set
\providecommand{\rhtaPctRidgeAllBar}{}\renewcommand{\rhtaPctRidgeAllBar}{\ensuremath{-}4.9}  % QLIKE % change, full set vs HAR benchmark: Ridge, ALL48 (all 48 intervals, 3,000 days), every bar
\providecommand{\rhtaDmRidgeAllBar}{}\renewcommand{\rhtaDmRidgeAllBar}{\ensuremath{-}6.71}  % ... its Diebold-Mariano statistic (negative: full more accurate)
\providecommand{\rhtaPRidgeAllBar}{}\renewcommand{\rhtaPRidgeAllBar}{\ensuremath{p < 0.001}}  % ... its p statement
\providecommand{\rhtaQHarRidgeAllBar}{}\renewcommand{\rhtaQHarRidgeAllBar}{0.1919}  % ... mean QLIKE, HAR benchmark
\providecommand{\rhtaQFullRidgeAllBar}{}\renewcommand{\rhtaQFullRidgeAllBar}{0.1825}  % ... mean QLIKE, full set
\providecommand{\rhtaPctLassoAllBar}{}\renewcommand{\rhtaPctLassoAllBar}{\ensuremath{-}5.1}  % QLIKE % change, full set vs HAR benchmark: Lasso, ALL48 (all 48 intervals, 3,000 days), every bar
\providecommand{\rhtaDmLassoAllBar}{}\renewcommand{\rhtaDmLassoAllBar}{\ensuremath{-}7.14}  % ... its Diebold-Mariano statistic (negative: full more accurate)
\providecommand{\rhtaPLassoAllBar}{}\renewcommand{\rhtaPLassoAllBar}{\ensuremath{p < 0.001}}  % ... its p statement
\providecommand{\rhtaQHarLassoAllBar}{}\renewcommand{\rhtaQHarLassoAllBar}{0.1917}  % ... mean QLIKE, HAR benchmark
\providecommand{\rhtaQFullLassoAllBar}{}\renewcommand{\rhtaQFullLassoAllBar}{0.1820}  % ... mean QLIKE, full set
\providecommand{\rhtaPctEnetAllBar}{}\renewcommand{\rhtaPctEnetAllBar}{\ensuremath{-}5.2}  % QLIKE % change, full set vs HAR benchmark: Elastic net, ALL48 (all 48 intervals, 3,000 days), every bar
\providecommand{\rhtaDmEnetAllBar}{}\renewcommand{\rhtaDmEnetAllBar}{\ensuremath{-}6.68}  % ... its Diebold-Mariano statistic (negative: full more accurate)
\providecommand{\rhtaPEnetAllBar}{}\renewcommand{\rhtaPEnetAllBar}{\ensuremath{p < 0.001}}  % ... its p statement
\providecommand{\rhtaQHarEnetAllBar}{}\renewcommand{\rhtaQHarEnetAllBar}{0.1916}  % ... mean QLIKE, HAR benchmark
\providecommand{\rhtaQFullEnetAllBar}{}\renewcommand{\rhtaQFullEnetAllBar}{0.1817}  % ... mean QLIKE, full set
\providecommand{\rhtaPctLgbmAllBar}{}\renewcommand{\rhtaPctLgbmAllBar}{\ensuremath{-}11.2}  % QLIKE % change, full set vs HAR benchmark: LightGBM, ALL48 (all 48 intervals, 3,000 days), every bar
\providecommand{\rhtaDmLgbmAllBar}{}\renewcommand{\rhtaDmLgbmAllBar}{\ensuremath{-}9.84}  % ... its Diebold-Mariano statistic (negative: full more accurate)
\providecommand{\rhtaPLgbmAllBar}{}\renewcommand{\rhtaPLgbmAllBar}{\ensuremath{p < 0.001}}  % ... its p statement
\providecommand{\rhtaQHarLgbmAllBar}{}\renewcommand{\rhtaQHarLgbmAllBar}{0.1898}  % ... mean QLIKE, HAR benchmark
\providecommand{\rhtaQFullLgbmAllBar}{}\renewcommand{\rhtaQFullLgbmAllBar}{0.1685}  % ... mean QLIKE, full set
\providecommand{\rhtaPctXgbAllBar}{}\renewcommand{\rhtaPctXgbAllBar}{\ensuremath{-}12.0}  % QLIKE % change, full set vs HAR benchmark: XGBoost, ALL48 (all 48 intervals, 3,000 days), every bar
\providecommand{\rhtaDmXgbAllBar}{}\renewcommand{\rhtaDmXgbAllBar}{\ensuremath{-}9.70}  % ... its Diebold-Mariano statistic (negative: full more accurate)
\providecommand{\rhtaPXgbAllBar}{}\renewcommand{\rhtaPXgbAllBar}{\ensuremath{p < 0.001}}  % ... its p statement
\providecommand{\rhtaQHarXgbAllBar}{}\renewcommand{\rhtaQHarXgbAllBar}{0.1908}  % ... mean QLIKE, HAR benchmark
\providecommand{\rhtaQFullXgbAllBar}{}\renewcommand{\rhtaQFullXgbAllBar}{0.1679}  % ... mean QLIKE, full set
\providecommand{\rhtaPctRfAllBar}{}\renewcommand{\rhtaPctRfAllBar}{\ensuremath{-}7.0}  % QLIKE % change, full set vs HAR benchmark: Random forest, ALL48 (all 48 intervals, 3,000 days), every bar
\providecommand{\rhtaDmRfAllBar}{}\renewcommand{\rhtaDmRfAllBar}{\ensuremath{-}9.54}  % ... its Diebold-Mariano statistic (negative: full more accurate)
\providecommand{\rhtaPRfAllBar}{}\renewcommand{\rhtaPRfAllBar}{\ensuremath{p < 0.001}}  % ... its p statement
\providecommand{\rhtaQHarRfAllBar}{}\renewcommand{\rhtaQHarRfAllBar}{0.1905}  % ... mean QLIKE, HAR benchmark
\providecommand{\rhtaQFullRfAllBar}{}\renewcommand{\rhtaQFullRfAllBar}{0.1773}  % ... mean QLIKE, full set
\providecommand{\rhtaPctRidgeClClose}{}\renewcommand{\rhtaPctRidgeClClose}{\ensuremath{-}6.1}  % QLIKE % change, full set vs HAR benchmark: Ridge, CLOSE (the closing half hour only, 3,000 closes), 16:00 bar
\providecommand{\rhtaDmRidgeClClose}{}\renewcommand{\rhtaDmRidgeClClose}{\ensuremath{-}3.62}  % ... its Diebold-Mariano statistic (negative: full more accurate)
\providecommand{\rhtaPRidgeClClose}{}\renewcommand{\rhtaPRidgeClClose}{\ensuremath{p < 0.001}}  % ... its p statement
\providecommand{\rhtaQHarRidgeClClose}{}\renewcommand{\rhtaQHarRidgeClClose}{0.1223}  % ... mean QLIKE, HAR benchmark
\providecommand{\rhtaQFullRidgeClClose}{}\renewcommand{\rhtaQFullRidgeClClose}{0.1148}  % ... mean QLIKE, full set
\providecommand{\rhtaPctLassoClClose}{}\renewcommand{\rhtaPctLassoClClose}{\ensuremath{-}5.4}  % QLIKE % change, full set vs HAR benchmark: Lasso, CLOSE (the closing half hour only, 3,000 closes), 16:00 bar
\providecommand{\rhtaDmLassoClClose}{}\renewcommand{\rhtaDmLassoClClose}{\ensuremath{-}5.19}  % ... its Diebold-Mariano statistic (negative: full more accurate)
\providecommand{\rhtaPLassoClClose}{}\renewcommand{\rhtaPLassoClClose}{\ensuremath{p < 0.001}}  % ... its p statement
\providecommand{\rhtaQHarLassoClClose}{}\renewcommand{\rhtaQHarLassoClClose}{0.1223}  % ... mean QLIKE, HAR benchmark
\providecommand{\rhtaQFullLassoClClose}{}\renewcommand{\rhtaQFullLassoClClose}{0.1156}  % ... mean QLIKE, full set
\providecommand{\rhtaPctEnetClClose}{}\renewcommand{\rhtaPctEnetClClose}{\ensuremath{-}6.5}  % QLIKE % change, full set vs HAR benchmark: Elastic net, CLOSE (the closing half hour only, 3,000 closes), 16:00 bar
\providecommand{\rhtaDmEnetClClose}{}\renewcommand{\rhtaDmEnetClClose}{\ensuremath{-}4.59}  % ... its Diebold-Mariano statistic (negative: full more accurate)
\providecommand{\rhtaPEnetClClose}{}\renewcommand{\rhtaPEnetClClose}{\ensuremath{p < 0.001}}  % ... its p statement
\providecommand{\rhtaQHarEnetClClose}{}\renewcommand{\rhtaQHarEnetClClose}{0.1223}  % ... mean QLIKE, HAR benchmark
\providecommand{\rhtaQFullEnetClClose}{}\renewcommand{\rhtaQFullEnetClClose}{0.1144}  % ... mean QLIKE, full set
\providecommand{\rhtaPctLgbmClClose}{}\renewcommand{\rhtaPctLgbmClClose}{\ensuremath{-}7.0}  % QLIKE % change, full set vs HAR benchmark: LightGBM, CLOSE (the closing half hour only, 3,000 closes), 16:00 bar
\providecommand{\rhtaDmLgbmClClose}{}\renewcommand{\rhtaDmLgbmClClose}{\ensuremath{-}4.50}  % ... its Diebold-Mariano statistic (negative: full more accurate)
\providecommand{\rhtaPLgbmClClose}{}\renewcommand{\rhtaPLgbmClClose}{\ensuremath{p < 0.001}}  % ... its p statement
\providecommand{\rhtaQHarLgbmClClose}{}\renewcommand{\rhtaQHarLgbmClClose}{0.1313}  % ... mean QLIKE, HAR benchmark
\providecommand{\rhtaQFullLgbmClClose}{}\renewcommand{\rhtaQFullLgbmClClose}{0.1221}  % ... mean QLIKE, full set
\providecommand{\rhtaPctXgbClClose}{}\renewcommand{\rhtaPctXgbClClose}{\ensuremath{-}8.5}  % QLIKE % change, full set vs HAR benchmark: XGBoost, CLOSE (the closing half hour only, 3,000 closes), 16:00 bar
\providecommand{\rhtaDmXgbClClose}{}\renewcommand{\rhtaDmXgbClClose}{\ensuremath{-}5.95}  % ... its Diebold-Mariano statistic (negative: full more accurate)
\providecommand{\rhtaPXgbClClose}{}\renewcommand{\rhtaPXgbClClose}{\ensuremath{p < 0.001}}  % ... its p statement
\providecommand{\rhtaQHarXgbClClose}{}\renewcommand{\rhtaQHarXgbClClose}{0.1290}  % ... mean QLIKE, HAR benchmark
\providecommand{\rhtaQFullXgbClClose}{}\renewcommand{\rhtaQFullXgbClClose}{0.1180}  % ... mean QLIKE, full set
\providecommand{\rhtaPctRfClClose}{}\renewcommand{\rhtaPctRfClClose}{\ensuremath{-}3.7}  % QLIKE % change, full set vs HAR benchmark: Random forest, CLOSE (the closing half hour only, 3,000 closes), 16:00 bar
\providecommand{\rhtaDmRfClClose}{}\renewcommand{\rhtaDmRfClClose}{\ensuremath{-}3.66}  % ... its Diebold-Mariano statistic (negative: full more accurate)
\providecommand{\rhtaPRfClClose}{}\renewcommand{\rhtaPRfClClose}{\ensuremath{p < 0.001}}  % ... its p statement
\providecommand{\rhtaQHarRfClClose}{}\renewcommand{\rhtaQHarRfClClose}{0.1267}  % ... mean QLIKE, HAR benchmark
\providecommand{\rhtaQFullRfClClose}{}\renewcommand{\rhtaQFullRfClClose}{0.1220}  % ... mean QLIKE, full set
\providecommand{\rhtaPctRidgeClNatClose}{}\renewcommand{\rhtaPctRidgeClNatClose}{\ensuremath{-}6.5}  % QLIKE % change, full set vs HAR benchmark: Ridge, CLOSE_NAT (closing rows, design built natively on them; linear only), 16:00 bar
\providecommand{\rhtaDmRidgeClNatClose}{}\renewcommand{\rhtaDmRidgeClNatClose}{\ensuremath{-}3.76}  % ... its Diebold-Mariano statistic (negative: full more accurate)
\providecommand{\rhtaPRidgeClNatClose}{}\renewcommand{\rhtaPRidgeClNatClose}{\ensuremath{p < 0.001}}  % ... its p statement
\providecommand{\rhtaQHarRidgeClNatClose}{}\renewcommand{\rhtaQHarRidgeClNatClose}{0.1226}  % ... mean QLIKE, HAR benchmark
\providecommand{\rhtaQFullRidgeClNatClose}{}\renewcommand{\rhtaQFullRidgeClNatClose}{0.1146}  % ... mean QLIKE, full set
\providecommand{\rhtaPctLassoClNatClose}{}\renewcommand{\rhtaPctLassoClNatClose}{\ensuremath{-}5.4}  % QLIKE % change, full set vs HAR benchmark: Lasso, CLOSE_NAT (closing rows, design built natively on them; linear only), 16:00 bar
\providecommand{\rhtaDmLassoClNatClose}{}\renewcommand{\rhtaDmLassoClNatClose}{\ensuremath{-}4.95}  % ... its Diebold-Mariano statistic (negative: full more accurate)
\providecommand{\rhtaPLassoClNatClose}{}\renewcommand{\rhtaPLassoClNatClose}{\ensuremath{p < 0.001}}  % ... its p statement
\providecommand{\rhtaQHarLassoClNatClose}{}\renewcommand{\rhtaQHarLassoClNatClose}{0.1226}  % ... mean QLIKE, HAR benchmark
\providecommand{\rhtaQFullLassoClNatClose}{}\renewcommand{\rhtaQFullLassoClNatClose}{0.1159}  % ... mean QLIKE, full set
\providecommand{\rhtaPctEnetClNatClose}{}\renewcommand{\rhtaPctEnetClNatClose}{\ensuremath{-}6.6}  % QLIKE % change, full set vs HAR benchmark: Elastic net, CLOSE_NAT (closing rows, design built natively on them; linear only), 16:00 bar
\providecommand{\rhtaDmEnetClNatClose}{}\renewcommand{\rhtaDmEnetClNatClose}{\ensuremath{-}4.48}  % ... its Diebold-Mariano statistic (negative: full more accurate)
\providecommand{\rhtaPEnetClNatClose}{}\renewcommand{\rhtaPEnetClNatClose}{\ensuremath{p < 0.001}}  % ... its p statement
\providecommand{\rhtaQHarEnetClNatClose}{}\renewcommand{\rhtaQHarEnetClNatClose}{0.1226}  % ... mean QLIKE, HAR benchmark
\providecommand{\rhtaQFullEnetClNatClose}{}\renewcommand{\rhtaQFullEnetClNatClose}{0.1145}  % ... mean QLIKE, full set
\providecommand{\rhtaPctRidgeAllCtClose}{}\renewcommand{\rhtaPctRidgeAllCtClose}{\ensuremath{-}1.9}  % QLIKE % change, full set vs HAR benchmark: Ridge, ALL48@16 (ALL48, penalties chosen on the 16:00 validation rows; linear only), 16:00 bar
\providecommand{\rhtaDmRidgeAllCtClose}{}\renewcommand{\rhtaDmRidgeAllCtClose}{\ensuremath{-}0.97}  % ... its Diebold-Mariano statistic (negative: full more accurate)
\providecommand{\rhtaPRidgeAllCtClose}{}\renewcommand{\rhtaPRidgeAllCtClose}{\ensuremath{p = 0.33}}  % ... its p statement
\providecommand{\rhtaQHarRidgeAllCtClose}{}\renewcommand{\rhtaQHarRidgeAllCtClose}{0.1262}  % ... mean QLIKE, HAR benchmark
\providecommand{\rhtaQFullRidgeAllCtClose}{}\renewcommand{\rhtaQFullRidgeAllCtClose}{0.1238}  % ... mean QLIKE, full set
\providecommand{\rhtaPctLassoAllCtClose}{}\renewcommand{\rhtaPctLassoAllCtClose}{\ensuremath{-}1.1}  % QLIKE % change, full set vs HAR benchmark: Lasso, ALL48@16 (ALL48, penalties chosen on the 16:00 validation rows; linear only), 16:00 bar
\providecommand{\rhtaDmLassoAllCtClose}{}\renewcommand{\rhtaDmLassoAllCtClose}{\ensuremath{-}0.53}  % ... its Diebold-Mariano statistic (negative: full more accurate)
\providecommand{\rhtaPLassoAllCtClose}{}\renewcommand{\rhtaPLassoAllCtClose}{\ensuremath{p = 0.59}}  % ... its p statement
\providecommand{\rhtaQHarLassoAllCtClose}{}\renewcommand{\rhtaQHarLassoAllCtClose}{0.1253}  % ... mean QLIKE, HAR benchmark
\providecommand{\rhtaQFullLassoAllCtClose}{}\renewcommand{\rhtaQFullLassoAllCtClose}{0.1239}  % ... mean QLIKE, full set
\providecommand{\rhtaPctEnetAllCtClose}{}\renewcommand{\rhtaPctEnetAllCtClose}{\ensuremath{-}0.8}  % QLIKE % change, full set vs HAR benchmark: Elastic net, ALL48@16 (ALL48, penalties chosen on the 16:00 validation rows; linear only), 16:00 bar
\providecommand{\rhtaDmEnetAllCtClose}{}\renewcommand{\rhtaDmEnetAllCtClose}{\ensuremath{-}0.40}  % ... its Diebold-Mariano statistic (negative: full more accurate)
\providecommand{\rhtaPEnetAllCtClose}{}\renewcommand{\rhtaPEnetAllCtClose}{\ensuremath{p = 0.69}}  % ... its p statement
\providecommand{\rhtaQHarEnetAllCtClose}{}\renewcommand{\rhtaQHarEnetAllCtClose}{0.1246}  % ... mean QLIKE, HAR benchmark
\providecommand{\rhtaQFullEnetAllCtClose}{}\renewcommand{\rhtaQFullEnetAllCtClose}{0.1236}  % ... mean QLIKE, full set
\providecommand{\rhtaPctRidgeAllCtBar}{}\renewcommand{\rhtaPctRidgeAllCtBar}{\ensuremath{-}5.2}  % QLIKE % change, full set vs HAR benchmark: Ridge, ALL48@16 (ALL48, penalties chosen on the 16:00 validation rows; linear only), every bar
\providecommand{\rhtaDmRidgeAllCtBar}{}\renewcommand{\rhtaDmRidgeAllCtBar}{\ensuremath{-}6.22}  % ... its Diebold-Mariano statistic (negative: full more accurate)
\providecommand{\rhtaPRidgeAllCtBar}{}\renewcommand{\rhtaPRidgeAllCtBar}{\ensuremath{p < 0.001}}  % ... its p statement
\providecommand{\rhtaQHarRidgeAllCtBar}{}\renewcommand{\rhtaQHarRidgeAllCtBar}{0.1919}  % ... mean QLIKE, HAR benchmark
\providecommand{\rhtaQFullRidgeAllCtBar}{}\renewcommand{\rhtaQFullRidgeAllCtBar}{0.1820}  % ... mean QLIKE, full set
\providecommand{\rhtaPctLassoAllCtBar}{}\renewcommand{\rhtaPctLassoAllCtBar}{\ensuremath{-}5.3}  % QLIKE % change, full set vs HAR benchmark: Lasso, ALL48@16 (ALL48, penalties chosen on the 16:00 validation rows; linear only), every bar
\providecommand{\rhtaDmLassoAllCtBar}{}\renewcommand{\rhtaDmLassoAllCtBar}{\ensuremath{-}6.18}  % ... its Diebold-Mariano statistic (negative: full more accurate)
\providecommand{\rhtaPLassoAllCtBar}{}\renewcommand{\rhtaPLassoAllCtBar}{\ensuremath{p < 0.001}}  % ... its p statement
\providecommand{\rhtaQHarLassoAllCtBar}{}\renewcommand{\rhtaQHarLassoAllCtBar}{0.1915}  % ... mean QLIKE, HAR benchmark
\providecommand{\rhtaQFullLassoAllCtBar}{}\renewcommand{\rhtaQFullLassoAllCtBar}{0.1814}  % ... mean QLIKE, full set
\providecommand{\rhtaPctEnetAllCtBar}{}\renewcommand{\rhtaPctEnetAllCtBar}{\ensuremath{-}5.5}  % QLIKE % change, full set vs HAR benchmark: Elastic net, ALL48@16 (ALL48, penalties chosen on the 16:00 validation rows; linear only), every bar
\providecommand{\rhtaDmEnetAllCtBar}{}\renewcommand{\rhtaDmEnetAllCtBar}{\ensuremath{-}7.11}  % ... its Diebold-Mariano statistic (negative: full more accurate)
\providecommand{\rhtaPEnetAllCtBar}{}\renewcommand{\rhtaPEnetAllCtBar}{\ensuremath{p < 0.001}}  % ... its p statement
\providecommand{\rhtaQHarEnetAllCtBar}{}\renewcommand{\rhtaQHarEnetAllCtBar}{0.1922}  % ... mean QLIKE, HAR benchmark
\providecommand{\rhtaQFullEnetAllCtBar}{}\renewcommand{\rhtaQFullEnetAllCtBar}{0.1817}  % ... mean QLIKE, full set
\providecommand{\rhtaPctRidgeRthCtClose}{}\renewcommand{\rhtaPctRidgeRthCtClose}{0.8}  % QLIKE % change, full set vs HAR benchmark: Ridge, RTH13@16 (RTH13, penalties chosen on the 16:00 validation rows; linear only), 16:00 bar
\providecommand{\rhtaDmRidgeRthCtClose}{}\renewcommand{\rhtaDmRidgeRthCtClose}{0.31}  % ... its Diebold-Mariano statistic (negative: full more accurate)
\providecommand{\rhtaPRidgeRthCtClose}{}\renewcommand{\rhtaPRidgeRthCtClose}{\ensuremath{p = 0.76}}  % ... its p statement
\providecommand{\rhtaQHarRidgeRthCtClose}{}\renewcommand{\rhtaQHarRidgeRthCtClose}{0.1269}  % ... mean QLIKE, HAR benchmark
\providecommand{\rhtaQFullRidgeRthCtClose}{}\renewcommand{\rhtaQFullRidgeRthCtClose}{0.1279}  % ... mean QLIKE, full set
\providecommand{\rhtaPctLassoRthCtClose}{}\renewcommand{\rhtaPctLassoRthCtClose}{\ensuremath{-}1.8}  % QLIKE % change, full set vs HAR benchmark: Lasso, RTH13@16 (RTH13, penalties chosen on the 16:00 validation rows; linear only), 16:00 bar
\providecommand{\rhtaDmLassoRthCtClose}{}\renewcommand{\rhtaDmLassoRthCtClose}{\ensuremath{-}0.87}  % ... its Diebold-Mariano statistic (negative: full more accurate)
\providecommand{\rhtaPLassoRthCtClose}{}\renewcommand{\rhtaPLassoRthCtClose}{\ensuremath{p = 0.38}}  % ... its p statement
\providecommand{\rhtaQHarLassoRthCtClose}{}\renewcommand{\rhtaQHarLassoRthCtClose}{0.1262}  % ... mean QLIKE, HAR benchmark
\providecommand{\rhtaQFullLassoRthCtClose}{}\renewcommand{\rhtaQFullLassoRthCtClose}{0.1239}  % ... mean QLIKE, full set
\providecommand{\rhtaPctEnetRthCtClose}{}\renewcommand{\rhtaPctEnetRthCtClose}{1.5}  % QLIKE % change, full set vs HAR benchmark: Elastic net, RTH13@16 (RTH13, penalties chosen on the 16:00 validation rows; linear only), 16:00 bar
\providecommand{\rhtaDmEnetRthCtClose}{}\renewcommand{\rhtaDmEnetRthCtClose}{0.59}  % ... its Diebold-Mariano statistic (negative: full more accurate)
\providecommand{\rhtaPEnetRthCtClose}{}\renewcommand{\rhtaPEnetRthCtClose}{\ensuremath{p = 0.56}}  % ... its p statement
\providecommand{\rhtaQHarEnetRthCtClose}{}\renewcommand{\rhtaQHarEnetRthCtClose}{0.1243}  % ... mean QLIKE, HAR benchmark
\providecommand{\rhtaQFullEnetRthCtClose}{}\renewcommand{\rhtaQFullEnetRthCtClose}{0.1262}  % ... mean QLIKE, full set
\providecommand{\rhtaPctRidgeRthCtBar}{}\renewcommand{\rhtaPctRidgeRthCtBar}{\ensuremath{-}5.4}  % QLIKE % change, full set vs HAR benchmark: Ridge, RTH13@16 (RTH13, penalties chosen on the 16:00 validation rows; linear only), every bar
\providecommand{\rhtaDmRidgeRthCtBar}{}\renewcommand{\rhtaDmRidgeRthCtBar}{\ensuremath{-}5.63}  % ... its Diebold-Mariano statistic (negative: full more accurate)
\providecommand{\rhtaPRidgeRthCtBar}{}\renewcommand{\rhtaPRidgeRthCtBar}{\ensuremath{p < 0.001}}  % ... its p statement
\providecommand{\rhtaQHarRidgeRthCtBar}{}\renewcommand{\rhtaQHarRidgeRthCtBar}{0.1893}  % ... mean QLIKE, HAR benchmark
\providecommand{\rhtaQFullRidgeRthCtBar}{}\renewcommand{\rhtaQFullRidgeRthCtBar}{0.1791}  % ... mean QLIKE, full set
\providecommand{\rhtaPctLassoRthCtBar}{}\renewcommand{\rhtaPctLassoRthCtBar}{\ensuremath{-}6.2}  % QLIKE % change, full set vs HAR benchmark: Lasso, RTH13@16 (RTH13, penalties chosen on the 16:00 validation rows; linear only), every bar
\providecommand{\rhtaDmLassoRthCtBar}{}\renewcommand{\rhtaDmLassoRthCtBar}{\ensuremath{-}8.49}  % ... its Diebold-Mariano statistic (negative: full more accurate)
\providecommand{\rhtaPLassoRthCtBar}{}\renewcommand{\rhtaPLassoRthCtBar}{\ensuremath{p < 0.001}}  % ... its p statement
\providecommand{\rhtaQHarLassoRthCtBar}{}\renewcommand{\rhtaQHarLassoRthCtBar}{0.1896}  % ... mean QLIKE, HAR benchmark
\providecommand{\rhtaQFullLassoRthCtBar}{}\renewcommand{\rhtaQFullLassoRthCtBar}{0.1777}  % ... mean QLIKE, full set
\providecommand{\rhtaPctEnetRthCtBar}{}\renewcommand{\rhtaPctEnetRthCtBar}{\ensuremath{-}6.9}  % QLIKE % change, full set vs HAR benchmark: Elastic net, RTH13@16 (RTH13, penalties chosen on the 16:00 validation rows; linear only), every bar
\providecommand{\rhtaDmEnetRthCtBar}{}\renewcommand{\rhtaDmEnetRthCtBar}{\ensuremath{-}7.94}  % ... its Diebold-Mariano statistic (negative: full more accurate)
\providecommand{\rhtaPEnetRthCtBar}{}\renewcommand{\rhtaPEnetRthCtBar}{\ensuremath{p < 0.001}}  % ... its p statement
\providecommand{\rhtaQHarEnetRthCtBar}{}\renewcommand{\rhtaQHarEnetRthCtBar}{0.1911}  % ... mean QLIKE, HAR benchmark
\providecommand{\rhtaQFullEnetRthCtBar}{}\renewcommand{\rhtaQFullEnetRthCtBar}{0.1778}  % ... mean QLIKE, full set
\providecommand{\rhtaPctRidgeAllThreeKClose}{}\renewcommand{\rhtaPctRidgeAllThreeKClose}{\ensuremath{-}4.7}  % QLIKE % change, full set vs HAR benchmark: Ridge, ALL48_N3k (ALL48, the last 3,000 rows), 16:00 bar
\providecommand{\rhtaDmRidgeAllThreeKClose}{}\renewcommand{\rhtaDmRidgeAllThreeKClose}{\ensuremath{-}4.44}  % ... its Diebold-Mariano statistic (negative: full more accurate)
\providecommand{\rhtaPRidgeAllThreeKClose}{}\renewcommand{\rhtaPRidgeAllThreeKClose}{\ensuremath{p < 0.001}}  % ... its p statement
\providecommand{\rhtaQHarRidgeAllThreeKClose}{}\renewcommand{\rhtaQHarRidgeAllThreeKClose}{0.1341}  % ... mean QLIKE, HAR benchmark
\providecommand{\rhtaQFullRidgeAllThreeKClose}{}\renewcommand{\rhtaQFullRidgeAllThreeKClose}{0.1278}  % ... mean QLIKE, full set
\providecommand{\rhtaPctLassoAllThreeKClose}{}\renewcommand{\rhtaPctLassoAllThreeKClose}{\ensuremath{-}1.6}  % QLIKE % change, full set vs HAR benchmark: Lasso, ALL48_N3k (ALL48, the last 3,000 rows), 16:00 bar
\providecommand{\rhtaDmLassoAllThreeKClose}{}\renewcommand{\rhtaDmLassoAllThreeKClose}{\ensuremath{-}1.25}  % ... its Diebold-Mariano statistic (negative: full more accurate)
\providecommand{\rhtaPLassoAllThreeKClose}{}\renewcommand{\rhtaPLassoAllThreeKClose}{\ensuremath{p = 0.21}}  % ... its p statement
\providecommand{\rhtaQHarLassoAllThreeKClose}{}\renewcommand{\rhtaQHarLassoAllThreeKClose}{0.1322}  % ... mean QLIKE, HAR benchmark
\providecommand{\rhtaQFullLassoAllThreeKClose}{}\renewcommand{\rhtaQFullLassoAllThreeKClose}{0.1301}  % ... mean QLIKE, full set
\providecommand{\rhtaPctEnetAllThreeKClose}{}\renewcommand{\rhtaPctEnetAllThreeKClose}{0.5}  % QLIKE % change, full set vs HAR benchmark: Elastic net, ALL48_N3k (ALL48, the last 3,000 rows), 16:00 bar
\providecommand{\rhtaDmEnetAllThreeKClose}{}\renewcommand{\rhtaDmEnetAllThreeKClose}{0.27}  % ... its Diebold-Mariano statistic (negative: full more accurate)
\providecommand{\rhtaPEnetAllThreeKClose}{}\renewcommand{\rhtaPEnetAllThreeKClose}{\ensuremath{p = 0.78}}  % ... its p statement
\providecommand{\rhtaQHarEnetAllThreeKClose}{}\renewcommand{\rhtaQHarEnetAllThreeKClose}{0.1331}  % ... mean QLIKE, HAR benchmark
\providecommand{\rhtaQFullEnetAllThreeKClose}{}\renewcommand{\rhtaQFullEnetAllThreeKClose}{0.1338}  % ... mean QLIKE, full set
\providecommand{\rhtaPctLgbmAllThreeKClose}{}\renewcommand{\rhtaPctLgbmAllThreeKClose}{\ensuremath{-}8.0}  % QLIKE % change, full set vs HAR benchmark: LightGBM, ALL48_N3k (ALL48, the last 3,000 rows), 16:00 bar
\providecommand{\rhtaDmLgbmAllThreeKClose}{}\renewcommand{\rhtaDmLgbmAllThreeKClose}{\ensuremath{-}3.83}  % ... its Diebold-Mariano statistic (negative: full more accurate)
\providecommand{\rhtaPLgbmAllThreeKClose}{}\renewcommand{\rhtaPLgbmAllThreeKClose}{\ensuremath{p < 0.001}}  % ... its p statement
\providecommand{\rhtaQHarLgbmAllThreeKClose}{}\renewcommand{\rhtaQHarLgbmAllThreeKClose}{0.1409}  % ... mean QLIKE, HAR benchmark
\providecommand{\rhtaQFullLgbmAllThreeKClose}{}\renewcommand{\rhtaQFullLgbmAllThreeKClose}{0.1295}  % ... mean QLIKE, full set
\providecommand{\rhtaPctXgbAllThreeKClose}{}\renewcommand{\rhtaPctXgbAllThreeKClose}{\ensuremath{-}6.6}  % QLIKE % change, full set vs HAR benchmark: XGBoost, ALL48_N3k (ALL48, the last 3,000 rows), 16:00 bar
\providecommand{\rhtaDmXgbAllThreeKClose}{}\renewcommand{\rhtaDmXgbAllThreeKClose}{\ensuremath{-}4.14}  % ... its Diebold-Mariano statistic (negative: full more accurate)
\providecommand{\rhtaPXgbAllThreeKClose}{}\renewcommand{\rhtaPXgbAllThreeKClose}{\ensuremath{p < 0.001}}  % ... its p statement
\providecommand{\rhtaQHarXgbAllThreeKClose}{}\renewcommand{\rhtaQHarXgbAllThreeKClose}{0.1348}  % ... mean QLIKE, HAR benchmark
\providecommand{\rhtaQFullXgbAllThreeKClose}{}\renewcommand{\rhtaQFullXgbAllThreeKClose}{0.1260}  % ... mean QLIKE, full set
\providecommand{\rhtaPctRfAllThreeKClose}{}\renewcommand{\rhtaPctRfAllThreeKClose}{\ensuremath{-}6.9}  % QLIKE % change, full set vs HAR benchmark: Random forest, ALL48_N3k (ALL48, the last 3,000 rows), 16:00 bar
\providecommand{\rhtaDmRfAllThreeKClose}{}\renewcommand{\rhtaDmRfAllThreeKClose}{\ensuremath{-}6.46}  % ... its Diebold-Mariano statistic (negative: full more accurate)
\providecommand{\rhtaPRfAllThreeKClose}{}\renewcommand{\rhtaPRfAllThreeKClose}{\ensuremath{p < 0.001}}  % ... its p statement
\providecommand{\rhtaQHarRfAllThreeKClose}{}\renewcommand{\rhtaQHarRfAllThreeKClose}{0.1355}  % ... mean QLIKE, HAR benchmark
\providecommand{\rhtaQFullRfAllThreeKClose}{}\renewcommand{\rhtaQFullRfAllThreeKClose}{0.1260}  % ... mean QLIKE, full set
\providecommand{\rhtaPctRidgeAllThreeKBar}{}\renewcommand{\rhtaPctRidgeAllThreeKBar}{\ensuremath{-}2.7}  % QLIKE % change, full set vs HAR benchmark: Ridge, ALL48_N3k (ALL48, the last 3,000 rows), every bar
\providecommand{\rhtaDmRidgeAllThreeKBar}{}\renewcommand{\rhtaDmRidgeAllThreeKBar}{\ensuremath{-}7.15}  % ... its Diebold-Mariano statistic (negative: full more accurate)
\providecommand{\rhtaPRidgeAllThreeKBar}{}\renewcommand{\rhtaPRidgeAllThreeKBar}{\ensuremath{p < 0.001}}  % ... its p statement
\providecommand{\rhtaQHarRidgeAllThreeKBar}{}\renewcommand{\rhtaQHarRidgeAllThreeKBar}{0.1976}  % ... mean QLIKE, HAR benchmark
\providecommand{\rhtaQFullRidgeAllThreeKBar}{}\renewcommand{\rhtaQFullRidgeAllThreeKBar}{0.1924}  % ... mean QLIKE, full set
\providecommand{\rhtaPctLassoAllThreeKBar}{}\renewcommand{\rhtaPctLassoAllThreeKBar}{\ensuremath{-}2.1}  % QLIKE % change, full set vs HAR benchmark: Lasso, ALL48_N3k (ALL48, the last 3,000 rows), every bar
\providecommand{\rhtaDmLassoAllThreeKBar}{}\renewcommand{\rhtaDmLassoAllThreeKBar}{\ensuremath{-}4.17}  % ... its Diebold-Mariano statistic (negative: full more accurate)
\providecommand{\rhtaPLassoAllThreeKBar}{}\renewcommand{\rhtaPLassoAllThreeKBar}{\ensuremath{p < 0.001}}  % ... its p statement
\providecommand{\rhtaQHarLassoAllThreeKBar}{}\renewcommand{\rhtaQHarLassoAllThreeKBar}{0.1967}  % ... mean QLIKE, HAR benchmark
\providecommand{\rhtaQFullLassoAllThreeKBar}{}\renewcommand{\rhtaQFullLassoAllThreeKBar}{0.1927}  % ... mean QLIKE, full set
\providecommand{\rhtaPctEnetAllThreeKBar}{}\renewcommand{\rhtaPctEnetAllThreeKBar}{\ensuremath{-}1.8}  % QLIKE % change, full set vs HAR benchmark: Elastic net, ALL48_N3k (ALL48, the last 3,000 rows), every bar
\providecommand{\rhtaDmEnetAllThreeKBar}{}\renewcommand{\rhtaDmEnetAllThreeKBar}{\ensuremath{-}2.57}  % ... its Diebold-Mariano statistic (negative: full more accurate)
\providecommand{\rhtaPEnetAllThreeKBar}{}\renewcommand{\rhtaPEnetAllThreeKBar}{\ensuremath{p = 0.0101}}  % ... its p statement
\providecommand{\rhtaQHarEnetAllThreeKBar}{}\renewcommand{\rhtaQHarEnetAllThreeKBar}{0.1972}  % ... mean QLIKE, HAR benchmark
\providecommand{\rhtaQFullEnetAllThreeKBar}{}\renewcommand{\rhtaQFullEnetAllThreeKBar}{0.1936}  % ... mean QLIKE, full set
\providecommand{\rhtaPctLgbmAllThreeKBar}{}\renewcommand{\rhtaPctLgbmAllThreeKBar}{\ensuremath{-}5.9}  % QLIKE % change, full set vs HAR benchmark: LightGBM, ALL48_N3k (ALL48, the last 3,000 rows), every bar
\providecommand{\rhtaDmLgbmAllThreeKBar}{}\renewcommand{\rhtaDmLgbmAllThreeKBar}{\ensuremath{-}8.86}  % ... its Diebold-Mariano statistic (negative: full more accurate)
\providecommand{\rhtaPLgbmAllThreeKBar}{}\renewcommand{\rhtaPLgbmAllThreeKBar}{\ensuremath{p < 0.001}}  % ... its p statement
\providecommand{\rhtaQHarLgbmAllThreeKBar}{}\renewcommand{\rhtaQHarLgbmAllThreeKBar}{0.2062}  % ... mean QLIKE, HAR benchmark
\providecommand{\rhtaQFullLgbmAllThreeKBar}{}\renewcommand{\rhtaQFullLgbmAllThreeKBar}{0.1941}  % ... mean QLIKE, full set
\providecommand{\rhtaPctXgbAllThreeKBar}{}\renewcommand{\rhtaPctXgbAllThreeKBar}{\ensuremath{-}5.9}  % QLIKE % change, full set vs HAR benchmark: XGBoost, ALL48_N3k (ALL48, the last 3,000 rows), every bar
\providecommand{\rhtaDmXgbAllThreeKBar}{}\renewcommand{\rhtaDmXgbAllThreeKBar}{\ensuremath{-}8.32}  % ... its Diebold-Mariano statistic (negative: full more accurate)
\providecommand{\rhtaPXgbAllThreeKBar}{}\renewcommand{\rhtaPXgbAllThreeKBar}{\ensuremath{p < 0.001}}  % ... its p statement
\providecommand{\rhtaQHarXgbAllThreeKBar}{}\renewcommand{\rhtaQHarXgbAllThreeKBar}{0.2011}  % ... mean QLIKE, HAR benchmark
\providecommand{\rhtaQFullXgbAllThreeKBar}{}\renewcommand{\rhtaQFullXgbAllThreeKBar}{0.1893}  % ... mean QLIKE, full set
\providecommand{\rhtaPctRfAllThreeKBar}{}\renewcommand{\rhtaPctRfAllThreeKBar}{\ensuremath{-}5.0}  % QLIKE % change, full set vs HAR benchmark: Random forest, ALL48_N3k (ALL48, the last 3,000 rows), every bar
\providecommand{\rhtaDmRfAllThreeKBar}{}\renewcommand{\rhtaDmRfAllThreeKBar}{\ensuremath{-}8.83}  % ... its Diebold-Mariano statistic (negative: full more accurate)
\providecommand{\rhtaPRfAllThreeKBar}{}\renewcommand{\rhtaPRfAllThreeKBar}{\ensuremath{p < 0.001}}  % ... its p statement
\providecommand{\rhtaQHarRfAllThreeKBar}{}\renewcommand{\rhtaQHarRfAllThreeKBar}{0.1994}  % ... mean QLIKE, HAR benchmark
\providecommand{\rhtaQFullRfAllThreeKBar}{}\renewcommand{\rhtaQFullRfAllThreeKBar}{0.1894}  % ... mean QLIKE, full set
\providecommand{\rhtaPctRidgeRthThreeKClose}{}\renewcommand{\rhtaPctRidgeRthThreeKClose}{\ensuremath{-}1.9}  % QLIKE % change, full set vs HAR benchmark: Ridge, RTH13_N3k (RTH13, the last 3,000 rows), 16:00 bar
\providecommand{\rhtaDmRidgeRthThreeKClose}{}\renewcommand{\rhtaDmRidgeRthThreeKClose}{\ensuremath{-}1.55}  % ... its Diebold-Mariano statistic (negative: full more accurate)
\providecommand{\rhtaPRidgeRthThreeKClose}{}\renewcommand{\rhtaPRidgeRthThreeKClose}{\ensuremath{p = 0.12}}  % ... its p statement
\providecommand{\rhtaQHarRidgeRthThreeKClose}{}\renewcommand{\rhtaQHarRidgeRthThreeKClose}{0.1418}  % ... mean QLIKE, HAR benchmark
\providecommand{\rhtaQFullRidgeRthThreeKClose}{}\renewcommand{\rhtaQFullRidgeRthThreeKClose}{0.1391}  % ... mean QLIKE, full set
\providecommand{\rhtaPctLassoRthThreeKClose}{}\renewcommand{\rhtaPctLassoRthThreeKClose}{\ensuremath{-}1.9}  % QLIKE % change, full set vs HAR benchmark: Lasso, RTH13_N3k (RTH13, the last 3,000 rows), 16:00 bar
\providecommand{\rhtaDmLassoRthThreeKClose}{}\renewcommand{\rhtaDmLassoRthThreeKClose}{\ensuremath{-}1.29}  % ... its Diebold-Mariano statistic (negative: full more accurate)
\providecommand{\rhtaPLassoRthThreeKClose}{}\renewcommand{\rhtaPLassoRthThreeKClose}{\ensuremath{p = 0.20}}  % ... its p statement
\providecommand{\rhtaQHarLassoRthThreeKClose}{}\renewcommand{\rhtaQHarLassoRthThreeKClose}{0.1418}  % ... mean QLIKE, HAR benchmark
\providecommand{\rhtaQFullLassoRthThreeKClose}{}\renewcommand{\rhtaQFullLassoRthThreeKClose}{0.1392}  % ... mean QLIKE, full set
\providecommand{\rhtaPctEnetRthThreeKClose}{}\renewcommand{\rhtaPctEnetRthThreeKClose}{\ensuremath{-}1.6}  % QLIKE % change, full set vs HAR benchmark: Elastic net, RTH13_N3k (RTH13, the last 3,000 rows), 16:00 bar
\providecommand{\rhtaDmEnetRthThreeKClose}{}\renewcommand{\rhtaDmEnetRthThreeKClose}{\ensuremath{-}0.85}  % ... its Diebold-Mariano statistic (negative: full more accurate)
\providecommand{\rhtaPEnetRthThreeKClose}{}\renewcommand{\rhtaPEnetRthThreeKClose}{\ensuremath{p = 0.39}}  % ... its p statement
\providecommand{\rhtaQHarEnetRthThreeKClose}{}\renewcommand{\rhtaQHarEnetRthThreeKClose}{0.1418}  % ... mean QLIKE, HAR benchmark
\providecommand{\rhtaQFullEnetRthThreeKClose}{}\renewcommand{\rhtaQFullEnetRthThreeKClose}{0.1396}  % ... mean QLIKE, full set
\providecommand{\rhtaPctLgbmRthThreeKClose}{}\renewcommand{\rhtaPctLgbmRthThreeKClose}{\ensuremath{-}3.7}  % QLIKE % change, full set vs HAR benchmark: LightGBM, RTH13_N3k (RTH13, the last 3,000 rows), 16:00 bar
\providecommand{\rhtaDmLgbmRthThreeKClose}{}\renewcommand{\rhtaDmLgbmRthThreeKClose}{\ensuremath{-}2.15}  % ... its Diebold-Mariano statistic (negative: full more accurate)
\providecommand{\rhtaPLgbmRthThreeKClose}{}\renewcommand{\rhtaPLgbmRthThreeKClose}{\ensuremath{p = 0.031}}  % ... its p statement
\providecommand{\rhtaQHarLgbmRthThreeKClose}{}\renewcommand{\rhtaQHarLgbmRthThreeKClose}{0.1401}  % ... mean QLIKE, HAR benchmark
\providecommand{\rhtaQFullLgbmRthThreeKClose}{}\renewcommand{\rhtaQFullLgbmRthThreeKClose}{0.1349}  % ... mean QLIKE, full set
\providecommand{\rhtaPctXgbRthThreeKClose}{}\renewcommand{\rhtaPctXgbRthThreeKClose}{\ensuremath{-}2.3}  % QLIKE % change, full set vs HAR benchmark: XGBoost, RTH13_N3k (RTH13, the last 3,000 rows), 16:00 bar
\providecommand{\rhtaDmXgbRthThreeKClose}{}\renewcommand{\rhtaDmXgbRthThreeKClose}{\ensuremath{-}1.26}  % ... its Diebold-Mariano statistic (negative: full more accurate)
\providecommand{\rhtaPXgbRthThreeKClose}{}\renewcommand{\rhtaPXgbRthThreeKClose}{\ensuremath{p = 0.21}}  % ... its p statement
\providecommand{\rhtaQHarXgbRthThreeKClose}{}\renewcommand{\rhtaQHarXgbRthThreeKClose}{0.1334}  % ... mean QLIKE, HAR benchmark
\providecommand{\rhtaQFullXgbRthThreeKClose}{}\renewcommand{\rhtaQFullXgbRthThreeKClose}{0.1303}  % ... mean QLIKE, full set
\providecommand{\rhtaPctRfRthThreeKClose}{}\renewcommand{\rhtaPctRfRthThreeKClose}{\ensuremath{-}4.5}  % QLIKE % change, full set vs HAR benchmark: Random forest, RTH13_N3k (RTH13, the last 3,000 rows), 16:00 bar
\providecommand{\rhtaDmRfRthThreeKClose}{}\renewcommand{\rhtaDmRfRthThreeKClose}{\ensuremath{-}3.54}  % ... its Diebold-Mariano statistic (negative: full more accurate)
\providecommand{\rhtaPRfRthThreeKClose}{}\renewcommand{\rhtaPRfRthThreeKClose}{\ensuremath{p < 0.001}}  % ... its p statement
\providecommand{\rhtaQHarRfRthThreeKClose}{}\renewcommand{\rhtaQHarRfRthThreeKClose}{0.1400}  % ... mean QLIKE, HAR benchmark
\providecommand{\rhtaQFullRfRthThreeKClose}{}\renewcommand{\rhtaQFullRfRthThreeKClose}{0.1336}  % ... mean QLIKE, full set
\providecommand{\rhtaPctRidgeRthThreeKBar}{}\renewcommand{\rhtaPctRidgeRthThreeKBar}{\ensuremath{-}3.6}  % QLIKE % change, full set vs HAR benchmark: Ridge, RTH13_N3k (RTH13, the last 3,000 rows), every bar
\providecommand{\rhtaDmRidgeRthThreeKBar}{}\renewcommand{\rhtaDmRidgeRthThreeKBar}{\ensuremath{-}9.44}  % ... its Diebold-Mariano statistic (negative: full more accurate)
\providecommand{\rhtaPRidgeRthThreeKBar}{}\renewcommand{\rhtaPRidgeRthThreeKBar}{\ensuremath{p < 0.001}}  % ... its p statement
\providecommand{\rhtaQHarRidgeRthThreeKBar}{}\renewcommand{\rhtaQHarRidgeRthThreeKBar}{0.1910}  % ... mean QLIKE, HAR benchmark
\providecommand{\rhtaQFullRidgeRthThreeKBar}{}\renewcommand{\rhtaQFullRidgeRthThreeKBar}{0.1840}  % ... mean QLIKE, full set
\providecommand{\rhtaPctLassoRthThreeKBar}{}\renewcommand{\rhtaPctLassoRthThreeKBar}{\ensuremath{-}3.8}  % QLIKE % change, full set vs HAR benchmark: Lasso, RTH13_N3k (RTH13, the last 3,000 rows), every bar
\providecommand{\rhtaDmLassoRthThreeKBar}{}\renewcommand{\rhtaDmLassoRthThreeKBar}{\ensuremath{-}7.70}  % ... its Diebold-Mariano statistic (negative: full more accurate)
\providecommand{\rhtaPLassoRthThreeKBar}{}\renewcommand{\rhtaPLassoRthThreeKBar}{\ensuremath{p < 0.001}}  % ... its p statement
\providecommand{\rhtaQHarLassoRthThreeKBar}{}\renewcommand{\rhtaQHarLassoRthThreeKBar}{0.1910}  % ... mean QLIKE, HAR benchmark
\providecommand{\rhtaQFullLassoRthThreeKBar}{}\renewcommand{\rhtaQFullLassoRthThreeKBar}{0.1837}  % ... mean QLIKE, full set
\providecommand{\rhtaPctEnetRthThreeKBar}{}\renewcommand{\rhtaPctEnetRthThreeKBar}{\ensuremath{-}4.2}  % QLIKE % change, full set vs HAR benchmark: Elastic net, RTH13_N3k (RTH13, the last 3,000 rows), every bar
\providecommand{\rhtaDmEnetRthThreeKBar}{}\renewcommand{\rhtaDmEnetRthThreeKBar}{\ensuremath{-}6.30}  % ... its Diebold-Mariano statistic (negative: full more accurate)
\providecommand{\rhtaPEnetRthThreeKBar}{}\renewcommand{\rhtaPEnetRthThreeKBar}{\ensuremath{p < 0.001}}  % ... its p statement
\providecommand{\rhtaQHarEnetRthThreeKBar}{}\renewcommand{\rhtaQHarEnetRthThreeKBar}{0.1910}  % ... mean QLIKE, HAR benchmark
\providecommand{\rhtaQFullEnetRthThreeKBar}{}\renewcommand{\rhtaQFullEnetRthThreeKBar}{0.1830}  % ... mean QLIKE, full set
\providecommand{\rhtaPctLgbmRthThreeKBar}{}\renewcommand{\rhtaPctLgbmRthThreeKBar}{\ensuremath{-}6.0}  % QLIKE % change, full set vs HAR benchmark: LightGBM, RTH13_N3k (RTH13, the last 3,000 rows), every bar
\providecommand{\rhtaDmLgbmRthThreeKBar}{}\renewcommand{\rhtaDmLgbmRthThreeKBar}{\ensuremath{-}7.85}  % ... its Diebold-Mariano statistic (negative: full more accurate)
\providecommand{\rhtaPLgbmRthThreeKBar}{}\renewcommand{\rhtaPLgbmRthThreeKBar}{\ensuremath{p < 0.001}}  % ... its p statement
\providecommand{\rhtaQHarLgbmRthThreeKBar}{}\renewcommand{\rhtaQHarLgbmRthThreeKBar}{0.1972}  % ... mean QLIKE, HAR benchmark
\providecommand{\rhtaQFullLgbmRthThreeKBar}{}\renewcommand{\rhtaQFullLgbmRthThreeKBar}{0.1854}  % ... mean QLIKE, full set
\providecommand{\rhtaPctXgbRthThreeKBar}{}\renewcommand{\rhtaPctXgbRthThreeKBar}{\ensuremath{-}6.8}  % QLIKE % change, full set vs HAR benchmark: XGBoost, RTH13_N3k (RTH13, the last 3,000 rows), every bar
\providecommand{\rhtaDmXgbRthThreeKBar}{}\renewcommand{\rhtaDmXgbRthThreeKBar}{\ensuremath{-}9.13}  % ... its Diebold-Mariano statistic (negative: full more accurate)
\providecommand{\rhtaPXgbRthThreeKBar}{}\renewcommand{\rhtaPXgbRthThreeKBar}{\ensuremath{p < 0.001}}  % ... its p statement
\providecommand{\rhtaQHarXgbRthThreeKBar}{}\renewcommand{\rhtaQHarXgbRthThreeKBar}{0.1931}  % ... mean QLIKE, HAR benchmark
\providecommand{\rhtaQFullXgbRthThreeKBar}{}\renewcommand{\rhtaQFullXgbRthThreeKBar}{0.1801}  % ... mean QLIKE, full set
\providecommand{\rhtaPctRfRthThreeKBar}{}\renewcommand{\rhtaPctRfRthThreeKBar}{\ensuremath{-}4.7}  % QLIKE % change, full set vs HAR benchmark: Random forest, RTH13_N3k (RTH13, the last 3,000 rows), every bar
\providecommand{\rhtaDmRfRthThreeKBar}{}\renewcommand{\rhtaDmRfRthThreeKBar}{\ensuremath{-}8.98}  % ... its Diebold-Mariano statistic (negative: full more accurate)
\providecommand{\rhtaPRfRthThreeKBar}{}\renewcommand{\rhtaPRfRthThreeKBar}{\ensuremath{p < 0.001}}  % ... its p statement
\providecommand{\rhtaQHarRfRthThreeKBar}{}\renewcommand{\rhtaQHarRfRthThreeKBar}{0.1935}  % ... mean QLIKE, HAR benchmark
\providecommand{\rhtaQFullRfRthThreeKBar}{}\renewcommand{\rhtaQFullRfRthThreeKBar}{0.1844}  % ... mean QLIKE, full set
\providecommand{\rhtaPctRidgeAllThirtyNineKClose}{}\renewcommand{\rhtaPctRidgeAllThirtyNineKClose}{\ensuremath{-}0.9}  % QLIKE % change, full set vs HAR benchmark: Ridge, ALL48_N39k (ALL48, the last 39,000 rows), 16:00 bar
\providecommand{\rhtaDmRidgeAllThirtyNineKClose}{}\renewcommand{\rhtaDmRidgeAllThirtyNineKClose}{\ensuremath{-}0.41}  % ... its Diebold-Mariano statistic (negative: full more accurate)
\providecommand{\rhtaPRidgeAllThirtyNineKClose}{}\renewcommand{\rhtaPRidgeAllThirtyNineKClose}{\ensuremath{p = 0.68}}  % ... its p statement
\providecommand{\rhtaQHarRidgeAllThirtyNineKClose}{}\renewcommand{\rhtaQHarRidgeAllThirtyNineKClose}{0.1296}  % ... mean QLIKE, HAR benchmark
\providecommand{\rhtaQFullRidgeAllThirtyNineKClose}{}\renewcommand{\rhtaQFullRidgeAllThirtyNineKClose}{0.1285}  % ... mean QLIKE, full set
\providecommand{\rhtaPctLassoAllThirtyNineKClose}{}\renewcommand{\rhtaPctLassoAllThirtyNineKClose}{\ensuremath{-}3.6}  % QLIKE % change, full set vs HAR benchmark: Lasso, ALL48_N39k (ALL48, the last 39,000 rows), 16:00 bar
\providecommand{\rhtaDmLassoAllThirtyNineKClose}{}\renewcommand{\rhtaDmLassoAllThirtyNineKClose}{\ensuremath{-}2.19}  % ... its Diebold-Mariano statistic (negative: full more accurate)
\providecommand{\rhtaPLassoAllThirtyNineKClose}{}\renewcommand{\rhtaPLassoAllThirtyNineKClose}{\ensuremath{p = 0.028}}  % ... its p statement
\providecommand{\rhtaQHarLassoAllThirtyNineKClose}{}\renewcommand{\rhtaQHarLassoAllThirtyNineKClose}{0.1311}  % ... mean QLIKE, HAR benchmark
\providecommand{\rhtaQFullLassoAllThirtyNineKClose}{}\renewcommand{\rhtaQFullLassoAllThirtyNineKClose}{0.1264}  % ... mean QLIKE, full set
\providecommand{\rhtaPctEnetAllThirtyNineKClose}{}\renewcommand{\rhtaPctEnetAllThirtyNineKClose}{\ensuremath{-}3.7}  % QLIKE % change, full set vs HAR benchmark: Elastic net, ALL48_N39k (ALL48, the last 39,000 rows), 16:00 bar
\providecommand{\rhtaDmEnetAllThirtyNineKClose}{}\renewcommand{\rhtaDmEnetAllThirtyNineKClose}{\ensuremath{-}3.22}  % ... its Diebold-Mariano statistic (negative: full more accurate)
\providecommand{\rhtaPEnetAllThirtyNineKClose}{}\renewcommand{\rhtaPEnetAllThirtyNineKClose}{\ensuremath{p = 0.001}}  % ... its p statement
\providecommand{\rhtaQHarEnetAllThirtyNineKClose}{}\renewcommand{\rhtaQHarEnetAllThirtyNineKClose}{0.1312}  % ... mean QLIKE, HAR benchmark
\providecommand{\rhtaQFullEnetAllThirtyNineKClose}{}\renewcommand{\rhtaQFullEnetAllThirtyNineKClose}{0.1263}  % ... mean QLIKE, full set
\providecommand{\rhtaPctLgbmAllThirtyNineKClose}{}\renewcommand{\rhtaPctLgbmAllThirtyNineKClose}{\ensuremath{-}5.5}  % QLIKE % change, full set vs HAR benchmark: LightGBM, ALL48_N39k (ALL48, the last 39,000 rows), 16:00 bar
\providecommand{\rhtaDmLgbmAllThirtyNineKClose}{}\renewcommand{\rhtaDmLgbmAllThirtyNineKClose}{\ensuremath{-}4.12}  % ... its Diebold-Mariano statistic (negative: full more accurate)
\providecommand{\rhtaPLgbmAllThirtyNineKClose}{}\renewcommand{\rhtaPLgbmAllThirtyNineKClose}{\ensuremath{p < 0.001}}  % ... its p statement
\providecommand{\rhtaQHarLgbmAllThirtyNineKClose}{}\renewcommand{\rhtaQHarLgbmAllThirtyNineKClose}{0.1276}  % ... mean QLIKE, HAR benchmark
\providecommand{\rhtaQFullLgbmAllThirtyNineKClose}{}\renewcommand{\rhtaQFullLgbmAllThirtyNineKClose}{0.1206}  % ... mean QLIKE, full set
\providecommand{\rhtaPctXgbAllThirtyNineKClose}{}\renewcommand{\rhtaPctXgbAllThirtyNineKClose}{\ensuremath{-}4.7}  % QLIKE % change, full set vs HAR benchmark: XGBoost, ALL48_N39k (ALL48, the last 39,000 rows), 16:00 bar
\providecommand{\rhtaDmXgbAllThirtyNineKClose}{}\renewcommand{\rhtaDmXgbAllThirtyNineKClose}{\ensuremath{-}3.66}  % ... its Diebold-Mariano statistic (negative: full more accurate)
\providecommand{\rhtaPXgbAllThirtyNineKClose}{}\renewcommand{\rhtaPXgbAllThirtyNineKClose}{\ensuremath{p < 0.001}}  % ... its p statement
\providecommand{\rhtaQHarXgbAllThirtyNineKClose}{}\renewcommand{\rhtaQHarXgbAllThirtyNineKClose}{0.1267}  % ... mean QLIKE, HAR benchmark
\providecommand{\rhtaQFullXgbAllThirtyNineKClose}{}\renewcommand{\rhtaQFullXgbAllThirtyNineKClose}{0.1208}  % ... mean QLIKE, full set
\providecommand{\rhtaPctRfAllThirtyNineKClose}{}\renewcommand{\rhtaPctRfAllThirtyNineKClose}{\ensuremath{-}4.9}  % QLIKE % change, full set vs HAR benchmark: Random forest, ALL48_N39k (ALL48, the last 39,000 rows), 16:00 bar
\providecommand{\rhtaDmRfAllThirtyNineKClose}{}\renewcommand{\rhtaDmRfAllThirtyNineKClose}{\ensuremath{-}3.89}  % ... its Diebold-Mariano statistic (negative: full more accurate)
\providecommand{\rhtaPRfAllThirtyNineKClose}{}\renewcommand{\rhtaPRfAllThirtyNineKClose}{\ensuremath{p < 0.001}}  % ... its p statement
\providecommand{\rhtaQHarRfAllThirtyNineKClose}{}\renewcommand{\rhtaQHarRfAllThirtyNineKClose}{0.1277}  % ... mean QLIKE, HAR benchmark
\providecommand{\rhtaQFullRfAllThirtyNineKClose}{}\renewcommand{\rhtaQFullRfAllThirtyNineKClose}{0.1214}  % ... mean QLIKE, full set
\providecommand{\rhtaPctRidgeAllThirtyNineKBar}{}\renewcommand{\rhtaPctRidgeAllThirtyNineKBar}{\ensuremath{-}4.4}  % QLIKE % change, full set vs HAR benchmark: Ridge, ALL48_N39k (ALL48, the last 39,000 rows), every bar
\providecommand{\rhtaDmRidgeAllThirtyNineKBar}{}\renewcommand{\rhtaDmRidgeAllThirtyNineKBar}{\ensuremath{-}4.51}  % ... its Diebold-Mariano statistic (negative: full more accurate)
\providecommand{\rhtaPRidgeAllThirtyNineKBar}{}\renewcommand{\rhtaPRidgeAllThirtyNineKBar}{\ensuremath{p < 0.001}}  % ... its p statement
\providecommand{\rhtaQHarRidgeAllThirtyNineKBar}{}\renewcommand{\rhtaQHarRidgeAllThirtyNineKBar}{0.1925}  % ... mean QLIKE, HAR benchmark
\providecommand{\rhtaQFullRidgeAllThirtyNineKBar}{}\renewcommand{\rhtaQFullRidgeAllThirtyNineKBar}{0.1840}  % ... mean QLIKE, full set
\providecommand{\rhtaPctLassoAllThirtyNineKBar}{}\renewcommand{\rhtaPctLassoAllThirtyNineKBar}{\ensuremath{-}3.8}  % QLIKE % change, full set vs HAR benchmark: Lasso, ALL48_N39k (ALL48, the last 39,000 rows), every bar
\providecommand{\rhtaDmLassoAllThirtyNineKBar}{}\renewcommand{\rhtaDmLassoAllThirtyNineKBar}{\ensuremath{-}7.03}  % ... its Diebold-Mariano statistic (negative: full more accurate)
\providecommand{\rhtaPLassoAllThirtyNineKBar}{}\renewcommand{\rhtaPLassoAllThirtyNineKBar}{\ensuremath{p < 0.001}}  % ... its p statement
\providecommand{\rhtaQHarLassoAllThirtyNineKBar}{}\renewcommand{\rhtaQHarLassoAllThirtyNineKBar}{0.1922}  % ... mean QLIKE, HAR benchmark
\providecommand{\rhtaQFullLassoAllThirtyNineKBar}{}\renewcommand{\rhtaQFullLassoAllThirtyNineKBar}{0.1849}  % ... mean QLIKE, full set
\providecommand{\rhtaPctEnetAllThirtyNineKBar}{}\renewcommand{\rhtaPctEnetAllThirtyNineKBar}{\ensuremath{-}3.0}  % QLIKE % change, full set vs HAR benchmark: Elastic net, ALL48_N39k (ALL48, the last 39,000 rows), every bar
\providecommand{\rhtaDmEnetAllThirtyNineKBar}{}\renewcommand{\rhtaDmEnetAllThirtyNineKBar}{\ensuremath{-}8.37}  % ... its Diebold-Mariano statistic (negative: full more accurate)
\providecommand{\rhtaPEnetAllThirtyNineKBar}{}\renewcommand{\rhtaPEnetAllThirtyNineKBar}{\ensuremath{p < 0.001}}  % ... its p statement
\providecommand{\rhtaQHarEnetAllThirtyNineKBar}{}\renewcommand{\rhtaQHarEnetAllThirtyNineKBar}{0.1922}  % ... mean QLIKE, HAR benchmark
\providecommand{\rhtaQFullEnetAllThirtyNineKBar}{}\renewcommand{\rhtaQFullEnetAllThirtyNineKBar}{0.1865}  % ... mean QLIKE, full set
\providecommand{\rhtaPctLgbmAllThirtyNineKBar}{}\renewcommand{\rhtaPctLgbmAllThirtyNineKBar}{\ensuremath{-}7.2}  % QLIKE % change, full set vs HAR benchmark: LightGBM, ALL48_N39k (ALL48, the last 39,000 rows), every bar
\providecommand{\rhtaDmLgbmAllThirtyNineKBar}{}\renewcommand{\rhtaDmLgbmAllThirtyNineKBar}{\ensuremath{-}10.73}  % ... its Diebold-Mariano statistic (negative: full more accurate)
\providecommand{\rhtaPLgbmAllThirtyNineKBar}{}\renewcommand{\rhtaPLgbmAllThirtyNineKBar}{\ensuremath{p < 0.001}}  % ... its p statement
\providecommand{\rhtaQHarLgbmAllThirtyNineKBar}{}\renewcommand{\rhtaQHarLgbmAllThirtyNineKBar}{0.1920}  % ... mean QLIKE, HAR benchmark
\providecommand{\rhtaQFullLgbmAllThirtyNineKBar}{}\renewcommand{\rhtaQFullLgbmAllThirtyNineKBar}{0.1781}  % ... mean QLIKE, full set
\providecommand{\rhtaPctXgbAllThirtyNineKBar}{}\renewcommand{\rhtaPctXgbAllThirtyNineKBar}{\ensuremath{-}10.1}  % QLIKE % change, full set vs HAR benchmark: XGBoost, ALL48_N39k (ALL48, the last 39,000 rows), every bar
\providecommand{\rhtaDmXgbAllThirtyNineKBar}{}\renewcommand{\rhtaDmXgbAllThirtyNineKBar}{\ensuremath{-}10.02}  % ... its Diebold-Mariano statistic (negative: full more accurate)
\providecommand{\rhtaPXgbAllThirtyNineKBar}{}\renewcommand{\rhtaPXgbAllThirtyNineKBar}{\ensuremath{p < 0.001}}  % ... its p statement
\providecommand{\rhtaQHarXgbAllThirtyNineKBar}{}\renewcommand{\rhtaQHarXgbAllThirtyNineKBar}{0.1927}  % ... mean QLIKE, HAR benchmark
\providecommand{\rhtaQFullXgbAllThirtyNineKBar}{}\renewcommand{\rhtaQFullXgbAllThirtyNineKBar}{0.1732}  % ... mean QLIKE, full set
\providecommand{\rhtaPctRfAllThirtyNineKBar}{}\renewcommand{\rhtaPctRfAllThirtyNineKBar}{\ensuremath{-}5.6}  % QLIKE % change, full set vs HAR benchmark: Random forest, ALL48_N39k (ALL48, the last 39,000 rows), every bar
\providecommand{\rhtaDmRfAllThirtyNineKBar}{}\renewcommand{\rhtaDmRfAllThirtyNineKBar}{\ensuremath{-}12.63}  % ... its Diebold-Mariano statistic (negative: full more accurate)
\providecommand{\rhtaPRfAllThirtyNineKBar}{}\renewcommand{\rhtaPRfAllThirtyNineKBar}{\ensuremath{p < 0.001}}  % ... its p statement
\providecommand{\rhtaQHarRfAllThirtyNineKBar}{}\renewcommand{\rhtaQHarRfAllThirtyNineKBar}{0.1926}  % ... mean QLIKE, HAR benchmark
\providecommand{\rhtaQFullRfAllThirtyNineKBar}{}\renewcommand{\rhtaQFullRfAllThirtyNineKBar}{0.1818}  % ... mean QLIKE, full set
\providecommand{\rhtaXPctRidgeAllVsRthClose}{}\renewcommand{\rhtaXPctRidgeAllVsRthClose}{\ensuremath{-}3.5}  % QLIKE % change, ALL48 vs RTH13 (full set; negative: ALL48 more accurate), Ridge, 16:00 bar
\providecommand{\rhtaXDmRidgeAllVsRthClose}{}\renewcommand{\rhtaXDmRidgeAllVsRthClose}{\ensuremath{-}2.72}  % ... its Diebold-Mariano statistic
\providecommand{\rhtaXPRidgeAllVsRthClose}{}\renewcommand{\rhtaXPRidgeAllVsRthClose}{\ensuremath{p = 0.006}}  % ... its p statement
\providecommand{\rhtaXPctLassoAllVsRthClose}{}\renewcommand{\rhtaXPctLassoAllVsRthClose}{\ensuremath{-}6.0}  % QLIKE % change, ALL48 vs RTH13 (full set; negative: ALL48 more accurate), Lasso, 16:00 bar
\providecommand{\rhtaXDmLassoAllVsRthClose}{}\renewcommand{\rhtaXDmLassoAllVsRthClose}{\ensuremath{-}3.31}  % ... its Diebold-Mariano statistic
\providecommand{\rhtaXPLassoAllVsRthClose}{}\renewcommand{\rhtaXPLassoAllVsRthClose}{\ensuremath{p < 0.001}}  % ... its p statement
\providecommand{\rhtaXPctEnetAllVsRthClose}{}\renewcommand{\rhtaXPctEnetAllVsRthClose}{\ensuremath{-}2.1}  % QLIKE % change, ALL48 vs RTH13 (full set; negative: ALL48 more accurate), Elastic net, 16:00 bar
\providecommand{\rhtaXDmEnetAllVsRthClose}{}\renewcommand{\rhtaXDmEnetAllVsRthClose}{\ensuremath{-}1.64}  % ... its Diebold-Mariano statistic
\providecommand{\rhtaXPEnetAllVsRthClose}{}\renewcommand{\rhtaXPEnetAllVsRthClose}{\ensuremath{p = 0.1006}}  % ... its p statement
\providecommand{\rhtaXPctLgbmAllVsRthClose}{}\renewcommand{\rhtaXPctLgbmAllVsRthClose}{\ensuremath{-}2.3}  % QLIKE % change, ALL48 vs RTH13 (full set; negative: ALL48 more accurate), LightGBM, 16:00 bar
\providecommand{\rhtaXDmLgbmAllVsRthClose}{}\renewcommand{\rhtaXDmLgbmAllVsRthClose}{\ensuremath{-}2.16}  % ... its Diebold-Mariano statistic
\providecommand{\rhtaXPLgbmAllVsRthClose}{}\renewcommand{\rhtaXPLgbmAllVsRthClose}{\ensuremath{p = 0.030}}  % ... its p statement
\providecommand{\rhtaXPctXgbAllVsRthClose}{}\renewcommand{\rhtaXPctXgbAllVsRthClose}{\ensuremath{-}1.9}  % QLIKE % change, ALL48 vs RTH13 (full set; negative: ALL48 more accurate), XGBoost, 16:00 bar
\providecommand{\rhtaXDmXgbAllVsRthClose}{}\renewcommand{\rhtaXDmXgbAllVsRthClose}{\ensuremath{-}2.21}  % ... its Diebold-Mariano statistic
\providecommand{\rhtaXPXgbAllVsRthClose}{}\renewcommand{\rhtaXPXgbAllVsRthClose}{\ensuremath{p = 0.027}}  % ... its p statement
\providecommand{\rhtaXPctRfAllVsRthClose}{}\renewcommand{\rhtaXPctRfAllVsRthClose}{\ensuremath{-}4.3}  % QLIKE % change, ALL48 vs RTH13 (full set; negative: ALL48 more accurate), Random forest, 16:00 bar
\providecommand{\rhtaXDmRfAllVsRthClose}{}\renewcommand{\rhtaXDmRfAllVsRthClose}{\ensuremath{-}4.26}  % ... its Diebold-Mariano statistic
\providecommand{\rhtaXPRfAllVsRthClose}{}\renewcommand{\rhtaXPRfAllVsRthClose}{\ensuremath{p < 0.001}}  % ... its p statement
\providecommand{\rhtaXPctRidgeAllVsRthBar}{}\renewcommand{\rhtaXPctRidgeAllVsRthBar}{1.9}  % QLIKE % change, ALL48 vs RTH13 (full set; negative: ALL48 more accurate), Ridge, every bar
\providecommand{\rhtaXDmRidgeAllVsRthBar}{}\renewcommand{\rhtaXDmRidgeAllVsRthBar}{3.53}  % ... its Diebold-Mariano statistic
\providecommand{\rhtaXPRidgeAllVsRthBar}{}\renewcommand{\rhtaXPRidgeAllVsRthBar}{\ensuremath{p < 0.001}}  % ... its p statement
\providecommand{\rhtaXPctLassoAllVsRthBar}{}\renewcommand{\rhtaXPctLassoAllVsRthBar}{2.5}  % QLIKE % change, ALL48 vs RTH13 (full set; negative: ALL48 more accurate), Lasso, every bar
\providecommand{\rhtaXDmLassoAllVsRthBar}{}\renewcommand{\rhtaXDmLassoAllVsRthBar}{2.64}  % ... its Diebold-Mariano statistic
\providecommand{\rhtaXPLassoAllVsRthBar}{}\renewcommand{\rhtaXPLassoAllVsRthBar}{\ensuremath{p = 0.008}}  % ... its p statement
\providecommand{\rhtaXPctEnetAllVsRthBar}{}\renewcommand{\rhtaXPctEnetAllVsRthBar}{2.2}  % QLIKE % change, ALL48 vs RTH13 (full set; negative: ALL48 more accurate), Elastic net, every bar
\providecommand{\rhtaXDmEnetAllVsRthBar}{}\renewcommand{\rhtaXDmEnetAllVsRthBar}{4.68}  % ... its Diebold-Mariano statistic
\providecommand{\rhtaXPEnetAllVsRthBar}{}\renewcommand{\rhtaXPEnetAllVsRthBar}{\ensuremath{p < 0.001}}  % ... its p statement
\providecommand{\rhtaXPctLgbmAllVsRthBar}{}\renewcommand{\rhtaXPctLgbmAllVsRthBar}{4.5}  % QLIKE % change, ALL48 vs RTH13 (full set; negative: ALL48 more accurate), LightGBM, every bar
\providecommand{\rhtaXDmLgbmAllVsRthBar}{}\renewcommand{\rhtaXDmLgbmAllVsRthBar}{7.39}  % ... its Diebold-Mariano statistic
\providecommand{\rhtaXPLgbmAllVsRthBar}{}\renewcommand{\rhtaXPLgbmAllVsRthBar}{\ensuremath{p < 0.001}}  % ... its p statement
\providecommand{\rhtaXPctXgbAllVsRthBar}{}\renewcommand{\rhtaXPctXgbAllVsRthBar}{3.2}  % QLIKE % change, ALL48 vs RTH13 (full set; negative: ALL48 more accurate), XGBoost, every bar
\providecommand{\rhtaXDmXgbAllVsRthBar}{}\renewcommand{\rhtaXDmXgbAllVsRthBar}{5.64}  % ... its Diebold-Mariano statistic
\providecommand{\rhtaXPXgbAllVsRthBar}{}\renewcommand{\rhtaXPXgbAllVsRthBar}{\ensuremath{p < 0.001}}  % ... its p statement
\providecommand{\rhtaXPctRfAllVsRthBar}{}\renewcommand{\rhtaXPctRfAllVsRthBar}{5.2}  % QLIKE % change, ALL48 vs RTH13 (full set; negative: ALL48 more accurate), Random forest, every bar
\providecommand{\rhtaXDmRfAllVsRthBar}{}\renewcommand{\rhtaXDmRfAllVsRthBar}{6.61}  % ... its Diebold-Mariano statistic
\providecommand{\rhtaXPRfAllVsRthBar}{}\renewcommand{\rhtaXPRfAllVsRthBar}{\ensuremath{p < 0.001}}  % ... its p statement
\providecommand{\rhtaXPctRidgeClVsRthClose}{}\renewcommand{\rhtaXPctRidgeClVsRthClose}{\ensuremath{-}10.2}  % QLIKE % change, CLOSE vs RTH13 (full set; negative: CLOSE more accurate), Ridge, 16:00 bar
\providecommand{\rhtaXDmRidgeClVsRthClose}{}\renewcommand{\rhtaXDmRidgeClVsRthClose}{\ensuremath{-}3.98}  % ... its Diebold-Mariano statistic
\providecommand{\rhtaXPRidgeClVsRthClose}{}\renewcommand{\rhtaXPRidgeClVsRthClose}{\ensuremath{p < 0.001}}  % ... its p statement
\providecommand{\rhtaXPctLassoClVsRthClose}{}\renewcommand{\rhtaXPctLassoClVsRthClose}{\ensuremath{-}11.8}  % QLIKE % change, CLOSE vs RTH13 (full set; negative: CLOSE more accurate), Lasso, 16:00 bar
\providecommand{\rhtaXDmLassoClVsRthClose}{}\renewcommand{\rhtaXDmLassoClVsRthClose}{\ensuremath{-}3.79}  % ... its Diebold-Mariano statistic
\providecommand{\rhtaXPLassoClVsRthClose}{}\renewcommand{\rhtaXPLassoClVsRthClose}{\ensuremath{p < 0.001}}  % ... its p statement
\providecommand{\rhtaXPctEnetClVsRthClose}{}\renewcommand{\rhtaXPctEnetClVsRthClose}{\ensuremath{-}9.4}  % QLIKE % change, CLOSE vs RTH13 (full set; negative: CLOSE more accurate), Elastic net, 16:00 bar
\providecommand{\rhtaXDmEnetClVsRthClose}{}\renewcommand{\rhtaXDmEnetClVsRthClose}{\ensuremath{-}3.86}  % ... its Diebold-Mariano statistic
\providecommand{\rhtaXPEnetClVsRthClose}{}\renewcommand{\rhtaXPEnetClVsRthClose}{\ensuremath{p < 0.001}}  % ... its p statement
\providecommand{\rhtaXPctLgbmClVsRthClose}{}\renewcommand{\rhtaXPctLgbmClVsRthClose}{1.6}  % QLIKE % change, CLOSE vs RTH13 (full set; negative: CLOSE more accurate), LightGBM, 16:00 bar
\providecommand{\rhtaXDmLgbmClVsRthClose}{}\renewcommand{\rhtaXDmLgbmClVsRthClose}{0.45}  % ... its Diebold-Mariano statistic
\providecommand{\rhtaXPLgbmClVsRthClose}{}\renewcommand{\rhtaXPLgbmClVsRthClose}{\ensuremath{p = 0.65}}  % ... its p statement
\providecommand{\rhtaXPctXgbClVsRthClose}{}\renewcommand{\rhtaXPctXgbClVsRthClose}{\ensuremath{-}1.1}  % QLIKE % change, CLOSE vs RTH13 (full set; negative: CLOSE more accurate), XGBoost, 16:00 bar
\providecommand{\rhtaXDmXgbClVsRthClose}{}\renewcommand{\rhtaXDmXgbClVsRthClose}{\ensuremath{-}0.40}  % ... its Diebold-Mariano statistic
\providecommand{\rhtaXPXgbClVsRthClose}{}\renewcommand{\rhtaXPXgbClVsRthClose}{\ensuremath{p = 0.69}}  % ... its p statement
\providecommand{\rhtaXPctRfClVsRthClose}{}\renewcommand{\rhtaXPctRfClVsRthClose}{\ensuremath{-}2.2}  % QLIKE % change, CLOSE vs RTH13 (full set; negative: CLOSE more accurate), Random forest, 16:00 bar
\providecommand{\rhtaXDmRfClVsRthClose}{}\renewcommand{\rhtaXDmRfClVsRthClose}{\ensuremath{-}0.81}  % ... its Diebold-Mariano statistic
\providecommand{\rhtaXPRfClVsRthClose}{}\renewcommand{\rhtaXPRfClVsRthClose}{\ensuremath{p = 0.42}}  % ... its p statement
\providecommand{\rhtaXPctRidgeClVsAllClose}{}\renewcommand{\rhtaXPctRidgeClVsAllClose}{\ensuremath{-}7.0}  % QLIKE % change, CLOSE vs ALL48 (full set; negative: CLOSE more accurate), Ridge, 16:00 bar
\providecommand{\rhtaXDmRidgeClVsAllClose}{}\renewcommand{\rhtaXDmRidgeClVsAllClose}{\ensuremath{-}3.31}  % ... its Diebold-Mariano statistic
\providecommand{\rhtaXPRidgeClVsAllClose}{}\renewcommand{\rhtaXPRidgeClVsAllClose}{\ensuremath{p < 0.001}}  % ... its p statement
\providecommand{\rhtaXPctLassoClVsAllClose}{}\renewcommand{\rhtaXPctLassoClVsAllClose}{\ensuremath{-}6.2}  % QLIKE % change, CLOSE vs ALL48 (full set; negative: CLOSE more accurate), Lasso, 16:00 bar
\providecommand{\rhtaXDmLassoClVsAllClose}{}\renewcommand{\rhtaXDmLassoClVsAllClose}{\ensuremath{-}3.16}  % ... its Diebold-Mariano statistic
\providecommand{\rhtaXPLassoClVsAllClose}{}\renewcommand{\rhtaXPLassoClVsAllClose}{\ensuremath{p = 0.002}}  % ... its p statement
\providecommand{\rhtaXPctEnetClVsAllClose}{}\renewcommand{\rhtaXPctEnetClVsAllClose}{\ensuremath{-}7.5}  % QLIKE % change, CLOSE vs ALL48 (full set; negative: CLOSE more accurate), Elastic net, 16:00 bar
\providecommand{\rhtaXDmEnetClVsAllClose}{}\renewcommand{\rhtaXDmEnetClVsAllClose}{\ensuremath{-}3.70}  % ... its Diebold-Mariano statistic
\providecommand{\rhtaXPEnetClVsAllClose}{}\renewcommand{\rhtaXPEnetClVsAllClose}{\ensuremath{p < 0.001}}  % ... its p statement
\providecommand{\rhtaXPctLgbmClVsAllClose}{}\renewcommand{\rhtaXPctLgbmClVsAllClose}{4.0}  % QLIKE % change, CLOSE vs ALL48 (full set; negative: CLOSE more accurate), LightGBM, 16:00 bar
\providecommand{\rhtaXDmLgbmClVsAllClose}{}\renewcommand{\rhtaXDmLgbmClVsAllClose}{0.94}  % ... its Diebold-Mariano statistic
\providecommand{\rhtaXPLgbmClVsAllClose}{}\renewcommand{\rhtaXPLgbmClVsAllClose}{\ensuremath{p = 0.35}}  % ... its p statement
\providecommand{\rhtaXPctXgbClVsAllClose}{}\renewcommand{\rhtaXPctXgbClVsAllClose}{0.8}  % QLIKE % change, CLOSE vs ALL48 (full set; negative: CLOSE more accurate), XGBoost, 16:00 bar
\providecommand{\rhtaXDmXgbClVsAllClose}{}\renewcommand{\rhtaXDmXgbClVsAllClose}{0.28}  % ... its Diebold-Mariano statistic
\providecommand{\rhtaXPXgbClVsAllClose}{}\renewcommand{\rhtaXPXgbClVsAllClose}{\ensuremath{p = 0.78}}  % ... its p statement
\providecommand{\rhtaXPctRfClVsAllClose}{}\renewcommand{\rhtaXPctRfClVsAllClose}{2.2}  % QLIKE % change, CLOSE vs ALL48 (full set; negative: CLOSE more accurate), Random forest, 16:00 bar
\providecommand{\rhtaXDmRfClVsAllClose}{}\renewcommand{\rhtaXDmRfClVsAllClose}{0.74}  % ... its Diebold-Mariano statistic
\providecommand{\rhtaXPRfClVsAllClose}{}\renewcommand{\rhtaXPRfClVsAllClose}{\ensuremath{p = 0.46}}  % ... its p statement
\providecommand{\rhtaXPctRidgeClNatVsRthClose}{}\renewcommand{\rhtaXPctRidgeClNatVsRthClose}{\ensuremath{-}10.4}  % QLIKE % change, CLOSE_NAT vs RTH13 (full set; negative: CLOSE_NAT more accurate), Ridge, 16:00 bar
\providecommand{\rhtaXDmRidgeClNatVsRthClose}{}\renewcommand{\rhtaXDmRidgeClNatVsRthClose}{\ensuremath{-}3.91}  % ... its Diebold-Mariano statistic
\providecommand{\rhtaXPRidgeClNatVsRthClose}{}\renewcommand{\rhtaXPRidgeClNatVsRthClose}{\ensuremath{p < 0.001}}  % ... its p statement
\providecommand{\rhtaXPctLassoClNatVsRthClose}{}\renewcommand{\rhtaXPctLassoClNatVsRthClose}{\ensuremath{-}11.5}  % QLIKE % change, CLOSE_NAT vs RTH13 (full set; negative: CLOSE_NAT more accurate), Lasso, 16:00 bar
\providecommand{\rhtaXDmLassoClNatVsRthClose}{}\renewcommand{\rhtaXDmLassoClNatVsRthClose}{\ensuremath{-}3.57}  % ... its Diebold-Mariano statistic
\providecommand{\rhtaXPLassoClNatVsRthClose}{}\renewcommand{\rhtaXPLassoClNatVsRthClose}{\ensuremath{p < 0.001}}  % ... its p statement
\providecommand{\rhtaXPctEnetClNatVsRthClose}{}\renewcommand{\rhtaXPctEnetClNatVsRthClose}{\ensuremath{-}9.3}  % QLIKE % change, CLOSE_NAT vs RTH13 (full set; negative: CLOSE_NAT more accurate), Elastic net, 16:00 bar
\providecommand{\rhtaXDmEnetClNatVsRthClose}{}\renewcommand{\rhtaXDmEnetClNatVsRthClose}{\ensuremath{-}3.58}  % ... its Diebold-Mariano statistic
\providecommand{\rhtaXPEnetClNatVsRthClose}{}\renewcommand{\rhtaXPEnetClNatVsRthClose}{\ensuremath{p < 0.001}}  % ... its p statement
\providecommand{\rhtaXPctRidgeAllCtVsRthClose}{}\renewcommand{\rhtaXPctRidgeAllCtVsRthClose}{\ensuremath{-}3.2}  % QLIKE % change, ALL48@16 vs RTH13 (full set; negative: ALL48@16 more accurate), Ridge, 16:00 bar
\providecommand{\rhtaXDmRidgeAllCtVsRthClose}{}\renewcommand{\rhtaXDmRidgeAllCtVsRthClose}{\ensuremath{-}2.35}  % ... its Diebold-Mariano statistic
\providecommand{\rhtaXPRidgeAllCtVsRthClose}{}\renewcommand{\rhtaXPRidgeAllCtVsRthClose}{\ensuremath{p = 0.019}}  % ... its p statement
\providecommand{\rhtaXPctLassoAllCtVsRthClose}{}\renewcommand{\rhtaXPctLassoAllCtVsRthClose}{\ensuremath{-}5.4}  % QLIKE % change, ALL48@16 vs RTH13 (full set; negative: ALL48@16 more accurate), Lasso, 16:00 bar
\providecommand{\rhtaXDmLassoAllCtVsRthClose}{}\renewcommand{\rhtaXDmLassoAllCtVsRthClose}{\ensuremath{-}3.04}  % ... its Diebold-Mariano statistic
\providecommand{\rhtaXPLassoAllCtVsRthClose}{}\renewcommand{\rhtaXPLassoAllCtVsRthClose}{\ensuremath{p = 0.002}}  % ... its p statement
\providecommand{\rhtaXPctEnetAllCtVsRthClose}{}\renewcommand{\rhtaXPctEnetAllCtVsRthClose}{\ensuremath{-}2.1}  % QLIKE % change, ALL48@16 vs RTH13 (full set; negative: ALL48@16 more accurate), Elastic net, 16:00 bar
\providecommand{\rhtaXDmEnetAllCtVsRthClose}{}\renewcommand{\rhtaXDmEnetAllCtVsRthClose}{\ensuremath{-}1.64}  % ... its Diebold-Mariano statistic
\providecommand{\rhtaXPEnetAllCtVsRthClose}{}\renewcommand{\rhtaXPEnetAllCtVsRthClose}{\ensuremath{p = 0.1006}}  % ... its p statement
\providecommand{\rhtaXPctRidgeRthCtVsRthClose}{}\renewcommand{\rhtaXPctRidgeRthCtVsRthClose}{0.0}  % QLIKE % change, RTH13@16 vs RTH13 (full set; negative: RTH13@16 more accurate), Ridge, 16:00 bar
\providecommand{\rhtaXDmRidgeRthCtVsRthClose}{}\renewcommand{\rhtaXDmRidgeRthCtVsRthClose}{0.00}  % ... its Diebold-Mariano statistic
\providecommand{\rhtaXPRidgeRthCtVsRthClose}{}\renewcommand{\rhtaXPRidgeRthCtVsRthClose}{\ensuremath{p = 1.00}}  % ... its p statement
\providecommand{\rhtaXPctLassoRthCtVsRthClose}{}\renewcommand{\rhtaXPctLassoRthCtVsRthClose}{\ensuremath{-}5.4}  % QLIKE % change, RTH13@16 vs RTH13 (full set; negative: RTH13@16 more accurate), Lasso, 16:00 bar
\providecommand{\rhtaXDmLassoRthCtVsRthClose}{}\renewcommand{\rhtaXDmLassoRthCtVsRthClose}{\ensuremath{-}3.40}  % ... its Diebold-Mariano statistic
\providecommand{\rhtaXPLassoRthCtVsRthClose}{}\renewcommand{\rhtaXPLassoRthCtVsRthClose}{\ensuremath{p < 0.001}}  % ... its p statement
\providecommand{\rhtaXPctEnetRthCtVsRthClose}{}\renewcommand{\rhtaXPctEnetRthCtVsRthClose}{0.0}  % QLIKE % change, RTH13@16 vs RTH13 (full set; negative: RTH13@16 more accurate), Elastic net, 16:00 bar
\providecommand{\rhtaXDmEnetRthCtVsRthClose}{}\renewcommand{\rhtaXDmEnetRthCtVsRthClose}{0.00}  % ... its Diebold-Mariano statistic
\providecommand{\rhtaXPEnetRthCtVsRthClose}{}\renewcommand{\rhtaXPEnetRthCtVsRthClose}{\ensuremath{p = 1.00}}  % ... its p statement
\providecommand{\rhtaXPctRidgeAllCtVsRthBar}{}\renewcommand{\rhtaXPctRidgeAllCtVsRthBar}{1.6}  % QLIKE % change, ALL48@16 vs RTH13 (full set; negative: ALL48@16 more accurate), Ridge, every bar
\providecommand{\rhtaXDmRidgeAllCtVsRthBar}{}\renewcommand{\rhtaXDmRidgeAllCtVsRthBar}{3.24}  % ... its Diebold-Mariano statistic
\providecommand{\rhtaXPRidgeAllCtVsRthBar}{}\renewcommand{\rhtaXPRidgeAllCtVsRthBar}{\ensuremath{p = 0.001}}  % ... its p statement
\providecommand{\rhtaXPctLassoAllCtVsRthBar}{}\renewcommand{\rhtaXPctLassoAllCtVsRthBar}{2.2}  % QLIKE % change, ALL48@16 vs RTH13 (full set; negative: ALL48@16 more accurate), Lasso, every bar
\providecommand{\rhtaXDmLassoAllCtVsRthBar}{}\renewcommand{\rhtaXDmLassoAllCtVsRthBar}{2.65}  % ... its Diebold-Mariano statistic
\providecommand{\rhtaXPLassoAllCtVsRthBar}{}\renewcommand{\rhtaXPLassoAllCtVsRthBar}{\ensuremath{p = 0.008}}  % ... its p statement
\providecommand{\rhtaXPctEnetAllCtVsRthBar}{}\renewcommand{\rhtaXPctEnetAllCtVsRthBar}{2.2}  % QLIKE % change, ALL48@16 vs RTH13 (full set; negative: ALL48@16 more accurate), Elastic net, every bar
\providecommand{\rhtaXDmEnetAllCtVsRthBar}{}\renewcommand{\rhtaXDmEnetAllCtVsRthBar}{4.68}  % ... its Diebold-Mariano statistic
\providecommand{\rhtaXPEnetAllCtVsRthBar}{}\renewcommand{\rhtaXPEnetAllCtVsRthBar}{\ensuremath{p < 0.001}}  % ... its p statement
\providecommand{\rhtaXPctRidgeRthCtVsRthBar}{}\renewcommand{\rhtaXPctRidgeRthCtVsRthBar}{0.0}  % QLIKE % change, RTH13@16 vs RTH13 (full set; negative: RTH13@16 more accurate), Ridge, every bar
\providecommand{\rhtaXDmRidgeRthCtVsRthBar}{}\renewcommand{\rhtaXDmRidgeRthCtVsRthBar}{0.00}  % ... its Diebold-Mariano statistic
\providecommand{\rhtaXPRidgeRthCtVsRthBar}{}\renewcommand{\rhtaXPRidgeRthCtVsRthBar}{\ensuremath{p = 1.00}}  % ... its p statement
\providecommand{\rhtaXPctLassoRthCtVsRthBar}{}\renewcommand{\rhtaXPctLassoRthCtVsRthBar}{0.2}  % QLIKE % change, RTH13@16 vs RTH13 (full set; negative: RTH13@16 more accurate), Lasso, every bar
\providecommand{\rhtaXDmLassoRthCtVsRthBar}{}\renewcommand{\rhtaXDmLassoRthCtVsRthBar}{0.19}  % ... its Diebold-Mariano statistic
\providecommand{\rhtaXPLassoRthCtVsRthBar}{}\renewcommand{\rhtaXPLassoRthCtVsRthBar}{\ensuremath{p = 0.85}}  % ... its p statement
\providecommand{\rhtaXPctEnetRthCtVsRthBar}{}\renewcommand{\rhtaXPctEnetRthCtVsRthBar}{0.0}  % QLIKE % change, RTH13@16 vs RTH13 (full set; negative: RTH13@16 more accurate), Elastic net, every bar
\providecommand{\rhtaXDmEnetRthCtVsRthBar}{}\renewcommand{\rhtaXDmEnetRthCtVsRthBar}{0.00}  % ... its Diebold-Mariano statistic
\providecommand{\rhtaXPEnetRthCtVsRthBar}{}\renewcommand{\rhtaXPEnetRthCtVsRthBar}{\ensuremath{p = 1.00}}  % ... its p statement
\providecommand{\rhtaXPctRidgeAllThreeKVsRthThreeKClose}{}\renewcommand{\rhtaXPctRidgeAllThreeKVsRthThreeKClose}{\ensuremath{-}8.1}  % QLIKE % change, ALL48_N3k vs RTH13_N3k (full set; negative: ALL48_N3k more accurate), Ridge, 16:00 bar
\providecommand{\rhtaXDmRidgeAllThreeKVsRthThreeKClose}{}\renewcommand{\rhtaXDmRidgeAllThreeKVsRthThreeKClose}{\ensuremath{-}5.43}  % ... its Diebold-Mariano statistic
\providecommand{\rhtaXPRidgeAllThreeKVsRthThreeKClose}{}\renewcommand{\rhtaXPRidgeAllThreeKVsRthThreeKClose}{\ensuremath{p < 0.001}}  % ... its p statement
\providecommand{\rhtaXPctLassoAllThreeKVsRthThreeKClose}{}\renewcommand{\rhtaXPctLassoAllThreeKVsRthThreeKClose}{\ensuremath{-}6.5}  % QLIKE % change, ALL48_N3k vs RTH13_N3k (full set; negative: ALL48_N3k more accurate), Lasso, 16:00 bar
\providecommand{\rhtaXDmLassoAllThreeKVsRthThreeKClose}{}\renewcommand{\rhtaXDmLassoAllThreeKVsRthThreeKClose}{\ensuremath{-}3.99}  % ... its Diebold-Mariano statistic
\providecommand{\rhtaXPLassoAllThreeKVsRthThreeKClose}{}\renewcommand{\rhtaXPLassoAllThreeKVsRthThreeKClose}{\ensuremath{p < 0.001}}  % ... its p statement
\providecommand{\rhtaXPctEnetAllThreeKVsRthThreeKClose}{}\renewcommand{\rhtaXPctEnetAllThreeKVsRthThreeKClose}{\ensuremath{-}4.2}  % QLIKE % change, ALL48_N3k vs RTH13_N3k (full set; negative: ALL48_N3k more accurate), Elastic net, 16:00 bar
\providecommand{\rhtaXDmEnetAllThreeKVsRthThreeKClose}{}\renewcommand{\rhtaXDmEnetAllThreeKVsRthThreeKClose}{\ensuremath{-}2.13}  % ... its Diebold-Mariano statistic
\providecommand{\rhtaXPEnetAllThreeKVsRthThreeKClose}{}\renewcommand{\rhtaXPEnetAllThreeKVsRthThreeKClose}{\ensuremath{p = 0.033}}  % ... its p statement
\providecommand{\rhtaXPctLgbmAllThreeKVsRthThreeKClose}{}\renewcommand{\rhtaXPctLgbmAllThreeKVsRthThreeKClose}{\ensuremath{-}3.9}  % QLIKE % change, ALL48_N3k vs RTH13_N3k (full set; negative: ALL48_N3k more accurate), LightGBM, 16:00 bar
\providecommand{\rhtaXDmLgbmAllThreeKVsRthThreeKClose}{}\renewcommand{\rhtaXDmLgbmAllThreeKVsRthThreeKClose}{\ensuremath{-}2.07}  % ... its Diebold-Mariano statistic
\providecommand{\rhtaXPLgbmAllThreeKVsRthThreeKClose}{}\renewcommand{\rhtaXPLgbmAllThreeKVsRthThreeKClose}{\ensuremath{p = 0.038}}  % ... its p statement
\providecommand{\rhtaXPctXgbAllThreeKVsRthThreeKClose}{}\renewcommand{\rhtaXPctXgbAllThreeKVsRthThreeKClose}{\ensuremath{-}3.3}  % QLIKE % change, ALL48_N3k vs RTH13_N3k (full set; negative: ALL48_N3k more accurate), XGBoost, 16:00 bar
\providecommand{\rhtaXDmXgbAllThreeKVsRthThreeKClose}{}\renewcommand{\rhtaXDmXgbAllThreeKVsRthThreeKClose}{\ensuremath{-}1.80}  % ... its Diebold-Mariano statistic
\providecommand{\rhtaXPXgbAllThreeKVsRthThreeKClose}{}\renewcommand{\rhtaXPXgbAllThreeKVsRthThreeKClose}{\ensuremath{p = 0.072}}  % ... its p statement
\providecommand{\rhtaXPctRfAllThreeKVsRthThreeKClose}{}\renewcommand{\rhtaXPctRfAllThreeKVsRthThreeKClose}{\ensuremath{-}5.7}  % QLIKE % change, ALL48_N3k vs RTH13_N3k (full set; negative: ALL48_N3k more accurate), Random forest, 16:00 bar
\providecommand{\rhtaXDmRfAllThreeKVsRthThreeKClose}{}\renewcommand{\rhtaXDmRfAllThreeKVsRthThreeKClose}{\ensuremath{-}3.60}  % ... its Diebold-Mariano statistic
\providecommand{\rhtaXPRfAllThreeKVsRthThreeKClose}{}\renewcommand{\rhtaXPRfAllThreeKVsRthThreeKClose}{\ensuremath{p < 0.001}}  % ... its p statement
\providecommand{\rhtaXPctRidgeAllThreeKVsRthThreeKBar}{}\renewcommand{\rhtaXPctRidgeAllThreeKVsRthThreeKBar}{4.5}  % QLIKE % change, ALL48_N3k vs RTH13_N3k (full set; negative: ALL48_N3k more accurate), Ridge, every bar
\providecommand{\rhtaXDmRidgeAllThreeKVsRthThreeKBar}{}\renewcommand{\rhtaXDmRidgeAllThreeKVsRthThreeKBar}{7.72}  % ... its Diebold-Mariano statistic
\providecommand{\rhtaXPRidgeAllThreeKVsRthThreeKBar}{}\renewcommand{\rhtaXPRidgeAllThreeKVsRthThreeKBar}{\ensuremath{p < 0.001}}  % ... its p statement
\providecommand{\rhtaXPctLassoAllThreeKVsRthThreeKBar}{}\renewcommand{\rhtaXPctLassoAllThreeKVsRthThreeKBar}{4.9}  % QLIKE % change, ALL48_N3k vs RTH13_N3k (full set; negative: ALL48_N3k more accurate), Lasso, every bar
\providecommand{\rhtaXDmLassoAllThreeKVsRthThreeKBar}{}\renewcommand{\rhtaXDmLassoAllThreeKVsRthThreeKBar}{7.96}  % ... its Diebold-Mariano statistic
\providecommand{\rhtaXPLassoAllThreeKVsRthThreeKBar}{}\renewcommand{\rhtaXPLassoAllThreeKVsRthThreeKBar}{\ensuremath{p < 0.001}}  % ... its p statement
\providecommand{\rhtaXPctEnetAllThreeKVsRthThreeKBar}{}\renewcommand{\rhtaXPctEnetAllThreeKVsRthThreeKBar}{5.8}  % QLIKE % change, ALL48_N3k vs RTH13_N3k (full set; negative: ALL48_N3k more accurate), Elastic net, every bar
\providecommand{\rhtaXDmEnetAllThreeKVsRthThreeKBar}{}\renewcommand{\rhtaXDmEnetAllThreeKVsRthThreeKBar}{7.58}  % ... its Diebold-Mariano statistic
\providecommand{\rhtaXPEnetAllThreeKVsRthThreeKBar}{}\renewcommand{\rhtaXPEnetAllThreeKVsRthThreeKBar}{\ensuremath{p < 0.001}}  % ... its p statement
\providecommand{\rhtaXPctLgbmAllThreeKVsRthThreeKBar}{}\renewcommand{\rhtaXPctLgbmAllThreeKVsRthThreeKBar}{4.7}  % QLIKE % change, ALL48_N3k vs RTH13_N3k (full set; negative: ALL48_N3k more accurate), LightGBM, every bar
\providecommand{\rhtaXDmLgbmAllThreeKVsRthThreeKBar}{}\renewcommand{\rhtaXDmLgbmAllThreeKVsRthThreeKBar}{5.00}  % ... its Diebold-Mariano statistic
\providecommand{\rhtaXPLgbmAllThreeKVsRthThreeKBar}{}\renewcommand{\rhtaXPLgbmAllThreeKVsRthThreeKBar}{\ensuremath{p < 0.001}}  % ... its p statement
\providecommand{\rhtaXPctXgbAllThreeKVsRthThreeKBar}{}\renewcommand{\rhtaXPctXgbAllThreeKVsRthThreeKBar}{5.1}  % QLIKE % change, ALL48_N3k vs RTH13_N3k (full set; negative: ALL48_N3k more accurate), XGBoost, every bar
\providecommand{\rhtaXDmXgbAllThreeKVsRthThreeKBar}{}\renewcommand{\rhtaXDmXgbAllThreeKVsRthThreeKBar}{5.92}  % ... its Diebold-Mariano statistic
\providecommand{\rhtaXPXgbAllThreeKVsRthThreeKBar}{}\renewcommand{\rhtaXPXgbAllThreeKVsRthThreeKBar}{\ensuremath{p < 0.001}}  % ... its p statement
\providecommand{\rhtaXPctRfAllThreeKVsRthThreeKBar}{}\renewcommand{\rhtaXPctRfAllThreeKVsRthThreeKBar}{2.7}  % QLIKE % change, ALL48_N3k vs RTH13_N3k (full set; negative: ALL48_N3k more accurate), Random forest, every bar
\providecommand{\rhtaXDmRfAllThreeKVsRthThreeKBar}{}\renewcommand{\rhtaXDmRfAllThreeKVsRthThreeKBar}{3.94}  % ... its Diebold-Mariano statistic
\providecommand{\rhtaXPRfAllThreeKVsRthThreeKBar}{}\renewcommand{\rhtaXPRfAllThreeKVsRthThreeKBar}{\ensuremath{p < 0.001}}  % ... its p statement
\providecommand{\rhtaXPctRidgeClVsRthThreeKClose}{}\renewcommand{\rhtaXPctRidgeClVsRthThreeKClose}{\ensuremath{-}17.4}  % QLIKE % change, CLOSE vs RTH13_N3k (full set; negative: CLOSE more accurate), Ridge, 16:00 bar
\providecommand{\rhtaXDmRidgeClVsRthThreeKClose}{}\renewcommand{\rhtaXDmRidgeClVsRthThreeKClose}{\ensuremath{-}6.84}  % ... its Diebold-Mariano statistic
\providecommand{\rhtaXPRidgeClVsRthThreeKClose}{}\renewcommand{\rhtaXPRidgeClVsRthThreeKClose}{\ensuremath{p < 0.001}}  % ... its p statement
\providecommand{\rhtaXPctLassoClVsRthThreeKClose}{}\renewcommand{\rhtaXPctLassoClVsRthThreeKClose}{\ensuremath{-}16.9}  % QLIKE % change, CLOSE vs RTH13_N3k (full set; negative: CLOSE more accurate), Lasso, 16:00 bar
\providecommand{\rhtaXDmLassoClVsRthThreeKClose}{}\renewcommand{\rhtaXDmLassoClVsRthThreeKClose}{\ensuremath{-}6.57}  % ... its Diebold-Mariano statistic
\providecommand{\rhtaXPLassoClVsRthThreeKClose}{}\renewcommand{\rhtaXPLassoClVsRthThreeKClose}{\ensuremath{p < 0.001}}  % ... its p statement
\providecommand{\rhtaXPctEnetClVsRthThreeKClose}{}\renewcommand{\rhtaXPctEnetClVsRthThreeKClose}{\ensuremath{-}18.0}  % QLIKE % change, CLOSE vs RTH13_N3k (full set; negative: CLOSE more accurate), Elastic net, 16:00 bar
\providecommand{\rhtaXDmEnetClVsRthThreeKClose}{}\renewcommand{\rhtaXDmEnetClVsRthThreeKClose}{\ensuremath{-}6.66}  % ... its Diebold-Mariano statistic
\providecommand{\rhtaXPEnetClVsRthThreeKClose}{}\renewcommand{\rhtaXPEnetClVsRthThreeKClose}{\ensuremath{p < 0.001}}  % ... its p statement
\providecommand{\rhtaXPctLgbmClVsRthThreeKClose}{}\renewcommand{\rhtaXPctLgbmClVsRthThreeKClose}{\ensuremath{-}9.5}  % QLIKE % change, CLOSE vs RTH13_N3k (full set; negative: CLOSE more accurate), LightGBM, 16:00 bar
\providecommand{\rhtaXDmLgbmClVsRthThreeKClose}{}\renewcommand{\rhtaXDmLgbmClVsRthThreeKClose}{\ensuremath{-}2.88}  % ... its Diebold-Mariano statistic
\providecommand{\rhtaXPLgbmClVsRthThreeKClose}{}\renewcommand{\rhtaXPLgbmClVsRthThreeKClose}{\ensuremath{p = 0.004}}  % ... its p statement
\providecommand{\rhtaXPctXgbClVsRthThreeKClose}{}\renewcommand{\rhtaXPctXgbClVsRthThreeKClose}{\ensuremath{-}9.5}  % QLIKE % change, CLOSE vs RTH13_N3k (full set; negative: CLOSE more accurate), XGBoost, 16:00 bar
\providecommand{\rhtaXDmXgbClVsRthThreeKClose}{}\renewcommand{\rhtaXDmXgbClVsRthThreeKClose}{\ensuremath{-}3.40}  % ... its Diebold-Mariano statistic
\providecommand{\rhtaXPXgbClVsRthThreeKClose}{}\renewcommand{\rhtaXPXgbClVsRthThreeKClose}{\ensuremath{p < 0.001}}  % ... its p statement
\providecommand{\rhtaXPctRfClVsRthThreeKClose}{}\renewcommand{\rhtaXPctRfClVsRthThreeKClose}{\ensuremath{-}8.7}  % QLIKE % change, CLOSE vs RTH13_N3k (full set; negative: CLOSE more accurate), Random forest, 16:00 bar
\providecommand{\rhtaXDmRfClVsRthThreeKClose}{}\renewcommand{\rhtaXDmRfClVsRthThreeKClose}{\ensuremath{-}3.08}  % ... its Diebold-Mariano statistic
\providecommand{\rhtaXPRfClVsRthThreeKClose}{}\renewcommand{\rhtaXPRfClVsRthThreeKClose}{\ensuremath{p = 0.002}}  % ... its p statement
\providecommand{\rhtaXPctRidgeClVsAllThreeKClose}{}\renewcommand{\rhtaXPctRidgeClVsAllThreeKClose}{\ensuremath{-}10.2}  % QLIKE % change, CLOSE vs ALL48_N3k (full set; negative: CLOSE more accurate), Ridge, 16:00 bar
\providecommand{\rhtaXDmRidgeClVsAllThreeKClose}{}\renewcommand{\rhtaXDmRidgeClVsAllThreeKClose}{\ensuremath{-}4.63}  % ... its Diebold-Mariano statistic
\providecommand{\rhtaXPRidgeClVsAllThreeKClose}{}\renewcommand{\rhtaXPRidgeClVsAllThreeKClose}{\ensuremath{p < 0.001}}  % ... its p statement
\providecommand{\rhtaXPctLassoClVsAllThreeKClose}{}\renewcommand{\rhtaXPctLassoClVsAllThreeKClose}{\ensuremath{-}11.1}  % QLIKE % change, CLOSE vs ALL48_N3k (full set; negative: CLOSE more accurate), Lasso, 16:00 bar
\providecommand{\rhtaXDmLassoClVsAllThreeKClose}{}\renewcommand{\rhtaXDmLassoClVsAllThreeKClose}{\ensuremath{-}4.62}  % ... its Diebold-Mariano statistic
\providecommand{\rhtaXPLassoClVsAllThreeKClose}{}\renewcommand{\rhtaXPLassoClVsAllThreeKClose}{\ensuremath{p < 0.001}}  % ... its p statement
\providecommand{\rhtaXPctEnetClVsAllThreeKClose}{}\renewcommand{\rhtaXPctEnetClVsAllThreeKClose}{\ensuremath{-}14.5}  % QLIKE % change, CLOSE vs ALL48_N3k (full set; negative: CLOSE more accurate), Elastic net, 16:00 bar
\providecommand{\rhtaXDmEnetClVsAllThreeKClose}{}\renewcommand{\rhtaXDmEnetClVsAllThreeKClose}{\ensuremath{-}5.41}  % ... its Diebold-Mariano statistic
\providecommand{\rhtaXPEnetClVsAllThreeKClose}{}\renewcommand{\rhtaXPEnetClVsAllThreeKClose}{\ensuremath{p < 0.001}}  % ... its p statement
\providecommand{\rhtaXPctLgbmClVsAllThreeKClose}{}\renewcommand{\rhtaXPctLgbmClVsAllThreeKClose}{\ensuremath{-}5.8}  % QLIKE % change, CLOSE vs ALL48_N3k (full set; negative: CLOSE more accurate), LightGBM, 16:00 bar
\providecommand{\rhtaXDmLgbmClVsAllThreeKClose}{}\renewcommand{\rhtaXDmLgbmClVsAllThreeKClose}{\ensuremath{-}1.68}  % ... its Diebold-Mariano statistic
\providecommand{\rhtaXPLgbmClVsAllThreeKClose}{}\renewcommand{\rhtaXPLgbmClVsAllThreeKClose}{\ensuremath{p = 0.093}}  % ... its p statement
\providecommand{\rhtaXPctXgbClVsAllThreeKClose}{}\renewcommand{\rhtaXPctXgbClVsAllThreeKClose}{\ensuremath{-}6.4}  % QLIKE % change, CLOSE vs ALL48_N3k (full set; negative: CLOSE more accurate), XGBoost, 16:00 bar
\providecommand{\rhtaXDmXgbClVsAllThreeKClose}{}\renewcommand{\rhtaXDmXgbClVsAllThreeKClose}{\ensuremath{-}2.23}  % ... its Diebold-Mariano statistic
\providecommand{\rhtaXPXgbClVsAllThreeKClose}{}\renewcommand{\rhtaXPXgbClVsAllThreeKClose}{\ensuremath{p = 0.026}}  % ... its p statement
\providecommand{\rhtaXPctRfClVsAllThreeKClose}{}\renewcommand{\rhtaXPctRfClVsAllThreeKClose}{\ensuremath{-}3.2}  % QLIKE % change, CLOSE vs ALL48_N3k (full set; negative: CLOSE more accurate), Random forest, 16:00 bar
\providecommand{\rhtaXDmRfClVsAllThreeKClose}{}\renewcommand{\rhtaXDmRfClVsAllThreeKClose}{\ensuremath{-}1.10}  % ... its Diebold-Mariano statistic
\providecommand{\rhtaXPRfClVsAllThreeKClose}{}\renewcommand{\rhtaXPRfClVsAllThreeKClose}{\ensuremath{p = 0.27}}  % ... its p statement
\providecommand{\rhtaXPctRidgeAllThirtyNineKVsRthClose}{}\renewcommand{\rhtaXPctRidgeAllThirtyNineKVsRthClose}{0.5}  % QLIKE % change, ALL48_N39k vs RTH13 (full set; negative: ALL48_N39k more accurate), Ridge, 16:00 bar
\providecommand{\rhtaXDmRidgeAllThirtyNineKVsRthClose}{}\renewcommand{\rhtaXDmRidgeAllThirtyNineKVsRthClose}{0.34}  % ... its Diebold-Mariano statistic
\providecommand{\rhtaXPRidgeAllThirtyNineKVsRthClose}{}\renewcommand{\rhtaXPRidgeAllThirtyNineKVsRthClose}{\ensuremath{p = 0.73}}  % ... its p statement
\providecommand{\rhtaXPctLassoAllThirtyNineKVsRthClose}{}\renewcommand{\rhtaXPctLassoAllThirtyNineKVsRthClose}{\ensuremath{-}3.6}  % QLIKE % change, ALL48_N39k vs RTH13 (full set; negative: ALL48_N39k more accurate), Lasso, 16:00 bar
\providecommand{\rhtaXDmLassoAllThirtyNineKVsRthClose}{}\renewcommand{\rhtaXDmLassoAllThirtyNineKVsRthClose}{\ensuremath{-}2.08}  % ... its Diebold-Mariano statistic
\providecommand{\rhtaXPLassoAllThirtyNineKVsRthClose}{}\renewcommand{\rhtaXPLassoAllThirtyNineKVsRthClose}{\ensuremath{p = 0.038}}  % ... its p statement
\providecommand{\rhtaXPctEnetAllThirtyNineKVsRthClose}{}\renewcommand{\rhtaXPctEnetAllThirtyNineKVsRthClose}{0.1}  % QLIKE % change, ALL48_N39k vs RTH13 (full set; negative: ALL48_N39k more accurate), Elastic net, 16:00 bar
\providecommand{\rhtaXDmEnetAllThirtyNineKVsRthClose}{}\renewcommand{\rhtaXDmEnetAllThirtyNineKVsRthClose}{0.07}  % ... its Diebold-Mariano statistic
\providecommand{\rhtaXPEnetAllThirtyNineKVsRthClose}{}\renewcommand{\rhtaXPEnetAllThirtyNineKVsRthClose}{\ensuremath{p = 0.95}}  % ... its p statement
\providecommand{\rhtaXPctLgbmAllThirtyNineKVsRthClose}{}\renewcommand{\rhtaXPctLgbmAllThirtyNineKVsRthClose}{0.4}  % QLIKE % change, ALL48_N39k vs RTH13 (full set; negative: ALL48_N39k more accurate), LightGBM, 16:00 bar
\providecommand{\rhtaXDmLgbmAllThirtyNineKVsRthClose}{}\renewcommand{\rhtaXDmLgbmAllThirtyNineKVsRthClose}{0.36}  % ... its Diebold-Mariano statistic
\providecommand{\rhtaXPLgbmAllThirtyNineKVsRthClose}{}\renewcommand{\rhtaXPLgbmAllThirtyNineKVsRthClose}{\ensuremath{p = 0.72}}  % ... its p statement
\providecommand{\rhtaXPctXgbAllThirtyNineKVsRthClose}{}\renewcommand{\rhtaXPctXgbAllThirtyNineKVsRthClose}{1.3}  % QLIKE % change, ALL48_N39k vs RTH13 (full set; negative: ALL48_N39k more accurate), XGBoost, 16:00 bar
\providecommand{\rhtaXDmXgbAllThirtyNineKVsRthClose}{}\renewcommand{\rhtaXDmXgbAllThirtyNineKVsRthClose}{1.45}  % ... its Diebold-Mariano statistic
\providecommand{\rhtaXPXgbAllThirtyNineKVsRthClose}{}\renewcommand{\rhtaXPXgbAllThirtyNineKVsRthClose}{\ensuremath{p = 0.15}}  % ... its p statement
\providecommand{\rhtaXPctRfAllThirtyNineKVsRthClose}{}\renewcommand{\rhtaXPctRfAllThirtyNineKVsRthClose}{\ensuremath{-}2.7}  % QLIKE % change, ALL48_N39k vs RTH13 (full set; negative: ALL48_N39k more accurate), Random forest, 16:00 bar
\providecommand{\rhtaXDmRfAllThirtyNineKVsRthClose}{}\renewcommand{\rhtaXDmRfAllThirtyNineKVsRthClose}{\ensuremath{-}2.72}  % ... its Diebold-Mariano statistic
\providecommand{\rhtaXPRfAllThirtyNineKVsRthClose}{}\renewcommand{\rhtaXPRfAllThirtyNineKVsRthClose}{\ensuremath{p = 0.007}}  % ... its p statement
\providecommand{\rhtaXPctRidgeAllThirtyNineKVsRthBar}{}\renewcommand{\rhtaXPctRidgeAllThirtyNineKVsRthBar}{2.7}  % QLIKE % change, ALL48_N39k vs RTH13 (full set; negative: ALL48_N39k more accurate), Ridge, every bar
\providecommand{\rhtaXDmRidgeAllThirtyNineKVsRthBar}{}\renewcommand{\rhtaXDmRidgeAllThirtyNineKVsRthBar}{4.54}  % ... its Diebold-Mariano statistic
\providecommand{\rhtaXPRidgeAllThirtyNineKVsRthBar}{}\renewcommand{\rhtaXPRidgeAllThirtyNineKVsRthBar}{\ensuremath{p < 0.001}}  % ... its p statement
\providecommand{\rhtaXPctLassoAllThirtyNineKVsRthBar}{}\renewcommand{\rhtaXPctLassoAllThirtyNineKVsRthBar}{4.2}  % QLIKE % change, ALL48_N39k vs RTH13 (full set; negative: ALL48_N39k more accurate), Lasso, every bar
\providecommand{\rhtaXDmLassoAllThirtyNineKVsRthBar}{}\renewcommand{\rhtaXDmLassoAllThirtyNineKVsRthBar}{3.44}  % ... its Diebold-Mariano statistic
\providecommand{\rhtaXPLassoAllThirtyNineKVsRthBar}{}\renewcommand{\rhtaXPLassoAllThirtyNineKVsRthBar}{\ensuremath{p < 0.001}}  % ... its p statement
\providecommand{\rhtaXPctEnetAllThirtyNineKVsRthBar}{}\renewcommand{\rhtaXPctEnetAllThirtyNineKVsRthBar}{4.9}  % QLIKE % change, ALL48_N39k vs RTH13 (full set; negative: ALL48_N39k more accurate), Elastic net, every bar
\providecommand{\rhtaXDmEnetAllThirtyNineKVsRthBar}{}\renewcommand{\rhtaXDmEnetAllThirtyNineKVsRthBar}{5.04}  % ... its Diebold-Mariano statistic
\providecommand{\rhtaXPEnetAllThirtyNineKVsRthBar}{}\renewcommand{\rhtaXPEnetAllThirtyNineKVsRthBar}{\ensuremath{p < 0.001}}  % ... its p statement
\providecommand{\rhtaXPctLgbmAllThirtyNineKVsRthBar}{}\renewcommand{\rhtaXPctLgbmAllThirtyNineKVsRthBar}{10.4}  % QLIKE % change, ALL48_N39k vs RTH13 (full set; negative: ALL48_N39k more accurate), LightGBM, every bar
\providecommand{\rhtaXDmLgbmAllThirtyNineKVsRthBar}{}\renewcommand{\rhtaXDmLgbmAllThirtyNineKVsRthBar}{8.35}  % ... its Diebold-Mariano statistic
\providecommand{\rhtaXPLgbmAllThirtyNineKVsRthBar}{}\renewcommand{\rhtaXPLgbmAllThirtyNineKVsRthBar}{\ensuremath{p < 0.001}}  % ... its p statement
\providecommand{\rhtaXPctXgbAllThirtyNineKVsRthBar}{}\renewcommand{\rhtaXPctXgbAllThirtyNineKVsRthBar}{6.4}  % QLIKE % change, ALL48_N39k vs RTH13 (full set; negative: ALL48_N39k more accurate), XGBoost, every bar
\providecommand{\rhtaXDmXgbAllThirtyNineKVsRthBar}{}\renewcommand{\rhtaXDmXgbAllThirtyNineKVsRthBar}{8.06}  % ... its Diebold-Mariano statistic
\providecommand{\rhtaXPXgbAllThirtyNineKVsRthBar}{}\renewcommand{\rhtaXPXgbAllThirtyNineKVsRthBar}{\ensuremath{p < 0.001}}  % ... its p statement
\providecommand{\rhtaXPctRfAllThirtyNineKVsRthBar}{}\renewcommand{\rhtaXPctRfAllThirtyNineKVsRthBar}{7.9}  % QLIKE % change, ALL48_N39k vs RTH13 (full set; negative: ALL48_N39k more accurate), Random forest, every bar
\providecommand{\rhtaXDmRfAllThirtyNineKVsRthBar}{}\renewcommand{\rhtaXDmRfAllThirtyNineKVsRthBar}{6.79}  % ... its Diebold-Mariano statistic
\providecommand{\rhtaXPRfAllThirtyNineKVsRthBar}{}\renewcommand{\rhtaXPRfAllThirtyNineKVsRthBar}{\ensuremath{p < 0.001}}  % ... its p statement
\providecommand{\rhtaXPctRidgeAllThreeKVsAllClose}{}\renewcommand{\rhtaXPctRidgeAllThreeKVsAllClose}{3.5}  % QLIKE % change, ALL48_N3k vs ALL48 (full set; negative: ALL48_N3k more accurate), Ridge, 16:00 bar
\providecommand{\rhtaXDmRidgeAllThreeKVsAllClose}{}\renewcommand{\rhtaXDmRidgeAllThreeKVsAllClose}{2.28}  % ... its Diebold-Mariano statistic
\providecommand{\rhtaXPRidgeAllThreeKVsAllClose}{}\renewcommand{\rhtaXPRidgeAllThreeKVsAllClose}{\ensuremath{p = 0.022}}  % ... its p statement
\providecommand{\rhtaXPctLassoAllThreeKVsAllClose}{}\renewcommand{\rhtaXPctLassoAllThreeKVsAllClose}{5.6}  % QLIKE % change, ALL48_N3k vs ALL48 (full set; negative: ALL48_N3k more accurate), Lasso, 16:00 bar
\providecommand{\rhtaXDmLassoAllThreeKVsAllClose}{}\renewcommand{\rhtaXDmLassoAllThreeKVsAllClose}{3.42}  % ... its Diebold-Mariano statistic
\providecommand{\rhtaXPLassoAllThreeKVsAllClose}{}\renewcommand{\rhtaXPLassoAllThreeKVsAllClose}{\ensuremath{p < 0.001}}  % ... its p statement
\providecommand{\rhtaXPctEnetAllThreeKVsAllClose}{}\renewcommand{\rhtaXPctEnetAllThreeKVsAllClose}{8.2}  % QLIKE % change, ALL48_N3k vs ALL48 (full set; negative: ALL48_N3k more accurate), Elastic net, 16:00 bar
\providecommand{\rhtaXDmEnetAllThreeKVsAllClose}{}\renewcommand{\rhtaXDmEnetAllThreeKVsAllClose}{4.01}  % ... its Diebold-Mariano statistic
\providecommand{\rhtaXPEnetAllThreeKVsAllClose}{}\renewcommand{\rhtaXPEnetAllThreeKVsAllClose}{\ensuremath{p < 0.001}}  % ... its p statement
\providecommand{\rhtaXPctLgbmAllThreeKVsAllClose}{}\renewcommand{\rhtaXPctLgbmAllThreeKVsAllClose}{10.4}  % QLIKE % change, ALL48_N3k vs ALL48 (full set; negative: ALL48_N3k more accurate), LightGBM, 16:00 bar
\providecommand{\rhtaXDmLgbmAllThreeKVsAllClose}{}\renewcommand{\rhtaXDmLgbmAllThreeKVsAllClose}{4.29}  % ... its Diebold-Mariano statistic
\providecommand{\rhtaXPLgbmAllThreeKVsAllClose}{}\renewcommand{\rhtaXPLgbmAllThreeKVsAllClose}{\ensuremath{p < 0.001}}  % ... its p statement
\providecommand{\rhtaXPctXgbAllThreeKVsAllClose}{}\renewcommand{\rhtaXPctXgbAllThreeKVsAllClose}{7.7}  % QLIKE % change, ALL48_N3k vs ALL48 (full set; negative: ALL48_N3k more accurate), XGBoost, 16:00 bar
\providecommand{\rhtaXDmXgbAllThreeKVsAllClose}{}\renewcommand{\rhtaXDmXgbAllThreeKVsAllClose}{4.20}  % ... its Diebold-Mariano statistic
\providecommand{\rhtaXPXgbAllThreeKVsAllClose}{}\renewcommand{\rhtaXPXgbAllThreeKVsAllClose}{\ensuremath{p < 0.001}}  % ... its p statement
\providecommand{\rhtaXPctRfAllThreeKVsAllClose}{}\renewcommand{\rhtaXPctRfAllThreeKVsAllClose}{5.6}  % QLIKE % change, ALL48_N3k vs ALL48 (full set; negative: ALL48_N3k more accurate), Random forest, 16:00 bar
\providecommand{\rhtaXDmRfAllThreeKVsAllClose}{}\renewcommand{\rhtaXDmRfAllThreeKVsAllClose}{3.47}  % ... its Diebold-Mariano statistic
\providecommand{\rhtaXPRfAllThreeKVsAllClose}{}\renewcommand{\rhtaXPRfAllThreeKVsAllClose}{\ensuremath{p < 0.001}}  % ... its p statement
\providecommand{\rhtaXPctRidgeAllThreeKVsAllBar}{}\renewcommand{\rhtaXPctRidgeAllThreeKVsAllBar}{5.4}  % QLIKE % change, ALL48_N3k vs ALL48 (full set; negative: ALL48_N3k more accurate), Ridge, every bar
\providecommand{\rhtaXDmRidgeAllThreeKVsAllBar}{}\renewcommand{\rhtaXDmRidgeAllThreeKVsAllBar}{6.67}  % ... its Diebold-Mariano statistic
\providecommand{\rhtaXPRidgeAllThreeKVsAllBar}{}\renewcommand{\rhtaXPRidgeAllThreeKVsAllBar}{\ensuremath{p < 0.001}}  % ... its p statement
\providecommand{\rhtaXPctLassoAllThreeKVsAllBar}{}\renewcommand{\rhtaXPctLassoAllThreeKVsAllBar}{5.9}  % QLIKE % change, ALL48_N3k vs ALL48 (full set; negative: ALL48_N3k more accurate), Lasso, every bar
\providecommand{\rhtaXDmLassoAllThreeKVsAllBar}{}\renewcommand{\rhtaXDmLassoAllThreeKVsAllBar}{7.47}  % ... its Diebold-Mariano statistic
\providecommand{\rhtaXPLassoAllThreeKVsAllBar}{}\renewcommand{\rhtaXPLassoAllThreeKVsAllBar}{\ensuremath{p < 0.001}}  % ... its p statement
\providecommand{\rhtaXPctEnetAllThreeKVsAllBar}{}\renewcommand{\rhtaXPctEnetAllThreeKVsAllBar}{6.6}  % QLIKE % change, ALL48_N3k vs ALL48 (full set; negative: ALL48_N3k more accurate), Elastic net, every bar
\providecommand{\rhtaXDmEnetAllThreeKVsAllBar}{}\renewcommand{\rhtaXDmEnetAllThreeKVsAllBar}{7.48}  % ... its Diebold-Mariano statistic
\providecommand{\rhtaXPEnetAllThreeKVsAllBar}{}\renewcommand{\rhtaXPEnetAllThreeKVsAllBar}{\ensuremath{p < 0.001}}  % ... its p statement
\providecommand{\rhtaXPctLgbmAllThreeKVsAllBar}{}\renewcommand{\rhtaXPctLgbmAllThreeKVsAllBar}{15.2}  % QLIKE % change, ALL48_N3k vs ALL48 (full set; negative: ALL48_N3k more accurate), LightGBM, every bar
\providecommand{\rhtaXDmLgbmAllThreeKVsAllBar}{}\renewcommand{\rhtaXDmLgbmAllThreeKVsAllBar}{8.62}  % ... its Diebold-Mariano statistic
\providecommand{\rhtaXPLgbmAllThreeKVsAllBar}{}\renewcommand{\rhtaXPLgbmAllThreeKVsAllBar}{\ensuremath{p < 0.001}}  % ... its p statement
\providecommand{\rhtaXPctXgbAllThreeKVsAllBar}{}\renewcommand{\rhtaXPctXgbAllThreeKVsAllBar}{12.7}  % QLIKE % change, ALL48_N3k vs ALL48 (full set; negative: ALL48_N3k more accurate), XGBoost, every bar
\providecommand{\rhtaXDmXgbAllThreeKVsAllBar}{}\renewcommand{\rhtaXDmXgbAllThreeKVsAllBar}{8.10}  % ... its Diebold-Mariano statistic
\providecommand{\rhtaXPXgbAllThreeKVsAllBar}{}\renewcommand{\rhtaXPXgbAllThreeKVsAllBar}{\ensuremath{p < 0.001}}  % ... its p statement
\providecommand{\rhtaXPctRfAllThreeKVsAllBar}{}\renewcommand{\rhtaXPctRfAllThreeKVsAllBar}{6.8}  % QLIKE % change, ALL48_N3k vs ALL48 (full set; negative: ALL48_N3k more accurate), Random forest, every bar
\providecommand{\rhtaXDmRfAllThreeKVsAllBar}{}\renewcommand{\rhtaXDmRfAllThreeKVsAllBar}{7.80}  % ... its Diebold-Mariano statistic
\providecommand{\rhtaXPRfAllThreeKVsAllBar}{}\renewcommand{\rhtaXPRfAllThreeKVsAllBar}{\ensuremath{p < 0.001}}  % ... its p statement
\providecommand{\rhtaXPctRidgeRthThreeKVsRthClose}{}\renewcommand{\rhtaXPctRidgeRthThreeKVsRthClose}{8.7}  % QLIKE % change, RTH13_N3k vs RTH13 (full set; negative: RTH13_N3k more accurate), Ridge, 16:00 bar
\providecommand{\rhtaXDmRidgeRthThreeKVsRthClose}{}\renewcommand{\rhtaXDmRidgeRthThreeKVsRthClose}{4.61}  % ... its Diebold-Mariano statistic
\providecommand{\rhtaXPRidgeRthThreeKVsRthClose}{}\renewcommand{\rhtaXPRidgeRthThreeKVsRthClose}{\ensuremath{p < 0.001}}  % ... its p statement
\providecommand{\rhtaXPctLassoRthThreeKVsRthClose}{}\renewcommand{\rhtaXPctLassoRthThreeKVsRthClose}{6.2}  % QLIKE % change, RTH13_N3k vs RTH13 (full set; negative: RTH13_N3k more accurate), Lasso, 16:00 bar
\providecommand{\rhtaXDmLassoRthThreeKVsRthClose}{}\renewcommand{\rhtaXDmLassoRthThreeKVsRthClose}{2.63}  % ... its Diebold-Mariano statistic
\providecommand{\rhtaXPLassoRthThreeKVsRthClose}{}\renewcommand{\rhtaXPLassoRthThreeKVsRthClose}{\ensuremath{p = 0.009}}  % ... its p statement
\providecommand{\rhtaXPctEnetRthThreeKVsRthClose}{}\renewcommand{\rhtaXPctEnetRthThreeKVsRthClose}{10.6}  % QLIKE % change, RTH13_N3k vs RTH13 (full set; negative: RTH13_N3k more accurate), Elastic net, 16:00 bar
\providecommand{\rhtaXDmEnetRthThreeKVsRthClose}{}\renewcommand{\rhtaXDmEnetRthThreeKVsRthClose}{5.13}  % ... its Diebold-Mariano statistic
\providecommand{\rhtaXPEnetRthThreeKVsRthClose}{}\renewcommand{\rhtaXPEnetRthThreeKVsRthClose}{\ensuremath{p < 0.001}}  % ... its p statement
\providecommand{\rhtaXPctLgbmRthThreeKVsRthClose}{}\renewcommand{\rhtaXPctLgbmRthThreeKVsRthClose}{12.3}  % QLIKE % change, RTH13_N3k vs RTH13 (full set; negative: RTH13_N3k more accurate), LightGBM, 16:00 bar
\providecommand{\rhtaXDmLgbmRthThreeKVsRthClose}{}\renewcommand{\rhtaXDmLgbmRthThreeKVsRthClose}{5.35}  % ... its Diebold-Mariano statistic
\providecommand{\rhtaXPLgbmRthThreeKVsRthClose}{}\renewcommand{\rhtaXPLgbmRthThreeKVsRthClose}{\ensuremath{p < 0.001}}  % ... its p statement
\providecommand{\rhtaXPctXgbRthThreeKVsRthClose}{}\renewcommand{\rhtaXPctXgbRthThreeKVsRthClose}{9.3}  % QLIKE % change, RTH13_N3k vs RTH13 (full set; negative: RTH13_N3k more accurate), XGBoost, 16:00 bar
\providecommand{\rhtaXDmXgbRthThreeKVsRthClose}{}\renewcommand{\rhtaXDmXgbRthThreeKVsRthClose}{5.11}  % ... its Diebold-Mariano statistic
\providecommand{\rhtaXPXgbRthThreeKVsRthClose}{}\renewcommand{\rhtaXPXgbRthThreeKVsRthClose}{\ensuremath{p < 0.001}}  % ... its p statement
\providecommand{\rhtaXPctRfRthThreeKVsRthClose}{}\renewcommand{\rhtaXPctRfRthThreeKVsRthClose}{7.1}  % QLIKE % change, RTH13_N3k vs RTH13 (full set; negative: RTH13_N3k more accurate), Random forest, 16:00 bar
\providecommand{\rhtaXDmRfRthThreeKVsRthClose}{}\renewcommand{\rhtaXDmRfRthThreeKVsRthClose}{5.16}  % ... its Diebold-Mariano statistic
\providecommand{\rhtaXPRfRthThreeKVsRthClose}{}\renewcommand{\rhtaXPRfRthThreeKVsRthClose}{\ensuremath{p < 0.001}}  % ... its p statement
\providecommand{\rhtaXPctRidgeRthThreeKVsRthBar}{}\renewcommand{\rhtaXPctRidgeRthThreeKVsRthBar}{2.7}  % QLIKE % change, RTH13_N3k vs RTH13 (full set; negative: RTH13_N3k more accurate), Ridge, every bar
\providecommand{\rhtaXDmRidgeRthThreeKVsRthBar}{}\renewcommand{\rhtaXDmRidgeRthThreeKVsRthBar}{3.01}  % ... its Diebold-Mariano statistic
\providecommand{\rhtaXPRidgeRthThreeKVsRthBar}{}\renewcommand{\rhtaXPRidgeRthThreeKVsRthBar}{\ensuremath{p = 0.003}}  % ... its p statement
\providecommand{\rhtaXPctLassoRthThreeKVsRthBar}{}\renewcommand{\rhtaXPctLassoRthThreeKVsRthBar}{3.5}  % QLIKE % change, RTH13_N3k vs RTH13 (full set; negative: RTH13_N3k more accurate), Lasso, every bar
\providecommand{\rhtaXDmLassoRthThreeKVsRthBar}{}\renewcommand{\rhtaXDmLassoRthThreeKVsRthBar}{2.71}  % ... its Diebold-Mariano statistic
\providecommand{\rhtaXPLassoRthThreeKVsRthBar}{}\renewcommand{\rhtaXPLassoRthThreeKVsRthBar}{\ensuremath{p = 0.007}}  % ... its p statement
\providecommand{\rhtaXPctEnetRthThreeKVsRthBar}{}\renewcommand{\rhtaXPctEnetRthThreeKVsRthBar}{2.9}  % QLIKE % change, RTH13_N3k vs RTH13 (full set; negative: RTH13_N3k more accurate), Elastic net, every bar
\providecommand{\rhtaXDmEnetRthThreeKVsRthBar}{}\renewcommand{\rhtaXDmEnetRthThreeKVsRthBar}{3.54}  % ... its Diebold-Mariano statistic
\providecommand{\rhtaXPEnetRthThreeKVsRthBar}{}\renewcommand{\rhtaXPEnetRthThreeKVsRthBar}{\ensuremath{p < 0.001}}  % ... its p statement
\providecommand{\rhtaXPctLgbmRthThreeKVsRthBar}{}\renewcommand{\rhtaXPctLgbmRthThreeKVsRthBar}{14.9}  % QLIKE % change, RTH13_N3k vs RTH13 (full set; negative: RTH13_N3k more accurate), LightGBM, every bar
\providecommand{\rhtaXDmLgbmRthThreeKVsRthBar}{}\renewcommand{\rhtaXDmLgbmRthThreeKVsRthBar}{8.05}  % ... its Diebold-Mariano statistic
\providecommand{\rhtaXPLgbmRthThreeKVsRthBar}{}\renewcommand{\rhtaXPLgbmRthThreeKVsRthBar}{\ensuremath{p < 0.001}}  % ... its p statement
\providecommand{\rhtaXPctXgbRthThreeKVsRthBar}{}\renewcommand{\rhtaXPctXgbRthThreeKVsRthBar}{10.7}  % QLIKE % change, RTH13_N3k vs RTH13 (full set; negative: RTH13_N3k more accurate), XGBoost, every bar
\providecommand{\rhtaXDmXgbRthThreeKVsRthBar}{}\renewcommand{\rhtaXDmXgbRthThreeKVsRthBar}{7.31}  % ... its Diebold-Mariano statistic
\providecommand{\rhtaXPXgbRthThreeKVsRthBar}{}\renewcommand{\rhtaXPXgbRthThreeKVsRthBar}{\ensuremath{p < 0.001}}  % ... its p statement
\providecommand{\rhtaXPctRfRthThreeKVsRthBar}{}\renewcommand{\rhtaXPctRfRthThreeKVsRthBar}{9.4}  % QLIKE % change, RTH13_N3k vs RTH13 (full set; negative: RTH13_N3k more accurate), Random forest, every bar
\providecommand{\rhtaXDmRfRthThreeKVsRthBar}{}\renewcommand{\rhtaXDmRfRthThreeKVsRthBar}{7.22}  % ... its Diebold-Mariano statistic
\providecommand{\rhtaXPRfRthThreeKVsRthBar}{}\renewcommand{\rhtaXPRfRthThreeKVsRthBar}{\ensuremath{p < 0.001}}  % ... its p statement
\providecommand{\rhtaXPctRidgeAllThirtyNineKVsAllClose}{}\renewcommand{\rhtaXPctRidgeAllThirtyNineKVsAllClose}{4.1}  % QLIKE % change, ALL48_N39k vs ALL48 (full set; negative: ALL48_N39k more accurate), Ridge, 16:00 bar
\providecommand{\rhtaXDmRidgeAllThirtyNineKVsAllClose}{}\renewcommand{\rhtaXDmRidgeAllThirtyNineKVsAllClose}{4.03}  % ... its Diebold-Mariano statistic
\providecommand{\rhtaXPRidgeAllThirtyNineKVsAllClose}{}\renewcommand{\rhtaXPRidgeAllThirtyNineKVsAllClose}{\ensuremath{p < 0.001}}  % ... its p statement
\providecommand{\rhtaXPctLassoAllThirtyNineKVsAllClose}{}\renewcommand{\rhtaXPctLassoAllThirtyNineKVsAllClose}{2.6}  % QLIKE % change, ALL48_N39k vs ALL48 (full set; negative: ALL48_N39k more accurate), Lasso, 16:00 bar
\providecommand{\rhtaXDmLassoAllThirtyNineKVsAllClose}{}\renewcommand{\rhtaXDmLassoAllThirtyNineKVsAllClose}{2.67}  % ... its Diebold-Mariano statistic
\providecommand{\rhtaXPLassoAllThirtyNineKVsAllClose}{}\renewcommand{\rhtaXPLassoAllThirtyNineKVsAllClose}{\ensuremath{p = 0.008}}  % ... its p statement
\providecommand{\rhtaXPctEnetAllThirtyNineKVsAllClose}{}\renewcommand{\rhtaXPctEnetAllThirtyNineKVsAllClose}{2.2}  % QLIKE % change, ALL48_N39k vs ALL48 (full set; negative: ALL48_N39k more accurate), Elastic net, 16:00 bar
\providecommand{\rhtaXDmEnetAllThirtyNineKVsAllClose}{}\renewcommand{\rhtaXDmEnetAllThirtyNineKVsAllClose}{1.89}  % ... its Diebold-Mariano statistic
\providecommand{\rhtaXPEnetAllThirtyNineKVsAllClose}{}\renewcommand{\rhtaXPEnetAllThirtyNineKVsAllClose}{\ensuremath{p = 0.059}}  % ... its p statement
\providecommand{\rhtaXPctLgbmAllThirtyNineKVsAllClose}{}\renewcommand{\rhtaXPctLgbmAllThirtyNineKVsAllClose}{2.7}  % QLIKE % change, ALL48_N39k vs ALL48 (full set; negative: ALL48_N39k more accurate), LightGBM, 16:00 bar
\providecommand{\rhtaXDmLgbmAllThirtyNineKVsAllClose}{}\renewcommand{\rhtaXDmLgbmAllThirtyNineKVsAllClose}{2.87}  % ... its Diebold-Mariano statistic
\providecommand{\rhtaXPLgbmAllThirtyNineKVsAllClose}{}\renewcommand{\rhtaXPLgbmAllThirtyNineKVsAllClose}{\ensuremath{p = 0.004}}  % ... its p statement
\providecommand{\rhtaXPctXgbAllThirtyNineKVsAllClose}{}\renewcommand{\rhtaXPctXgbAllThirtyNineKVsAllClose}{3.3}  % QLIKE % change, ALL48_N39k vs ALL48 (full set; negative: ALL48_N39k more accurate), XGBoost, 16:00 bar
\providecommand{\rhtaXDmXgbAllThirtyNineKVsAllClose}{}\renewcommand{\rhtaXDmXgbAllThirtyNineKVsAllClose}{4.05}  % ... its Diebold-Mariano statistic
\providecommand{\rhtaXPXgbAllThirtyNineKVsAllClose}{}\renewcommand{\rhtaXPXgbAllThirtyNineKVsAllClose}{\ensuremath{p < 0.001}}  % ... its p statement
\providecommand{\rhtaXPctRfAllThirtyNineKVsAllClose}{}\renewcommand{\rhtaXPctRfAllThirtyNineKVsAllClose}{1.7}  % QLIKE % change, ALL48_N39k vs ALL48 (full set; negative: ALL48_N39k more accurate), Random forest, 16:00 bar
\providecommand{\rhtaXDmRfAllThirtyNineKVsAllClose}{}\renewcommand{\rhtaXDmRfAllThirtyNineKVsAllClose}{2.38}  % ... its Diebold-Mariano statistic
\providecommand{\rhtaXPRfAllThirtyNineKVsAllClose}{}\renewcommand{\rhtaXPRfAllThirtyNineKVsAllClose}{\ensuremath{p = 0.017}}  % ... its p statement
\providecommand{\rhtaXPctRidgeAllThirtyNineKVsAllBar}{}\renewcommand{\rhtaXPctRidgeAllThirtyNineKVsAllBar}{0.8}  % QLIKE % change, ALL48_N39k vs ALL48 (full set; negative: ALL48_N39k more accurate), Ridge, every bar
\providecommand{\rhtaXDmRidgeAllThirtyNineKVsAllBar}{}\renewcommand{\rhtaXDmRidgeAllThirtyNineKVsAllBar}{1.56}  % ... its Diebold-Mariano statistic
\providecommand{\rhtaXPRidgeAllThirtyNineKVsAllBar}{}\renewcommand{\rhtaXPRidgeAllThirtyNineKVsAllBar}{\ensuremath{p = 0.12}}  % ... its p statement
\providecommand{\rhtaXPctLassoAllThirtyNineKVsAllBar}{}\renewcommand{\rhtaXPctLassoAllThirtyNineKVsAllBar}{1.6}  % QLIKE % change, ALL48_N39k vs ALL48 (full set; negative: ALL48_N39k more accurate), Lasso, every bar
\providecommand{\rhtaXDmLassoAllThirtyNineKVsAllBar}{}\renewcommand{\rhtaXDmLassoAllThirtyNineKVsAllBar}{3.26}  % ... its Diebold-Mariano statistic
\providecommand{\rhtaXPLassoAllThirtyNineKVsAllBar}{}\renewcommand{\rhtaXPLassoAllThirtyNineKVsAllBar}{\ensuremath{p = 0.001}}  % ... its p statement
\providecommand{\rhtaXPctEnetAllThirtyNineKVsAllBar}{}\renewcommand{\rhtaXPctEnetAllThirtyNineKVsAllBar}{2.7}  % QLIKE % change, ALL48_N39k vs ALL48 (full set; negative: ALL48_N39k more accurate), Elastic net, every bar
\providecommand{\rhtaXDmEnetAllThirtyNineKVsAllBar}{}\renewcommand{\rhtaXDmEnetAllThirtyNineKVsAllBar}{3.42}  % ... its Diebold-Mariano statistic
\providecommand{\rhtaXPEnetAllThirtyNineKVsAllBar}{}\renewcommand{\rhtaXPEnetAllThirtyNineKVsAllBar}{\ensuremath{p < 0.001}}  % ... its p statement
\providecommand{\rhtaXPctLgbmAllThirtyNineKVsAllBar}{}\renewcommand{\rhtaXPctLgbmAllThirtyNineKVsAllBar}{5.7}  % QLIKE % change, ALL48_N39k vs ALL48 (full set; negative: ALL48_N39k more accurate), LightGBM, every bar
\providecommand{\rhtaXDmLgbmAllThirtyNineKVsAllBar}{}\renewcommand{\rhtaXDmLgbmAllThirtyNineKVsAllBar}{6.76}  % ... its Diebold-Mariano statistic
\providecommand{\rhtaXPLgbmAllThirtyNineKVsAllBar}{}\renewcommand{\rhtaXPLgbmAllThirtyNineKVsAllBar}{\ensuremath{p < 0.001}}  % ... its p statement
\providecommand{\rhtaXPctXgbAllThirtyNineKVsAllBar}{}\renewcommand{\rhtaXPctXgbAllThirtyNineKVsAllBar}{3.2}  % QLIKE % change, ALL48_N39k vs ALL48 (full set; negative: ALL48_N39k more accurate), XGBoost, every bar
\providecommand{\rhtaXDmXgbAllThirtyNineKVsAllBar}{}\renewcommand{\rhtaXDmXgbAllThirtyNineKVsAllBar}{6.86}  % ... its Diebold-Mariano statistic
\providecommand{\rhtaXPXgbAllThirtyNineKVsAllBar}{}\renewcommand{\rhtaXPXgbAllThirtyNineKVsAllBar}{\ensuremath{p < 0.001}}  % ... its p statement
\providecommand{\rhtaXPctRfAllThirtyNineKVsAllBar}{}\renewcommand{\rhtaXPctRfAllThirtyNineKVsAllBar}{2.6}  % QLIKE % change, ALL48_N39k vs ALL48 (full set; negative: ALL48_N39k more accurate), Random forest, every bar
\providecommand{\rhtaXDmRfAllThirtyNineKVsAllBar}{}\renewcommand{\rhtaXDmRfAllThirtyNineKVsAllBar}{4.91}  % ... its Diebold-Mariano statistic
\providecommand{\rhtaXPRfAllThirtyNineKVsAllBar}{}\renewcommand{\rhtaXPRfAllThirtyNineKVsAllBar}{\ensuremath{p < 0.001}}  % ... its p statement
\providecommand{\rhtaShRidgeHarRth}{}\renewcommand{\rhtaShRidgeHarRth}{0.95}  % Sharpe ratio (D19), mid-quotes: Ridge, baseline, RTH13 (main design: the 13 regular-hours bars, 3,000 days)
\providecommand{\rhtaShXRidgeHarRth}{}\renewcommand{\rhtaShXRidgeHarRth}{0.60}  % ... at quoted bid/ask prices
\providecommand{\rhtaShRidgeFullRth}{}\renewcommand{\rhtaShRidgeFullRth}{1.77}  % Sharpe ratio (D19), mid-quotes: Ridge, all_features, RTH13 (main design: the 13 regular-hours bars, 3,000 days)
\providecommand{\rhtaShXRidgeFullRth}{}\renewcommand{\rhtaShXRidgeFullRth}{1.41}  % ... at quoted bid/ask prices
\providecommand{\rhtaPShRidgeRth}{}\renewcommand{\rhtaPShRidgeRth}{\ensuremath{p = 0.11}}  % Ledoit-Wolf p, Sharpe ratio full minus HAR (mid): Ridge, RTH13
\providecommand{\rhtaShLassoHarRth}{}\renewcommand{\rhtaShLassoHarRth}{0.95}  % Sharpe ratio (D19), mid-quotes: Lasso, baseline, RTH13 (main design: the 13 regular-hours bars, 3,000 days)
\providecommand{\rhtaShXLassoHarRth}{}\renewcommand{\rhtaShXLassoHarRth}{0.60}  % ... at quoted bid/ask prices
\providecommand{\rhtaShLassoFullRth}{}\renewcommand{\rhtaShLassoFullRth}{1.72}  % Sharpe ratio (D19), mid-quotes: Lasso, all_features, RTH13 (main design: the 13 regular-hours bars, 3,000 days)
\providecommand{\rhtaShXLassoFullRth}{}\renewcommand{\rhtaShXLassoFullRth}{1.37}  % ... at quoted bid/ask prices
\providecommand{\rhtaPShLassoRth}{}\renewcommand{\rhtaPShLassoRth}{\ensuremath{p = 0.12}}  % Ledoit-Wolf p, Sharpe ratio full minus HAR (mid): Lasso, RTH13
\providecommand{\rhtaShEnetHarRth}{}\renewcommand{\rhtaShEnetHarRth}{0.95}  % Sharpe ratio (D19), mid-quotes: Elastic net, baseline, RTH13 (main design: the 13 regular-hours bars, 3,000 days)
\providecommand{\rhtaShXEnetHarRth}{}\renewcommand{\rhtaShXEnetHarRth}{0.60}  % ... at quoted bid/ask prices
\providecommand{\rhtaShEnetFullRth}{}\renewcommand{\rhtaShEnetFullRth}{1.86}  % Sharpe ratio (D19), mid-quotes: Elastic net, all_features, RTH13 (main design: the 13 regular-hours bars, 3,000 days)
\providecommand{\rhtaShXEnetFullRth}{}\renewcommand{\rhtaShXEnetFullRth}{1.50}  % ... at quoted bid/ask prices
\providecommand{\rhtaPShEnetRth}{}\renewcommand{\rhtaPShEnetRth}{\ensuremath{p = 0.077}}  % Ledoit-Wolf p, Sharpe ratio full minus HAR (mid): Elastic net, RTH13
\providecommand{\rhtaShLgbmHarRth}{}\renewcommand{\rhtaShLgbmHarRth}{0.65}  % Sharpe ratio (D19), mid-quotes: LightGBM, baseline, RTH13 (main design: the 13 regular-hours bars, 3,000 days)
\providecommand{\rhtaShXLgbmHarRth}{}\renewcommand{\rhtaShXLgbmHarRth}{0.31}  % ... at quoted bid/ask prices
\providecommand{\rhtaShLgbmFullRth}{}\renewcommand{\rhtaShLgbmFullRth}{1.59}  % Sharpe ratio (D19), mid-quotes: LightGBM, all_features, RTH13 (main design: the 13 regular-hours bars, 3,000 days)
\providecommand{\rhtaShXLgbmFullRth}{}\renewcommand{\rhtaShXLgbmFullRth}{1.23}  % ... at quoted bid/ask prices
\providecommand{\rhtaPShLgbmRth}{}\renewcommand{\rhtaPShLgbmRth}{\ensuremath{p = 0.040}}  % Ledoit-Wolf p, Sharpe ratio full minus HAR (mid): LightGBM, RTH13
\providecommand{\rhtaShXgbHarRth}{}\renewcommand{\rhtaShXgbHarRth}{0.98}  % Sharpe ratio (D19), mid-quotes: XGBoost, baseline, RTH13 (main design: the 13 regular-hours bars, 3,000 days)
\providecommand{\rhtaShXXgbHarRth}{}\renewcommand{\rhtaShXXgbHarRth}{0.63}  % ... at quoted bid/ask prices
\providecommand{\rhtaShXgbFullRth}{}\renewcommand{\rhtaShXgbFullRth}{1.47}  % Sharpe ratio (D19), mid-quotes: XGBoost, all_features, RTH13 (main design: the 13 regular-hours bars, 3,000 days)
\providecommand{\rhtaShXXgbFullRth}{}\renewcommand{\rhtaShXXgbFullRth}{1.12}  % ... at quoted bid/ask prices
\providecommand{\rhtaPShXgbRth}{}\renewcommand{\rhtaPShXgbRth}{\ensuremath{p = 0.27}}  % Ledoit-Wolf p, Sharpe ratio full minus HAR (mid): XGBoost, RTH13
\providecommand{\rhtaShRfHarRth}{}\renewcommand{\rhtaShRfHarRth}{0.32}  % Sharpe ratio (D19), mid-quotes: Random forest, baseline, RTH13 (main design: the 13 regular-hours bars, 3,000 days)
\providecommand{\rhtaShXRfHarRth}{}\renewcommand{\rhtaShXRfHarRth}{\ensuremath{-}0.03}  % ... at quoted bid/ask prices
\providecommand{\rhtaShRfFullRth}{}\renewcommand{\rhtaShRfFullRth}{1.88}  % Sharpe ratio (D19), mid-quotes: Random forest, all_features, RTH13 (main design: the 13 regular-hours bars, 3,000 days)
\providecommand{\rhtaShXRfFullRth}{}\renewcommand{\rhtaShXRfFullRth}{1.52}  % ... at quoted bid/ask prices
\providecommand{\rhtaPShRfRth}{}\renewcommand{\rhtaPShRfRth}{\ensuremath{p = 0.020}}  % Ledoit-Wolf p, Sharpe ratio full minus HAR (mid): Random forest, RTH13
\providecommand{\rhtaShRidgeHarAll}{}\renewcommand{\rhtaShRidgeHarAll}{0.95}  % Sharpe ratio (D19), mid-quotes: Ridge, baseline, ALL48 (all 48 intervals, 3,000 days)
\providecommand{\rhtaShXRidgeHarAll}{}\renewcommand{\rhtaShXRidgeHarAll}{0.59}  % ... at quoted bid/ask prices
\providecommand{\rhtaShRidgeFullAll}{}\renewcommand{\rhtaShRidgeFullAll}{1.61}  % Sharpe ratio (D19), mid-quotes: Ridge, all_features, ALL48 (all 48 intervals, 3,000 days)
\providecommand{\rhtaShXRidgeFullAll}{}\renewcommand{\rhtaShXRidgeFullAll}{1.26}  % ... at quoted bid/ask prices
\providecommand{\rhtaShLassoHarAll}{}\renewcommand{\rhtaShLassoHarAll}{0.93}  % Sharpe ratio (D19), mid-quotes: Lasso, baseline, ALL48 (all 48 intervals, 3,000 days)
\providecommand{\rhtaShXLassoHarAll}{}\renewcommand{\rhtaShXLassoHarAll}{0.58}  % ... at quoted bid/ask prices
\providecommand{\rhtaShLassoFullAll}{}\renewcommand{\rhtaShLassoFullAll}{1.50}  % Sharpe ratio (D19), mid-quotes: Lasso, all_features, ALL48 (all 48 intervals, 3,000 days)
\providecommand{\rhtaShXLassoFullAll}{}\renewcommand{\rhtaShXLassoFullAll}{1.14}  % ... at quoted bid/ask prices
\providecommand{\rhtaShEnetHarAll}{}\renewcommand{\rhtaShEnetHarAll}{0.89}  % Sharpe ratio (D19), mid-quotes: Elastic net, baseline, ALL48 (all 48 intervals, 3,000 days)
\providecommand{\rhtaShXEnetHarAll}{}\renewcommand{\rhtaShXEnetHarAll}{0.54}  % ... at quoted bid/ask prices
\providecommand{\rhtaShEnetFullAll}{}\renewcommand{\rhtaShEnetFullAll}{1.52}  % Sharpe ratio (D19), mid-quotes: Elastic net, all_features, ALL48 (all 48 intervals, 3,000 days)
\providecommand{\rhtaShXEnetFullAll}{}\renewcommand{\rhtaShXEnetFullAll}{1.17}  % ... at quoted bid/ask prices
\providecommand{\rhtaShLgbmHarAll}{}\renewcommand{\rhtaShLgbmHarAll}{0.21}  % Sharpe ratio (D19), mid-quotes: LightGBM, baseline, ALL48 (all 48 intervals, 3,000 days)
\providecommand{\rhtaShXLgbmHarAll}{}\renewcommand{\rhtaShXLgbmHarAll}{\ensuremath{-}0.14}  % ... at quoted bid/ask prices
\providecommand{\rhtaShLgbmFullAll}{}\renewcommand{\rhtaShLgbmFullAll}{1.25}  % Sharpe ratio (D19), mid-quotes: LightGBM, all_features, ALL48 (all 48 intervals, 3,000 days)
\providecommand{\rhtaShXLgbmFullAll}{}\renewcommand{\rhtaShXLgbmFullAll}{0.90}  % ... at quoted bid/ask prices
\providecommand{\rhtaPShLgbmAll}{}\renewcommand{\rhtaPShLgbmAll}{\ensuremath{p = 0.11}}  % Ledoit-Wolf p, Sharpe ratio full minus HAR (mid): LightGBM, ALL48
\providecommand{\rhtaShXgbHarAll}{}\renewcommand{\rhtaShXgbHarAll}{0.73}  % Sharpe ratio (D19), mid-quotes: XGBoost, baseline, ALL48 (all 48 intervals, 3,000 days)
\providecommand{\rhtaShXXgbHarAll}{}\renewcommand{\rhtaShXXgbHarAll}{0.38}  % ... at quoted bid/ask prices
\providecommand{\rhtaShXgbFullAll}{}\renewcommand{\rhtaShXgbFullAll}{1.50}  % Sharpe ratio (D19), mid-quotes: XGBoost, all_features, ALL48 (all 48 intervals, 3,000 days)
\providecommand{\rhtaShXXgbFullAll}{}\renewcommand{\rhtaShXXgbFullAll}{1.15}  % ... at quoted bid/ask prices
\providecommand{\rhtaPShXgbAll}{}\renewcommand{\rhtaPShXgbAll}{\ensuremath{p = 0.11}}  % Ledoit-Wolf p, Sharpe ratio full minus HAR (mid): XGBoost, ALL48
\providecommand{\rhtaShRfHarAll}{}\renewcommand{\rhtaShRfHarAll}{0.65}  % Sharpe ratio (D19), mid-quotes: Random forest, baseline, ALL48 (all 48 intervals, 3,000 days)
\providecommand{\rhtaShXRfHarAll}{}\renewcommand{\rhtaShXRfHarAll}{0.29}  % ... at quoted bid/ask prices
\providecommand{\rhtaShRfFullAll}{}\renewcommand{\rhtaShRfFullAll}{1.82}  % Sharpe ratio (D19), mid-quotes: Random forest, all_features, ALL48 (all 48 intervals, 3,000 days)
\providecommand{\rhtaShXRfFullAll}{}\renewcommand{\rhtaShXRfFullAll}{1.46}  % ... at quoted bid/ask prices
\providecommand{\rhtaPShRfAll}{}\renewcommand{\rhtaPShRfAll}{\ensuremath{p = 0.066}}  % Ledoit-Wolf p, Sharpe ratio full minus HAR (mid): Random forest, ALL48
\providecommand{\rhtaShRidgeHarCl}{}\renewcommand{\rhtaShRidgeHarCl}{1.03}  % Sharpe ratio (D19), mid-quotes: Ridge, baseline, CLOSE (the closing half hour only, 3,000 closes)
\providecommand{\rhtaShXRidgeHarCl}{}\renewcommand{\rhtaShXRidgeHarCl}{0.68}  % ... at quoted bid/ask prices
\providecommand{\rhtaShRidgeFullCl}{}\renewcommand{\rhtaShRidgeFullCl}{1.26}  % Sharpe ratio (D19), mid-quotes: Ridge, all_features, CLOSE (the closing half hour only, 3,000 closes)
\providecommand{\rhtaShXRidgeFullCl}{}\renewcommand{\rhtaShXRidgeFullCl}{0.91}  % ... at quoted bid/ask prices
\providecommand{\rhtaShLassoHarCl}{}\renewcommand{\rhtaShLassoHarCl}{1.03}  % Sharpe ratio (D19), mid-quotes: Lasso, baseline, CLOSE (the closing half hour only, 3,000 closes)
\providecommand{\rhtaShXLassoHarCl}{}\renewcommand{\rhtaShXLassoHarCl}{0.68}  % ... at quoted bid/ask prices
\providecommand{\rhtaShLassoFullCl}{}\renewcommand{\rhtaShLassoFullCl}{1.58}  % Sharpe ratio (D19), mid-quotes: Lasso, all_features, CLOSE (the closing half hour only, 3,000 closes)
\providecommand{\rhtaShXLassoFullCl}{}\renewcommand{\rhtaShXLassoFullCl}{1.23}  % ... at quoted bid/ask prices
\providecommand{\rhtaShEnetHarCl}{}\renewcommand{\rhtaShEnetHarCl}{1.03}  % Sharpe ratio (D19), mid-quotes: Elastic net, baseline, CLOSE (the closing half hour only, 3,000 closes)
\providecommand{\rhtaShXEnetHarCl}{}\renewcommand{\rhtaShXEnetHarCl}{0.68}  % ... at quoted bid/ask prices
\providecommand{\rhtaShEnetFullCl}{}\renewcommand{\rhtaShEnetFullCl}{1.71}  % Sharpe ratio (D19), mid-quotes: Elastic net, all_features, CLOSE (the closing half hour only, 3,000 closes)
\providecommand{\rhtaShXEnetFullCl}{}\renewcommand{\rhtaShXEnetFullCl}{1.36}  % ... at quoted bid/ask prices
\providecommand{\rhtaShLgbmHarCl}{}\renewcommand{\rhtaShLgbmHarCl}{\ensuremath{-}0.03}  % Sharpe ratio (D19), mid-quotes: LightGBM, baseline, CLOSE (the closing half hour only, 3,000 closes)
\providecommand{\rhtaShXLgbmHarCl}{}\renewcommand{\rhtaShXLgbmHarCl}{\ensuremath{-}0.37}  % ... at quoted bid/ask prices
\providecommand{\rhtaShLgbmFullCl}{}\renewcommand{\rhtaShLgbmFullCl}{1.94}  % Sharpe ratio (D19), mid-quotes: LightGBM, all_features, CLOSE (the closing half hour only, 3,000 closes)
\providecommand{\rhtaShXLgbmFullCl}{}\renewcommand{\rhtaShXLgbmFullCl}{1.60}  % ... at quoted bid/ask prices
\providecommand{\rhtaPShLgbmCl}{}\renewcommand{\rhtaPShLgbmCl}{\ensuremath{p = 0.003}}  % Ledoit-Wolf p, Sharpe ratio full minus HAR (mid): LightGBM, CLOSE
\providecommand{\rhtaShXgbHarCl}{}\renewcommand{\rhtaShXgbHarCl}{0.40}  % Sharpe ratio (D19), mid-quotes: XGBoost, baseline, CLOSE (the closing half hour only, 3,000 closes)
\providecommand{\rhtaShXXgbHarCl}{}\renewcommand{\rhtaShXXgbHarCl}{0.05}  % ... at quoted bid/ask prices
\providecommand{\rhtaShXgbFullCl}{}\renewcommand{\rhtaShXgbFullCl}{1.62}  % Sharpe ratio (D19), mid-quotes: XGBoost, all_features, CLOSE (the closing half hour only, 3,000 closes)
\providecommand{\rhtaShXXgbFullCl}{}\renewcommand{\rhtaShXXgbFullCl}{1.27}  % ... at quoted bid/ask prices
\providecommand{\rhtaPShXgbCl}{}\renewcommand{\rhtaPShXgbCl}{\ensuremath{p = 0.055}}  % Ledoit-Wolf p, Sharpe ratio full minus HAR (mid): XGBoost, CLOSE
\providecommand{\rhtaShRfHarCl}{}\renewcommand{\rhtaShRfHarCl}{0.29}  % Sharpe ratio (D19), mid-quotes: Random forest, baseline, CLOSE (the closing half hour only, 3,000 closes)
\providecommand{\rhtaShXRfHarCl}{}\renewcommand{\rhtaShXRfHarCl}{\ensuremath{-}0.05}  % ... at quoted bid/ask prices
\providecommand{\rhtaShRfFullCl}{}\renewcommand{\rhtaShRfFullCl}{1.29}  % Sharpe ratio (D19), mid-quotes: Random forest, all_features, CLOSE (the closing half hour only, 3,000 closes)
\providecommand{\rhtaShXRfFullCl}{}\renewcommand{\rhtaShXRfFullCl}{0.94}  % ... at quoted bid/ask prices
\providecommand{\rhtaPShRfCl}{}\renewcommand{\rhtaPShRfCl}{\ensuremath{p = 0.14}}  % Ledoit-Wolf p, Sharpe ratio full minus HAR (mid): Random forest, CLOSE
\providecommand{\rhtaShRidgeHarClNat}{}\renewcommand{\rhtaShRidgeHarClNat}{0.94}  % Sharpe ratio (D19), mid-quotes: Ridge, baseline, CLOSE_NAT (closing rows, design built natively on them; linear only)
\providecommand{\rhtaShXRidgeHarClNat}{}\renewcommand{\rhtaShXRidgeHarClNat}{0.59}  % ... at quoted bid/ask prices
\providecommand{\rhtaShRidgeFullClNat}{}\renewcommand{\rhtaShRidgeFullClNat}{1.24}  % Sharpe ratio (D19), mid-quotes: Ridge, all_features, CLOSE_NAT (closing rows, design built natively on them; linear only)
\providecommand{\rhtaShXRidgeFullClNat}{}\renewcommand{\rhtaShXRidgeFullClNat}{0.89}  % ... at quoted bid/ask prices
\providecommand{\rhtaShLassoHarClNat}{}\renewcommand{\rhtaShLassoHarClNat}{0.94}  % Sharpe ratio (D19), mid-quotes: Lasso, baseline, CLOSE_NAT (closing rows, design built natively on them; linear only)
\providecommand{\rhtaShXLassoHarClNat}{}\renewcommand{\rhtaShXLassoHarClNat}{0.59}  % ... at quoted bid/ask prices
\providecommand{\rhtaShLassoFullClNat}{}\renewcommand{\rhtaShLassoFullClNat}{1.23}  % Sharpe ratio (D19), mid-quotes: Lasso, all_features, CLOSE_NAT (closing rows, design built natively on them; linear only)
\providecommand{\rhtaShXLassoFullClNat}{}\renewcommand{\rhtaShXLassoFullClNat}{0.89}  % ... at quoted bid/ask prices
\providecommand{\rhtaShEnetHarClNat}{}\renewcommand{\rhtaShEnetHarClNat}{0.94}  % Sharpe ratio (D19), mid-quotes: Elastic net, baseline, CLOSE_NAT (closing rows, design built natively on them; linear only)
\providecommand{\rhtaShXEnetHarClNat}{}\renewcommand{\rhtaShXEnetHarClNat}{0.59}  % ... at quoted bid/ask prices
\providecommand{\rhtaShEnetFullClNat}{}\renewcommand{\rhtaShEnetFullClNat}{1.71}  % Sharpe ratio (D19), mid-quotes: Elastic net, all_features, CLOSE_NAT (closing rows, design built natively on them; linear only)
\providecommand{\rhtaShXEnetFullClNat}{}\renewcommand{\rhtaShXEnetFullClNat}{1.37}  % ... at quoted bid/ask prices
\providecommand{\rhtaShRidgeHarAllCt}{}\renewcommand{\rhtaShRidgeHarAllCt}{0.95}  % Sharpe ratio (D19), mid-quotes: Ridge, baseline, ALL48@16 (ALL48, penalties chosen on the 16:00 validation rows; linear only)
\providecommand{\rhtaShXRidgeHarAllCt}{}\renewcommand{\rhtaShXRidgeHarAllCt}{0.59}  % ... at quoted bid/ask prices
\providecommand{\rhtaShRidgeFullAllCt}{}\renewcommand{\rhtaShRidgeFullAllCt}{1.40}  % Sharpe ratio (D19), mid-quotes: Ridge, all_features, ALL48@16 (ALL48, penalties chosen on the 16:00 validation rows; linear only)
\providecommand{\rhtaShXRidgeFullAllCt}{}\renewcommand{\rhtaShXRidgeFullAllCt}{1.05}  % ... at quoted bid/ask prices
\providecommand{\rhtaShLassoHarAllCt}{}\renewcommand{\rhtaShLassoHarAllCt}{0.96}  % Sharpe ratio (D19), mid-quotes: Lasso, baseline, ALL48@16 (ALL48, penalties chosen on the 16:00 validation rows; linear only)
\providecommand{\rhtaShXLassoHarAllCt}{}\renewcommand{\rhtaShXLassoHarAllCt}{0.60}  % ... at quoted bid/ask prices
\providecommand{\rhtaShLassoFullAllCt}{}\renewcommand{\rhtaShLassoFullAllCt}{1.35}  % Sharpe ratio (D19), mid-quotes: Lasso, all_features, ALL48@16 (ALL48, penalties chosen on the 16:00 validation rows; linear only)
\providecommand{\rhtaShXLassoFullAllCt}{}\renewcommand{\rhtaShXLassoFullAllCt}{1.00}  % ... at quoted bid/ask prices
\providecommand{\rhtaShEnetHarAllCt}{}\renewcommand{\rhtaShEnetHarAllCt}{0.80}  % Sharpe ratio (D19), mid-quotes: Elastic net, baseline, ALL48@16 (ALL48, penalties chosen on the 16:00 validation rows; linear only)
\providecommand{\rhtaShXEnetHarAllCt}{}\renewcommand{\rhtaShXEnetHarAllCt}{0.45}  % ... at quoted bid/ask prices
\providecommand{\rhtaShEnetFullAllCt}{}\renewcommand{\rhtaShEnetFullAllCt}{1.52}  % Sharpe ratio (D19), mid-quotes: Elastic net, all_features, ALL48@16 (ALL48, penalties chosen on the 16:00 validation rows; linear only)
\providecommand{\rhtaShXEnetFullAllCt}{}\renewcommand{\rhtaShXEnetFullAllCt}{1.17}  % ... at quoted bid/ask prices
\providecommand{\rhtaShRidgeHarRthCt}{}\renewcommand{\rhtaShRidgeHarRthCt}{0.95}  % Sharpe ratio (D19), mid-quotes: Ridge, baseline, RTH13@16 (RTH13, penalties chosen on the 16:00 validation rows; linear only)
\providecommand{\rhtaShXRidgeHarRthCt}{}\renewcommand{\rhtaShXRidgeHarRthCt}{0.60}  % ... at quoted bid/ask prices
\providecommand{\rhtaShRidgeFullRthCt}{}\renewcommand{\rhtaShRidgeFullRthCt}{1.77}  % Sharpe ratio (D19), mid-quotes: Ridge, all_features, RTH13@16 (RTH13, penalties chosen on the 16:00 validation rows; linear only)
\providecommand{\rhtaShXRidgeFullRthCt}{}\renewcommand{\rhtaShXRidgeFullRthCt}{1.41}  % ... at quoted bid/ask prices
\providecommand{\rhtaShLassoHarRthCt}{}\renewcommand{\rhtaShLassoHarRthCt}{1.01}  % Sharpe ratio (D19), mid-quotes: Lasso, baseline, RTH13@16 (RTH13, penalties chosen on the 16:00 validation rows; linear only)
\providecommand{\rhtaShXLassoHarRthCt}{}\renewcommand{\rhtaShXLassoHarRthCt}{0.65}  % ... at quoted bid/ask prices
\providecommand{\rhtaShLassoFullRthCt}{}\renewcommand{\rhtaShLassoFullRthCt}{1.88}  % Sharpe ratio (D19), mid-quotes: Lasso, all_features, RTH13@16 (RTH13, penalties chosen on the 16:00 validation rows; linear only)
\providecommand{\rhtaShXLassoFullRthCt}{}\renewcommand{\rhtaShXLassoFullRthCt}{1.52}  % ... at quoted bid/ask prices
\providecommand{\rhtaShEnetHarRthCt}{}\renewcommand{\rhtaShEnetHarRthCt}{1.01}  % Sharpe ratio (D19), mid-quotes: Elastic net, baseline, RTH13@16 (RTH13, penalties chosen on the 16:00 validation rows; linear only)
\providecommand{\rhtaShXEnetHarRthCt}{}\renewcommand{\rhtaShXEnetHarRthCt}{0.66}  % ... at quoted bid/ask prices
\providecommand{\rhtaShEnetFullRthCt}{}\renewcommand{\rhtaShEnetFullRthCt}{1.86}  % Sharpe ratio (D19), mid-quotes: Elastic net, all_features, RTH13@16 (RTH13, penalties chosen on the 16:00 validation rows; linear only)
\providecommand{\rhtaShXEnetFullRthCt}{}\renewcommand{\rhtaShXEnetFullRthCt}{1.50}  % ... at quoted bid/ask prices
\providecommand{\rhtaShRidgeHarAllThreeK}{}\renewcommand{\rhtaShRidgeHarAllThreeK}{1.56}  % Sharpe ratio (D19), mid-quotes: Ridge, baseline, ALL48_N3k (ALL48, the last 3,000 rows)
\providecommand{\rhtaShXRidgeHarAllThreeK}{}\renewcommand{\rhtaShXRidgeHarAllThreeK}{1.21}  % ... at quoted bid/ask prices
\providecommand{\rhtaShRidgeFullAllThreeK}{}\renewcommand{\rhtaShRidgeFullAllThreeK}{1.45}  % Sharpe ratio (D19), mid-quotes: Ridge, all_features, ALL48_N3k (ALL48, the last 3,000 rows)
\providecommand{\rhtaShXRidgeFullAllThreeK}{}\renewcommand{\rhtaShXRidgeFullAllThreeK}{1.10}  % ... at quoted bid/ask prices
\providecommand{\rhtaShLassoHarAllThreeK}{}\renewcommand{\rhtaShLassoHarAllThreeK}{1.57}  % Sharpe ratio (D19), mid-quotes: Lasso, baseline, ALL48_N3k (ALL48, the last 3,000 rows)
\providecommand{\rhtaShXLassoHarAllThreeK}{}\renewcommand{\rhtaShXLassoHarAllThreeK}{1.21}  % ... at quoted bid/ask prices
\providecommand{\rhtaShLassoFullAllThreeK}{}\renewcommand{\rhtaShLassoFullAllThreeK}{1.41}  % Sharpe ratio (D19), mid-quotes: Lasso, all_features, ALL48_N3k (ALL48, the last 3,000 rows)
\providecommand{\rhtaShXLassoFullAllThreeK}{}\renewcommand{\rhtaShXLassoFullAllThreeK}{1.06}  % ... at quoted bid/ask prices
\providecommand{\rhtaShEnetHarAllThreeK}{}\renewcommand{\rhtaShEnetHarAllThreeK}{1.52}  % Sharpe ratio (D19), mid-quotes: Elastic net, baseline, ALL48_N3k (ALL48, the last 3,000 rows)
\providecommand{\rhtaShXEnetHarAllThreeK}{}\renewcommand{\rhtaShXEnetHarAllThreeK}{1.16}  % ... at quoted bid/ask prices
\providecommand{\rhtaShEnetFullAllThreeK}{}\renewcommand{\rhtaShEnetFullAllThreeK}{1.63}  % Sharpe ratio (D19), mid-quotes: Elastic net, all_features, ALL48_N3k (ALL48, the last 3,000 rows)
\providecommand{\rhtaShXEnetFullAllThreeK}{}\renewcommand{\rhtaShXEnetFullAllThreeK}{1.27}  % ... at quoted bid/ask prices
\providecommand{\rhtaShLgbmHarAllThreeK}{}\renewcommand{\rhtaShLgbmHarAllThreeK}{0.92}  % Sharpe ratio (D19), mid-quotes: LightGBM, baseline, ALL48_N3k (ALL48, the last 3,000 rows)
\providecommand{\rhtaShXLgbmHarAllThreeK}{}\renewcommand{\rhtaShXLgbmHarAllThreeK}{0.57}  % ... at quoted bid/ask prices
\providecommand{\rhtaShLgbmFullAllThreeK}{}\renewcommand{\rhtaShLgbmFullAllThreeK}{1.35}  % Sharpe ratio (D19), mid-quotes: LightGBM, all_features, ALL48_N3k (ALL48, the last 3,000 rows)
\providecommand{\rhtaShXLgbmFullAllThreeK}{}\renewcommand{\rhtaShXLgbmFullAllThreeK}{0.99}  % ... at quoted bid/ask prices
\providecommand{\rhtaPShLgbmAllThreeK}{}\renewcommand{\rhtaPShLgbmAllThreeK}{\ensuremath{p = 0.20}}  % Ledoit-Wolf p, Sharpe ratio full minus HAR (mid): LightGBM, ALL48_N3k
\providecommand{\rhtaShXgbHarAllThreeK}{}\renewcommand{\rhtaShXgbHarAllThreeK}{0.93}  % Sharpe ratio (D19), mid-quotes: XGBoost, baseline, ALL48_N3k (ALL48, the last 3,000 rows)
\providecommand{\rhtaShXXgbHarAllThreeK}{}\renewcommand{\rhtaShXXgbHarAllThreeK}{0.58}  % ... at quoted bid/ask prices
\providecommand{\rhtaShXgbFullAllThreeK}{}\renewcommand{\rhtaShXgbFullAllThreeK}{1.24}  % Sharpe ratio (D19), mid-quotes: XGBoost, all_features, ALL48_N3k (ALL48, the last 3,000 rows)
\providecommand{\rhtaShXXgbFullAllThreeK}{}\renewcommand{\rhtaShXXgbFullAllThreeK}{0.88}  % ... at quoted bid/ask prices
\providecommand{\rhtaPShXgbAllThreeK}{}\renewcommand{\rhtaPShXgbAllThreeK}{\ensuremath{p = 0.41}}  % Ledoit-Wolf p, Sharpe ratio full minus HAR (mid): XGBoost, ALL48_N3k
\providecommand{\rhtaShRfHarAllThreeK}{}\renewcommand{\rhtaShRfHarAllThreeK}{0.87}  % Sharpe ratio (D19), mid-quotes: Random forest, baseline, ALL48_N3k (ALL48, the last 3,000 rows)
\providecommand{\rhtaShXRfHarAllThreeK}{}\renewcommand{\rhtaShXRfHarAllThreeK}{0.52}  % ... at quoted bid/ask prices
\providecommand{\rhtaShRfFullAllThreeK}{}\renewcommand{\rhtaShRfFullAllThreeK}{1.30}  % Sharpe ratio (D19), mid-quotes: Random forest, all_features, ALL48_N3k (ALL48, the last 3,000 rows)
\providecommand{\rhtaShXRfFullAllThreeK}{}\renewcommand{\rhtaShXRfFullAllThreeK}{0.94}  % ... at quoted bid/ask prices
\providecommand{\rhtaPShRfAllThreeK}{}\renewcommand{\rhtaPShRfAllThreeK}{\ensuremath{p = 0.21}}  % Ledoit-Wolf p, Sharpe ratio full minus HAR (mid): Random forest, ALL48_N3k
\providecommand{\rhtaShRidgeHarRthThreeK}{}\renewcommand{\rhtaShRidgeHarRthThreeK}{1.16}  % Sharpe ratio (D19), mid-quotes: Ridge, baseline, RTH13_N3k (RTH13, the last 3,000 rows)
\providecommand{\rhtaShXRidgeHarRthThreeK}{}\renewcommand{\rhtaShXRidgeHarRthThreeK}{0.80}  % ... at quoted bid/ask prices
\providecommand{\rhtaShRidgeFullRthThreeK}{}\renewcommand{\rhtaShRidgeFullRthThreeK}{1.12}  % Sharpe ratio (D19), mid-quotes: Ridge, all_features, RTH13_N3k (RTH13, the last 3,000 rows)
\providecommand{\rhtaShXRidgeFullRthThreeK}{}\renewcommand{\rhtaShXRidgeFullRthThreeK}{0.77}  % ... at quoted bid/ask prices
\providecommand{\rhtaShLassoHarRthThreeK}{}\renewcommand{\rhtaShLassoHarRthThreeK}{1.16}  % Sharpe ratio (D19), mid-quotes: Lasso, baseline, RTH13_N3k (RTH13, the last 3,000 rows)
\providecommand{\rhtaShXLassoHarRthThreeK}{}\renewcommand{\rhtaShXLassoHarRthThreeK}{0.80}  % ... at quoted bid/ask prices
\providecommand{\rhtaShLassoFullRthThreeK}{}\renewcommand{\rhtaShLassoFullRthThreeK}{1.55}  % Sharpe ratio (D19), mid-quotes: Lasso, all_features, RTH13_N3k (RTH13, the last 3,000 rows)
\providecommand{\rhtaShXLassoFullRthThreeK}{}\renewcommand{\rhtaShXLassoFullRthThreeK}{1.19}  % ... at quoted bid/ask prices
\providecommand{\rhtaShEnetHarRthThreeK}{}\renewcommand{\rhtaShEnetHarRthThreeK}{1.16}  % Sharpe ratio (D19), mid-quotes: Elastic net, baseline, RTH13_N3k (RTH13, the last 3,000 rows)
\providecommand{\rhtaShXEnetHarRthThreeK}{}\renewcommand{\rhtaShXEnetHarRthThreeK}{0.80}  % ... at quoted bid/ask prices
\providecommand{\rhtaShEnetFullRthThreeK}{}\renewcommand{\rhtaShEnetFullRthThreeK}{1.25}  % Sharpe ratio (D19), mid-quotes: Elastic net, all_features, RTH13_N3k (RTH13, the last 3,000 rows)
\providecommand{\rhtaShXEnetFullRthThreeK}{}\renewcommand{\rhtaShXEnetFullRthThreeK}{0.89}  % ... at quoted bid/ask prices
\providecommand{\rhtaShLgbmHarRthThreeK}{}\renewcommand{\rhtaShLgbmHarRthThreeK}{0.25}  % Sharpe ratio (D19), mid-quotes: LightGBM, baseline, RTH13_N3k (RTH13, the last 3,000 rows)
\providecommand{\rhtaShXLgbmHarRthThreeK}{}\renewcommand{\rhtaShXLgbmHarRthThreeK}{\ensuremath{-}0.10}  % ... at quoted bid/ask prices
\providecommand{\rhtaShLgbmFullRthThreeK}{}\renewcommand{\rhtaShLgbmFullRthThreeK}{0.83}  % Sharpe ratio (D19), mid-quotes: LightGBM, all_features, RTH13_N3k (RTH13, the last 3,000 rows)
\providecommand{\rhtaShXLgbmFullRthThreeK}{}\renewcommand{\rhtaShXLgbmFullRthThreeK}{0.47}  % ... at quoted bid/ask prices
\providecommand{\rhtaPShLgbmRthThreeK}{}\renewcommand{\rhtaPShLgbmRthThreeK}{\ensuremath{p = 0.22}}  % Ledoit-Wolf p, Sharpe ratio full minus HAR (mid): LightGBM, RTH13_N3k
\providecommand{\rhtaShXgbHarRthThreeK}{}\renewcommand{\rhtaShXgbHarRthThreeK}{0.87}  % Sharpe ratio (D19), mid-quotes: XGBoost, baseline, RTH13_N3k (RTH13, the last 3,000 rows)
\providecommand{\rhtaShXXgbHarRthThreeK}{}\renewcommand{\rhtaShXXgbHarRthThreeK}{0.52}  % ... at quoted bid/ask prices
\providecommand{\rhtaShXgbFullRthThreeK}{}\renewcommand{\rhtaShXgbFullRthThreeK}{0.86}  % Sharpe ratio (D19), mid-quotes: XGBoost, all_features, RTH13_N3k (RTH13, the last 3,000 rows)
\providecommand{\rhtaShXXgbFullRthThreeK}{}\renewcommand{\rhtaShXXgbFullRthThreeK}{0.51}  % ... at quoted bid/ask prices
\providecommand{\rhtaPShXgbRthThreeK}{}\renewcommand{\rhtaPShXgbRthThreeK}{\ensuremath{p = 0.98}}  % Ledoit-Wolf p, Sharpe ratio full minus HAR (mid): XGBoost, RTH13_N3k
\providecommand{\rhtaShRfHarRthThreeK}{}\renewcommand{\rhtaShRfHarRthThreeK}{0.67}  % Sharpe ratio (D19), mid-quotes: Random forest, baseline, RTH13_N3k (RTH13, the last 3,000 rows)
\providecommand{\rhtaShXRfHarRthThreeK}{}\renewcommand{\rhtaShXRfHarRthThreeK}{0.31}  % ... at quoted bid/ask prices
\providecommand{\rhtaShRfFullRthThreeK}{}\renewcommand{\rhtaShRfFullRthThreeK}{0.91}  % Sharpe ratio (D19), mid-quotes: Random forest, all_features, RTH13_N3k (RTH13, the last 3,000 rows)
\providecommand{\rhtaShXRfFullRthThreeK}{}\renewcommand{\rhtaShXRfFullRthThreeK}{0.56}  % ... at quoted bid/ask prices
\providecommand{\rhtaPShRfRthThreeK}{}\renewcommand{\rhtaPShRfRthThreeK}{\ensuremath{p = 0.53}}  % Ledoit-Wolf p, Sharpe ratio full minus HAR (mid): Random forest, RTH13_N3k
\providecommand{\rhtaShRidgeHarAllThirtyNineK}{}\renewcommand{\rhtaShRidgeHarAllThirtyNineK}{1.23}  % Sharpe ratio (D19), mid-quotes: Ridge, baseline, ALL48_N39k (ALL48, the last 39,000 rows)
\providecommand{\rhtaShXRidgeHarAllThirtyNineK}{}\renewcommand{\rhtaShXRidgeHarAllThirtyNineK}{0.87}  % ... at quoted bid/ask prices
\providecommand{\rhtaShRidgeFullAllThirtyNineK}{}\renewcommand{\rhtaShRidgeFullAllThirtyNineK}{0.76}  % Sharpe ratio (D19), mid-quotes: Ridge, all_features, ALL48_N39k (ALL48, the last 39,000 rows)
\providecommand{\rhtaShXRidgeFullAllThirtyNineK}{}\renewcommand{\rhtaShXRidgeFullAllThirtyNineK}{0.40}  % ... at quoted bid/ask prices
\providecommand{\rhtaShLassoHarAllThirtyNineK}{}\renewcommand{\rhtaShLassoHarAllThirtyNineK}{1.20}  % Sharpe ratio (D19), mid-quotes: Lasso, baseline, ALL48_N39k (ALL48, the last 39,000 rows)
\providecommand{\rhtaShXLassoHarAllThirtyNineK}{}\renewcommand{\rhtaShXLassoHarAllThirtyNineK}{0.84}  % ... at quoted bid/ask prices
\providecommand{\rhtaShLassoFullAllThirtyNineK}{}\renewcommand{\rhtaShLassoFullAllThirtyNineK}{1.33}  % Sharpe ratio (D19), mid-quotes: Lasso, all_features, ALL48_N39k (ALL48, the last 39,000 rows)
\providecommand{\rhtaShXLassoFullAllThirtyNineK}{}\renewcommand{\rhtaShXLassoFullAllThirtyNineK}{0.98}  % ... at quoted bid/ask prices
\providecommand{\rhtaShEnetHarAllThirtyNineK}{}\renewcommand{\rhtaShEnetHarAllThirtyNineK}{1.22}  % Sharpe ratio (D19), mid-quotes: Elastic net, baseline, ALL48_N39k (ALL48, the last 39,000 rows)
\providecommand{\rhtaShXEnetHarAllThirtyNineK}{}\renewcommand{\rhtaShXEnetHarAllThirtyNineK}{0.87}  % ... at quoted bid/ask prices
\providecommand{\rhtaShEnetFullAllThirtyNineK}{}\renewcommand{\rhtaShEnetFullAllThirtyNineK}{1.63}  % Sharpe ratio (D19), mid-quotes: Elastic net, all_features, ALL48_N39k (ALL48, the last 39,000 rows)
\providecommand{\rhtaShXEnetFullAllThirtyNineK}{}\renewcommand{\rhtaShXEnetFullAllThirtyNineK}{1.27}  % ... at quoted bid/ask prices
\providecommand{\rhtaShLgbmHarAllThirtyNineK}{}\renewcommand{\rhtaShLgbmHarAllThirtyNineK}{0.48}  % Sharpe ratio (D19), mid-quotes: LightGBM, baseline, ALL48_N39k (ALL48, the last 39,000 rows)
\providecommand{\rhtaShXLgbmHarAllThirtyNineK}{}\renewcommand{\rhtaShXLgbmHarAllThirtyNineK}{0.13}  % ... at quoted bid/ask prices
\providecommand{\rhtaShLgbmFullAllThirtyNineK}{}\renewcommand{\rhtaShLgbmFullAllThirtyNineK}{1.73}  % Sharpe ratio (D19), mid-quotes: LightGBM, all_features, ALL48_N39k (ALL48, the last 39,000 rows)
\providecommand{\rhtaShXLgbmFullAllThirtyNineK}{}\renewcommand{\rhtaShXLgbmFullAllThirtyNineK}{1.37}  % ... at quoted bid/ask prices
\providecommand{\rhtaPShLgbmAllThirtyNineK}{}\renewcommand{\rhtaPShLgbmAllThirtyNineK}{\ensuremath{p = 0.060}}  % Ledoit-Wolf p, Sharpe ratio full minus HAR (mid): LightGBM, ALL48_N39k
\providecommand{\rhtaShXgbHarAllThirtyNineK}{}\renewcommand{\rhtaShXgbHarAllThirtyNineK}{0.62}  % Sharpe ratio (D19), mid-quotes: XGBoost, baseline, ALL48_N39k (ALL48, the last 39,000 rows)
\providecommand{\rhtaShXXgbHarAllThirtyNineK}{}\renewcommand{\rhtaShXXgbHarAllThirtyNineK}{0.27}  % ... at quoted bid/ask prices
\providecommand{\rhtaShXgbFullAllThirtyNineK}{}\renewcommand{\rhtaShXgbFullAllThirtyNineK}{1.64}  % Sharpe ratio (D19), mid-quotes: XGBoost, all_features, ALL48_N39k (ALL48, the last 39,000 rows)
\providecommand{\rhtaShXXgbFullAllThirtyNineK}{}\renewcommand{\rhtaShXXgbFullAllThirtyNineK}{1.28}  % ... at quoted bid/ask prices
\providecommand{\rhtaPShXgbAllThirtyNineK}{}\renewcommand{\rhtaPShXgbAllThirtyNineK}{\ensuremath{p = 0.12}}  % Ledoit-Wolf p, Sharpe ratio full minus HAR (mid): XGBoost, ALL48_N39k
\providecommand{\rhtaShRfHarAllThirtyNineK}{}\renewcommand{\rhtaShRfHarAllThirtyNineK}{0.87}  % Sharpe ratio (D19), mid-quotes: Random forest, baseline, ALL48_N39k (ALL48, the last 39,000 rows)
\providecommand{\rhtaShXRfHarAllThirtyNineK}{}\renewcommand{\rhtaShXRfHarAllThirtyNineK}{0.50}  % ... at quoted bid/ask prices
\providecommand{\rhtaShRfFullAllThirtyNineK}{}\renewcommand{\rhtaShRfFullAllThirtyNineK}{0.95}  % Sharpe ratio (D19), mid-quotes: Random forest, all_features, ALL48_N39k (ALL48, the last 39,000 rows)
\providecommand{\rhtaShXRfFullAllThirtyNineK}{}\renewcommand{\rhtaShXRfFullAllThirtyNineK}{0.58}  % ... at quoted bid/ask prices
\providecommand{\rhtaPShRfAllThirtyNineK}{}\renewcommand{\rhtaPShRfAllThirtyNineK}{\ensuremath{p = 0.83}}  % Ledoit-Wolf p, Sharpe ratio full minus HAR (mid): Random forest, ALL48_N39k
\providecommand{\rhtaDShRidgeAllVsRth}{}\renewcommand{\rhtaDShRidgeAllVsRth}{\ensuremath{-}0.16}  % Sharpe ratio (D19, mid), ALL48 minus RTH13, full set, Ridge
\providecommand{\rhtaPDShRidgeAllVsRth}{}\renewcommand{\rhtaPDShRidgeAllVsRth}{\ensuremath{p = 0.63}}  % ... its Ledoit-Wolf p
\providecommand{\rhtaDShLassoAllVsRth}{}\renewcommand{\rhtaDShLassoAllVsRth}{\ensuremath{-}0.23}  % Sharpe ratio (D19, mid), ALL48 minus RTH13, full set, Lasso
\providecommand{\rhtaPDShLassoAllVsRth}{}\renewcommand{\rhtaPDShLassoAllVsRth}{\ensuremath{p = 0.51}}  % ... its Ledoit-Wolf p
\providecommand{\rhtaDShEnetAllVsRth}{}\renewcommand{\rhtaDShEnetAllVsRth}{\ensuremath{-}0.33}  % Sharpe ratio (D19, mid), ALL48 minus RTH13, full set, Elastic net
\providecommand{\rhtaPDShEnetAllVsRth}{}\renewcommand{\rhtaPDShEnetAllVsRth}{\ensuremath{p = 0.32}}  % ... its Ledoit-Wolf p
\providecommand{\rhtaDShLgbmAllVsRth}{}\renewcommand{\rhtaDShLgbmAllVsRth}{\ensuremath{-}0.34}  % Sharpe ratio (D19, mid), ALL48 minus RTH13, full set, LightGBM
\providecommand{\rhtaPDShLgbmAllVsRth}{}\renewcommand{\rhtaPDShLgbmAllVsRth}{\ensuremath{p = 0.17}}  % ... its Ledoit-Wolf p
\providecommand{\rhtaDShXgbAllVsRth}{}\renewcommand{\rhtaDShXgbAllVsRth}{0.03}  % Sharpe ratio (D19, mid), ALL48 minus RTH13, full set, XGBoost
\providecommand{\rhtaPDShXgbAllVsRth}{}\renewcommand{\rhtaPDShXgbAllVsRth}{\ensuremath{p = 0.90}}  % ... its Ledoit-Wolf p
\providecommand{\rhtaDShRfAllVsRth}{}\renewcommand{\rhtaDShRfAllVsRth}{\ensuremath{-}0.06}  % Sharpe ratio (D19, mid), ALL48 minus RTH13, full set, Random forest
\providecommand{\rhtaPDShRfAllVsRth}{}\renewcommand{\rhtaPDShRfAllVsRth}{\ensuremath{p = 0.85}}  % ... its Ledoit-Wolf p
\providecommand{\rhtaDShRidgeClVsRth}{}\renewcommand{\rhtaDShRidgeClVsRth}{\ensuremath{-}0.51}  % Sharpe ratio (D19, mid), CLOSE minus RTH13, full set, Ridge
\providecommand{\rhtaPDShRidgeClVsRth}{}\renewcommand{\rhtaPDShRidgeClVsRth}{\ensuremath{p = 0.22}}  % ... its Ledoit-Wolf p
\providecommand{\rhtaDShLassoClVsRth}{}\renewcommand{\rhtaDShLassoClVsRth}{\ensuremath{-}0.15}  % Sharpe ratio (D19, mid), CLOSE minus RTH13, full set, Lasso
\providecommand{\rhtaPDShLassoClVsRth}{}\renewcommand{\rhtaPDShLassoClVsRth}{\ensuremath{p = 0.72}}  % ... its Ledoit-Wolf p
\providecommand{\rhtaDShEnetClVsRth}{}\renewcommand{\rhtaDShEnetClVsRth}{\ensuremath{-}0.15}  % Sharpe ratio (D19, mid), CLOSE minus RTH13, full set, Elastic net
\providecommand{\rhtaPDShEnetClVsRth}{}\renewcommand{\rhtaPDShEnetClVsRth}{\ensuremath{p = 0.70}}  % ... its Ledoit-Wolf p
\providecommand{\rhtaDShLgbmClVsRth}{}\renewcommand{\rhtaDShLgbmClVsRth}{0.36}  % Sharpe ratio (D19, mid), CLOSE minus RTH13, full set, LightGBM
\providecommand{\rhtaPDShLgbmClVsRth}{}\renewcommand{\rhtaPDShLgbmClVsRth}{\ensuremath{p = 0.41}}  % ... its Ledoit-Wolf p
\providecommand{\rhtaDShXgbClVsRth}{}\renewcommand{\rhtaDShXgbClVsRth}{0.15}  % Sharpe ratio (D19, mid), CLOSE minus RTH13, full set, XGBoost
\providecommand{\rhtaPDShXgbClVsRth}{}\renewcommand{\rhtaPDShXgbClVsRth}{\ensuremath{p = 0.60}}  % ... its Ledoit-Wolf p
\providecommand{\rhtaDShRfClVsRth}{}\renewcommand{\rhtaDShRfClVsRth}{\ensuremath{-}0.59}  % Sharpe ratio (D19, mid), CLOSE minus RTH13, full set, Random forest
\providecommand{\rhtaPDShRfClVsRth}{}\renewcommand{\rhtaPDShRfClVsRth}{\ensuremath{p = 0.22}}  % ... its Ledoit-Wolf p
\providecommand{\rhtaWsPctRidgeWOne}{}\renewcommand{\rhtaWsPctRidgeWOne}{\ensuremath{-}4.8}  % CLOSE at W = 1,000 closes, 16:00, QLIKE % change full vs HAR, Ridge
\providecommand{\rhtaWsPRidgeWOne}{}\renewcommand{\rhtaWsPRidgeWOne}{\ensuremath{p < 0.001}}  % ... its p statement
\providecommand{\rhtaWsPctLassoWOne}{}\renewcommand{\rhtaWsPctLassoWOne}{\ensuremath{-}2.6}  % CLOSE at W = 1,000 closes, 16:00, QLIKE % change full vs HAR, Lasso
\providecommand{\rhtaWsPLassoWOne}{}\renewcommand{\rhtaWsPLassoWOne}{\ensuremath{p = 0.12}}  % ... its p statement
\providecommand{\rhtaWsPctEnetWOne}{}\renewcommand{\rhtaWsPctEnetWOne}{\ensuremath{-}1.5}  % CLOSE at W = 1,000 closes, 16:00, QLIKE % change full vs HAR, Elastic net
\providecommand{\rhtaWsPEnetWOne}{}\renewcommand{\rhtaWsPEnetWOne}{\ensuremath{p = 0.51}}  % ... its p statement
\providecommand{\rhtaWsPctLgbmWOne}{}\renewcommand{\rhtaWsPctLgbmWOne}{\ensuremath{-}7.8}  % CLOSE at W = 1,000 closes, 16:00, QLIKE % change full vs HAR, LightGBM
\providecommand{\rhtaWsPLgbmWOne}{}\renewcommand{\rhtaWsPLgbmWOne}{\ensuremath{p < 0.001}}  % ... its p statement
\providecommand{\rhtaWsPctXgbWOne}{}\renewcommand{\rhtaWsPctXgbWOne}{\ensuremath{-}6.8}  % CLOSE at W = 1,000 closes, 16:00, QLIKE % change full vs HAR, XGBoost
\providecommand{\rhtaWsPXgbWOne}{}\renewcommand{\rhtaWsPXgbWOne}{\ensuremath{p < 0.001}}  % ... its p statement
\providecommand{\rhtaWsPctRfWOne}{}\renewcommand{\rhtaWsPctRfWOne}{\ensuremath{-}2.1}  % CLOSE at W = 1,000 closes, 16:00, QLIKE % change full vs HAR, Random forest
\providecommand{\rhtaWsPRfWOne}{}\renewcommand{\rhtaWsPRfWOne}{\ensuremath{p = 0.069}}  % ... its p statement
\providecommand{\rhtaWsPctRidgeWOneHalf}{}\renewcommand{\rhtaWsPctRidgeWOneHalf}{\ensuremath{-}4.1}  % CLOSE at W = 1,500 closes, 16:00, QLIKE % change full vs HAR, Ridge
\providecommand{\rhtaWsPRidgeWOneHalf}{}\renewcommand{\rhtaWsPRidgeWOneHalf}{\ensuremath{p = 0.036}}  % ... its p statement
\providecommand{\rhtaWsPctLassoWOneHalf}{}\renewcommand{\rhtaWsPctLassoWOneHalf}{\ensuremath{-}5.0}  % CLOSE at W = 1,500 closes, 16:00, QLIKE % change full vs HAR, Lasso
\providecommand{\rhtaWsPLassoWOneHalf}{}\renewcommand{\rhtaWsPLassoWOneHalf}{\ensuremath{p < 0.001}}  % ... its p statement
\providecommand{\rhtaWsPctEnetWOneHalf}{}\renewcommand{\rhtaWsPctEnetWOneHalf}{\ensuremath{-}4.3}  % CLOSE at W = 1,500 closes, 16:00, QLIKE % change full vs HAR, Elastic net
\providecommand{\rhtaWsPEnetWOneHalf}{}\renewcommand{\rhtaWsPEnetWOneHalf}{\ensuremath{p = 0.035}}  % ... its p statement
\providecommand{\rhtaWsPctLgbmWOneHalf}{}\renewcommand{\rhtaWsPctLgbmWOneHalf}{\ensuremath{-}5.8}  % CLOSE at W = 1,500 closes, 16:00, QLIKE % change full vs HAR, LightGBM
\providecommand{\rhtaWsPLgbmWOneHalf}{}\renewcommand{\rhtaWsPLgbmWOneHalf}{\ensuremath{p = 0.001}}  % ... its p statement
\providecommand{\rhtaWsPctXgbWOneHalf}{}\renewcommand{\rhtaWsPctXgbWOneHalf}{\ensuremath{-}7.2}  % CLOSE at W = 1,500 closes, 16:00, QLIKE % change full vs HAR, XGBoost
\providecommand{\rhtaWsPXgbWOneHalf}{}\renewcommand{\rhtaWsPXgbWOneHalf}{\ensuremath{p < 0.001}}  % ... its p statement
\providecommand{\rhtaWsPctRfWOneHalf}{}\renewcommand{\rhtaWsPctRfWOneHalf}{\ensuremath{-}2.2}  % CLOSE at W = 1,500 closes, 16:00, QLIKE % change full vs HAR, Random forest
\providecommand{\rhtaWsPRfWOneHalf}{}\renewcommand{\rhtaWsPRfWOneHalf}{\ensuremath{p = 0.052}}  % ... its p statement
\providecommand{\rhtaWsPctRidgeWTwo}{}\renewcommand{\rhtaWsPctRidgeWTwo}{\ensuremath{-}3.9}  % CLOSE at W = 2,000 closes, 16:00, QLIKE % change full vs HAR, Ridge
\providecommand{\rhtaWsPRidgeWTwo}{}\renewcommand{\rhtaWsPRidgeWTwo}{\ensuremath{p = 0.049}}  % ... its p statement
\providecommand{\rhtaWsPctLassoWTwo}{}\renewcommand{\rhtaWsPctLassoWTwo}{\ensuremath{-}4.6}  % CLOSE at W = 2,000 closes, 16:00, QLIKE % change full vs HAR, Lasso
\providecommand{\rhtaWsPLassoWTwo}{}\renewcommand{\rhtaWsPLassoWTwo}{\ensuremath{p < 0.001}}  % ... its p statement
\providecommand{\rhtaWsPctEnetWTwo}{}\renewcommand{\rhtaWsPctEnetWTwo}{\ensuremath{-}4.6}  % CLOSE at W = 2,000 closes, 16:00, QLIKE % change full vs HAR, Elastic net
\providecommand{\rhtaWsPEnetWTwo}{}\renewcommand{\rhtaWsPEnetWTwo}{\ensuremath{p = 0.007}}  % ... its p statement
\providecommand{\rhtaWsPctLgbmWTwo}{}\renewcommand{\rhtaWsPctLgbmWTwo}{\ensuremath{-}7.4}  % CLOSE at W = 2,000 closes, 16:00, QLIKE % change full vs HAR, LightGBM
\providecommand{\rhtaWsPLgbmWTwo}{}\renewcommand{\rhtaWsPLgbmWTwo}{\ensuremath{p < 0.001}}  % ... its p statement
\providecommand{\rhtaWsPctXgbWTwo}{}\renewcommand{\rhtaWsPctXgbWTwo}{\ensuremath{-}7.8}  % CLOSE at W = 2,000 closes, 16:00, QLIKE % change full vs HAR, XGBoost
\providecommand{\rhtaWsPXgbWTwo}{}\renewcommand{\rhtaWsPXgbWTwo}{\ensuremath{p < 0.001}}  % ... its p statement
\providecommand{\rhtaWsPctRfWTwo}{}\renewcommand{\rhtaWsPctRfWTwo}{\ensuremath{-}3.2}  % CLOSE at W = 2,000 closes, 16:00, QLIKE % change full vs HAR, Random forest
\providecommand{\rhtaWsPRfWTwo}{}\renewcommand{\rhtaWsPRfWTwo}{\ensuremath{p = 0.002}}  % ... its p statement
\providecommand{\rhtaWsPctRidgeWTwoHalf}{}\renewcommand{\rhtaWsPctRidgeWTwoHalf}{\ensuremath{-}4.8}  % CLOSE at W = 2,500 closes, 16:00, QLIKE % change full vs HAR, Ridge
\providecommand{\rhtaWsPRidgeWTwoHalf}{}\renewcommand{\rhtaWsPRidgeWTwoHalf}{\ensuremath{p = 0.008}}  % ... its p statement
\providecommand{\rhtaWsPctLassoWTwoHalf}{}\renewcommand{\rhtaWsPctLassoWTwoHalf}{\ensuremath{-}5.3}  % CLOSE at W = 2,500 closes, 16:00, QLIKE % change full vs HAR, Lasso
\providecommand{\rhtaWsPLassoWTwoHalf}{}\renewcommand{\rhtaWsPLassoWTwoHalf}{\ensuremath{p < 0.001}}  % ... its p statement
\providecommand{\rhtaWsPctEnetWTwoHalf}{}\renewcommand{\rhtaWsPctEnetWTwoHalf}{\ensuremath{-}5.8}  % CLOSE at W = 2,500 closes, 16:00, QLIKE % change full vs HAR, Elastic net
\providecommand{\rhtaWsPEnetWTwoHalf}{}\renewcommand{\rhtaWsPEnetWTwoHalf}{\ensuremath{p < 0.001}}  % ... its p statement
\providecommand{\rhtaWsPctLgbmWTwoHalf}{}\renewcommand{\rhtaWsPctLgbmWTwoHalf}{\ensuremath{-}7.0}  % CLOSE at W = 2,500 closes, 16:00, QLIKE % change full vs HAR, LightGBM
\providecommand{\rhtaWsPLgbmWTwoHalf}{}\renewcommand{\rhtaWsPLgbmWTwoHalf}{\ensuremath{p < 0.001}}  % ... its p statement
\providecommand{\rhtaWsPctXgbWTwoHalf}{}\renewcommand{\rhtaWsPctXgbWTwoHalf}{\ensuremath{-}7.6}  % CLOSE at W = 2,500 closes, 16:00, QLIKE % change full vs HAR, XGBoost
\providecommand{\rhtaWsPXgbWTwoHalf}{}\renewcommand{\rhtaWsPXgbWTwoHalf}{\ensuremath{p < 0.001}}  % ... its p statement
\providecommand{\rhtaWsPctRfWTwoHalf}{}\renewcommand{\rhtaWsPctRfWTwoHalf}{\ensuremath{-}3.7}  % CLOSE at W = 2,500 closes, 16:00, QLIKE % change full vs HAR, Random forest
\providecommand{\rhtaWsPRfWTwoHalf}{}\renewcommand{\rhtaWsPRfWTwoHalf}{\ensuremath{p < 0.001}}  % ... its p statement
\providecommand{\rhtaWsPctRidgeWThree}{}\renewcommand{\rhtaWsPctRidgeWThree}{\ensuremath{-}6.1}  % CLOSE at W = 3,000 closes, 16:00, QLIKE % change full vs HAR, Ridge
\providecommand{\rhtaWsPRidgeWThree}{}\renewcommand{\rhtaWsPRidgeWThree}{\ensuremath{p < 0.001}}  % ... its p statement
\providecommand{\rhtaWsPctLassoWThree}{}\renewcommand{\rhtaWsPctLassoWThree}{\ensuremath{-}5.4}  % CLOSE at W = 3,000 closes, 16:00, QLIKE % change full vs HAR, Lasso
\providecommand{\rhtaWsPLassoWThree}{}\renewcommand{\rhtaWsPLassoWThree}{\ensuremath{p < 0.001}}  % ... its p statement
\providecommand{\rhtaWsPctEnetWThree}{}\renewcommand{\rhtaWsPctEnetWThree}{\ensuremath{-}6.5}  % CLOSE at W = 3,000 closes, 16:00, QLIKE % change full vs HAR, Elastic net
\providecommand{\rhtaWsPEnetWThree}{}\renewcommand{\rhtaWsPEnetWThree}{\ensuremath{p < 0.001}}  % ... its p statement
\providecommand{\rhtaWsPctLgbmWThree}{}\renewcommand{\rhtaWsPctLgbmWThree}{\ensuremath{-}7.0}  % CLOSE at W = 3,000 closes, 16:00, QLIKE % change full vs HAR, LightGBM
\providecommand{\rhtaWsPLgbmWThree}{}\renewcommand{\rhtaWsPLgbmWThree}{\ensuremath{p < 0.001}}  % ... its p statement
\providecommand{\rhtaWsPctXgbWThree}{}\renewcommand{\rhtaWsPctXgbWThree}{\ensuremath{-}8.5}  % CLOSE at W = 3,000 closes, 16:00, QLIKE % change full vs HAR, XGBoost
\providecommand{\rhtaWsPXgbWThree}{}\renewcommand{\rhtaWsPXgbWThree}{\ensuremath{p < 0.001}}  % ... its p statement
\providecommand{\rhtaWsPctRfWThree}{}\renewcommand{\rhtaWsPctRfWThree}{\ensuremath{-}3.7}  % CLOSE at W = 3,000 closes, 16:00, QLIKE % change full vs HAR, Random forest
\providecommand{\rhtaWsPRfWThree}{}\renewcommand{\rhtaWsPRfWThree}{\ensuremath{p < 0.001}}  % ... its p statement
\providecommand{\rhtaWsPctRidgeWThreeHalf}{}\renewcommand{\rhtaWsPctRidgeWThreeHalf}{\ensuremath{-}6.5}  % CLOSE at W = 3,500 closes, 16:00, QLIKE % change full vs HAR, Ridge
\providecommand{\rhtaWsPRidgeWThreeHalf}{}\renewcommand{\rhtaWsPRidgeWThreeHalf}{\ensuremath{p < 0.001}}  % ... its p statement
\providecommand{\rhtaWsPctLassoWThreeHalf}{}\renewcommand{\rhtaWsPctLassoWThreeHalf}{\ensuremath{-}5.1}  % CLOSE at W = 3,500 closes, 16:00, QLIKE % change full vs HAR, Lasso
\providecommand{\rhtaWsPLassoWThreeHalf}{}\renewcommand{\rhtaWsPLassoWThreeHalf}{\ensuremath{p = 0.011}}  % ... its p statement
\providecommand{\rhtaWsPctEnetWThreeHalf}{}\renewcommand{\rhtaWsPctEnetWThreeHalf}{\ensuremath{-}6.7}  % CLOSE at W = 3,500 closes, 16:00, QLIKE % change full vs HAR, Elastic net
\providecommand{\rhtaWsPEnetWThreeHalf}{}\renewcommand{\rhtaWsPEnetWThreeHalf}{\ensuremath{p < 0.001}}  % ... its p statement
\providecommand{\rhtaWsPctLgbmWThreeHalf}{}\renewcommand{\rhtaWsPctLgbmWThreeHalf}{\ensuremath{-}7.2}  % CLOSE at W = 3,500 closes, 16:00, QLIKE % change full vs HAR, LightGBM
\providecommand{\rhtaWsPLgbmWThreeHalf}{}\renewcommand{\rhtaWsPLgbmWThreeHalf}{\ensuremath{p < 0.001}}  % ... its p statement
\providecommand{\rhtaWsPctXgbWThreeHalf}{}\renewcommand{\rhtaWsPctXgbWThreeHalf}{\ensuremath{-}7.3}  % CLOSE at W = 3,500 closes, 16:00, QLIKE % change full vs HAR, XGBoost
\providecommand{\rhtaWsPXgbWThreeHalf}{}\renewcommand{\rhtaWsPXgbWThreeHalf}{\ensuremath{p < 0.001}}  % ... its p statement
\providecommand{\rhtaWsPctRfWThreeHalf}{}\renewcommand{\rhtaWsPctRfWThreeHalf}{\ensuremath{-}3.4}  % CLOSE at W = 3,500 closes, 16:00, QLIKE % change full vs HAR, Random forest
\providecommand{\rhtaWsPRfWThreeHalf}{}\renewcommand{\rhtaWsPRfWThreeHalf}{\ensuremath{p = 0.001}}  % ... its p statement
\providecommand{\rhtaLinPctAllCloseMin}{}\renewcommand{\rhtaLinPctAllCloseMin}{\ensuremath{-}2.6}  % 16:00, full vs HAR, ALL48, penalized linear models: lowest of the range, %
\providecommand{\rhtaLinPctAllCloseMax}{}\renewcommand{\rhtaLinPctAllCloseMax}{\ensuremath{-}1.8}  % 16:00, full vs HAR, ALL48, penalized linear models: highest of the range, %
\providecommand{\rhtaTreePctAllCloseMin}{}\renewcommand{\rhtaTreePctAllCloseMin}{\ensuremath{-}6.5}  % 16:00, full vs HAR, ALL48, tree ensembles: lowest of the range, %
\providecommand{\rhtaTreePctAllCloseMax}{}\renewcommand{\rhtaTreePctAllCloseMax}{\ensuremath{-}5.6}  % 16:00, full vs HAR, ALL48, tree ensembles: highest of the range, %
\providecommand{\rhtaPctAllCloseMin}{}\renewcommand{\rhtaPctAllCloseMin}{\ensuremath{-}6.5}  % 16:00, full vs HAR, ALL48, all six classes: lowest of the range, %
\providecommand{\rhtaPctAllCloseMax}{}\renewcommand{\rhtaPctAllCloseMax}{\ensuremath{-}1.8}  % 16:00, full vs HAR, ALL48, all six classes: highest of the range, %
\providecommand{\rhtaLinPctClCloseMin}{}\renewcommand{\rhtaLinPctClCloseMin}{\ensuremath{-}6.5}  % 16:00, full vs HAR, CLOSE, penalized linear models: lowest of the range, %
\providecommand{\rhtaLinPctClCloseMax}{}\renewcommand{\rhtaLinPctClCloseMax}{\ensuremath{-}5.4}  % 16:00, full vs HAR, CLOSE, penalized linear models: highest of the range, %
\providecommand{\rhtaTreePctClCloseMin}{}\renewcommand{\rhtaTreePctClCloseMin}{\ensuremath{-}8.5}  % 16:00, full vs HAR, CLOSE, tree ensembles: lowest of the range, %
\providecommand{\rhtaTreePctClCloseMax}{}\renewcommand{\rhtaTreePctClCloseMax}{\ensuremath{-}3.7}  % 16:00, full vs HAR, CLOSE, tree ensembles: highest of the range, %
\providecommand{\rhtaPctClCloseMin}{}\renewcommand{\rhtaPctClCloseMin}{\ensuremath{-}8.5}  % 16:00, full vs HAR, CLOSE, all six classes: lowest of the range, %
\providecommand{\rhtaPctClCloseMax}{}\renewcommand{\rhtaPctClCloseMax}{\ensuremath{-}3.7}  % 16:00, full vs HAR, CLOSE, all six classes: highest of the range, %
\providecommand{\rhtaLinPctRthCloseMin}{}\renewcommand{\rhtaLinPctRthCloseMin}{\ensuremath{-}0.5}  % 16:00, full vs HAR, RTH13, penalized linear models: lowest of the range, %
\providecommand{\rhtaLinPctRthCloseMax}{}\renewcommand{\rhtaLinPctRthCloseMax}{3.3}  % 16:00, full vs HAR, RTH13, penalized linear models: highest of the range, %
\providecommand{\rhtaTreePctRthCloseMin}{}\renewcommand{\rhtaTreePctRthCloseMin}{\ensuremath{-}3.4}  % 16:00, full vs HAR, RTH13, tree ensembles: lowest of the range, %
\providecommand{\rhtaTreePctRthCloseMax}{}\renewcommand{\rhtaTreePctRthCloseMax}{\ensuremath{-}1.4}  % 16:00, full vs HAR, RTH13, tree ensembles: highest of the range, %
\providecommand{\rhtaPctRthCloseMin}{}\renewcommand{\rhtaPctRthCloseMin}{\ensuremath{-}3.4}  % 16:00, full vs HAR, RTH13, all six classes: lowest of the range, %
\providecommand{\rhtaPctRthCloseMax}{}\renewcommand{\rhtaPctRthCloseMax}{3.3}  % 16:00, full vs HAR, RTH13, all six classes: highest of the range, %
\providecommand{\rhtaLinPctClNatCloseMin}{}\renewcommand{\rhtaLinPctClNatCloseMin}{\ensuremath{-}6.6}  % 16:00, full vs HAR, CLOSE_NAT, penalized linear models: lowest of the range, %
\providecommand{\rhtaLinPctClNatCloseMax}{}\renewcommand{\rhtaLinPctClNatCloseMax}{\ensuremath{-}5.4}  % 16:00, full vs HAR, CLOSE_NAT, penalized linear models: highest of the range, %
\providecommand{\rhtaLinAllClosePMin}{}\renewcommand{\rhtaLinAllClosePMin}{\ensuremath{p = 0.15}}  % 16:00, ALL48, linear: smallest p
\providecommand{\rhtaLinAllCtClosePMin}{}\renewcommand{\rhtaLinAllCtClosePMin}{\ensuremath{p = 0.33}}  % 16:00, ALL48@16, linear: smallest p
\providecommand{\rhtaLinRthCtClosePMin}{}\renewcommand{\rhtaLinRthCtClosePMin}{\ensuremath{p = 0.38}}  % 16:00, RTH13@16, linear: smallest p
\providecommand{\rhtaLinPctRthBarMin}{}\renewcommand{\rhtaLinPctRthBarMin}{\ensuremath{-}6.2}  % every bar, full vs HAR, RTH13, penalized linear models: lowest of the range, %
\providecommand{\rhtaLinPctRthBarMax}{}\renewcommand{\rhtaLinPctRthBarMax}{\ensuremath{-}5.4}  % every bar, full vs HAR, RTH13, penalized linear models: highest of the range, %
\providecommand{\rhtaTreePctRthBarMin}{}\renewcommand{\rhtaTreePctRthBarMin}{\ensuremath{-}14.0}  % every bar, full vs HAR, RTH13, tree ensembles: lowest of the range, %
\providecommand{\rhtaTreePctRthBarMax}{}\renewcommand{\rhtaTreePctRthBarMax}{\ensuremath{-}11.1}  % every bar, full vs HAR, RTH13, tree ensembles: highest of the range, %
\providecommand{\rhtaLinPctAllBarMin}{}\renewcommand{\rhtaLinPctAllBarMin}{\ensuremath{-}5.2}  % every bar, full vs HAR, ALL48, penalized linear models: lowest of the range, %
\providecommand{\rhtaLinPctAllBarMax}{}\renewcommand{\rhtaLinPctAllBarMax}{\ensuremath{-}4.9}  % every bar, full vs HAR, ALL48, penalized linear models: highest of the range, %
\providecommand{\rhtaTreePctAllBarMin}{}\renewcommand{\rhtaTreePctAllBarMin}{\ensuremath{-}12.0}  % every bar, full vs HAR, ALL48, tree ensembles: lowest of the range, %
\providecommand{\rhtaTreePctAllBarMax}{}\renewcommand{\rhtaTreePctAllBarMax}{\ensuremath{-}7.0}  % every bar, full vs HAR, ALL48, tree ensembles: highest of the range, %
\providecommand{\rhtaAllVsRthClosePctMin}{}\renewcommand{\rhtaAllVsRthClosePctMin}{\ensuremath{-}6.0}  % 16:00, ALL48 vs RTH13, full set, all six classes: lowest of the range, %
\providecommand{\rhtaAllVsRthClosePctMax}{}\renewcommand{\rhtaAllVsRthClosePctMax}{\ensuremath{-}1.9}  % 16:00, ALL48 vs RTH13, full set, all six classes: highest of the range, %
\providecommand{\rhtaNAllVsRthCloseSig}{}\renewcommand{\rhtaNAllVsRthCloseSig}{five}  % 16:00, ALL48 vs RTH13: classes with p < 0.05 (lower)
\providecommand{\rhtaAllVsRthBarPctMin}{}\renewcommand{\rhtaAllVsRthBarPctMin}{1.9}  % every bar, ALL48 vs RTH13, full set, all six classes (RTH13 more accurate): lowest of the range, %
\providecommand{\rhtaAllVsRthBarPctMax}{}\renewcommand{\rhtaAllVsRthBarPctMax}{5.2}  % every bar, ALL48 vs RTH13, full set, all six classes (RTH13 more accurate): highest of the range, %
\providecommand{\rhtaAllVsRthBarPMax}{}\renewcommand{\rhtaAllVsRthBarPMax}{\ensuremath{p = 0.008}}  % every bar, ALL48 vs RTH13: largest p of the six
\providecommand{\rhtaLinClVsRthClosePctMin}{}\renewcommand{\rhtaLinClVsRthClosePctMin}{\ensuremath{-}11.8}  % 16:00, CLOSE vs RTH13, full set, penalized linear models: lowest of the range, %
\providecommand{\rhtaLinClVsRthClosePctMax}{}\renewcommand{\rhtaLinClVsRthClosePctMax}{\ensuremath{-}9.4}  % 16:00, CLOSE vs RTH13, full set, penalized linear models: highest of the range, %
\providecommand{\rhtaTreeClVsRthClosePctMin}{}\renewcommand{\rhtaTreeClVsRthClosePctMin}{\ensuremath{-}2.2}  % 16:00, CLOSE vs RTH13, full set, tree ensembles: lowest of the range, %
\providecommand{\rhtaTreeClVsRthClosePctMax}{}\renewcommand{\rhtaTreeClVsRthClosePctMax}{1.6}  % 16:00, CLOSE vs RTH13, full set, tree ensembles: highest of the range, %
\providecommand{\rhtaTreeClVsRthClosePMin}{}\renewcommand{\rhtaTreeClVsRthClosePMin}{\ensuremath{p = 0.42}}  % ... tree ensembles: smallest p
\providecommand{\rhtaTreeClVsRthClosePMax}{}\renewcommand{\rhtaTreeClVsRthClosePMax}{\ensuremath{p = 0.69}}  % ... tree ensembles: largest p
\providecommand{\rhtaAllThirtyNineKVsRthClosePctMin}{}\renewcommand{\rhtaAllThirtyNineKVsRthClosePctMin}{\ensuremath{-}3.6}  % 16:00, ALL48_N39k vs RTH13, full set, all six classes: lowest of the range, %
\providecommand{\rhtaAllThirtyNineKVsRthClosePctMax}{}\renewcommand{\rhtaAllThirtyNineKVsRthClosePctMax}{1.3}  % 16:00, ALL48_N39k vs RTH13, full set, all six classes: highest of the range, %
\providecommand{\rhtaNAllThirtyNineKVsRthCloseSig}{}\renewcommand{\rhtaNAllThirtyNineKVsRthCloseSig}{two}  % ... classes with p < 0.05
\providecommand{\rhtaAllThirtyNineKVsRthBarPctMin}{}\renewcommand{\rhtaAllThirtyNineKVsRthBarPctMin}{2.7}  % every bar, ALL48_N39k vs RTH13, full set, all six classes: lowest of the range, %
\providecommand{\rhtaAllThirtyNineKVsRthBarPctMax}{}\renewcommand{\rhtaAllThirtyNineKVsRthBarPctMax}{10.4}  % every bar, ALL48_N39k vs RTH13, full set, all six classes: highest of the range, %
\providecommand{\rhtaAllThirtyNineKVsAllClosePctMin}{}\renewcommand{\rhtaAllThirtyNineKVsAllClosePctMin}{1.7}  % 16:00, ALL48_N39k vs ALL48 (39,000 vs 144,000 rows), full set: lowest of the range, %
\providecommand{\rhtaAllThirtyNineKVsAllClosePctMax}{}\renewcommand{\rhtaAllThirtyNineKVsAllClosePctMax}{4.1}  % 16:00, ALL48_N39k vs ALL48 (39,000 vs 144,000 rows), full set: highest of the range, %
\providecommand{\rhtaNAllThirtyNineKVsAllCloseSig}{}\renewcommand{\rhtaNAllThirtyNineKVsAllCloseSig}{five}  % ... classes with p < 0.05 (higher loss)
\providecommand{\rhtaAllThreeKVsRthThreeKClosePctMin}{}\renewcommand{\rhtaAllThreeKVsRthThreeKClosePctMin}{\ensuremath{-}8.1}  % 16:00, ALL48_N3k vs RTH13_N3k, full set, all six classes: lowest of the range, %
\providecommand{\rhtaAllThreeKVsRthThreeKClosePctMax}{}\renewcommand{\rhtaAllThreeKVsRthThreeKClosePctMax}{\ensuremath{-}3.3}  % 16:00, ALL48_N3k vs RTH13_N3k, full set, all six classes: highest of the range, %
\providecommand{\rhtaNAllThreeKVsRthThreeKCloseSig}{}\renewcommand{\rhtaNAllThreeKVsRthThreeKCloseSig}{five}  % ... classes with p < 0.05 (lower)
\providecommand{\rhtaAllThreeKVsRthThreeKBarPctMin}{}\renewcommand{\rhtaAllThreeKVsRthThreeKBarPctMin}{2.7}  % every bar, ALL48_N3k vs RTH13_N3k, full set, all six classes: lowest of the range, %
\providecommand{\rhtaAllThreeKVsRthThreeKBarPctMax}{}\renewcommand{\rhtaAllThreeKVsRthThreeKBarPctMax}{5.8}  % every bar, ALL48_N3k vs RTH13_N3k, full set, all six classes: highest of the range, %
\providecommand{\rhtaClVsRthThreeKClosePctMin}{}\renewcommand{\rhtaClVsRthThreeKClosePctMin}{\ensuremath{-}18.0}  % 16:00, CLOSE vs RTH13_N3k (3,000 rows each), full set: lowest of the range, %
\providecommand{\rhtaClVsRthThreeKClosePctMax}{}\renewcommand{\rhtaClVsRthThreeKClosePctMax}{\ensuremath{-}8.7}  % 16:00, CLOSE vs RTH13_N3k (3,000 rows each), full set: highest of the range, %
\providecommand{\rhtaClVsRthThreeKClosePMax}{}\renewcommand{\rhtaClVsRthThreeKClosePMax}{\ensuremath{p = 0.004}}  % ... largest p of the six
\providecommand{\rhtaRthThreeKVsRthClosePctMin}{}\renewcommand{\rhtaRthThreeKVsRthClosePctMin}{6.2}  % 16:00, RTH13_N3k vs RTH13 (3,000 vs 39,000 rows), full set: lowest of the range, %
\providecommand{\rhtaRthThreeKVsRthClosePctMax}{}\renewcommand{\rhtaRthThreeKVsRthClosePctMax}{12.3}  % 16:00, RTH13_N3k vs RTH13 (3,000 vs 39,000 rows), full set: highest of the range, %
\providecommand{\rhtaRthThreeKVsRthBarPctMin}{}\renewcommand{\rhtaRthThreeKVsRthBarPctMin}{2.7}  % every bar, RTH13_N3k vs RTH13 (3,000 vs 39,000 rows), full set: lowest of the range, %
\providecommand{\rhtaRthThreeKVsRthBarPctMax}{}\renewcommand{\rhtaRthThreeKVsRthBarPctMax}{14.9}  % every bar, RTH13_N3k vs RTH13 (3,000 vs 39,000 rows), full set: highest of the range, %
\providecommand{\rhtaShFullRthMin}{}\renewcommand{\rhtaShFullRthMin}{1.47}  % Sharpe ratio (D19, mid), full set, RTH13: lowest of the six
\providecommand{\rhtaShFullRthMax}{}\renewcommand{\rhtaShFullRthMax}{1.88}  % ... highest of the six
\providecommand{\rhtaNTreeShSigRth}{}\renewcommand{\rhtaNTreeShSigRth}{two}  % tree classes whose full-vs-HAR Sharpe ratio difference has Ledoit-Wolf p < 0.05, RTH13
\providecommand{\rhtaShFullAllMin}{}\renewcommand{\rhtaShFullAllMin}{1.25}  % Sharpe ratio (D19, mid), full set, ALL48: lowest of the six
\providecommand{\rhtaShFullAllMax}{}\renewcommand{\rhtaShFullAllMax}{1.82}  % ... highest of the six
\providecommand{\rhtaNTreeShSigAll}{}\renewcommand{\rhtaNTreeShSigAll}{none}  % tree classes whose full-vs-HAR Sharpe ratio difference has Ledoit-Wolf p < 0.05, ALL48
\providecommand{\rhtaShFullClMin}{}\renewcommand{\rhtaShFullClMin}{1.26}  % Sharpe ratio (D19, mid), full set, CLOSE: lowest of the six
\providecommand{\rhtaShFullClMax}{}\renewcommand{\rhtaShFullClMax}{1.94}  % ... highest of the six
\providecommand{\rhtaNTreeShSigCl}{}\renewcommand{\rhtaNTreeShSigCl}{one}  % tree classes whose full-vs-HAR Sharpe ratio difference has Ledoit-Wolf p < 0.05, CLOSE
\providecommand{\rhtaShTreeHarClMin}{}\renewcommand{\rhtaShTreeHarClMin}{\ensuremath{-}0.03}  % Sharpe ratio (D19, mid), HAR benchmark, CLOSE, tree ensembles: lowest
\providecommand{\rhtaShTreeHarClMax}{}\renewcommand{\rhtaShTreeHarClMax}{0.40}  % ... highest
\providecommand{\rhtaNDShTests}{}\renewcommand{\rhtaNDShTests}{12}  % Sharpe ratio tests ALL48 - RTH13 and CLOSE - RTH13, full set (six classes each)
\providecommand{\rhtaNDShSig}{}\renewcommand{\rhtaNDShSig}{none}  % ... of them with Ledoit-Wolf p < 0.05
\providecommand{\rhtaDShPMin}{}\renewcommand{\rhtaDShPMin}{\ensuremath{p = 0.17}}  % ... the smallest Ledoit-Wolf p
\providecommand{\rhtaWsNCells}{}\renewcommand{\rhtaWsNCells}{36}  % window x class cells of the close-only sweep (16:00, full vs HAR)
\providecommand{\rhtaWsNLower}{}\renewcommand{\rhtaWsNLower}{36}  % ... with the lower QLIKE on the full set
\providecommand{\rhtaWsNSig}{}\renewcommand{\rhtaWsNSig}{32}  % ... with p < 0.05
\providecommand{\rhtaWsSigFrom}{}\renewcommand{\rhtaWsSigFrom}{2{,}000}  % smallest W from which every cell has p < 0.05
\providecommand{\rhtaWsLinPctMin}{}\renewcommand{\rhtaWsLinPctMin}{\ensuremath{-}6.7}  % close-only sweep, 16:00, full vs HAR, linear, all W: lowest of the range, %
\providecommand{\rhtaWsLinPctMax}{}\renewcommand{\rhtaWsLinPctMax}{\ensuremath{-}1.5}  % close-only sweep, 16:00, full vs HAR, linear, all W: highest of the range, %
\providecommand{\rhtaWsTreePctMin}{}\renewcommand{\rhtaWsTreePctMin}{\ensuremath{-}8.5}  % close-only sweep, 16:00, full vs HAR, trees, all W: lowest of the range, %
\providecommand{\rhtaWsTreePctMax}{}\renewcommand{\rhtaWsTreePctMax}{\ensuremath{-}2.1}  % close-only sweep, 16:00, full vs HAR, trees, all W: highest of the range, %
\fi
\ifdefined\rhrbWinMain\else% AUTO-GENERATED by writeup/make_rehaul_robust_window.py -- do not edit; re-run the script.
% \rhrb... macros of the appendix robustness table tab:rh_robust_w3500 (rehaul_robust_w3500_table.tex): the final
% pooled models on a rolling window of 3,500 against 3,000 trading days; every bar scored on the 2,579 days 2014-01-02 .. 2024-04-30
% (33,527 bars), Sharpe ratios (D19 pooled) on the 614 post-warmup option days. Names: Q (QLIKE), SR (Sharpe ratio), Pct / DM / P
% (3,500 vs 3,000: % change, DM, p statement), DSR / PSR (Sharpe change 3,500 - 3,000, Ledoit-Wolf p), PctF / DMF / PF / DSRF / PSRF
% (full set vs HAR benchmark at a window); Model Ols Ridge Lasso Enet Lgbm Xgb Rf; Set Har Full; Bar (every bar) / Close (16:00);
% Mid / Ba (bid/ask); Main (W = 3,000) / Long (W = 3,500). Pending: none.
\providecommand{\rhPending}{\textit{pending}}
\providecommand{\rhrbWinMain}{}\renewcommand{\rhrbWinMain}{3{,}000}  % main rolling window, trading days
\providecommand{\rhrbWinLong}{}\renewcommand{\rhrbWinLong}{3{,}500}  % robustness rolling window, trading days
\providecommand{\rhrbWinMainBars}{}\renewcommand{\rhrbWinMainBars}{39{,}000}  % main rolling window, 13-bar observations
\providecommand{\rhrbWinLongBars}{}\renewcommand{\rhrbWinLongBars}{45{,}500}  % robustness rolling window, 13-bar observations
\providecommand{\rhrbNdays}{}\renewcommand{\rhrbNdays}{2{,}579}  % evaluation days (every bar and 16:00)
\providecommand{\rhrbNbars}{}\renewcommand{\rhrbNbars}{33{,}527}  % scored bars, the same at both windows (asserted)
\providecommand{\rhrbSpanStart}{}\renewcommand{\rhrbSpanStart}{2 January 2014}  % first evaluation day
\providecommand{\rhrbSpanEnd}{}\renewcommand{\rhrbSpanEnd}{30 April 2024}  % last evaluation day
\providecommand{\rhrbDMLag}{}\renewcommand{\rhrbDMLag}{8}  % Newey-West lag of every DM test
\providecommand{\rhrbNopt}{}\renewcommand{\rhrbNopt}{866}  % option days of the sign strategy
\providecommand{\rhrbNpw}{}\renewcommand{\rhrbNpw}{614}  % post-warmup option days, the Sharpe-ratio sample
\providecommand{\rhrbBootN}{}\renewcommand{\rhrbBootN}{4{,}999}  % bootstrap resamples of the Sharpe-ratio test (Ledoit-Wolf since D19)
\providecommand{\rhrbBootBlock}{}\renewcommand{\rhrbBootBlock}{10}  % block length of the Sharpe-ratio bootstrap
\providecommand{\rhrbPendingNote}{}\renewcommand{\rhrbPendingNote}{}  % table-note sentence on the pending forest (empty when nothing is pending)
\providecommand{\rhrbQOlsHarBarMain}{}\renewcommand{\rhrbQOlsHarBarMain}{0.1893}  % QLIKE, every bar, OLS, HAR benchmark, W = 3,000
\providecommand{\rhrbQOlsHarCloseMain}{}\renewcommand{\rhrbQOlsHarCloseMain}{0.1269}  % QLIKE, 16:00 bar, OLS, HAR benchmark, W = 3,000
\providecommand{\rhrbQOlsHarBarLong}{}\renewcommand{\rhrbQOlsHarBarLong}{0.1897}  % QLIKE, every bar, OLS, HAR benchmark, W = 3,500
\providecommand{\rhrbQOlsHarCloseLong}{}\renewcommand{\rhrbQOlsHarCloseLong}{0.1272}  % QLIKE, 16:00 bar, OLS, HAR benchmark, W = 3,500
\providecommand{\rhrbPctOlsHarBar}{}\renewcommand{\rhrbPctOlsHarBar}{0.2}  % QLIKE change, %, W = 3,500 vs 3,000, every bar, OLS, HAR benchmark
\providecommand{\rhrbDMOlsHarBar}{}\renewcommand{\rhrbDMOlsHarBar}{2.48}  % DM, W = 3,500 vs 3,000 (positive: 3,500 has the higher loss), every bar, OLS, HAR benchmark
\providecommand{\rhrbPOlsHarBar}{}\renewcommand{\rhrbPOlsHarBar}{\ensuremath{p = 0.013}}  % ... its p statement
\providecommand{\rhrbPctOlsHarClose}{}\renewcommand{\rhrbPctOlsHarClose}{0.3}  % QLIKE change, %, W = 3,500 vs 3,000, 16:00 bar, OLS, HAR benchmark
\providecommand{\rhrbDMOlsHarClose}{}\renewcommand{\rhrbDMOlsHarClose}{1.31}  % DM, W = 3,500 vs 3,000 (positive: 3,500 has the higher loss), 16:00 bar, OLS, HAR benchmark
\providecommand{\rhrbPOlsHarClose}{}\renewcommand{\rhrbPOlsHarClose}{\ensuremath{p = 0.19}}  % ... its p statement
\providecommand{\rhrbQRidgeHarBarMain}{}\renewcommand{\rhrbQRidgeHarBarMain}{0.1893}  % QLIKE, every bar, Ridge, HAR benchmark, W = 3,000
\providecommand{\rhrbQRidgeHarCloseMain}{}\renewcommand{\rhrbQRidgeHarCloseMain}{0.1269}  % QLIKE, 16:00 bar, Ridge, HAR benchmark, W = 3,000
\providecommand{\rhrbSRRidgeHarMidMain}{}\renewcommand{\rhrbSRRidgeHarMidMain}{0.95}  % Sharpe ratio (D19 pooled), sign strategy at mid-quotes, Ridge, HAR benchmark, W = 3,000
\providecommand{\rhrbSRRidgeHarBaMain}{}\renewcommand{\rhrbSRRidgeHarBaMain}{0.60}  % Sharpe ratio (D19 pooled), sign strategy at bid/ask prices, Ridge, HAR benchmark, W = 3,000
\providecommand{\rhrbQRidgeHarBarLong}{}\renewcommand{\rhrbQRidgeHarBarLong}{0.1897}  % QLIKE, every bar, Ridge, HAR benchmark, W = 3,500
\providecommand{\rhrbQRidgeHarCloseLong}{}\renewcommand{\rhrbQRidgeHarCloseLong}{0.1272}  % QLIKE, 16:00 bar, Ridge, HAR benchmark, W = 3,500
\providecommand{\rhrbSRRidgeHarMidLong}{}\renewcommand{\rhrbSRRidgeHarMidLong}{1.01}  % Sharpe ratio (D19 pooled), sign strategy at mid-quotes, Ridge, HAR benchmark, W = 3,500
\providecommand{\rhrbSRRidgeHarBaLong}{}\renewcommand{\rhrbSRRidgeHarBaLong}{0.65}  % Sharpe ratio (D19 pooled), sign strategy at bid/ask prices, Ridge, HAR benchmark, W = 3,500
\providecommand{\rhrbPctRidgeHarBar}{}\renewcommand{\rhrbPctRidgeHarBar}{0.2}  % QLIKE change, %, W = 3,500 vs 3,000, every bar, Ridge, HAR benchmark
\providecommand{\rhrbDMRidgeHarBar}{}\renewcommand{\rhrbDMRidgeHarBar}{2.48}  % DM, W = 3,500 vs 3,000 (positive: 3,500 has the higher loss), every bar, Ridge, HAR benchmark
\providecommand{\rhrbPRidgeHarBar}{}\renewcommand{\rhrbPRidgeHarBar}{\ensuremath{p = 0.013}}  % ... its p statement
\providecommand{\rhrbPctRidgeHarClose}{}\renewcommand{\rhrbPctRidgeHarClose}{0.3}  % QLIKE change, %, W = 3,500 vs 3,000, 16:00 bar, Ridge, HAR benchmark
\providecommand{\rhrbDMRidgeHarClose}{}\renewcommand{\rhrbDMRidgeHarClose}{1.31}  % DM, W = 3,500 vs 3,000 (positive: 3,500 has the higher loss), 16:00 bar, Ridge, HAR benchmark
\providecommand{\rhrbPRidgeHarClose}{}\renewcommand{\rhrbPRidgeHarClose}{\ensuremath{p = 0.19}}  % ... its p statement
\providecommand{\rhrbDSRRidgeHarMid}{}\renewcommand{\rhrbDSRRidgeHarMid}{0.05}  % Sharpe ratio change, W = 3,500 minus 3,000, mid-quotes, Ridge, HAR benchmark
\providecommand{\rhrbPSRRidgeHarMid}{}\renewcommand{\rhrbPSRRidgeHarMid}{\ensuremath{p = 0.58}}  % ... its p statement (Ledoit-Wolf, D19; was the D6 bootstrap)
\providecommand{\rhrbDSRRidgeHarBa}{}\renewcommand{\rhrbDSRRidgeHarBa}{0.05}  % Sharpe ratio change, W = 3,500 minus 3,000, bid/ask prices, Ridge, HAR benchmark
\providecommand{\rhrbPSRRidgeHarBa}{}\renewcommand{\rhrbPSRRidgeHarBa}{\ensuremath{p = 0.58}}  % ... its p statement (Ledoit-Wolf, D19; was the D6 bootstrap)
\providecommand{\rhrbQRidgeFullBarMain}{}\renewcommand{\rhrbQRidgeFullBarMain}{0.1791}  % QLIKE, every bar, Ridge, full predictor set, W = 3,000
\providecommand{\rhrbQRidgeFullCloseMain}{}\renewcommand{\rhrbQRidgeFullCloseMain}{0.1279}  % QLIKE, 16:00 bar, Ridge, full predictor set, W = 3,000
\providecommand{\rhrbSRRidgeFullMidMain}{}\renewcommand{\rhrbSRRidgeFullMidMain}{1.77}  % Sharpe ratio (D19 pooled), sign strategy at mid-quotes, Ridge, full predictor set, W = 3,000
\providecommand{\rhrbSRRidgeFullBaMain}{}\renewcommand{\rhrbSRRidgeFullBaMain}{1.41}  % Sharpe ratio (D19 pooled), sign strategy at bid/ask prices, Ridge, full predictor set, W = 3,000
\providecommand{\rhrbQRidgeFullBarLong}{}\renewcommand{\rhrbQRidgeFullBarLong}{0.1777}  % QLIKE, every bar, Ridge, full predictor set, W = 3,500
\providecommand{\rhrbQRidgeFullCloseLong}{}\renewcommand{\rhrbQRidgeFullCloseLong}{0.1233}  % QLIKE, 16:00 bar, Ridge, full predictor set, W = 3,500
\providecommand{\rhrbSRRidgeFullMidLong}{}\renewcommand{\rhrbSRRidgeFullMidLong}{1.84}  % Sharpe ratio (D19 pooled), sign strategy at mid-quotes, Ridge, full predictor set, W = 3,500
\providecommand{\rhrbSRRidgeFullBaLong}{}\renewcommand{\rhrbSRRidgeFullBaLong}{1.48}  % Sharpe ratio (D19 pooled), sign strategy at bid/ask prices, Ridge, full predictor set, W = 3,500
\providecommand{\rhrbPctRidgeFullBar}{}\renewcommand{\rhrbPctRidgeFullBar}{\ensuremath{-}0.8}  % QLIKE change, %, W = 3,500 vs 3,000, every bar, Ridge, full predictor set
\providecommand{\rhrbDMRidgeFullBar}{}\renewcommand{\rhrbDMRidgeFullBar}{\ensuremath{-}2.26}  % DM, W = 3,500 vs 3,000 (positive: 3,500 has the higher loss), every bar, Ridge, full predictor set
\providecommand{\rhrbPRidgeFullBar}{}\renewcommand{\rhrbPRidgeFullBar}{\ensuremath{p = 0.024}}  % ... its p statement
\providecommand{\rhrbPctRidgeFullClose}{}\renewcommand{\rhrbPctRidgeFullClose}{\ensuremath{-}3.6}  % QLIKE change, %, W = 3,500 vs 3,000, 16:00 bar, Ridge, full predictor set
\providecommand{\rhrbDMRidgeFullClose}{}\renewcommand{\rhrbDMRidgeFullClose}{\ensuremath{-}3.10}  % DM, W = 3,500 vs 3,000 (positive: 3,500 has the higher loss), 16:00 bar, Ridge, full predictor set
\providecommand{\rhrbPRidgeFullClose}{}\renewcommand{\rhrbPRidgeFullClose}{\ensuremath{p = 0.002}}  % ... its p statement
\providecommand{\rhrbDSRRidgeFullMid}{}\renewcommand{\rhrbDSRRidgeFullMid}{0.07}  % Sharpe ratio change, W = 3,500 minus 3,000, mid-quotes, Ridge, full predictor set
\providecommand{\rhrbPSRRidgeFullMid}{}\renewcommand{\rhrbPSRRidgeFullMid}{\ensuremath{p = 0.81}}  % ... its p statement (Ledoit-Wolf, D19; was the D6 bootstrap)
\providecommand{\rhrbDSRRidgeFullBa}{}\renewcommand{\rhrbDSRRidgeFullBa}{0.07}  % Sharpe ratio change, W = 3,500 minus 3,000, bid/ask prices, Ridge, full predictor set
\providecommand{\rhrbPSRRidgeFullBa}{}\renewcommand{\rhrbPSRRidgeFullBa}{\ensuremath{p = 0.81}}  % ... its p statement (Ledoit-Wolf, D19; was the D6 bootstrap)
\providecommand{\rhrbQLassoHarBarMain}{}\renewcommand{\rhrbQLassoHarBarMain}{0.1893}  % QLIKE, every bar, Lasso, HAR benchmark, W = 3,000
\providecommand{\rhrbQLassoHarCloseMain}{}\renewcommand{\rhrbQLassoHarCloseMain}{0.1269}  % QLIKE, 16:00 bar, Lasso, HAR benchmark, W = 3,000
\providecommand{\rhrbSRLassoHarMidMain}{}\renewcommand{\rhrbSRLassoHarMidMain}{0.95}  % Sharpe ratio (D19 pooled), sign strategy at mid-quotes, Lasso, HAR benchmark, W = 3,000
\providecommand{\rhrbSRLassoHarBaMain}{}\renewcommand{\rhrbSRLassoHarBaMain}{0.60}  % Sharpe ratio (D19 pooled), sign strategy at bid/ask prices, Lasso, HAR benchmark, W = 3,000
\providecommand{\rhrbQLassoHarBarLong}{}\renewcommand{\rhrbQLassoHarBarLong}{0.1897}  % QLIKE, every bar, Lasso, HAR benchmark, W = 3,500
\providecommand{\rhrbQLassoHarCloseLong}{}\renewcommand{\rhrbQLassoHarCloseLong}{0.1272}  % QLIKE, 16:00 bar, Lasso, HAR benchmark, W = 3,500
\providecommand{\rhrbSRLassoHarMidLong}{}\renewcommand{\rhrbSRLassoHarMidLong}{1.01}  % Sharpe ratio (D19 pooled), sign strategy at mid-quotes, Lasso, HAR benchmark, W = 3,500
\providecommand{\rhrbSRLassoHarBaLong}{}\renewcommand{\rhrbSRLassoHarBaLong}{0.65}  % Sharpe ratio (D19 pooled), sign strategy at bid/ask prices, Lasso, HAR benchmark, W = 3,500
\providecommand{\rhrbPctLassoHarBar}{}\renewcommand{\rhrbPctLassoHarBar}{0.2}  % QLIKE change, %, W = 3,500 vs 3,000, every bar, Lasso, HAR benchmark
\providecommand{\rhrbDMLassoHarBar}{}\renewcommand{\rhrbDMLassoHarBar}{2.48}  % DM, W = 3,500 vs 3,000 (positive: 3,500 has the higher loss), every bar, Lasso, HAR benchmark
\providecommand{\rhrbPLassoHarBar}{}\renewcommand{\rhrbPLassoHarBar}{\ensuremath{p = 0.013}}  % ... its p statement
\providecommand{\rhrbPctLassoHarClose}{}\renewcommand{\rhrbPctLassoHarClose}{0.3}  % QLIKE change, %, W = 3,500 vs 3,000, 16:00 bar, Lasso, HAR benchmark
\providecommand{\rhrbDMLassoHarClose}{}\renewcommand{\rhrbDMLassoHarClose}{1.31}  % DM, W = 3,500 vs 3,000 (positive: 3,500 has the higher loss), 16:00 bar, Lasso, HAR benchmark
\providecommand{\rhrbPLassoHarClose}{}\renewcommand{\rhrbPLassoHarClose}{\ensuremath{p = 0.19}}  % ... its p statement
\providecommand{\rhrbDSRLassoHarMid}{}\renewcommand{\rhrbDSRLassoHarMid}{0.05}  % Sharpe ratio change, W = 3,500 minus 3,000, mid-quotes, Lasso, HAR benchmark
\providecommand{\rhrbPSRLassoHarMid}{}\renewcommand{\rhrbPSRLassoHarMid}{\ensuremath{p = 0.58}}  % ... its p statement (Ledoit-Wolf, D19; was the D6 bootstrap)
\providecommand{\rhrbDSRLassoHarBa}{}\renewcommand{\rhrbDSRLassoHarBa}{0.05}  % Sharpe ratio change, W = 3,500 minus 3,000, bid/ask prices, Lasso, HAR benchmark
\providecommand{\rhrbPSRLassoHarBa}{}\renewcommand{\rhrbPSRLassoHarBa}{\ensuremath{p = 0.58}}  % ... its p statement (Ledoit-Wolf, D19; was the D6 bootstrap)
\providecommand{\rhrbQLassoFullBarMain}{}\renewcommand{\rhrbQLassoFullBarMain}{0.1775}  % QLIKE, every bar, Lasso, full predictor set, W = 3,000
\providecommand{\rhrbQLassoFullCloseMain}{}\renewcommand{\rhrbQLassoFullCloseMain}{0.1310}  % QLIKE, 16:00 bar, Lasso, full predictor set, W = 3,000
\providecommand{\rhrbSRLassoFullMidMain}{}\renewcommand{\rhrbSRLassoFullMidMain}{1.72}  % Sharpe ratio (D19 pooled), sign strategy at mid-quotes, Lasso, full predictor set, W = 3,000
\providecommand{\rhrbSRLassoFullBaMain}{}\renewcommand{\rhrbSRLassoFullBaMain}{1.37}  % Sharpe ratio (D19 pooled), sign strategy at bid/ask prices, Lasso, full predictor set, W = 3,000
\providecommand{\rhrbQLassoFullBarLong}{}\renewcommand{\rhrbQLassoFullBarLong}{0.1752}  % QLIKE, every bar, Lasso, full predictor set, W = 3,500
\providecommand{\rhrbQLassoFullCloseLong}{}\renewcommand{\rhrbQLassoFullCloseLong}{0.1253}  % QLIKE, 16:00 bar, Lasso, full predictor set, W = 3,500
\providecommand{\rhrbSRLassoFullMidLong}{}\renewcommand{\rhrbSRLassoFullMidLong}{1.80}  % Sharpe ratio (D19 pooled), sign strategy at mid-quotes, Lasso, full predictor set, W = 3,500
\providecommand{\rhrbSRLassoFullBaLong}{}\renewcommand{\rhrbSRLassoFullBaLong}{1.44}  % Sharpe ratio (D19 pooled), sign strategy at bid/ask prices, Lasso, full predictor set, W = 3,500
\providecommand{\rhrbPctLassoFullBar}{}\renewcommand{\rhrbPctLassoFullBar}{\ensuremath{-}1.3}  % QLIKE change, %, W = 3,500 vs 3,000, every bar, Lasso, full predictor set
\providecommand{\rhrbDMLassoFullBar}{}\renewcommand{\rhrbDMLassoFullBar}{\ensuremath{-}2.04}  % DM, W = 3,500 vs 3,000 (positive: 3,500 has the higher loss), every bar, Lasso, full predictor set
\providecommand{\rhrbPLassoFullBar}{}\renewcommand{\rhrbPLassoFullBar}{\ensuremath{p = 0.042}}  % ... its p statement
\providecommand{\rhrbPctLassoFullClose}{}\renewcommand{\rhrbPctLassoFullClose}{\ensuremath{-}4.4}  % QLIKE change, %, W = 3,500 vs 3,000, 16:00 bar, Lasso, full predictor set
\providecommand{\rhrbDMLassoFullClose}{}\renewcommand{\rhrbDMLassoFullClose}{\ensuremath{-}2.89}  % DM, W = 3,500 vs 3,000 (positive: 3,500 has the higher loss), 16:00 bar, Lasso, full predictor set
\providecommand{\rhrbPLassoFullClose}{}\renewcommand{\rhrbPLassoFullClose}{\ensuremath{p = 0.004}}  % ... its p statement
\providecommand{\rhrbDSRLassoFullMid}{}\renewcommand{\rhrbDSRLassoFullMid}{0.07}  % Sharpe ratio change, W = 3,500 minus 3,000, mid-quotes, Lasso, full predictor set
\providecommand{\rhrbPSRLassoFullMid}{}\renewcommand{\rhrbPSRLassoFullMid}{\ensuremath{p = 0.83}}  % ... its p statement (Ledoit-Wolf, D19; was the D6 bootstrap)
\providecommand{\rhrbDSRLassoFullBa}{}\renewcommand{\rhrbDSRLassoFullBa}{0.07}  % Sharpe ratio change, W = 3,500 minus 3,000, bid/ask prices, Lasso, full predictor set
\providecommand{\rhrbPSRLassoFullBa}{}\renewcommand{\rhrbPSRLassoFullBa}{\ensuremath{p = 0.83}}  % ... its p statement (Ledoit-Wolf, D19; was the D6 bootstrap)
\providecommand{\rhrbQEnetHarBarMain}{}\renewcommand{\rhrbQEnetHarBarMain}{0.1893}  % QLIKE, every bar, Elastic net, HAR benchmark, W = 3,000
\providecommand{\rhrbQEnetHarCloseMain}{}\renewcommand{\rhrbQEnetHarCloseMain}{0.1269}  % QLIKE, 16:00 bar, Elastic net, HAR benchmark, W = 3,000
\providecommand{\rhrbSREnetHarMidMain}{}\renewcommand{\rhrbSREnetHarMidMain}{0.95}  % Sharpe ratio (D19 pooled), sign strategy at mid-quotes, Elastic net, HAR benchmark, W = 3,000
\providecommand{\rhrbSREnetHarBaMain}{}\renewcommand{\rhrbSREnetHarBaMain}{0.60}  % Sharpe ratio (D19 pooled), sign strategy at bid/ask prices, Elastic net, HAR benchmark, W = 3,000
\providecommand{\rhrbQEnetHarBarLong}{}\renewcommand{\rhrbQEnetHarBarLong}{0.1897}  % QLIKE, every bar, Elastic net, HAR benchmark, W = 3,500
\providecommand{\rhrbQEnetHarCloseLong}{}\renewcommand{\rhrbQEnetHarCloseLong}{0.1272}  % QLIKE, 16:00 bar, Elastic net, HAR benchmark, W = 3,500
\providecommand{\rhrbSREnetHarMidLong}{}\renewcommand{\rhrbSREnetHarMidLong}{1.01}  % Sharpe ratio (D19 pooled), sign strategy at mid-quotes, Elastic net, HAR benchmark, W = 3,500
\providecommand{\rhrbSREnetHarBaLong}{}\renewcommand{\rhrbSREnetHarBaLong}{0.65}  % Sharpe ratio (D19 pooled), sign strategy at bid/ask prices, Elastic net, HAR benchmark, W = 3,500
\providecommand{\rhrbPctEnetHarBar}{}\renewcommand{\rhrbPctEnetHarBar}{0.2}  % QLIKE change, %, W = 3,500 vs 3,000, every bar, Elastic net, HAR benchmark
\providecommand{\rhrbDMEnetHarBar}{}\renewcommand{\rhrbDMEnetHarBar}{2.48}  % DM, W = 3,500 vs 3,000 (positive: 3,500 has the higher loss), every bar, Elastic net, HAR benchmark
\providecommand{\rhrbPEnetHarBar}{}\renewcommand{\rhrbPEnetHarBar}{\ensuremath{p = 0.013}}  % ... its p statement
\providecommand{\rhrbPctEnetHarClose}{}\renewcommand{\rhrbPctEnetHarClose}{0.3}  % QLIKE change, %, W = 3,500 vs 3,000, 16:00 bar, Elastic net, HAR benchmark
\providecommand{\rhrbDMEnetHarClose}{}\renewcommand{\rhrbDMEnetHarClose}{1.31}  % DM, W = 3,500 vs 3,000 (positive: 3,500 has the higher loss), 16:00 bar, Elastic net, HAR benchmark
\providecommand{\rhrbPEnetHarClose}{}\renewcommand{\rhrbPEnetHarClose}{\ensuremath{p = 0.19}}  % ... its p statement
\providecommand{\rhrbDSREnetHarMid}{}\renewcommand{\rhrbDSREnetHarMid}{0.05}  % Sharpe ratio change, W = 3,500 minus 3,000, mid-quotes, Elastic net, HAR benchmark
\providecommand{\rhrbPSREnetHarMid}{}\renewcommand{\rhrbPSREnetHarMid}{\ensuremath{p = 0.58}}  % ... its p statement (Ledoit-Wolf, D19; was the D6 bootstrap)
\providecommand{\rhrbDSREnetHarBa}{}\renewcommand{\rhrbDSREnetHarBa}{0.05}  % Sharpe ratio change, W = 3,500 minus 3,000, bid/ask prices, Elastic net, HAR benchmark
\providecommand{\rhrbPSREnetHarBa}{}\renewcommand{\rhrbPSREnetHarBa}{\ensuremath{p = 0.58}}  % ... its p statement (Ledoit-Wolf, D19; was the D6 bootstrap)
\providecommand{\rhrbQEnetFullBarMain}{}\renewcommand{\rhrbQEnetFullBarMain}{0.1778}  % QLIKE, every bar, Elastic net, full predictor set, W = 3,000
\providecommand{\rhrbQEnetFullCloseMain}{}\renewcommand{\rhrbQEnetFullCloseMain}{0.1262}  % QLIKE, 16:00 bar, Elastic net, full predictor set, W = 3,000
\providecommand{\rhrbSREnetFullMidMain}{}\renewcommand{\rhrbSREnetFullMidMain}{1.86}  % Sharpe ratio (D19 pooled), sign strategy at mid-quotes, Elastic net, full predictor set, W = 3,000
\providecommand{\rhrbSREnetFullBaMain}{}\renewcommand{\rhrbSREnetFullBaMain}{1.50}  % Sharpe ratio (D19 pooled), sign strategy at bid/ask prices, Elastic net, full predictor set, W = 3,000
\providecommand{\rhrbQEnetFullBarLong}{}\renewcommand{\rhrbQEnetFullBarLong}{0.1770}  % QLIKE, every bar, Elastic net, full predictor set, W = 3,500
\providecommand{\rhrbQEnetFullCloseLong}{}\renewcommand{\rhrbQEnetFullCloseLong}{0.1223}  % QLIKE, 16:00 bar, Elastic net, full predictor set, W = 3,500
\providecommand{\rhrbSREnetFullMidLong}{}\renewcommand{\rhrbSREnetFullMidLong}{2.00}  % Sharpe ratio (D19 pooled), sign strategy at mid-quotes, Elastic net, full predictor set, W = 3,500
\providecommand{\rhrbSREnetFullBaLong}{}\renewcommand{\rhrbSREnetFullBaLong}{1.65}  % Sharpe ratio (D19 pooled), sign strategy at bid/ask prices, Elastic net, full predictor set, W = 3,500
\providecommand{\rhrbPctEnetFullBar}{}\renewcommand{\rhrbPctEnetFullBar}{\ensuremath{-}0.5}  % QLIKE change, %, W = 3,500 vs 3,000, every bar, Elastic net, full predictor set
\providecommand{\rhrbDMEnetFullBar}{}\renewcommand{\rhrbDMEnetFullBar}{\ensuremath{-}1.47}  % DM, W = 3,500 vs 3,000 (positive: 3,500 has the higher loss), every bar, Elastic net, full predictor set
\providecommand{\rhrbPEnetFullBar}{}\renewcommand{\rhrbPEnetFullBar}{\ensuremath{p = 0.14}}  % ... its p statement
\providecommand{\rhrbPctEnetFullClose}{}\renewcommand{\rhrbPctEnetFullClose}{\ensuremath{-}3.1}  % QLIKE change, %, W = 3,500 vs 3,000, 16:00 bar, Elastic net, full predictor set
\providecommand{\rhrbDMEnetFullClose}{}\renewcommand{\rhrbDMEnetFullClose}{\ensuremath{-}2.93}  % DM, W = 3,500 vs 3,000 (positive: 3,500 has the higher loss), 16:00 bar, Elastic net, full predictor set
\providecommand{\rhrbPEnetFullClose}{}\renewcommand{\rhrbPEnetFullClose}{\ensuremath{p = 0.003}}  % ... its p statement
\providecommand{\rhrbDSREnetFullMid}{}\renewcommand{\rhrbDSREnetFullMid}{0.15}  % Sharpe ratio change, W = 3,500 minus 3,000, mid-quotes, Elastic net, full predictor set
\providecommand{\rhrbPSREnetFullMid}{}\renewcommand{\rhrbPSREnetFullMid}{\ensuremath{p = 0.56}}  % ... its p statement (Ledoit-Wolf, D19; was the D6 bootstrap)
\providecommand{\rhrbDSREnetFullBa}{}\renewcommand{\rhrbDSREnetFullBa}{0.15}  % Sharpe ratio change, W = 3,500 minus 3,000, bid/ask prices, Elastic net, full predictor set
\providecommand{\rhrbPSREnetFullBa}{}\renewcommand{\rhrbPSREnetFullBa}{\ensuremath{p = 0.55}}  % ... its p statement (Ledoit-Wolf, D19; was the D6 bootstrap)
\providecommand{\rhrbQLgbmHarBarMain}{}\renewcommand{\rhrbQLgbmHarBarMain}{0.1875}  % QLIKE, every bar, LightGBM, HAR benchmark, W = 3,000
\providecommand{\rhrbQLgbmHarCloseMain}{}\renewcommand{\rhrbQLgbmHarCloseMain}{0.1235}  % QLIKE, 16:00 bar, LightGBM, HAR benchmark, W = 3,000
\providecommand{\rhrbSRLgbmHarMidMain}{}\renewcommand{\rhrbSRLgbmHarMidMain}{0.65}  % Sharpe ratio (D19 pooled), sign strategy at mid-quotes, LightGBM, HAR benchmark, W = 3,000
\providecommand{\rhrbSRLgbmHarBaMain}{}\renewcommand{\rhrbSRLgbmHarBaMain}{0.31}  % Sharpe ratio (D19 pooled), sign strategy at bid/ask prices, LightGBM, HAR benchmark, W = 3,000
\providecommand{\rhrbQLgbmHarBarLong}{}\renewcommand{\rhrbQLgbmHarBarLong}{0.1876}  % QLIKE, every bar, LightGBM, HAR benchmark, W = 3,500
\providecommand{\rhrbQLgbmHarCloseLong}{}\renewcommand{\rhrbQLgbmHarCloseLong}{0.1233}  % QLIKE, 16:00 bar, LightGBM, HAR benchmark, W = 3,500
\providecommand{\rhrbSRLgbmHarMidLong}{}\renewcommand{\rhrbSRLgbmHarMidLong}{0.59}  % Sharpe ratio (D19 pooled), sign strategy at mid-quotes, LightGBM, HAR benchmark, W = 3,500
\providecommand{\rhrbSRLgbmHarBaLong}{}\renewcommand{\rhrbSRLgbmHarBaLong}{0.24}  % Sharpe ratio (D19 pooled), sign strategy at bid/ask prices, LightGBM, HAR benchmark, W = 3,500
\providecommand{\rhrbPctLgbmHarBar}{}\renewcommand{\rhrbPctLgbmHarBar}{0.0}  % QLIKE change, %, W = 3,500 vs 3,000, every bar, LightGBM, HAR benchmark
\providecommand{\rhrbDMLgbmHarBar}{}\renewcommand{\rhrbDMLgbmHarBar}{0.20}  % DM, W = 3,500 vs 3,000 (positive: 3,500 has the higher loss), every bar, LightGBM, HAR benchmark
\providecommand{\rhrbPLgbmHarBar}{}\renewcommand{\rhrbPLgbmHarBar}{\ensuremath{p = 0.84}}  % ... its p statement
\providecommand{\rhrbPctLgbmHarClose}{}\renewcommand{\rhrbPctLgbmHarClose}{\ensuremath{-}0.2}  % QLIKE change, %, W = 3,500 vs 3,000, 16:00 bar, LightGBM, HAR benchmark
\providecommand{\rhrbDMLgbmHarClose}{}\renewcommand{\rhrbDMLgbmHarClose}{\ensuremath{-}0.45}  % DM, W = 3,500 vs 3,000 (positive: 3,500 has the higher loss), 16:00 bar, LightGBM, HAR benchmark
\providecommand{\rhrbPLgbmHarClose}{}\renewcommand{\rhrbPLgbmHarClose}{\ensuremath{p = 0.65}}  % ... its p statement
\providecommand{\rhrbDSRLgbmHarMid}{}\renewcommand{\rhrbDSRLgbmHarMid}{\ensuremath{-}0.06}  % Sharpe ratio change, W = 3,500 minus 3,000, mid-quotes, LightGBM, HAR benchmark
\providecommand{\rhrbPSRLgbmHarMid}{}\renewcommand{\rhrbPSRLgbmHarMid}{\ensuremath{p = 0.57}}  % ... its p statement (Ledoit-Wolf, D19; was the D6 bootstrap)
\providecommand{\rhrbDSRLgbmHarBa}{}\renewcommand{\rhrbDSRLgbmHarBa}{\ensuremath{-}0.06}  % Sharpe ratio change, W = 3,500 minus 3,000, bid/ask prices, LightGBM, HAR benchmark
\providecommand{\rhrbPSRLgbmHarBa}{}\renewcommand{\rhrbPSRLgbmHarBa}{\ensuremath{p = 0.57}}  % ... its p statement (Ledoit-Wolf, D19; was the D6 bootstrap)
\providecommand{\rhrbQLgbmFullBarMain}{}\renewcommand{\rhrbQLgbmFullBarMain}{0.1613}  % QLIKE, every bar, LightGBM, full predictor set, W = 3,000
\providecommand{\rhrbQLgbmFullCloseMain}{}\renewcommand{\rhrbQLgbmFullCloseMain}{0.1201}  % QLIKE, 16:00 bar, LightGBM, full predictor set, W = 3,000
\providecommand{\rhrbSRLgbmFullMidMain}{}\renewcommand{\rhrbSRLgbmFullMidMain}{1.59}  % Sharpe ratio (D19 pooled), sign strategy at mid-quotes, LightGBM, full predictor set, W = 3,000
\providecommand{\rhrbSRLgbmFullBaMain}{}\renewcommand{\rhrbSRLgbmFullBaMain}{1.23}  % Sharpe ratio (D19 pooled), sign strategy at bid/ask prices, LightGBM, full predictor set, W = 3,000
\providecommand{\rhrbQLgbmFullBarLong}{}\renewcommand{\rhrbQLgbmFullBarLong}{0.1613}  % QLIKE, every bar, LightGBM, full predictor set, W = 3,500
\providecommand{\rhrbQLgbmFullCloseLong}{}\renewcommand{\rhrbQLgbmFullCloseLong}{0.1204}  % QLIKE, 16:00 bar, LightGBM, full predictor set, W = 3,500
\providecommand{\rhrbSRLgbmFullMidLong}{}\renewcommand{\rhrbSRLgbmFullMidLong}{1.54}  % Sharpe ratio (D19 pooled), sign strategy at mid-quotes, LightGBM, full predictor set, W = 3,500
\providecommand{\rhrbSRLgbmFullBaLong}{}\renewcommand{\rhrbSRLgbmFullBaLong}{1.19}  % Sharpe ratio (D19 pooled), sign strategy at bid/ask prices, LightGBM, full predictor set, W = 3,500
\providecommand{\rhrbPctLgbmFullBar}{}\renewcommand{\rhrbPctLgbmFullBar}{0.0}  % QLIKE change, %, W = 3,500 vs 3,000, every bar, LightGBM, full predictor set
\providecommand{\rhrbDMLgbmFullBar}{}\renewcommand{\rhrbDMLgbmFullBar}{0.20}  % DM, W = 3,500 vs 3,000 (positive: 3,500 has the higher loss), every bar, LightGBM, full predictor set
\providecommand{\rhrbPLgbmFullBar}{}\renewcommand{\rhrbPLgbmFullBar}{\ensuremath{p = 0.84}}  % ... its p statement
\providecommand{\rhrbPctLgbmFullClose}{}\renewcommand{\rhrbPctLgbmFullClose}{0.2}  % QLIKE change, %, W = 3,500 vs 3,000, 16:00 bar, LightGBM, full predictor set
\providecommand{\rhrbDMLgbmFullClose}{}\renewcommand{\rhrbDMLgbmFullClose}{0.55}  % DM, W = 3,500 vs 3,000 (positive: 3,500 has the higher loss), 16:00 bar, LightGBM, full predictor set
\providecommand{\rhrbPLgbmFullClose}{}\renewcommand{\rhrbPLgbmFullClose}{\ensuremath{p = 0.58}}  % ... its p statement
\providecommand{\rhrbDSRLgbmFullMid}{}\renewcommand{\rhrbDSRLgbmFullMid}{\ensuremath{-}0.05}  % Sharpe ratio change, W = 3,500 minus 3,000, mid-quotes, LightGBM, full predictor set
\providecommand{\rhrbPSRLgbmFullMid}{}\renewcommand{\rhrbPSRLgbmFullMid}{\ensuremath{p = 0.76}}  % ... its p statement (Ledoit-Wolf, D19; was the D6 bootstrap)
\providecommand{\rhrbDSRLgbmFullBa}{}\renewcommand{\rhrbDSRLgbmFullBa}{\ensuremath{-}0.05}  % Sharpe ratio change, W = 3,500 minus 3,000, bid/ask prices, LightGBM, full predictor set
\providecommand{\rhrbPSRLgbmFullBa}{}\renewcommand{\rhrbPSRLgbmFullBa}{\ensuremath{p = 0.76}}  % ... its p statement (Ledoit-Wolf, D19; was the D6 bootstrap)
\providecommand{\rhrbQXgbHarBarMain}{}\renewcommand{\rhrbQXgbHarBarMain}{0.1888}  % QLIKE, every bar, XGBoost, HAR benchmark, W = 3,000
\providecommand{\rhrbQXgbHarCloseMain}{}\renewcommand{\rhrbQXgbHarCloseMain}{0.1234}  % QLIKE, 16:00 bar, XGBoost, HAR benchmark, W = 3,000
\providecommand{\rhrbSRXgbHarMidMain}{}\renewcommand{\rhrbSRXgbHarMidMain}{0.98}  % Sharpe ratio (D19 pooled), sign strategy at mid-quotes, XGBoost, HAR benchmark, W = 3,000
\providecommand{\rhrbSRXgbHarBaMain}{}\renewcommand{\rhrbSRXgbHarBaMain}{0.63}  % Sharpe ratio (D19 pooled), sign strategy at bid/ask prices, XGBoost, HAR benchmark, W = 3,000
\providecommand{\rhrbQXgbHarBarLong}{}\renewcommand{\rhrbQXgbHarBarLong}{0.1890}  % QLIKE, every bar, XGBoost, HAR benchmark, W = 3,500
\providecommand{\rhrbQXgbHarCloseLong}{}\renewcommand{\rhrbQXgbHarCloseLong}{0.1232}  % QLIKE, 16:00 bar, XGBoost, HAR benchmark, W = 3,500
\providecommand{\rhrbSRXgbHarMidLong}{}\renewcommand{\rhrbSRXgbHarMidLong}{0.78}  % Sharpe ratio (D19 pooled), sign strategy at mid-quotes, XGBoost, HAR benchmark, W = 3,500
\providecommand{\rhrbSRXgbHarBaLong}{}\renewcommand{\rhrbSRXgbHarBaLong}{0.43}  % Sharpe ratio (D19 pooled), sign strategy at bid/ask prices, XGBoost, HAR benchmark, W = 3,500
\providecommand{\rhrbPctXgbHarBar}{}\renewcommand{\rhrbPctXgbHarBar}{0.1}  % QLIKE change, %, W = 3,500 vs 3,000, every bar, XGBoost, HAR benchmark
\providecommand{\rhrbDMXgbHarBar}{}\renewcommand{\rhrbDMXgbHarBar}{1.13}  % DM, W = 3,500 vs 3,000 (positive: 3,500 has the higher loss), every bar, XGBoost, HAR benchmark
\providecommand{\rhrbPXgbHarBar}{}\renewcommand{\rhrbPXgbHarBar}{\ensuremath{p = 0.26}}  % ... its p statement
\providecommand{\rhrbPctXgbHarClose}{}\renewcommand{\rhrbPctXgbHarClose}{\ensuremath{-}0.2}  % QLIKE change, %, W = 3,500 vs 3,000, 16:00 bar, XGBoost, HAR benchmark
\providecommand{\rhrbDMXgbHarClose}{}\renewcommand{\rhrbDMXgbHarClose}{\ensuremath{-}0.55}  % DM, W = 3,500 vs 3,000 (positive: 3,500 has the higher loss), 16:00 bar, XGBoost, HAR benchmark
\providecommand{\rhrbPXgbHarClose}{}\renewcommand{\rhrbPXgbHarClose}{\ensuremath{p = 0.58}}  % ... its p statement
\providecommand{\rhrbDSRXgbHarMid}{}\renewcommand{\rhrbDSRXgbHarMid}{\ensuremath{-}0.20}  % Sharpe ratio change, W = 3,500 minus 3,000, mid-quotes, XGBoost, HAR benchmark
\providecommand{\rhrbPSRXgbHarMid}{}\renewcommand{\rhrbPSRXgbHarMid}{\ensuremath{p = 0.40}}  % ... its p statement (Ledoit-Wolf, D19; was the D6 bootstrap)
\providecommand{\rhrbDSRXgbHarBa}{}\renewcommand{\rhrbDSRXgbHarBa}{\ensuremath{-}0.20}  % Sharpe ratio change, W = 3,500 minus 3,000, bid/ask prices, XGBoost, HAR benchmark
\providecommand{\rhrbPSRXgbHarBa}{}\renewcommand{\rhrbPSRXgbHarBa}{\ensuremath{p = 0.40}}  % ... its p statement (Ledoit-Wolf, D19; was the D6 bootstrap)
\providecommand{\rhrbQXgbFullBarMain}{}\renewcommand{\rhrbQXgbFullBarMain}{0.1627}  % QLIKE, every bar, XGBoost, full predictor set, W = 3,000
\providecommand{\rhrbQXgbFullCloseMain}{}\renewcommand{\rhrbQXgbFullCloseMain}{0.1192}  % QLIKE, 16:00 bar, XGBoost, full predictor set, W = 3,000
\providecommand{\rhrbSRXgbFullMidMain}{}\renewcommand{\rhrbSRXgbFullMidMain}{1.47}  % Sharpe ratio (D19 pooled), sign strategy at mid-quotes, XGBoost, full predictor set, W = 3,000
\providecommand{\rhrbSRXgbFullBaMain}{}\renewcommand{\rhrbSRXgbFullBaMain}{1.12}  % Sharpe ratio (D19 pooled), sign strategy at bid/ask prices, XGBoost, full predictor set, W = 3,000
\providecommand{\rhrbQXgbFullBarLong}{}\renewcommand{\rhrbQXgbFullBarLong}{0.1636}  % QLIKE, every bar, XGBoost, full predictor set, W = 3,500
\providecommand{\rhrbQXgbFullCloseLong}{}\renewcommand{\rhrbQXgbFullCloseLong}{0.1192}  % QLIKE, 16:00 bar, XGBoost, full predictor set, W = 3,500
\providecommand{\rhrbSRXgbFullMidLong}{}\renewcommand{\rhrbSRXgbFullMidLong}{1.35}  % Sharpe ratio (D19 pooled), sign strategy at mid-quotes, XGBoost, full predictor set, W = 3,500
\providecommand{\rhrbSRXgbFullBaLong}{}\renewcommand{\rhrbSRXgbFullBaLong}{0.99}  % Sharpe ratio (D19 pooled), sign strategy at bid/ask prices, XGBoost, full predictor set, W = 3,500
\providecommand{\rhrbPctXgbFullBar}{}\renewcommand{\rhrbPctXgbFullBar}{0.5}  % QLIKE change, %, W = 3,500 vs 3,000, every bar, XGBoost, full predictor set
\providecommand{\rhrbDMXgbFullBar}{}\renewcommand{\rhrbDMXgbFullBar}{1.87}  % DM, W = 3,500 vs 3,000 (positive: 3,500 has the higher loss), every bar, XGBoost, full predictor set
\providecommand{\rhrbPXgbFullBar}{}\renewcommand{\rhrbPXgbFullBar}{\ensuremath{p = 0.062}}  % ... its p statement
\providecommand{\rhrbPctXgbFullClose}{}\renewcommand{\rhrbPctXgbFullClose}{\ensuremath{-}0.1}  % QLIKE change, %, W = 3,500 vs 3,000, 16:00 bar, XGBoost, full predictor set
\providecommand{\rhrbDMXgbFullClose}{}\renewcommand{\rhrbDMXgbFullClose}{\ensuremath{-}0.15}  % DM, W = 3,500 vs 3,000 (positive: 3,500 has the higher loss), 16:00 bar, XGBoost, full predictor set
\providecommand{\rhrbPXgbFullClose}{}\renewcommand{\rhrbPXgbFullClose}{\ensuremath{p = 0.88}}  % ... its p statement
\providecommand{\rhrbDSRXgbFullMid}{}\renewcommand{\rhrbDSRXgbFullMid}{\ensuremath{-}0.12}  % Sharpe ratio change, W = 3,500 minus 3,000, mid-quotes, XGBoost, full predictor set
\providecommand{\rhrbPSRXgbFullMid}{}\renewcommand{\rhrbPSRXgbFullMid}{\ensuremath{p = 0.62}}  % ... its p statement (Ledoit-Wolf, D19; was the D6 bootstrap)
\providecommand{\rhrbDSRXgbFullBa}{}\renewcommand{\rhrbDSRXgbFullBa}{\ensuremath{-}0.12}  % Sharpe ratio change, W = 3,500 minus 3,000, bid/ask prices, XGBoost, full predictor set
\providecommand{\rhrbPSRXgbFullBa}{}\renewcommand{\rhrbPSRXgbFullBa}{\ensuremath{p = 0.62}}  % ... its p statement (Ledoit-Wolf, D19; was the D6 bootstrap)
\providecommand{\rhrbQRfHarBarMain}{}\renewcommand{\rhrbQRfHarBarMain}{0.1895}  % QLIKE, every bar, Random forest, HAR benchmark, W = 3,000
\providecommand{\rhrbQRfHarCloseMain}{}\renewcommand{\rhrbQRfHarCloseMain}{0.1266}  % QLIKE, 16:00 bar, Random forest, HAR benchmark, W = 3,000
\providecommand{\rhrbSRRfHarMidMain}{}\renewcommand{\rhrbSRRfHarMidMain}{0.32}  % Sharpe ratio (D19 pooled), sign strategy at mid-quotes, Random forest, HAR benchmark, W = 3,000
\providecommand{\rhrbSRRfHarBaMain}{}\renewcommand{\rhrbSRRfHarBaMain}{\ensuremath{-}0.03}  % Sharpe ratio (D19 pooled), sign strategy at bid/ask prices, Random forest, HAR benchmark, W = 3,000
\providecommand{\rhrbQRfHarBarLong}{}\renewcommand{\rhrbQRfHarBarLong}{0.1898}  % QLIKE, every bar, Random forest, HAR benchmark, W = 3,500
\providecommand{\rhrbQRfHarCloseLong}{}\renewcommand{\rhrbQRfHarCloseLong}{0.1264}  % QLIKE, 16:00 bar, Random forest, HAR benchmark, W = 3,500
\providecommand{\rhrbSRRfHarMidLong}{}\renewcommand{\rhrbSRRfHarMidLong}{0.80}  % Sharpe ratio (D19 pooled), sign strategy at mid-quotes, Random forest, HAR benchmark, W = 3,500
\providecommand{\rhrbSRRfHarBaLong}{}\renewcommand{\rhrbSRRfHarBaLong}{0.45}  % Sharpe ratio (D19 pooled), sign strategy at bid/ask prices, Random forest, HAR benchmark, W = 3,500
\providecommand{\rhrbPctRfHarBar}{}\renewcommand{\rhrbPctRfHarBar}{0.1}  % QLIKE change, %, W = 3,500 vs 3,000, every bar, Random forest, HAR benchmark
\providecommand{\rhrbDMRfHarBar}{}\renewcommand{\rhrbDMRfHarBar}{1.15}  % DM, W = 3,500 vs 3,000 (positive: 3,500 has the higher loss), every bar, Random forest, HAR benchmark
\providecommand{\rhrbPRfHarBar}{}\renewcommand{\rhrbPRfHarBar}{\ensuremath{p = 0.25}}  % ... its p statement
\providecommand{\rhrbPctRfHarClose}{}\renewcommand{\rhrbPctRfHarClose}{\ensuremath{-}0.1}  % QLIKE change, %, W = 3,500 vs 3,000, 16:00 bar, Random forest, HAR benchmark
\providecommand{\rhrbDMRfHarClose}{}\renewcommand{\rhrbDMRfHarClose}{\ensuremath{-}0.36}  % DM, W = 3,500 vs 3,000 (positive: 3,500 has the higher loss), 16:00 bar, Random forest, HAR benchmark
\providecommand{\rhrbPRfHarClose}{}\renewcommand{\rhrbPRfHarClose}{\ensuremath{p = 0.72}}  % ... its p statement
\providecommand{\rhrbDSRRfHarMid}{}\renewcommand{\rhrbDSRRfHarMid}{0.48}  % Sharpe ratio change, W = 3,500 minus 3,000, mid-quotes, Random forest, HAR benchmark
\providecommand{\rhrbPSRRfHarMid}{}\renewcommand{\rhrbPSRRfHarMid}{\ensuremath{p = 0.42}}  % ... its p statement (Ledoit-Wolf, D19; was the D6 bootstrap)
\providecommand{\rhrbDSRRfHarBa}{}\renewcommand{\rhrbDSRRfHarBa}{0.48}  % Sharpe ratio change, W = 3,500 minus 3,000, bid/ask prices, Random forest, HAR benchmark
\providecommand{\rhrbPSRRfHarBa}{}\renewcommand{\rhrbPSRRfHarBa}{\ensuremath{p = 0.42}}  % ... its p statement (Ledoit-Wolf, D19; was the D6 bootstrap)
\providecommand{\rhrbQRfFullBarMain}{}\renewcommand{\rhrbQRfFullBarMain}{0.1685}  % QLIKE, every bar, Random forest, full predictor set, W = 3,000
\providecommand{\rhrbQRfFullCloseMain}{}\renewcommand{\rhrbQRfFullCloseMain}{0.1248}  % QLIKE, 16:00 bar, Random forest, full predictor set, W = 3,000
\providecommand{\rhrbSRRfFullMidMain}{}\renewcommand{\rhrbSRRfFullMidMain}{1.88}  % Sharpe ratio (D19 pooled), sign strategy at mid-quotes, Random forest, full predictor set, W = 3,000
\providecommand{\rhrbSRRfFullBaMain}{}\renewcommand{\rhrbSRRfFullBaMain}{1.52}  % Sharpe ratio (D19 pooled), sign strategy at bid/ask prices, Random forest, full predictor set, W = 3,000
\providecommand{\rhrbQRfFullBarLong}{}\renewcommand{\rhrbQRfFullBarLong}{0.1685}  % QLIKE, every bar, Random forest, full predictor set, W = 3,500
\providecommand{\rhrbQRfFullCloseLong}{}\renewcommand{\rhrbQRfFullCloseLong}{0.1241}  % QLIKE, 16:00 bar, Random forest, full predictor set, W = 3,500
\providecommand{\rhrbSRRfFullMidLong}{}\renewcommand{\rhrbSRRfFullMidLong}{1.83}  % Sharpe ratio (D19 pooled), sign strategy at mid-quotes, Random forest, full predictor set, W = 3,500
\providecommand{\rhrbSRRfFullBaLong}{}\renewcommand{\rhrbSRRfFullBaLong}{1.47}  % Sharpe ratio (D19 pooled), sign strategy at bid/ask prices, Random forest, full predictor set, W = 3,500
\providecommand{\rhrbPctRfFullBar}{}\renewcommand{\rhrbPctRfFullBar}{0.0}  % QLIKE change, %, W = 3,500 vs 3,000, every bar, Random forest, full predictor set
\providecommand{\rhrbDMRfFullBar}{}\renewcommand{\rhrbDMRfFullBar}{0.02}  % DM, W = 3,500 vs 3,000 (positive: 3,500 has the higher loss), every bar, Random forest, full predictor set
\providecommand{\rhrbPRfFullBar}{}\renewcommand{\rhrbPRfFullBar}{\ensuremath{p = 0.99}}  % ... its p statement
\providecommand{\rhrbPctRfFullClose}{}\renewcommand{\rhrbPctRfFullClose}{\ensuremath{-}0.5}  % QLIKE change, %, W = 3,500 vs 3,000, 16:00 bar, Random forest, full predictor set
\providecommand{\rhrbDMRfFullClose}{}\renewcommand{\rhrbDMRfFullClose}{\ensuremath{-}1.43}  % DM, W = 3,500 vs 3,000 (positive: 3,500 has the higher loss), 16:00 bar, Random forest, full predictor set
\providecommand{\rhrbPRfFullClose}{}\renewcommand{\rhrbPRfFullClose}{\ensuremath{p = 0.15}}  % ... its p statement
\providecommand{\rhrbDSRRfFullMid}{}\renewcommand{\rhrbDSRRfFullMid}{\ensuremath{-}0.05}  % Sharpe ratio change, W = 3,500 minus 3,000, mid-quotes, Random forest, full predictor set
\providecommand{\rhrbPSRRfFullMid}{}\renewcommand{\rhrbPSRRfFullMid}{\ensuremath{p = 0.58}}  % ... its p statement (Ledoit-Wolf, D19; was the D6 bootstrap)
\providecommand{\rhrbDSRRfFullBa}{}\renewcommand{\rhrbDSRRfFullBa}{\ensuremath{-}0.05}  % Sharpe ratio change, W = 3,500 minus 3,000, bid/ask prices, Random forest, full predictor set
\providecommand{\rhrbPSRRfFullBa}{}\renewcommand{\rhrbPSRRfFullBa}{\ensuremath{p = 0.59}}  % ... its p statement (Ledoit-Wolf, D19; was the D6 bootstrap)
\providecommand{\rhrbPctFRidgeBarMain}{}\renewcommand{\rhrbPctFRidgeBarMain}{\ensuremath{-}5.4}  % QLIKE, full set vs HAR benchmark, %, every bar, Ridge, W = 3,000
\providecommand{\rhrbDMFRidgeBarMain}{}\renewcommand{\rhrbDMFRidgeBarMain}{\ensuremath{-}5.63}  % DM, full set vs HAR benchmark (negative: full more accurate), every bar, Ridge, W = 3,000
\providecommand{\rhrbPFRidgeBarMain}{}\renewcommand{\rhrbPFRidgeBarMain}{\ensuremath{p < 0.001}}  % ... its p statement
\providecommand{\rhrbPctFRidgeCloseMain}{}\renewcommand{\rhrbPctFRidgeCloseMain}{0.8}  % QLIKE, full set vs HAR benchmark, %, 16:00 bar, Ridge, W = 3,000
\providecommand{\rhrbDMFRidgeCloseMain}{}\renewcommand{\rhrbDMFRidgeCloseMain}{0.31}  % DM, full set vs HAR benchmark (negative: full more accurate), 16:00 bar, Ridge, W = 3,000
\providecommand{\rhrbPFRidgeCloseMain}{}\renewcommand{\rhrbPFRidgeCloseMain}{\ensuremath{p = 0.76}}  % ... its p statement
\providecommand{\rhrbDSRFRidgeMidMain}{}\renewcommand{\rhrbDSRFRidgeMidMain}{0.82}  % Sharpe ratio, full set minus HAR benchmark, mid-quotes, Ridge, W = 3,000
\providecommand{\rhrbPSRFRidgeMidMain}{}\renewcommand{\rhrbPSRFRidgeMidMain}{\ensuremath{p = 0.11}}  % ... its p statement (Ledoit-Wolf, D19; was the D6 bootstrap)
\providecommand{\rhrbDSRFRidgeBaMain}{}\renewcommand{\rhrbDSRFRidgeBaMain}{0.81}  % Sharpe ratio, full set minus HAR benchmark, bid/ask prices, Ridge, W = 3,000
\providecommand{\rhrbPSRFRidgeBaMain}{}\renewcommand{\rhrbPSRFRidgeBaMain}{\ensuremath{p = 0.11}}  % ... its p statement (Ledoit-Wolf, D19; was the D6 bootstrap)
\providecommand{\rhrbPctFRidgeBarLong}{}\renewcommand{\rhrbPctFRidgeBarLong}{\ensuremath{-}6.3}  % QLIKE, full set vs HAR benchmark, %, every bar, Ridge, W = 3,500
\providecommand{\rhrbDMFRidgeBarLong}{}\renewcommand{\rhrbDMFRidgeBarLong}{\ensuremath{-}7.06}  % DM, full set vs HAR benchmark (negative: full more accurate), every bar, Ridge, W = 3,500
\providecommand{\rhrbPFRidgeBarLong}{}\renewcommand{\rhrbPFRidgeBarLong}{\ensuremath{p < 0.001}}  % ... its p statement
\providecommand{\rhrbPctFRidgeCloseLong}{}\renewcommand{\rhrbPctFRidgeCloseLong}{\ensuremath{-}3.1}  % QLIKE, full set vs HAR benchmark, %, 16:00 bar, Ridge, W = 3,500
\providecommand{\rhrbDMFRidgeCloseLong}{}\renewcommand{\rhrbDMFRidgeCloseLong}{\ensuremath{-}1.55}  % DM, full set vs HAR benchmark (negative: full more accurate), 16:00 bar, Ridge, W = 3,500
\providecommand{\rhrbPFRidgeCloseLong}{}\renewcommand{\rhrbPFRidgeCloseLong}{\ensuremath{p = 0.12}}  % ... its p statement
\providecommand{\rhrbDSRFRidgeMidLong}{}\renewcommand{\rhrbDSRFRidgeMidLong}{0.83}  % Sharpe ratio, full set minus HAR benchmark, mid-quotes, Ridge, W = 3,500
\providecommand{\rhrbPSRFRidgeMidLong}{}\renewcommand{\rhrbPSRFRidgeMidLong}{\ensuremath{p = 0.069}}  % ... its p statement (Ledoit-Wolf, D19; was the D6 bootstrap)
\providecommand{\rhrbDSRFRidgeBaLong}{}\renewcommand{\rhrbDSRFRidgeBaLong}{0.83}  % Sharpe ratio, full set minus HAR benchmark, bid/ask prices, Ridge, W = 3,500
\providecommand{\rhrbPSRFRidgeBaLong}{}\renewcommand{\rhrbPSRFRidgeBaLong}{\ensuremath{p = 0.068}}  % ... its p statement (Ledoit-Wolf, D19; was the D6 bootstrap)
\providecommand{\rhrbPctFLassoBarMain}{}\renewcommand{\rhrbPctFLassoBarMain}{\ensuremath{-}6.2}  % QLIKE, full set vs HAR benchmark, %, every bar, Lasso, W = 3,000
\providecommand{\rhrbDMFLassoBarMain}{}\renewcommand{\rhrbDMFLassoBarMain}{\ensuremath{-}4.72}  % DM, full set vs HAR benchmark (negative: full more accurate), every bar, Lasso, W = 3,000
\providecommand{\rhrbPFLassoBarMain}{}\renewcommand{\rhrbPFLassoBarMain}{\ensuremath{p < 0.001}}  % ... its p statement
\providecommand{\rhrbPctFLassoCloseMain}{}\renewcommand{\rhrbPctFLassoCloseMain}{3.3}  % QLIKE, full set vs HAR benchmark, %, 16:00 bar, Lasso, W = 3,000
\providecommand{\rhrbDMFLassoCloseMain}{}\renewcommand{\rhrbDMFLassoCloseMain}{1.04}  % DM, full set vs HAR benchmark (negative: full more accurate), 16:00 bar, Lasso, W = 3,000
\providecommand{\rhrbPFLassoCloseMain}{}\renewcommand{\rhrbPFLassoCloseMain}{\ensuremath{p = 0.30}}  % ... its p statement
\providecommand{\rhrbDSRFLassoMidMain}{}\renewcommand{\rhrbDSRFLassoMidMain}{0.77}  % Sharpe ratio, full set minus HAR benchmark, mid-quotes, Lasso, W = 3,000
\providecommand{\rhrbPSRFLassoMidMain}{}\renewcommand{\rhrbPSRFLassoMidMain}{\ensuremath{p = 0.12}}  % ... its p statement (Ledoit-Wolf, D19; was the D6 bootstrap)
\providecommand{\rhrbDSRFLassoBaMain}{}\renewcommand{\rhrbDSRFLassoBaMain}{0.77}  % Sharpe ratio, full set minus HAR benchmark, bid/ask prices, Lasso, W = 3,000
\providecommand{\rhrbPSRFLassoBaMain}{}\renewcommand{\rhrbPSRFLassoBaMain}{\ensuremath{p = 0.12}}  % ... its p statement (Ledoit-Wolf, D19; was the D6 bootstrap)
\providecommand{\rhrbPctFLassoBarLong}{}\renewcommand{\rhrbPctFLassoBarLong}{\ensuremath{-}7.6}  % QLIKE, full set vs HAR benchmark, %, every bar, Lasso, W = 3,500
\providecommand{\rhrbDMFLassoBarLong}{}\renewcommand{\rhrbDMFLassoBarLong}{\ensuremath{-}6.39}  % DM, full set vs HAR benchmark (negative: full more accurate), every bar, Lasso, W = 3,500
\providecommand{\rhrbPFLassoBarLong}{}\renewcommand{\rhrbPFLassoBarLong}{\ensuremath{p < 0.001}}  % ... its p statement
\providecommand{\rhrbPctFLassoCloseLong}{}\renewcommand{\rhrbPctFLassoCloseLong}{\ensuremath{-}1.6}  % QLIKE, full set vs HAR benchmark, %, 16:00 bar, Lasso, W = 3,500
\providecommand{\rhrbDMFLassoCloseLong}{}\renewcommand{\rhrbDMFLassoCloseLong}{\ensuremath{-}0.69}  % DM, full set vs HAR benchmark (negative: full more accurate), 16:00 bar, Lasso, W = 3,500
\providecommand{\rhrbPFLassoCloseLong}{}\renewcommand{\rhrbPFLassoCloseLong}{\ensuremath{p = 0.49}}  % ... its p statement
\providecommand{\rhrbDSRFLassoMidLong}{}\renewcommand{\rhrbDSRFLassoMidLong}{0.79}  % Sharpe ratio, full set minus HAR benchmark, mid-quotes, Lasso, W = 3,500
\providecommand{\rhrbPSRFLassoMidLong}{}\renewcommand{\rhrbPSRFLassoMidLong}{\ensuremath{p = 0.1006}}  % ... its p statement (Ledoit-Wolf, D19; was the D6 bootstrap)
\providecommand{\rhrbDSRFLassoBaLong}{}\renewcommand{\rhrbDSRFLassoBaLong}{0.79}  % Sharpe ratio, full set minus HAR benchmark, bid/ask prices, Lasso, W = 3,500
\providecommand{\rhrbPSRFLassoBaLong}{}\renewcommand{\rhrbPSRFLassoBaLong}{\ensuremath{p = 0.1008}}  % ... its p statement (Ledoit-Wolf, D19; was the D6 bootstrap)
\providecommand{\rhrbPctFEnetBarMain}{}\renewcommand{\rhrbPctFEnetBarMain}{\ensuremath{-}6.0}  % QLIKE, full set vs HAR benchmark, %, every bar, Elastic net, W = 3,000
\providecommand{\rhrbDMFEnetBarMain}{}\renewcommand{\rhrbDMFEnetBarMain}{\ensuremath{-}6.98}  % DM, full set vs HAR benchmark (negative: full more accurate), every bar, Elastic net, W = 3,000
\providecommand{\rhrbPFEnetBarMain}{}\renewcommand{\rhrbPFEnetBarMain}{\ensuremath{p < 0.001}}  % ... its p statement
\providecommand{\rhrbPctFEnetCloseMain}{}\renewcommand{\rhrbPctFEnetCloseMain}{\ensuremath{-}0.5}  % QLIKE, full set vs HAR benchmark, %, 16:00 bar, Elastic net, W = 3,000
\providecommand{\rhrbDMFEnetCloseMain}{}\renewcommand{\rhrbDMFEnetCloseMain}{\ensuremath{-}0.21}  % DM, full set vs HAR benchmark (negative: full more accurate), 16:00 bar, Elastic net, W = 3,000
\providecommand{\rhrbPFEnetCloseMain}{}\renewcommand{\rhrbPFEnetCloseMain}{\ensuremath{p = 0.83}}  % ... its p statement
\providecommand{\rhrbDSRFEnetMidMain}{}\renewcommand{\rhrbDSRFEnetMidMain}{0.90}  % Sharpe ratio, full set minus HAR benchmark, mid-quotes, Elastic net, W = 3,000
\providecommand{\rhrbPSRFEnetMidMain}{}\renewcommand{\rhrbPSRFEnetMidMain}{\ensuremath{p = 0.077}}  % ... its p statement (Ledoit-Wolf, D19; was the D6 bootstrap)
\providecommand{\rhrbDSRFEnetBaMain}{}\renewcommand{\rhrbDSRFEnetBaMain}{0.90}  % Sharpe ratio, full set minus HAR benchmark, bid/ask prices, Elastic net, W = 3,000
\providecommand{\rhrbPSRFEnetBaMain}{}\renewcommand{\rhrbPSRFEnetBaMain}{\ensuremath{p = 0.076}}  % ... its p statement (Ledoit-Wolf, D19; was the D6 bootstrap)
\providecommand{\rhrbPctFEnetBarLong}{}\renewcommand{\rhrbPctFEnetBarLong}{\ensuremath{-}6.7}  % QLIKE, full set vs HAR benchmark, %, every bar, Elastic net, W = 3,500
\providecommand{\rhrbDMFEnetBarLong}{}\renewcommand{\rhrbDMFEnetBarLong}{\ensuremath{-}8.45}  % DM, full set vs HAR benchmark (negative: full more accurate), every bar, Elastic net, W = 3,500
\providecommand{\rhrbPFEnetBarLong}{}\renewcommand{\rhrbPFEnetBarLong}{\ensuremath{p < 0.001}}  % ... its p statement
\providecommand{\rhrbPctFEnetCloseLong}{}\renewcommand{\rhrbPctFEnetCloseLong}{\ensuremath{-}3.9}  % QLIKE, full set vs HAR benchmark, %, 16:00 bar, Elastic net, W = 3,500
\providecommand{\rhrbDMFEnetCloseLong}{}\renewcommand{\rhrbDMFEnetCloseLong}{\ensuremath{-}2.15}  % DM, full set vs HAR benchmark (negative: full more accurate), 16:00 bar, Elastic net, W = 3,500
\providecommand{\rhrbPFEnetCloseLong}{}\renewcommand{\rhrbPFEnetCloseLong}{\ensuremath{p = 0.032}}  % ... its p statement
\providecommand{\rhrbDSRFEnetMidLong}{}\renewcommand{\rhrbDSRFEnetMidLong}{1.00}  % Sharpe ratio, full set minus HAR benchmark, mid-quotes, Elastic net, W = 3,500
\providecommand{\rhrbPSRFEnetMidLong}{}\renewcommand{\rhrbPSRFEnetMidLong}{\ensuremath{p = 0.027}}  % ... its p statement (Ledoit-Wolf, D19; was the D6 bootstrap)
\providecommand{\rhrbDSRFEnetBaLong}{}\renewcommand{\rhrbDSRFEnetBaLong}{0.99}  % Sharpe ratio, full set minus HAR benchmark, bid/ask prices, Elastic net, W = 3,500
\providecommand{\rhrbPSRFEnetBaLong}{}\renewcommand{\rhrbPSRFEnetBaLong}{\ensuremath{p = 0.027}}  % ... its p statement (Ledoit-Wolf, D19; was the D6 bootstrap)
\providecommand{\rhrbPctFLgbmBarMain}{}\renewcommand{\rhrbPctFLgbmBarMain}{\ensuremath{-}14.0}  % QLIKE, full set vs HAR benchmark, %, every bar, LightGBM, W = 3,000
\providecommand{\rhrbDMFLgbmBarMain}{}\renewcommand{\rhrbDMFLgbmBarMain}{\ensuremath{-}10.32}  % DM, full set vs HAR benchmark (negative: full more accurate), every bar, LightGBM, W = 3,000
\providecommand{\rhrbPFLgbmBarMain}{}\renewcommand{\rhrbPFLgbmBarMain}{\ensuremath{p < 0.001}}  % ... its p statement
\providecommand{\rhrbPctFLgbmCloseMain}{}\renewcommand{\rhrbPctFLgbmCloseMain}{\ensuremath{-}2.7}  % QLIKE, full set vs HAR benchmark, %, 16:00 bar, LightGBM, W = 3,000
\providecommand{\rhrbDMFLgbmCloseMain}{}\renewcommand{\rhrbDMFLgbmCloseMain}{\ensuremath{-}1.81}  % DM, full set vs HAR benchmark (negative: full more accurate), 16:00 bar, LightGBM, W = 3,000
\providecommand{\rhrbPFLgbmCloseMain}{}\renewcommand{\rhrbPFLgbmCloseMain}{\ensuremath{p = 0.071}}  % ... its p statement
\providecommand{\rhrbDSRFLgbmMidMain}{}\renewcommand{\rhrbDSRFLgbmMidMain}{0.93}  % Sharpe ratio, full set minus HAR benchmark, mid-quotes, LightGBM, W = 3,000
\providecommand{\rhrbPSRFLgbmMidMain}{}\renewcommand{\rhrbPSRFLgbmMidMain}{\ensuremath{p = 0.040}}  % ... its p statement (Ledoit-Wolf, D19; was the D6 bootstrap)
\providecommand{\rhrbDSRFLgbmBaMain}{}\renewcommand{\rhrbDSRFLgbmBaMain}{0.93}  % Sharpe ratio, full set minus HAR benchmark, bid/ask prices, LightGBM, W = 3,000
\providecommand{\rhrbPSRFLgbmBaMain}{}\renewcommand{\rhrbPSRFLgbmBaMain}{\ensuremath{p = 0.042}}  % ... its p statement (Ledoit-Wolf, D19; was the D6 bootstrap)
\providecommand{\rhrbPctFLgbmBarLong}{}\renewcommand{\rhrbPctFLgbmBarLong}{\ensuremath{-}14.0}  % QLIKE, full set vs HAR benchmark, %, every bar, LightGBM, W = 3,500
\providecommand{\rhrbDMFLgbmBarLong}{}\renewcommand{\rhrbDMFLgbmBarLong}{\ensuremath{-}10.02}  % DM, full set vs HAR benchmark (negative: full more accurate), every bar, LightGBM, W = 3,500
\providecommand{\rhrbPFLgbmBarLong}{}\renewcommand{\rhrbPFLgbmBarLong}{\ensuremath{p < 0.001}}  % ... its p statement
\providecommand{\rhrbPctFLgbmCloseLong}{}\renewcommand{\rhrbPctFLgbmCloseLong}{\ensuremath{-}2.4}  % QLIKE, full set vs HAR benchmark, %, 16:00 bar, LightGBM, W = 3,500
\providecommand{\rhrbDMFLgbmCloseLong}{}\renewcommand{\rhrbDMFLgbmCloseLong}{\ensuremath{-}1.62}  % DM, full set vs HAR benchmark (negative: full more accurate), 16:00 bar, LightGBM, W = 3,500
\providecommand{\rhrbPFLgbmCloseLong}{}\renewcommand{\rhrbPFLgbmCloseLong}{\ensuremath{p = 0.11}}  % ... its p statement
\providecommand{\rhrbDSRFLgbmMidLong}{}\renewcommand{\rhrbDSRFLgbmMidLong}{0.95}  % Sharpe ratio, full set minus HAR benchmark, mid-quotes, LightGBM, W = 3,500
\providecommand{\rhrbPSRFLgbmMidLong}{}\renewcommand{\rhrbPSRFLgbmMidLong}{\ensuremath{p = 0.035}}  % ... its p statement (Ledoit-Wolf, D19; was the D6 bootstrap)
\providecommand{\rhrbDSRFLgbmBaLong}{}\renewcommand{\rhrbDSRFLgbmBaLong}{0.94}  % Sharpe ratio, full set minus HAR benchmark, bid/ask prices, LightGBM, W = 3,500
\providecommand{\rhrbPSRFLgbmBaLong}{}\renewcommand{\rhrbPSRFLgbmBaLong}{\ensuremath{p = 0.036}}  % ... its p statement (Ledoit-Wolf, D19; was the D6 bootstrap)
\providecommand{\rhrbPctFXgbBarMain}{}\renewcommand{\rhrbPctFXgbBarMain}{\ensuremath{-}13.8}  % QLIKE, full set vs HAR benchmark, %, every bar, XGBoost, W = 3,000
\providecommand{\rhrbDMFXgbBarMain}{}\renewcommand{\rhrbDMFXgbBarMain}{\ensuremath{-}9.94}  % DM, full set vs HAR benchmark (negative: full more accurate), every bar, XGBoost, W = 3,000
\providecommand{\rhrbPFXgbBarMain}{}\renewcommand{\rhrbPFXgbBarMain}{\ensuremath{p < 0.001}}  % ... its p statement
\providecommand{\rhrbPctFXgbCloseMain}{}\renewcommand{\rhrbPctFXgbCloseMain}{\ensuremath{-}3.4}  % QLIKE, full set vs HAR benchmark, %, 16:00 bar, XGBoost, W = 3,000
\providecommand{\rhrbDMFXgbCloseMain}{}\renewcommand{\rhrbDMFXgbCloseMain}{\ensuremath{-}2.18}  % DM, full set vs HAR benchmark (negative: full more accurate), 16:00 bar, XGBoost, W = 3,000
\providecommand{\rhrbPFXgbCloseMain}{}\renewcommand{\rhrbPFXgbCloseMain}{\ensuremath{p = 0.030}}  % ... its p statement
\providecommand{\rhrbDSRFXgbMidMain}{}\renewcommand{\rhrbDSRFXgbMidMain}{0.49}  % Sharpe ratio, full set minus HAR benchmark, mid-quotes, XGBoost, W = 3,000
\providecommand{\rhrbPSRFXgbMidMain}{}\renewcommand{\rhrbPSRFXgbMidMain}{\ensuremath{p = 0.27}}  % ... its p statement (Ledoit-Wolf, D19; was the D6 bootstrap)
\providecommand{\rhrbDSRFXgbBaMain}{}\renewcommand{\rhrbDSRFXgbBaMain}{0.49}  % Sharpe ratio, full set minus HAR benchmark, bid/ask prices, XGBoost, W = 3,000
\providecommand{\rhrbPSRFXgbBaMain}{}\renewcommand{\rhrbPSRFXgbBaMain}{\ensuremath{p = 0.27}}  % ... its p statement (Ledoit-Wolf, D19; was the D6 bootstrap)
\providecommand{\rhrbPctFXgbBarLong}{}\renewcommand{\rhrbPctFXgbBarLong}{\ensuremath{-}13.5}  % QLIKE, full set vs HAR benchmark, %, every bar, XGBoost, W = 3,500
\providecommand{\rhrbDMFXgbBarLong}{}\renewcommand{\rhrbDMFXgbBarLong}{\ensuremath{-}9.97}  % DM, full set vs HAR benchmark (negative: full more accurate), every bar, XGBoost, W = 3,500
\providecommand{\rhrbPFXgbBarLong}{}\renewcommand{\rhrbPFXgbBarLong}{\ensuremath{p < 0.001}}  % ... its p statement
\providecommand{\rhrbPctFXgbCloseLong}{}\renewcommand{\rhrbPctFXgbCloseLong}{\ensuremath{-}3.2}  % QLIKE, full set vs HAR benchmark, %, 16:00 bar, XGBoost, W = 3,500
\providecommand{\rhrbDMFXgbCloseLong}{}\renewcommand{\rhrbDMFXgbCloseLong}{\ensuremath{-}2.19}  % DM, full set vs HAR benchmark (negative: full more accurate), 16:00 bar, XGBoost, W = 3,500
\providecommand{\rhrbPFXgbCloseLong}{}\renewcommand{\rhrbPFXgbCloseLong}{\ensuremath{p = 0.029}}  % ... its p statement
\providecommand{\rhrbDSRFXgbMidLong}{}\renewcommand{\rhrbDSRFXgbMidLong}{0.57}  % Sharpe ratio, full set minus HAR benchmark, mid-quotes, XGBoost, W = 3,500
\providecommand{\rhrbPSRFXgbMidLong}{}\renewcommand{\rhrbPSRFXgbMidLong}{\ensuremath{p = 0.15}}  % ... its p statement (Ledoit-Wolf, D19; was the D6 bootstrap)
\providecommand{\rhrbDSRFXgbBaLong}{}\renewcommand{\rhrbDSRFXgbBaLong}{0.56}  % Sharpe ratio, full set minus HAR benchmark, bid/ask prices, XGBoost, W = 3,500
\providecommand{\rhrbPSRFXgbBaLong}{}\renewcommand{\rhrbPSRFXgbBaLong}{\ensuremath{p = 0.15}}  % ... its p statement (Ledoit-Wolf, D19; was the D6 bootstrap)
\providecommand{\rhrbPctFRfBarMain}{}\renewcommand{\rhrbPctFRfBarMain}{\ensuremath{-}11.1}  % QLIKE, full set vs HAR benchmark, %, every bar, Random forest, W = 3,000
\providecommand{\rhrbDMFRfBarMain}{}\renewcommand{\rhrbDMFRfBarMain}{\ensuremath{-}9.85}  % DM, full set vs HAR benchmark (negative: full more accurate), every bar, Random forest, W = 3,000
\providecommand{\rhrbPFRfBarMain}{}\renewcommand{\rhrbPFRfBarMain}{\ensuremath{p < 0.001}}  % ... its p statement
\providecommand{\rhrbPctFRfCloseMain}{}\renewcommand{\rhrbPctFRfCloseMain}{\ensuremath{-}1.4}  % QLIKE, full set vs HAR benchmark, %, 16:00 bar, Random forest, W = 3,000
\providecommand{\rhrbDMFRfCloseMain}{}\renewcommand{\rhrbDMFRfCloseMain}{\ensuremath{-}0.92}  % DM, full set vs HAR benchmark (negative: full more accurate), 16:00 bar, Random forest, W = 3,000
\providecommand{\rhrbPFRfCloseMain}{}\renewcommand{\rhrbPFRfCloseMain}{\ensuremath{p = 0.36}}  % ... its p statement
\providecommand{\rhrbDSRFRfMidMain}{}\renewcommand{\rhrbDSRFRfMidMain}{1.56}  % Sharpe ratio, full set minus HAR benchmark, mid-quotes, Random forest, W = 3,000
\providecommand{\rhrbPSRFRfMidMain}{}\renewcommand{\rhrbPSRFRfMidMain}{\ensuremath{p = 0.020}}  % ... its p statement (Ledoit-Wolf, D19; was the D6 bootstrap)
\providecommand{\rhrbDSRFRfBaMain}{}\renewcommand{\rhrbDSRFRfBaMain}{1.55}  % Sharpe ratio, full set minus HAR benchmark, bid/ask prices, Random forest, W = 3,000
\providecommand{\rhrbPSRFRfBaMain}{}\renewcommand{\rhrbPSRFRfBaMain}{\ensuremath{p = 0.023}}  % ... its p statement (Ledoit-Wolf, D19; was the D6 bootstrap)
\providecommand{\rhrbPctFRfBarLong}{}\renewcommand{\rhrbPctFRfBarLong}{\ensuremath{-}11.2}  % QLIKE, full set vs HAR benchmark, %, every bar, Random forest, W = 3,500
\providecommand{\rhrbDMFRfBarLong}{}\renewcommand{\rhrbDMFRfBarLong}{\ensuremath{-}9.40}  % DM, full set vs HAR benchmark (negative: full more accurate), every bar, Random forest, W = 3,500
\providecommand{\rhrbPFRfBarLong}{}\renewcommand{\rhrbPFRfBarLong}{\ensuremath{p < 0.001}}  % ... its p statement
\providecommand{\rhrbPctFRfCloseLong}{}\renewcommand{\rhrbPctFRfCloseLong}{\ensuremath{-}1.8}  % QLIKE, full set vs HAR benchmark, %, 16:00 bar, Random forest, W = 3,500
\providecommand{\rhrbDMFRfCloseLong}{}\renewcommand{\rhrbDMFRfCloseLong}{\ensuremath{-}1.15}  % DM, full set vs HAR benchmark (negative: full more accurate), 16:00 bar, Random forest, W = 3,500
\providecommand{\rhrbPFRfCloseLong}{}\renewcommand{\rhrbPFRfCloseLong}{\ensuremath{p = 0.25}}  % ... its p statement
\providecommand{\rhrbDSRFRfMidLong}{}\renewcommand{\rhrbDSRFRfMidLong}{1.03}  % Sharpe ratio, full set minus HAR benchmark, mid-quotes, Random forest, W = 3,500
\providecommand{\rhrbPSRFRfMidLong}{}\renewcommand{\rhrbPSRFRfMidLong}{\ensuremath{p = 0.0098}}  % ... its p statement (Ledoit-Wolf, D19; was the D6 bootstrap)
\providecommand{\rhrbDSRFRfBaLong}{}\renewcommand{\rhrbDSRFRfBaLong}{1.02}  % Sharpe ratio, full set minus HAR benchmark, bid/ask prices, Random forest, W = 3,500
\providecommand{\rhrbPSRFRfBaLong}{}\renewcommand{\rhrbPSRFRfBaLong}{\ensuremath{p = 0.0100}}  % ... its p statement (Ledoit-Wolf, D19; was the D6 bootstrap)
\providecommand{\rhrbNArms}{}\renewcommand{\rhrbNArms}{12}  % pooled forecasts compared at both windows (finished; OLS reference row out)
\providecommand{\rhrbNArmsAll}{}\renewcommand{\rhrbNArmsAll}{12}  % pooled forecasts in all (finished or pending)
\providecommand{\rhrbNArmsPending}{}\renewcommand{\rhrbNArmsPending}{0}  % pooled forecasts pending at a window
\providecommand{\rhrbNSigBar}{}\renewcommand{\rhrbNSigBar}{5}  % of rhrbNArms: QLIKE differs between the windows at 5 % (DM), every bar
\providecommand{\rhrbNLowerBar}{}\renewcommand{\rhrbNLowerBar}{3}  % of rhrbNArms: lower QLIKE at 3,500 (DM < 0), every bar
\providecommand{\rhrbNHigherSigBar}{}\renewcommand{\rhrbNHigherSigBar}{3}  % of rhrbNArms: significantly higher QLIKE at 3,500 (DM > 0, 5 %), every bar
\providecommand{\rhrbNLowerSigBar}{}\renewcommand{\rhrbNLowerSigBar}{2}  % of rhrbNArms: significantly lower QLIKE at 3,500 (DM < 0, 5 %), every bar
\providecommand{\rhrbMaxAbsPctBar}{}\renewcommand{\rhrbMaxAbsPctBar}{1.3}  % largest absolute QLIKE change between the windows, %, every bar (finished arms)
\providecommand{\rhrbNSigClose}{}\renewcommand{\rhrbNSigClose}{3}  % of rhrbNArms: QLIKE differs between the windows at 5 % (DM), 16:00 bar
\providecommand{\rhrbNLowerClose}{}\renewcommand{\rhrbNLowerClose}{8}  % of rhrbNArms: lower QLIKE at 3,500 (DM < 0), 16:00 bar
\providecommand{\rhrbNHigherSigClose}{}\renewcommand{\rhrbNHigherSigClose}{0}  % of rhrbNArms: significantly higher QLIKE at 3,500 (DM > 0, 5 %), 16:00 bar
\providecommand{\rhrbNLowerSigClose}{}\renewcommand{\rhrbNLowerSigClose}{3}  % of rhrbNArms: significantly lower QLIKE at 3,500 (DM < 0, 5 %), 16:00 bar
\providecommand{\rhrbMaxAbsPctClose}{}\renewcommand{\rhrbMaxAbsPctClose}{4.4}  % largest absolute QLIKE change between the windows, %, 16:00 bar (finished arms)
\providecommand{\rhrbNModelsMain}{}\renewcommand{\rhrbNModelsMain}{6}  % model classes with both predictor sets finished, W = 3,000
\providecommand{\rhrbNFullLowerSigBarMain}{}\renewcommand{\rhrbNFullLowerSigBarMain}{6}  % of rhrbNModelsMain: full set significantly more accurate than HAR (DM < 0, 5 %), every bar
\providecommand{\rhrbNFullLowerSigCloseMain}{}\renewcommand{\rhrbNFullLowerSigCloseMain}{1}  % of rhrbNModelsMain: full set significantly more accurate than HAR (DM < 0, 5 %), 16:00 bar
\providecommand{\rhrbNDSRFPosMidMain}{}\renewcommand{\rhrbNDSRFPosMidMain}{6}  % of rhrbNModelsMain: full-set Sharpe ratio above HAR, mid-quotes
\providecommand{\rhrbNDSRFSigMidMain}{}\renewcommand{\rhrbNDSRFSigMidMain}{2}  % of rhrbNModelsMain: ... and Ledoit-Wolf p < 0.05 (D19), mid-quotes
\providecommand{\rhrbNDSRFPosBaMain}{}\renewcommand{\rhrbNDSRFPosBaMain}{6}  % of rhrbNModelsMain: full-set Sharpe ratio above HAR, bid/ask prices
\providecommand{\rhrbNDSRFSigBaMain}{}\renewcommand{\rhrbNDSRFSigBaMain}{2}  % of rhrbNModelsMain: ... and Ledoit-Wolf p < 0.05 (D19), bid/ask prices
\providecommand{\rhrbNModelsLong}{}\renewcommand{\rhrbNModelsLong}{6}  % model classes with both predictor sets finished, W = 3,500
\providecommand{\rhrbNFullLowerSigBarLong}{}\renewcommand{\rhrbNFullLowerSigBarLong}{6}  % of rhrbNModelsLong: full set significantly more accurate than HAR (DM < 0, 5 %), every bar
\providecommand{\rhrbNFullLowerSigCloseLong}{}\renewcommand{\rhrbNFullLowerSigCloseLong}{2}  % of rhrbNModelsLong: full set significantly more accurate than HAR (DM < 0, 5 %), 16:00 bar
\providecommand{\rhrbNDSRFPosMidLong}{}\renewcommand{\rhrbNDSRFPosMidLong}{6}  % of rhrbNModelsLong: full-set Sharpe ratio above HAR, mid-quotes
\providecommand{\rhrbNDSRFSigMidLong}{}\renewcommand{\rhrbNDSRFSigMidLong}{3}  % of rhrbNModelsLong: ... and Ledoit-Wolf p < 0.05 (D19), mid-quotes
\providecommand{\rhrbNDSRFPosBaLong}{}\renewcommand{\rhrbNDSRFPosBaLong}{6}  % of rhrbNModelsLong: full-set Sharpe ratio above HAR, bid/ask prices
\providecommand{\rhrbNDSRFSigBaLong}{}\renewcommand{\rhrbNDSRFSigBaLong}{3}  % of rhrbNModelsLong: ... and Ledoit-Wolf p < 0.05 (D19), bid/ask prices
\providecommand{\rhrbNModelsAll}{}\renewcommand{\rhrbNModelsAll}{6}  % model classes in all
\providecommand{\rhrbNSRSigMid}{}\renewcommand{\rhrbNSRSigMid}{0}  % of rhrbNArms: Sharpe ratio differs between the windows (Ledoit-Wolf p < 0.05, D19), mid-quotes
\providecommand{\rhrbMaxAbsDSRMid}{}\renewcommand{\rhrbMaxAbsDSRMid}{0.48}  % largest absolute Sharpe-ratio change between the windows, mid-quotes (finished arms)
\providecommand{\rhrbMinPSRMid}{}\renewcommand{\rhrbMinPSRMid}{\ensuremath{p = 0.40}}  % smallest Ledoit-Wolf p (D19) of a Sharpe-ratio change between the windows, mid-quotes
\providecommand{\rhrbNSRSigBa}{}\renewcommand{\rhrbNSRSigBa}{0}  % of rhrbNArms: Sharpe ratio differs between the windows (Ledoit-Wolf p < 0.05, D19), bid/ask prices
\providecommand{\rhrbMaxAbsDSRBa}{}\renewcommand{\rhrbMaxAbsDSRBa}{0.48}  % largest absolute Sharpe-ratio change between the windows, bid/ask prices (finished arms)
\providecommand{\rhrbMinPSRBa}{}\renewcommand{\rhrbMinPSRBa}{\ensuremath{p = 0.40}}  % smallest Ledoit-Wolf p (D19) of a Sharpe-ratio change between the windows, bid/ask prices
\fi
% 2026-10-26 (coordinator): the file now exists; plain guard form, which the Overleaf flattener drops. Was:
% \ifdefined\rhholNSessions\else\IfFileExists{generated/rehaul_holiday_numbers.tex}{% AUTO-GENERATED by writeup/make_rehaul_holiday.py -- do not edit; re-run the script.
% \rhhol... macros of the NYSE-holiday disclosure (D15; Table tab:rh_holiday): results/fomcfix_2026-10-26,
% RULE.txt ADDENDUM A (cells A-D; Eff CA = C - A, Eff DB = D - B, the holiday effect without / with the FOMC fix).
% Sh... (2026-10-28, D19): the paper's pooled Sharpe ratio, recomputed from each cell's 16:00 forecasts (was D6).
% NOTE FOR OFFICIAL PUBLICATION (D15): rerun with DROP_HOLIDAY_SESSIONS on (all models) before the final version.
\providecommand{\rhholNDatesAll}{}\renewcommand{\rhholNDatesAll}{123}  % holiday dates with a shortened E-mini session in the sample
\providecommand{\rhholFirstYearAll}{}\renewcommand{\rhholFirstYearAll}{2004}  % first year of such a date
\providecommand{\rhholNSessions}{}\renewcommand{\rhholNSessions}{64}  % holiday sessions in the evaluation span 2014-01-02 .. 2024-04-30
\providecommand{\rhholNBars}{}\renewcommand{\rhholNBars}{443}  % their regular-hours bars (forecast, not scored)
\providecommand{\rhholNForecast}{}\renewcommand{\rhholNForecast}{33{,}970}  % bars forecast in the span (33,527 scored + the holiday bars)
\providecommand{\rhholNScored}{}\renewcommand{\rhholNScored}{33{,}527}  % bars scored (Table 1)
\providecommand{\rhholLastBar}{}\renewcommand{\rhholLastBar}{13:00}  % the usual last regular-hours bar end of a holiday session (ET)
\providecommand{\rhholEffCABarPpMin}{}\renewcommand{\rhholEffCABarPpMin}{+0.13}  % holiday effect C-A, linear: lowest change of pct_vs_har_bar
\providecommand{\rhholEffCABarPpMax}{}\renewcommand{\rhholEffCABarPpMax}{+0.14}  % holiday effect C-A, linear: highest change of pct_vs_har_bar
\providecommand{\rhholEffCABarPpRng}{}\renewcommand{\rhholEffCABarPpRng}{0.13--0.14}  % holiday effect C-A: range of pct_vs_har_bar
\providecommand{\rhholEffCARidgeBarPp}{}\renewcommand{\rhholEffCARidgeBarPp}{+0.14}  % holiday effect C-A, ridge: pct_vs_har_bar
\providecommand{\rhholEffCALassoBarPp}{}\renewcommand{\rhholEffCALassoBarPp}{+0.14}  % holiday effect C-A, lasso: pct_vs_har_bar
\providecommand{\rhholEffCAEnetBarPp}{}\renewcommand{\rhholEffCAEnetBarPp}{+0.13}  % holiday effect C-A, enet: pct_vs_har_bar
\providecommand{\rhholEffCAClosePpMin}{}\renewcommand{\rhholEffCAClosePpMin}{\ensuremath{-}0.98}  % holiday effect C-A, linear: lowest change of pct_vs_har_c16
\providecommand{\rhholEffCAClosePpMax}{}\renewcommand{\rhholEffCAClosePpMax}{\ensuremath{-}0.49}  % holiday effect C-A, linear: highest change of pct_vs_har_c16
\providecommand{\rhholEffCAClosePpRng}{}\renewcommand{\rhholEffCAClosePpRng}{\ensuremath{-}0.98 to \ensuremath{-}0.49}  % holiday effect C-A: range of pct_vs_har_c16
\providecommand{\rhholEffCARidgeClosePp}{}\renewcommand{\rhholEffCARidgeClosePp}{\ensuremath{-}0.98}  % holiday effect C-A, ridge: pct_vs_har_c16
\providecommand{\rhholEffCALassoClosePp}{}\renewcommand{\rhholEffCALassoClosePp}{\ensuremath{-}0.49}  % holiday effect C-A, lasso: pct_vs_har_c16
\providecommand{\rhholEffCAEnetClosePp}{}\renewcommand{\rhholEffCAEnetClosePp}{\ensuremath{-}0.75}  % holiday effect C-A, enet: pct_vs_har_c16
\providecommand{\rhholEffCAShFullMin}{}\renewcommand{\rhholEffCAShFullMin}{\ensuremath{-}0.15}  % holiday effect C-A, linear: lowest change of sr_mid_full
\providecommand{\rhholEffCAShFullMax}{}\renewcommand{\rhholEffCAShFullMax}{+0.39}  % holiday effect C-A, linear: highest change of sr_mid_full
\providecommand{\rhholEffCAShFullRng}{}\renewcommand{\rhholEffCAShFullRng}{\ensuremath{-}0.15 to +0.39}  % holiday effect C-A: range of sr_mid_full
\providecommand{\rhholEffCARidgeShFull}{}\renewcommand{\rhholEffCARidgeShFull}{+0.17}  % holiday effect C-A, ridge: sr_mid_full
\providecommand{\rhholEffCALassoShFull}{}\renewcommand{\rhholEffCALassoShFull}{+0.39}  % holiday effect C-A, lasso: sr_mid_full
\providecommand{\rhholEffCAEnetShFull}{}\renewcommand{\rhholEffCAEnetShFull}{\ensuremath{-}0.15}  % holiday effect C-A, enet: sr_mid_full
\providecommand{\rhholEffCADShMin}{}\renewcommand{\rhholEffCADShMin}{\ensuremath{-}0.18}  % holiday effect C-A, linear: lowest change of dsr_mid
\providecommand{\rhholEffCADShMax}{}\renewcommand{\rhholEffCADShMax}{+0.36}  % holiday effect C-A, linear: highest change of dsr_mid
\providecommand{\rhholEffCADShRng}{}\renewcommand{\rhholEffCADShRng}{\ensuremath{-}0.18 to +0.36}  % holiday effect C-A: range of dsr_mid
\providecommand{\rhholEffCARidgeDSh}{}\renewcommand{\rhholEffCARidgeDSh}{+0.13}  % holiday effect C-A, ridge: dsr_mid
\providecommand{\rhholEffCALassoDSh}{}\renewcommand{\rhholEffCALassoDSh}{+0.36}  % holiday effect C-A, lasso: dsr_mid
\providecommand{\rhholEffCAEnetDSh}{}\renewcommand{\rhholEffCAEnetDSh}{\ensuremath{-}0.18}  % holiday effect C-A, enet: dsr_mid
\providecommand{\rhholEffDBBarPpMin}{}\renewcommand{\rhholEffDBBarPpMin}{+0.11}  % holiday effect D-B, linear: lowest change of pct_vs_har_bar
\providecommand{\rhholEffDBBarPpMax}{}\renewcommand{\rhholEffDBBarPpMax}{+0.30}  % holiday effect D-B, linear: highest change of pct_vs_har_bar
\providecommand{\rhholEffDBBarPpRng}{}\renewcommand{\rhholEffDBBarPpRng}{0.11--0.30}  % holiday effect D-B: range of pct_vs_har_bar
\providecommand{\rhholEffDBRidgeBarPp}{}\renewcommand{\rhholEffDBRidgeBarPp}{+0.12}  % holiday effect D-B, ridge: pct_vs_har_bar
\providecommand{\rhholEffDBLassoBarPp}{}\renewcommand{\rhholEffDBLassoBarPp}{+0.30}  % holiday effect D-B, lasso: pct_vs_har_bar
\providecommand{\rhholEffDBEnetBarPp}{}\renewcommand{\rhholEffDBEnetBarPp}{+0.11}  % holiday effect D-B, enet: pct_vs_har_bar
\providecommand{\rhholEffDBClosePpMin}{}\renewcommand{\rhholEffDBClosePpMin}{\ensuremath{-}0.86}  % holiday effect D-B, linear: lowest change of pct_vs_har_c16
\providecommand{\rhholEffDBClosePpMax}{}\renewcommand{\rhholEffDBClosePpMax}{\ensuremath{-}0.20}  % holiday effect D-B, linear: highest change of pct_vs_har_c16
\providecommand{\rhholEffDBClosePpRng}{}\renewcommand{\rhholEffDBClosePpRng}{\ensuremath{-}0.86 to \ensuremath{-}0.20}  % holiday effect D-B: range of pct_vs_har_c16
\providecommand{\rhholEffDBRidgeClosePp}{}\renewcommand{\rhholEffDBRidgeClosePp}{\ensuremath{-}0.86}  % holiday effect D-B, ridge: pct_vs_har_c16
\providecommand{\rhholEffDBLassoClosePp}{}\renewcommand{\rhholEffDBLassoClosePp}{\ensuremath{-}0.20}  % holiday effect D-B, lasso: pct_vs_har_c16
\providecommand{\rhholEffDBEnetClosePp}{}\renewcommand{\rhholEffDBEnetClosePp}{\ensuremath{-}0.65}  % holiday effect D-B, enet: pct_vs_har_c16
\providecommand{\rhholEffDBShFullMin}{}\renewcommand{\rhholEffDBShFullMin}{\ensuremath{-}0.20}  % holiday effect D-B, linear: lowest change of sr_mid_full
\providecommand{\rhholEffDBShFullMax}{}\renewcommand{\rhholEffDBShFullMax}{+0.38}  % holiday effect D-B, linear: highest change of sr_mid_full
\providecommand{\rhholEffDBShFullRng}{}\renewcommand{\rhholEffDBShFullRng}{\ensuremath{-}0.20 to +0.38}  % holiday effect D-B: range of sr_mid_full
\providecommand{\rhholEffDBRidgeShFull}{}\renewcommand{\rhholEffDBRidgeShFull}{+0.20}  % holiday effect D-B, ridge: sr_mid_full
\providecommand{\rhholEffDBLassoShFull}{}\renewcommand{\rhholEffDBLassoShFull}{+0.38}  % holiday effect D-B, lasso: sr_mid_full
\providecommand{\rhholEffDBEnetShFull}{}\renewcommand{\rhholEffDBEnetShFull}{\ensuremath{-}0.20}  % holiday effect D-B, enet: sr_mid_full
\providecommand{\rhholEffDBDShMin}{}\renewcommand{\rhholEffDBDShMin}{\ensuremath{-}0.24}  % holiday effect D-B, linear: lowest change of dsr_mid
\providecommand{\rhholEffDBDShMax}{}\renewcommand{\rhholEffDBDShMax}{+0.35}  % holiday effect D-B, linear: highest change of dsr_mid
\providecommand{\rhholEffDBDShRng}{}\renewcommand{\rhholEffDBDShRng}{\ensuremath{-}0.24 to +0.35}  % holiday effect D-B: range of dsr_mid
\providecommand{\rhholEffDBRidgeDSh}{}\renewcommand{\rhholEffDBRidgeDSh}{+0.17}  % holiday effect D-B, ridge: dsr_mid
\providecommand{\rhholEffDBLassoDSh}{}\renewcommand{\rhholEffDBLassoDSh}{+0.35}  % holiday effect D-B, lasso: dsr_mid
\providecommand{\rhholEffDBEnetDSh}{}\renewcommand{\rhholEffDBEnetDSh}{\ensuremath{-}0.24}  % holiday effect D-B, enet: dsr_mid
\providecommand{\rhholRidgeAPctBar}{}\renewcommand{\rhholRidgeAPctBar}{\ensuremath{-}5.4}  % cell A, ridge: every-bar % vs HAR
\providecommand{\rhholRidgeAPBar}{}\renewcommand{\rhholRidgeAPBar}{\ensuremath{<}0.001}  % cell A, ridge: every-bar DM p
\providecommand{\rhholRidgeAPctClose}{}\renewcommand{\rhholRidgeAPctClose}{0.8}  % cell A, ridge: 16:00 % vs HAR
\providecommand{\rhholRidgeAShFull}{}\renewcommand{\rhholRidgeAShFull}{1.77}  % cell A, ridge: Sharpe ratio (D19), full set, mid
\providecommand{\rhholLassoAPctBar}{}\renewcommand{\rhholLassoAPctBar}{\ensuremath{-}6.2}  % cell A, lasso: every-bar % vs HAR
\providecommand{\rhholLassoAPBar}{}\renewcommand{\rhholLassoAPBar}{\ensuremath{<}0.001}  % cell A, lasso: every-bar DM p
\providecommand{\rhholLassoAPctClose}{}\renewcommand{\rhholLassoAPctClose}{3.3}  % cell A, lasso: 16:00 % vs HAR
\providecommand{\rhholLassoAShFull}{}\renewcommand{\rhholLassoAShFull}{1.72}  % cell A, lasso: Sharpe ratio (D19), full set, mid
\providecommand{\rhholEnetAPctBar}{}\renewcommand{\rhholEnetAPctBar}{\ensuremath{-}6.0}  % cell A, enet: every-bar % vs HAR
\providecommand{\rhholEnetAPBar}{}\renewcommand{\rhholEnetAPBar}{\ensuremath{<}0.001}  % cell A, enet: every-bar DM p
\providecommand{\rhholEnetAPctClose}{}\renewcommand{\rhholEnetAPctClose}{\ensuremath{-}0.5}  % cell A, enet: 16:00 % vs HAR
\providecommand{\rhholEnetAShFull}{}\renewcommand{\rhholEnetAShFull}{1.86}  % cell A, enet: Sharpe ratio (D19), full set, mid
\providecommand{\rhholRidgeBPctBar}{}\renewcommand{\rhholRidgeBPctBar}{\ensuremath{-}6.0}  % cell B, ridge: every-bar % vs HAR
\providecommand{\rhholRidgeBPBar}{}\renewcommand{\rhholRidgeBPBar}{\ensuremath{<}0.001}  % cell B, ridge: every-bar DM p
\providecommand{\rhholRidgeBPctClose}{}\renewcommand{\rhholRidgeBPctClose}{1.8}  % cell B, ridge: 16:00 % vs HAR
\providecommand{\rhholRidgeBShFull}{}\renewcommand{\rhholRidgeBShFull}{1.64}  % cell B, ridge: Sharpe ratio (D19), full set, mid
\providecommand{\rhholLassoBPctBar}{}\renewcommand{\rhholLassoBPctBar}{\ensuremath{-}9.0}  % cell B, lasso: every-bar % vs HAR
\providecommand{\rhholLassoBPBar}{}\renewcommand{\rhholLassoBPBar}{\ensuremath{<}0.001}  % cell B, lasso: every-bar DM p
\providecommand{\rhholLassoBPctClose}{}\renewcommand{\rhholLassoBPctClose}{4.2}  % cell B, lasso: 16:00 % vs HAR
\providecommand{\rhholLassoBShFull}{}\renewcommand{\rhholLassoBShFull}{1.73}  % cell B, lasso: Sharpe ratio (D19), full set, mid
\providecommand{\rhholEnetBPctBar}{}\renewcommand{\rhholEnetBPctBar}{\ensuremath{-}6.7}  % cell B, enet: every-bar % vs HAR
\providecommand{\rhholEnetBPBar}{}\renewcommand{\rhholEnetBPBar}{\ensuremath{<}0.001}  % cell B, enet: every-bar DM p
\providecommand{\rhholEnetBPctClose}{}\renewcommand{\rhholEnetBPctClose}{0.6}  % cell B, enet: 16:00 % vs HAR
\providecommand{\rhholEnetBShFull}{}\renewcommand{\rhholEnetBShFull}{1.85}  % cell B, enet: Sharpe ratio (D19), full set, mid
\providecommand{\rhholRidgeCPctBar}{}\renewcommand{\rhholRidgeCPctBar}{\ensuremath{-}5.2}  % cell C, ridge: every-bar % vs HAR
\providecommand{\rhholRidgeCPBar}{}\renewcommand{\rhholRidgeCPBar}{\ensuremath{<}0.001}  % cell C, ridge: every-bar DM p
\providecommand{\rhholRidgeCPctClose}{}\renewcommand{\rhholRidgeCPctClose}{\ensuremath{-}0.2}  % cell C, ridge: 16:00 % vs HAR
\providecommand{\rhholRidgeCShFull}{}\renewcommand{\rhholRidgeCShFull}{1.94}  % cell C, ridge: Sharpe ratio (D19), full set, mid
\providecommand{\rhholLassoCPctBar}{}\renewcommand{\rhholLassoCPctBar}{\ensuremath{-}6.1}  % cell C, lasso: every-bar % vs HAR
\providecommand{\rhholLassoCPBar}{}\renewcommand{\rhholLassoCPBar}{\ensuremath{<}0.001}  % cell C, lasso: every-bar DM p
\providecommand{\rhholLassoCPctClose}{}\renewcommand{\rhholLassoCPctClose}{2.8}  % cell C, lasso: 16:00 % vs HAR
\providecommand{\rhholLassoCShFull}{}\renewcommand{\rhholLassoCShFull}{2.12}  % cell C, lasso: Sharpe ratio (D19), full set, mid
\providecommand{\rhholEnetCPctBar}{}\renewcommand{\rhholEnetCPctBar}{\ensuremath{-}5.9}  % cell C, enet: every-bar % vs HAR
\providecommand{\rhholEnetCPBar}{}\renewcommand{\rhholEnetCPBar}{\ensuremath{<}0.001}  % cell C, enet: every-bar DM p
\providecommand{\rhholEnetCPctClose}{}\renewcommand{\rhholEnetCPctClose}{\ensuremath{-}1.3}  % cell C, enet: 16:00 % vs HAR
\providecommand{\rhholEnetCShFull}{}\renewcommand{\rhholEnetCShFull}{1.71}  % cell C, enet: Sharpe ratio (D19), full set, mid
\providecommand{\rhholRidgeDPctBar}{}\renewcommand{\rhholRidgeDPctBar}{\ensuremath{-}5.9}  % cell D, ridge: every-bar % vs HAR
\providecommand{\rhholRidgeDPBar}{}\renewcommand{\rhholRidgeDPBar}{\ensuremath{<}0.001}  % cell D, ridge: every-bar DM p
\providecommand{\rhholRidgeDPctClose}{}\renewcommand{\rhholRidgeDPctClose}{0.9}  % cell D, ridge: 16:00 % vs HAR
\providecommand{\rhholRidgeDShFull}{}\renewcommand{\rhholRidgeDShFull}{1.85}  % cell D, ridge: Sharpe ratio (D19), full set, mid
\providecommand{\rhholLassoDPctBar}{}\renewcommand{\rhholLassoDPctBar}{\ensuremath{-}8.7}  % cell D, lasso: every-bar % vs HAR
\providecommand{\rhholLassoDPBar}{}\renewcommand{\rhholLassoDPBar}{\ensuremath{<}0.001}  % cell D, lasso: every-bar DM p
\providecommand{\rhholLassoDPctClose}{}\renewcommand{\rhholLassoDPctClose}{4.0}  % cell D, lasso: 16:00 % vs HAR
\providecommand{\rhholLassoDShFull}{}\renewcommand{\rhholLassoDShFull}{2.11}  % cell D, lasso: Sharpe ratio (D19), full set, mid
\providecommand{\rhholEnetDPctBar}{}\renewcommand{\rhholEnetDPctBar}{\ensuremath{-}6.6}  % cell D, enet: every-bar % vs HAR
\providecommand{\rhholEnetDPBar}{}\renewcommand{\rhholEnetDPBar}{\ensuremath{<}0.001}  % cell D, enet: every-bar DM p
\providecommand{\rhholEnetDPctClose}{}\renewcommand{\rhholEnetDPctClose}{0.0}  % cell D, enet: 16:00 % vs HAR
\providecommand{\rhholEnetDShFull}{}\renewcommand{\rhholEnetDShFull}{1.65}  % cell D, enet: Sharpe ratio (D19), full set, mid
}{}\fi
\ifdefined\rhholNSessions\else% AUTO-GENERATED by writeup/make_rehaul_holiday.py -- do not edit; re-run the script.
% \rhhol... macros of the NYSE-holiday disclosure (D15; Table tab:rh_holiday): results/fomcfix_2026-10-26,
% RULE.txt ADDENDUM A (cells A-D; Eff CA = C - A, Eff DB = D - B, the holiday effect without / with the FOMC fix).
% Sh... (2026-10-28, D19): the paper's pooled Sharpe ratio, recomputed from each cell's 16:00 forecasts (was D6).
% NOTE FOR OFFICIAL PUBLICATION (D15): rerun with DROP_HOLIDAY_SESSIONS on (all models) before the final version.
\providecommand{\rhholNDatesAll}{}\renewcommand{\rhholNDatesAll}{123}  % holiday dates with a shortened E-mini session in the sample
\providecommand{\rhholFirstYearAll}{}\renewcommand{\rhholFirstYearAll}{2004}  % first year of such a date
\providecommand{\rhholNSessions}{}\renewcommand{\rhholNSessions}{64}  % holiday sessions in the evaluation span 2014-01-02 .. 2024-04-30
\providecommand{\rhholNBars}{}\renewcommand{\rhholNBars}{443}  % their regular-hours bars (forecast, not scored)
\providecommand{\rhholNForecast}{}\renewcommand{\rhholNForecast}{33{,}970}  % bars forecast in the span (33,527 scored + the holiday bars)
\providecommand{\rhholNScored}{}\renewcommand{\rhholNScored}{33{,}527}  % bars scored (Table 1)
\providecommand{\rhholLastBar}{}\renewcommand{\rhholLastBar}{13:00}  % the usual last regular-hours bar end of a holiday session (ET)
\providecommand{\rhholEffCABarPpMin}{}\renewcommand{\rhholEffCABarPpMin}{+0.13}  % holiday effect C-A, linear: lowest change of pct_vs_har_bar
\providecommand{\rhholEffCABarPpMax}{}\renewcommand{\rhholEffCABarPpMax}{+0.14}  % holiday effect C-A, linear: highest change of pct_vs_har_bar
\providecommand{\rhholEffCABarPpRng}{}\renewcommand{\rhholEffCABarPpRng}{0.13--0.14}  % holiday effect C-A: range of pct_vs_har_bar
\providecommand{\rhholEffCARidgeBarPp}{}\renewcommand{\rhholEffCARidgeBarPp}{+0.14}  % holiday effect C-A, ridge: pct_vs_har_bar
\providecommand{\rhholEffCALassoBarPp}{}\renewcommand{\rhholEffCALassoBarPp}{+0.14}  % holiday effect C-A, lasso: pct_vs_har_bar
\providecommand{\rhholEffCAEnetBarPp}{}\renewcommand{\rhholEffCAEnetBarPp}{+0.13}  % holiday effect C-A, enet: pct_vs_har_bar
\providecommand{\rhholEffCAClosePpMin}{}\renewcommand{\rhholEffCAClosePpMin}{\ensuremath{-}0.98}  % holiday effect C-A, linear: lowest change of pct_vs_har_c16
\providecommand{\rhholEffCAClosePpMax}{}\renewcommand{\rhholEffCAClosePpMax}{\ensuremath{-}0.49}  % holiday effect C-A, linear: highest change of pct_vs_har_c16
\providecommand{\rhholEffCAClosePpRng}{}\renewcommand{\rhholEffCAClosePpRng}{\ensuremath{-}0.98 to \ensuremath{-}0.49}  % holiday effect C-A: range of pct_vs_har_c16
\providecommand{\rhholEffCARidgeClosePp}{}\renewcommand{\rhholEffCARidgeClosePp}{\ensuremath{-}0.98}  % holiday effect C-A, ridge: pct_vs_har_c16
\providecommand{\rhholEffCALassoClosePp}{}\renewcommand{\rhholEffCALassoClosePp}{\ensuremath{-}0.49}  % holiday effect C-A, lasso: pct_vs_har_c16
\providecommand{\rhholEffCAEnetClosePp}{}\renewcommand{\rhholEffCAEnetClosePp}{\ensuremath{-}0.75}  % holiday effect C-A, enet: pct_vs_har_c16
\providecommand{\rhholEffCAShFullMin}{}\renewcommand{\rhholEffCAShFullMin}{\ensuremath{-}0.15}  % holiday effect C-A, linear: lowest change of sr_mid_full
\providecommand{\rhholEffCAShFullMax}{}\renewcommand{\rhholEffCAShFullMax}{+0.39}  % holiday effect C-A, linear: highest change of sr_mid_full
\providecommand{\rhholEffCAShFullRng}{}\renewcommand{\rhholEffCAShFullRng}{\ensuremath{-}0.15 to +0.39}  % holiday effect C-A: range of sr_mid_full
\providecommand{\rhholEffCARidgeShFull}{}\renewcommand{\rhholEffCARidgeShFull}{+0.17}  % holiday effect C-A, ridge: sr_mid_full
\providecommand{\rhholEffCALassoShFull}{}\renewcommand{\rhholEffCALassoShFull}{+0.39}  % holiday effect C-A, lasso: sr_mid_full
\providecommand{\rhholEffCAEnetShFull}{}\renewcommand{\rhholEffCAEnetShFull}{\ensuremath{-}0.15}  % holiday effect C-A, enet: sr_mid_full
\providecommand{\rhholEffCADShMin}{}\renewcommand{\rhholEffCADShMin}{\ensuremath{-}0.18}  % holiday effect C-A, linear: lowest change of dsr_mid
\providecommand{\rhholEffCADShMax}{}\renewcommand{\rhholEffCADShMax}{+0.36}  % holiday effect C-A, linear: highest change of dsr_mid
\providecommand{\rhholEffCADShRng}{}\renewcommand{\rhholEffCADShRng}{\ensuremath{-}0.18 to +0.36}  % holiday effect C-A: range of dsr_mid
\providecommand{\rhholEffCARidgeDSh}{}\renewcommand{\rhholEffCARidgeDSh}{+0.13}  % holiday effect C-A, ridge: dsr_mid
\providecommand{\rhholEffCALassoDSh}{}\renewcommand{\rhholEffCALassoDSh}{+0.36}  % holiday effect C-A, lasso: dsr_mid
\providecommand{\rhholEffCAEnetDSh}{}\renewcommand{\rhholEffCAEnetDSh}{\ensuremath{-}0.18}  % holiday effect C-A, enet: dsr_mid
\providecommand{\rhholEffDBBarPpMin}{}\renewcommand{\rhholEffDBBarPpMin}{+0.11}  % holiday effect D-B, linear: lowest change of pct_vs_har_bar
\providecommand{\rhholEffDBBarPpMax}{}\renewcommand{\rhholEffDBBarPpMax}{+0.30}  % holiday effect D-B, linear: highest change of pct_vs_har_bar
\providecommand{\rhholEffDBBarPpRng}{}\renewcommand{\rhholEffDBBarPpRng}{0.11--0.30}  % holiday effect D-B: range of pct_vs_har_bar
\providecommand{\rhholEffDBRidgeBarPp}{}\renewcommand{\rhholEffDBRidgeBarPp}{+0.12}  % holiday effect D-B, ridge: pct_vs_har_bar
\providecommand{\rhholEffDBLassoBarPp}{}\renewcommand{\rhholEffDBLassoBarPp}{+0.30}  % holiday effect D-B, lasso: pct_vs_har_bar
\providecommand{\rhholEffDBEnetBarPp}{}\renewcommand{\rhholEffDBEnetBarPp}{+0.11}  % holiday effect D-B, enet: pct_vs_har_bar
\providecommand{\rhholEffDBClosePpMin}{}\renewcommand{\rhholEffDBClosePpMin}{\ensuremath{-}0.86}  % holiday effect D-B, linear: lowest change of pct_vs_har_c16
\providecommand{\rhholEffDBClosePpMax}{}\renewcommand{\rhholEffDBClosePpMax}{\ensuremath{-}0.20}  % holiday effect D-B, linear: highest change of pct_vs_har_c16
\providecommand{\rhholEffDBClosePpRng}{}\renewcommand{\rhholEffDBClosePpRng}{\ensuremath{-}0.86 to \ensuremath{-}0.20}  % holiday effect D-B: range of pct_vs_har_c16
\providecommand{\rhholEffDBRidgeClosePp}{}\renewcommand{\rhholEffDBRidgeClosePp}{\ensuremath{-}0.86}  % holiday effect D-B, ridge: pct_vs_har_c16
\providecommand{\rhholEffDBLassoClosePp}{}\renewcommand{\rhholEffDBLassoClosePp}{\ensuremath{-}0.20}  % holiday effect D-B, lasso: pct_vs_har_c16
\providecommand{\rhholEffDBEnetClosePp}{}\renewcommand{\rhholEffDBEnetClosePp}{\ensuremath{-}0.65}  % holiday effect D-B, enet: pct_vs_har_c16
\providecommand{\rhholEffDBShFullMin}{}\renewcommand{\rhholEffDBShFullMin}{\ensuremath{-}0.20}  % holiday effect D-B, linear: lowest change of sr_mid_full
\providecommand{\rhholEffDBShFullMax}{}\renewcommand{\rhholEffDBShFullMax}{+0.38}  % holiday effect D-B, linear: highest change of sr_mid_full
\providecommand{\rhholEffDBShFullRng}{}\renewcommand{\rhholEffDBShFullRng}{\ensuremath{-}0.20 to +0.38}  % holiday effect D-B: range of sr_mid_full
\providecommand{\rhholEffDBRidgeShFull}{}\renewcommand{\rhholEffDBRidgeShFull}{+0.20}  % holiday effect D-B, ridge: sr_mid_full
\providecommand{\rhholEffDBLassoShFull}{}\renewcommand{\rhholEffDBLassoShFull}{+0.38}  % holiday effect D-B, lasso: sr_mid_full
\providecommand{\rhholEffDBEnetShFull}{}\renewcommand{\rhholEffDBEnetShFull}{\ensuremath{-}0.20}  % holiday effect D-B, enet: sr_mid_full
\providecommand{\rhholEffDBDShMin}{}\renewcommand{\rhholEffDBDShMin}{\ensuremath{-}0.24}  % holiday effect D-B, linear: lowest change of dsr_mid
\providecommand{\rhholEffDBDShMax}{}\renewcommand{\rhholEffDBDShMax}{+0.35}  % holiday effect D-B, linear: highest change of dsr_mid
\providecommand{\rhholEffDBDShRng}{}\renewcommand{\rhholEffDBDShRng}{\ensuremath{-}0.24 to +0.35}  % holiday effect D-B: range of dsr_mid
\providecommand{\rhholEffDBRidgeDSh}{}\renewcommand{\rhholEffDBRidgeDSh}{+0.17}  % holiday effect D-B, ridge: dsr_mid
\providecommand{\rhholEffDBLassoDSh}{}\renewcommand{\rhholEffDBLassoDSh}{+0.35}  % holiday effect D-B, lasso: dsr_mid
\providecommand{\rhholEffDBEnetDSh}{}\renewcommand{\rhholEffDBEnetDSh}{\ensuremath{-}0.24}  % holiday effect D-B, enet: dsr_mid
\providecommand{\rhholRidgeAPctBar}{}\renewcommand{\rhholRidgeAPctBar}{\ensuremath{-}5.4}  % cell A, ridge: every-bar % vs HAR
\providecommand{\rhholRidgeAPBar}{}\renewcommand{\rhholRidgeAPBar}{\ensuremath{<}0.001}  % cell A, ridge: every-bar DM p
\providecommand{\rhholRidgeAPctClose}{}\renewcommand{\rhholRidgeAPctClose}{0.8}  % cell A, ridge: 16:00 % vs HAR
\providecommand{\rhholRidgeAShFull}{}\renewcommand{\rhholRidgeAShFull}{1.77}  % cell A, ridge: Sharpe ratio (D19), full set, mid
\providecommand{\rhholLassoAPctBar}{}\renewcommand{\rhholLassoAPctBar}{\ensuremath{-}6.2}  % cell A, lasso: every-bar % vs HAR
\providecommand{\rhholLassoAPBar}{}\renewcommand{\rhholLassoAPBar}{\ensuremath{<}0.001}  % cell A, lasso: every-bar DM p
\providecommand{\rhholLassoAPctClose}{}\renewcommand{\rhholLassoAPctClose}{3.3}  % cell A, lasso: 16:00 % vs HAR
\providecommand{\rhholLassoAShFull}{}\renewcommand{\rhholLassoAShFull}{1.72}  % cell A, lasso: Sharpe ratio (D19), full set, mid
\providecommand{\rhholEnetAPctBar}{}\renewcommand{\rhholEnetAPctBar}{\ensuremath{-}6.0}  % cell A, enet: every-bar % vs HAR
\providecommand{\rhholEnetAPBar}{}\renewcommand{\rhholEnetAPBar}{\ensuremath{<}0.001}  % cell A, enet: every-bar DM p
\providecommand{\rhholEnetAPctClose}{}\renewcommand{\rhholEnetAPctClose}{\ensuremath{-}0.5}  % cell A, enet: 16:00 % vs HAR
\providecommand{\rhholEnetAShFull}{}\renewcommand{\rhholEnetAShFull}{1.86}  % cell A, enet: Sharpe ratio (D19), full set, mid
\providecommand{\rhholRidgeBPctBar}{}\renewcommand{\rhholRidgeBPctBar}{\ensuremath{-}6.0}  % cell B, ridge: every-bar % vs HAR
\providecommand{\rhholRidgeBPBar}{}\renewcommand{\rhholRidgeBPBar}{\ensuremath{<}0.001}  % cell B, ridge: every-bar DM p
\providecommand{\rhholRidgeBPctClose}{}\renewcommand{\rhholRidgeBPctClose}{1.8}  % cell B, ridge: 16:00 % vs HAR
\providecommand{\rhholRidgeBShFull}{}\renewcommand{\rhholRidgeBShFull}{1.64}  % cell B, ridge: Sharpe ratio (D19), full set, mid
\providecommand{\rhholLassoBPctBar}{}\renewcommand{\rhholLassoBPctBar}{\ensuremath{-}9.0}  % cell B, lasso: every-bar % vs HAR
\providecommand{\rhholLassoBPBar}{}\renewcommand{\rhholLassoBPBar}{\ensuremath{<}0.001}  % cell B, lasso: every-bar DM p
\providecommand{\rhholLassoBPctClose}{}\renewcommand{\rhholLassoBPctClose}{4.2}  % cell B, lasso: 16:00 % vs HAR
\providecommand{\rhholLassoBShFull}{}\renewcommand{\rhholLassoBShFull}{1.73}  % cell B, lasso: Sharpe ratio (D19), full set, mid
\providecommand{\rhholEnetBPctBar}{}\renewcommand{\rhholEnetBPctBar}{\ensuremath{-}6.7}  % cell B, enet: every-bar % vs HAR
\providecommand{\rhholEnetBPBar}{}\renewcommand{\rhholEnetBPBar}{\ensuremath{<}0.001}  % cell B, enet: every-bar DM p
\providecommand{\rhholEnetBPctClose}{}\renewcommand{\rhholEnetBPctClose}{0.6}  % cell B, enet: 16:00 % vs HAR
\providecommand{\rhholEnetBShFull}{}\renewcommand{\rhholEnetBShFull}{1.85}  % cell B, enet: Sharpe ratio (D19), full set, mid
\providecommand{\rhholRidgeCPctBar}{}\renewcommand{\rhholRidgeCPctBar}{\ensuremath{-}5.2}  % cell C, ridge: every-bar % vs HAR
\providecommand{\rhholRidgeCPBar}{}\renewcommand{\rhholRidgeCPBar}{\ensuremath{<}0.001}  % cell C, ridge: every-bar DM p
\providecommand{\rhholRidgeCPctClose}{}\renewcommand{\rhholRidgeCPctClose}{\ensuremath{-}0.2}  % cell C, ridge: 16:00 % vs HAR
\providecommand{\rhholRidgeCShFull}{}\renewcommand{\rhholRidgeCShFull}{1.94}  % cell C, ridge: Sharpe ratio (D19), full set, mid
\providecommand{\rhholLassoCPctBar}{}\renewcommand{\rhholLassoCPctBar}{\ensuremath{-}6.1}  % cell C, lasso: every-bar % vs HAR
\providecommand{\rhholLassoCPBar}{}\renewcommand{\rhholLassoCPBar}{\ensuremath{<}0.001}  % cell C, lasso: every-bar DM p
\providecommand{\rhholLassoCPctClose}{}\renewcommand{\rhholLassoCPctClose}{2.8}  % cell C, lasso: 16:00 % vs HAR
\providecommand{\rhholLassoCShFull}{}\renewcommand{\rhholLassoCShFull}{2.12}  % cell C, lasso: Sharpe ratio (D19), full set, mid
\providecommand{\rhholEnetCPctBar}{}\renewcommand{\rhholEnetCPctBar}{\ensuremath{-}5.9}  % cell C, enet: every-bar % vs HAR
\providecommand{\rhholEnetCPBar}{}\renewcommand{\rhholEnetCPBar}{\ensuremath{<}0.001}  % cell C, enet: every-bar DM p
\providecommand{\rhholEnetCPctClose}{}\renewcommand{\rhholEnetCPctClose}{\ensuremath{-}1.3}  % cell C, enet: 16:00 % vs HAR
\providecommand{\rhholEnetCShFull}{}\renewcommand{\rhholEnetCShFull}{1.71}  % cell C, enet: Sharpe ratio (D19), full set, mid
\providecommand{\rhholRidgeDPctBar}{}\renewcommand{\rhholRidgeDPctBar}{\ensuremath{-}5.9}  % cell D, ridge: every-bar % vs HAR
\providecommand{\rhholRidgeDPBar}{}\renewcommand{\rhholRidgeDPBar}{\ensuremath{<}0.001}  % cell D, ridge: every-bar DM p
\providecommand{\rhholRidgeDPctClose}{}\renewcommand{\rhholRidgeDPctClose}{0.9}  % cell D, ridge: 16:00 % vs HAR
\providecommand{\rhholRidgeDShFull}{}\renewcommand{\rhholRidgeDShFull}{1.85}  % cell D, ridge: Sharpe ratio (D19), full set, mid
\providecommand{\rhholLassoDPctBar}{}\renewcommand{\rhholLassoDPctBar}{\ensuremath{-}8.7}  % cell D, lasso: every-bar % vs HAR
\providecommand{\rhholLassoDPBar}{}\renewcommand{\rhholLassoDPBar}{\ensuremath{<}0.001}  % cell D, lasso: every-bar DM p
\providecommand{\rhholLassoDPctClose}{}\renewcommand{\rhholLassoDPctClose}{4.0}  % cell D, lasso: 16:00 % vs HAR
\providecommand{\rhholLassoDShFull}{}\renewcommand{\rhholLassoDShFull}{2.11}  % cell D, lasso: Sharpe ratio (D19), full set, mid
\providecommand{\rhholEnetDPctBar}{}\renewcommand{\rhholEnetDPctBar}{\ensuremath{-}6.6}  % cell D, enet: every-bar % vs HAR
\providecommand{\rhholEnetDPBar}{}\renewcommand{\rhholEnetDPBar}{\ensuremath{<}0.001}  % cell D, enet: every-bar DM p
\providecommand{\rhholEnetDPctClose}{}\renewcommand{\rhholEnetDPctClose}{0.0}  % cell D, enet: 16:00 % vs HAR
\providecommand{\rhholEnetDShFull}{}\renewcommand{\rhholEnetDShFull}{1.65}  % cell D, enet: Sharpe ratio (D19), full set, mid
\fi
% 2026-10-26 (coordinator): parked stop-gap, the macro is now generated -- \providecommand{\rhPending}{\textit{pending}}
% Window-sweep macros, not generated yet (NUMBERS: from results/window_sweep_2014start_2026-10-22/part2 and
% latestart/latestart_spans.txt). Stop-gaps at the values of the sweep's design (part2/tables.txt: windows 1,000,
% 1,500, ..., 3,500 trading days; latestart_spans.txt, fold-3 validation year: some late-start columns live at
% W = 1,000, none at W >= 2,000); \providecommand, so a generated macro wins once it exists.
% 2026-10-26 (coordinator): parked stop-gap, the macro is now generated -- \providecommand{\rhwsNWin}{six}  % windows scored out of sample
% 2026-10-26 (coordinator): parked stop-gap, the macro is now generated -- \providecommand{\rhwsWinMin}{1{,}000}  % shortest window, trading days
% 2026-10-26 (coordinator): parked stop-gap, the macro is now generated -- \providecommand{\rhwsWinMax}{3{,}500}  % longest window, trading days
% 2026-10-26 (coordinator): parked stop-gap, the macro is now generated -- \providecommand{\rhwsLateNoneFrom}{2{,}000}  % shortest window at which no late-start column is live in the fold

\subsection{Estimation window}
\label{app:rh_robust_window}
\begin{itemize}
\item \textbf{How the window was chosen.} The rolling window of the main design, $W = \rhrbWinMain$ trading days
  ($\rhrbWinMainBars$ bars), was fixed after the pooled linear models had been estimated and scored out of sample,
  over the evaluation span $\mathcal T_{\mathrm{oos}}$ of Table~\ref{tab:rh_oos}, at \rhwsNWin{} windows from
  \rhwsWinMin{} to \rhwsWinMax{} trading days, on the design before the FOMC fixes of
  Appendix~\ref{app:rh_dp_fomc}. The window is therefore not a pre-sample choice, and the items below report the
  windows that were scored.
  \begin{itemize}
  \item The pre-sample learning curve, on the fold on which the linear penalties are chosen
    (Appendix~\ref{app:rh_methods_splits}), scores every regular-hours bar, has no separate measure for the 16:00
    bar, and its validation MSE does not single out the main window.
  \item In that fold, some of the late-starting series (Appendix~\ref{app:rh_tuning}) are live, wholly or in part,
    at a window of \rhwsWinMin{} trading days, but none is live at a window of \rhwsLateNoneFrom{} trading days or
    more, so the learning curve compares different effective predictor sets.
  \end{itemize}
% 2026-10-26 (coordinator): unwrapped \IfFileExists{generated/rehaul_window_sweep_table.tex} (the file exists; the flat Overleaf build has no rehaul/ directory, so the wrapper skipped its body there). Its false branch was:
% | %
% | \item \textbf{Every window, penalized linear models.} The forecasts and scores at every window are in the
% |   replication package.
%
\item \textbf{Every window, penalized linear models.} Table~\ref{tab:rh_window_sweep} reports, at each window, the
  change in QLIKE from the full predictor set against the HAR benchmark, on every bar and on the 16:00 bar. Every bar:
  the elastic net's change lies between \rhwsEnetBarPctMin\% and \rhwsEnetBarPctMax\%, significant at each window.
  The 16:00 bar: comparisons are window-sensitive, and the same change ranges from \rhwsEnetClosePctMin\% to
  \rhwsEnetClosePctMax\%. % 2026-10-26 (coordinator): the 16:00 range moved here from Section 6 (author: no 16:00-bar results in the main text)
% AUTO-GENERATED by writeup/make_rehaul_window_sweep.py -- do not edit; re-run the script.
% Table tab:rh_window_sweep (appendix, D8): results/window_sweep_2014start_2026-10-22/part2 (gap.csv, econ.csv).
\begin{table}[htbp]
\centering
\small
\setlength{\tabcolsep}{4pt}
\begin{threeparttable}
\caption{The full predictor set against the HAR set across rolling windows}
\label{tab:rh_window_sweep}
\begin{tabular*}{\linewidth}{@{\extracolsep{\fill}}lrrrrrr@{}}
\toprule
 & & \multicolumn{2}{c}{Every bar} & \multicolumn{2}{c}{16:00} & \\
\cmidrule(lr){3-4}\cmidrule(lr){5-6}
Model & $W$ (days) & \% vs HAR & DM ($p$) & \% vs HAR & DM ($p$) & $\Delta\SR$ ($p$) \\
\midrule
Ridge & 1{,}000 & \ensuremath{-}5.5 & \ensuremath{-}7.59 (\ensuremath{<}0.001) & \ensuremath{-}3.1 & \ensuremath{-}1.80 (0.071) & 0.32 (0.553) \\
 & 1{,}500 & \ensuremath{-}5.4 & \ensuremath{-}7.35 (\ensuremath{<}0.001) & \ensuremath{-}1.9 & \ensuremath{-}1.00 (0.320) & 0.55 (0.299) \\
 & 2{,}000 & \ensuremath{-}4.9 & \ensuremath{-}4.53 (\ensuremath{<}0.001) & 3.1 & 0.99 (0.323) & 0.40 (0.593) \\
 & 2{,}500 & \ensuremath{-}4.7 & \ensuremath{-}4.03 (\ensuremath{<}0.001) & 2.0 & 0.78 (0.434) & 0.21 (0.791) \\
 & 3{,}000$^{\ast}$ & \ensuremath{-}5.4 & \ensuremath{-}5.63 (\ensuremath{<}0.001) & 0.8 & 0.31 (0.759) & 0.86 (0.161) \\
 & 3{,}500 & \ensuremath{-}6.3 & \ensuremath{-}7.06 (\ensuremath{<}0.001) & \ensuremath{-}3.1 & \ensuremath{-}1.55 (0.120) & 0.74 (0.186) \\
\addlinespace[3pt]
Lasso & 1{,}000 & \ensuremath{-}4.9 & \ensuremath{-}3.75 (\ensuremath{<}0.001) & \ensuremath{-}2.8 & \ensuremath{-}1.23 (0.218) & 0.45 (0.412) \\
 & 1{,}500 & \ensuremath{-}5.9 & \ensuremath{-}6.38 (\ensuremath{<}0.001) & \ensuremath{-}2.2 & \ensuremath{-}0.98 (0.326) & 0.67 (0.242) \\
 & 2{,}000 & \ensuremath{-}6.2 & \ensuremath{-}5.17 (\ensuremath{<}0.001) & 5.4 & 1.46 (0.145) & 0.15 (0.850) \\
 & 2{,}500 & \ensuremath{-}5.3 & \ensuremath{-}5.81 (\ensuremath{<}0.001) & \ensuremath{-}1.5 & \ensuremath{-}0.75 (0.455) & 1.04 (0.122) \\
 & 3{,}000$^{\ast}$ & \ensuremath{-}6.2 & \ensuremath{-}4.72 (\ensuremath{<}0.001) & 3.3 & 1.04 (0.300) & 0.74 (0.227) \\
 & 3{,}500 & \ensuremath{-}7.6 & \ensuremath{-}6.39 (\ensuremath{<}0.001) & \ensuremath{-}1.6 & \ensuremath{-}0.69 (0.488) & 0.74 (0.236) \\
\addlinespace[3pt]
Elastic net & 1{,}000 & \ensuremath{-}5.9 & \ensuremath{-}8.83 (\ensuremath{<}0.001) & \ensuremath{-}4.0 & \ensuremath{-}2.56 (0.011) & 0.72 (0.134) \\
 & 1{,}500 & \ensuremath{-}5.6 & \ensuremath{-}5.30 (\ensuremath{<}0.001) & \ensuremath{-}0.5 & \ensuremath{-}0.21 (0.835) & 0.63 (0.268) \\
 & 2{,}000 & \ensuremath{-}5.7 & \ensuremath{-}5.54 (\ensuremath{<}0.001) & 1.3 & 0.48 (0.633) & 0.44 (0.571) \\
 & 2{,}500 & \ensuremath{-}5.2 & \ensuremath{-}4.88 (\ensuremath{<}0.001) & 0.4 & 0.17 (0.862) & 1.00 (0.120) \\
 & 3{,}000$^{\ast}$ & \ensuremath{-}6.0 & \ensuremath{-}6.98 (\ensuremath{<}0.001) & \ensuremath{-}0.5 & \ensuremath{-}0.21 (0.832) & 0.96 (0.135) \\
 & 3{,}500 & \ensuremath{-}6.7 & \ensuremath{-}8.45 (\ensuremath{<}0.001) & \ensuremath{-}3.9 & \ensuremath{-}2.15 (0.032) & 0.98 (0.101) \\
\bottomrule
\end{tabular*}
\begin{tablenotes}[flushleft]
\footnotesize
\item \textit{Notes:} Each penalized linear class estimated on the HAR set and on the full set with a rolling window of $W$ trading days ($13W$ bars), its penalties chosen on the selection fold at that window and frozen; $^{\ast}$ the window of the main text. \% vs HAR: the full set's QLIKE relative to the HAR set's (negative: lower loss), on the \rhorNIntB{} bars of Table~\ref{tab:rh_oos} and on their 16:00 bar; DM: Diebold--Mariano statistic of the daily (average) QLIKE, HLN correction, Newey--West lag \rhorHacLag, $p$-value in parentheses. $\Delta\SR$: the full-set minus the HAR-set Sharpe ratio of the sign strategy at mid-quotes on the \rheSzDays{} post-warmup option days, with the bootstrap $p$-value. Design before the FOMC calendar corrections of Section~\ref{sec:rh_data}.
\end{tablenotes}
\end{threeparttable}
\end{table}

% 2026-10-26 (integration, JFEc): moved here from the item "Estimation window" of 10_appendix.tex (section Robustness
% of the statistical results), split into sub-items; the pre-patch design marked (D13), and the 16:00 counts at the
% two windows added (audit correctness-2: the 16:00 comparisons are window-sensitive). The 2026-10-23 check of that
% item still holds: the \rhrbNSigBar every-bar and \rhrbNSigClose 16:00 arms whose QLIKE differs between the windows
% at 5% are ridge, lasso and elastic net; the largest Sharpe-ratio change at mid-quotes is the random forest's on the
% HAR benchmark.
\item \textbf{A longer window, every model.} Table~\ref{tab:rh_robust_w3500} refits the models of
  Table~\ref{tab:rh_oos} on a rolling window of \rhrbWinLong{} trading days ($W = \rhrbWinLongBars$ bars) instead
  of \rhrbWinMain, with the same first evaluation day, the same frozen hyperparameters and a refit at every bar, on
  the design before the FOMC fixes.
  \begin{itemize}
  \item \emph{Every bar.} The full predictor set is significantly more accurate than the HAR benchmark for
    \rhrbNFullLowerSigBarLong{} of the \rhrbNModelsLong{} model classes (\rhrbNFullLowerSigBarMain{} of
    \rhrbNModelsMain{} at \rhrbWinMain{} days). Mean QLIKE moves by at most \rhrbMaxAbsPctBar\% on every bar and
    \rhrbMaxAbsPctClose\% on the 16:00 bar; the \rhrbNSigBar{} forecasts whose every-bar QLIKE differs significantly
    between the windows, and the \rhrbNSigClose{} at 16:00, are all linear models.
  \item \emph{The 16:00 bar is sensitive to the window.} There the full predictor set is significantly more accurate
    than the HAR benchmark for \rhrbNFullLowerSigCloseLong{} of the \rhrbNModelsLong{} classes at \rhrbWinLong{}
    days and for \rhrbNFullLowerSigCloseMain{} at \rhrbWinMain{} days.
  \item \emph{Sharpe ratio.} It changes significantly between the windows for \rhrbNSRSigMid{} of the \rhrbNArms{}
    forecasts at mid-quotes and \rhrbNSRSigBa{} at quoted bid and ask prices; the largest change at mid-quotes is
    that of the random forest on the HAR benchmark, \rhrbDSRRfHarMid{} (\rhrbPSRRfHarMid).
  \end{itemize}
% AUTO-GENERATED by writeup/make_rehaul_robust_window.py -- do not edit; re-run the script.
% tab:rh_robust_w3500 (appendix): the final pooled models on a rolling window of 3,500 trading days against
% 3,000 (same first day, frozen hyperparameters and refit at every bar); QLIKE on every bar and on the 16:00 bar of
% the 2,579 evaluation days, DM on the daily loss (HLN, Newey-West lag 8); Sharpe ratio (D6) of the sign
% strategy on the 614 post-warmup option days, differences tested by the paired block bootstrap (hac_inference).
% Runs: results/window_sweep_2014start_2026-10-22/w3000, results/pooled_trees_final_2014start_w3000, results/pooled_rf_skx_2014start_w3000 and their w3500 twins.
% Pending: none.
\providecommand{\rhPending}{\textit{pending}}
\ifdefined\rhrbNdays\else% AUTO-GENERATED by writeup/make_rehaul_robust_window.py -- do not edit; re-run the script.
% \rhrb... macros of the appendix robustness table tab:rh_robust_w3500 (rehaul_robust_w3500_table.tex): the final
% pooled models on a rolling window of 3,500 against 3,000 trading days; every bar scored on the 2,579 days 2014-01-02 .. 2024-04-30
% (33,527 bars), Sharpe ratios (D19 pooled) on the 614 post-warmup option days. Names: Q (QLIKE), SR (Sharpe ratio), Pct / DM / P
% (3,500 vs 3,000: % change, DM, p statement), DSR / PSR (Sharpe change 3,500 - 3,000, Ledoit-Wolf p), PctF / DMF / PF / DSRF / PSRF
% (full set vs HAR benchmark at a window); Model Ols Ridge Lasso Enet Lgbm Xgb Rf; Set Har Full; Bar (every bar) / Close (16:00);
% Mid / Ba (bid/ask); Main (W = 3,000) / Long (W = 3,500). Pending: none.
\providecommand{\rhPending}{\textit{pending}}
\providecommand{\rhrbWinMain}{}\renewcommand{\rhrbWinMain}{3{,}000}  % main rolling window, trading days
\providecommand{\rhrbWinLong}{}\renewcommand{\rhrbWinLong}{3{,}500}  % robustness rolling window, trading days
\providecommand{\rhrbWinMainBars}{}\renewcommand{\rhrbWinMainBars}{39{,}000}  % main rolling window, 13-bar observations
\providecommand{\rhrbWinLongBars}{}\renewcommand{\rhrbWinLongBars}{45{,}500}  % robustness rolling window, 13-bar observations
\providecommand{\rhrbNdays}{}\renewcommand{\rhrbNdays}{2{,}579}  % evaluation days (every bar and 16:00)
\providecommand{\rhrbNbars}{}\renewcommand{\rhrbNbars}{33{,}527}  % scored bars, the same at both windows (asserted)
\providecommand{\rhrbSpanStart}{}\renewcommand{\rhrbSpanStart}{2 January 2014}  % first evaluation day
\providecommand{\rhrbSpanEnd}{}\renewcommand{\rhrbSpanEnd}{30 April 2024}  % last evaluation day
\providecommand{\rhrbDMLag}{}\renewcommand{\rhrbDMLag}{8}  % Newey-West lag of every DM test
\providecommand{\rhrbNopt}{}\renewcommand{\rhrbNopt}{866}  % option days of the sign strategy
\providecommand{\rhrbNpw}{}\renewcommand{\rhrbNpw}{614}  % post-warmup option days, the Sharpe-ratio sample
\providecommand{\rhrbBootN}{}\renewcommand{\rhrbBootN}{4{,}999}  % bootstrap resamples of the Sharpe-ratio test (Ledoit-Wolf since D19)
\providecommand{\rhrbBootBlock}{}\renewcommand{\rhrbBootBlock}{10}  % block length of the Sharpe-ratio bootstrap
\providecommand{\rhrbPendingNote}{}\renewcommand{\rhrbPendingNote}{}  % table-note sentence on the pending forest (empty when nothing is pending)
\providecommand{\rhrbQOlsHarBarMain}{}\renewcommand{\rhrbQOlsHarBarMain}{0.1893}  % QLIKE, every bar, OLS, HAR benchmark, W = 3,000
\providecommand{\rhrbQOlsHarCloseMain}{}\renewcommand{\rhrbQOlsHarCloseMain}{0.1269}  % QLIKE, 16:00 bar, OLS, HAR benchmark, W = 3,000
\providecommand{\rhrbQOlsHarBarLong}{}\renewcommand{\rhrbQOlsHarBarLong}{0.1897}  % QLIKE, every bar, OLS, HAR benchmark, W = 3,500
\providecommand{\rhrbQOlsHarCloseLong}{}\renewcommand{\rhrbQOlsHarCloseLong}{0.1272}  % QLIKE, 16:00 bar, OLS, HAR benchmark, W = 3,500
\providecommand{\rhrbPctOlsHarBar}{}\renewcommand{\rhrbPctOlsHarBar}{0.2}  % QLIKE change, %, W = 3,500 vs 3,000, every bar, OLS, HAR benchmark
\providecommand{\rhrbDMOlsHarBar}{}\renewcommand{\rhrbDMOlsHarBar}{2.48}  % DM, W = 3,500 vs 3,000 (positive: 3,500 has the higher loss), every bar, OLS, HAR benchmark
\providecommand{\rhrbPOlsHarBar}{}\renewcommand{\rhrbPOlsHarBar}{\ensuremath{p = 0.013}}  % ... its p statement
\providecommand{\rhrbPctOlsHarClose}{}\renewcommand{\rhrbPctOlsHarClose}{0.3}  % QLIKE change, %, W = 3,500 vs 3,000, 16:00 bar, OLS, HAR benchmark
\providecommand{\rhrbDMOlsHarClose}{}\renewcommand{\rhrbDMOlsHarClose}{1.31}  % DM, W = 3,500 vs 3,000 (positive: 3,500 has the higher loss), 16:00 bar, OLS, HAR benchmark
\providecommand{\rhrbPOlsHarClose}{}\renewcommand{\rhrbPOlsHarClose}{\ensuremath{p = 0.19}}  % ... its p statement
\providecommand{\rhrbQRidgeHarBarMain}{}\renewcommand{\rhrbQRidgeHarBarMain}{0.1893}  % QLIKE, every bar, Ridge, HAR benchmark, W = 3,000
\providecommand{\rhrbQRidgeHarCloseMain}{}\renewcommand{\rhrbQRidgeHarCloseMain}{0.1269}  % QLIKE, 16:00 bar, Ridge, HAR benchmark, W = 3,000
\providecommand{\rhrbSRRidgeHarMidMain}{}\renewcommand{\rhrbSRRidgeHarMidMain}{0.95}  % Sharpe ratio (D19 pooled), sign strategy at mid-quotes, Ridge, HAR benchmark, W = 3,000
\providecommand{\rhrbSRRidgeHarBaMain}{}\renewcommand{\rhrbSRRidgeHarBaMain}{0.60}  % Sharpe ratio (D19 pooled), sign strategy at bid/ask prices, Ridge, HAR benchmark, W = 3,000
\providecommand{\rhrbQRidgeHarBarLong}{}\renewcommand{\rhrbQRidgeHarBarLong}{0.1897}  % QLIKE, every bar, Ridge, HAR benchmark, W = 3,500
\providecommand{\rhrbQRidgeHarCloseLong}{}\renewcommand{\rhrbQRidgeHarCloseLong}{0.1272}  % QLIKE, 16:00 bar, Ridge, HAR benchmark, W = 3,500
\providecommand{\rhrbSRRidgeHarMidLong}{}\renewcommand{\rhrbSRRidgeHarMidLong}{1.01}  % Sharpe ratio (D19 pooled), sign strategy at mid-quotes, Ridge, HAR benchmark, W = 3,500
\providecommand{\rhrbSRRidgeHarBaLong}{}\renewcommand{\rhrbSRRidgeHarBaLong}{0.65}  % Sharpe ratio (D19 pooled), sign strategy at bid/ask prices, Ridge, HAR benchmark, W = 3,500
\providecommand{\rhrbPctRidgeHarBar}{}\renewcommand{\rhrbPctRidgeHarBar}{0.2}  % QLIKE change, %, W = 3,500 vs 3,000, every bar, Ridge, HAR benchmark
\providecommand{\rhrbDMRidgeHarBar}{}\renewcommand{\rhrbDMRidgeHarBar}{2.48}  % DM, W = 3,500 vs 3,000 (positive: 3,500 has the higher loss), every bar, Ridge, HAR benchmark
\providecommand{\rhrbPRidgeHarBar}{}\renewcommand{\rhrbPRidgeHarBar}{\ensuremath{p = 0.013}}  % ... its p statement
\providecommand{\rhrbPctRidgeHarClose}{}\renewcommand{\rhrbPctRidgeHarClose}{0.3}  % QLIKE change, %, W = 3,500 vs 3,000, 16:00 bar, Ridge, HAR benchmark
\providecommand{\rhrbDMRidgeHarClose}{}\renewcommand{\rhrbDMRidgeHarClose}{1.31}  % DM, W = 3,500 vs 3,000 (positive: 3,500 has the higher loss), 16:00 bar, Ridge, HAR benchmark
\providecommand{\rhrbPRidgeHarClose}{}\renewcommand{\rhrbPRidgeHarClose}{\ensuremath{p = 0.19}}  % ... its p statement
\providecommand{\rhrbDSRRidgeHarMid}{}\renewcommand{\rhrbDSRRidgeHarMid}{0.05}  % Sharpe ratio change, W = 3,500 minus 3,000, mid-quotes, Ridge, HAR benchmark
\providecommand{\rhrbPSRRidgeHarMid}{}\renewcommand{\rhrbPSRRidgeHarMid}{\ensuremath{p = 0.58}}  % ... its p statement (Ledoit-Wolf, D19; was the D6 bootstrap)
\providecommand{\rhrbDSRRidgeHarBa}{}\renewcommand{\rhrbDSRRidgeHarBa}{0.05}  % Sharpe ratio change, W = 3,500 minus 3,000, bid/ask prices, Ridge, HAR benchmark
\providecommand{\rhrbPSRRidgeHarBa}{}\renewcommand{\rhrbPSRRidgeHarBa}{\ensuremath{p = 0.58}}  % ... its p statement (Ledoit-Wolf, D19; was the D6 bootstrap)
\providecommand{\rhrbQRidgeFullBarMain}{}\renewcommand{\rhrbQRidgeFullBarMain}{0.1791}  % QLIKE, every bar, Ridge, full predictor set, W = 3,000
\providecommand{\rhrbQRidgeFullCloseMain}{}\renewcommand{\rhrbQRidgeFullCloseMain}{0.1279}  % QLIKE, 16:00 bar, Ridge, full predictor set, W = 3,000
\providecommand{\rhrbSRRidgeFullMidMain}{}\renewcommand{\rhrbSRRidgeFullMidMain}{1.77}  % Sharpe ratio (D19 pooled), sign strategy at mid-quotes, Ridge, full predictor set, W = 3,000
\providecommand{\rhrbSRRidgeFullBaMain}{}\renewcommand{\rhrbSRRidgeFullBaMain}{1.41}  % Sharpe ratio (D19 pooled), sign strategy at bid/ask prices, Ridge, full predictor set, W = 3,000
\providecommand{\rhrbQRidgeFullBarLong}{}\renewcommand{\rhrbQRidgeFullBarLong}{0.1777}  % QLIKE, every bar, Ridge, full predictor set, W = 3,500
\providecommand{\rhrbQRidgeFullCloseLong}{}\renewcommand{\rhrbQRidgeFullCloseLong}{0.1233}  % QLIKE, 16:00 bar, Ridge, full predictor set, W = 3,500
\providecommand{\rhrbSRRidgeFullMidLong}{}\renewcommand{\rhrbSRRidgeFullMidLong}{1.84}  % Sharpe ratio (D19 pooled), sign strategy at mid-quotes, Ridge, full predictor set, W = 3,500
\providecommand{\rhrbSRRidgeFullBaLong}{}\renewcommand{\rhrbSRRidgeFullBaLong}{1.48}  % Sharpe ratio (D19 pooled), sign strategy at bid/ask prices, Ridge, full predictor set, W = 3,500
\providecommand{\rhrbPctRidgeFullBar}{}\renewcommand{\rhrbPctRidgeFullBar}{\ensuremath{-}0.8}  % QLIKE change, %, W = 3,500 vs 3,000, every bar, Ridge, full predictor set
\providecommand{\rhrbDMRidgeFullBar}{}\renewcommand{\rhrbDMRidgeFullBar}{\ensuremath{-}2.26}  % DM, W = 3,500 vs 3,000 (positive: 3,500 has the higher loss), every bar, Ridge, full predictor set
\providecommand{\rhrbPRidgeFullBar}{}\renewcommand{\rhrbPRidgeFullBar}{\ensuremath{p = 0.024}}  % ... its p statement
\providecommand{\rhrbPctRidgeFullClose}{}\renewcommand{\rhrbPctRidgeFullClose}{\ensuremath{-}3.6}  % QLIKE change, %, W = 3,500 vs 3,000, 16:00 bar, Ridge, full predictor set
\providecommand{\rhrbDMRidgeFullClose}{}\renewcommand{\rhrbDMRidgeFullClose}{\ensuremath{-}3.10}  % DM, W = 3,500 vs 3,000 (positive: 3,500 has the higher loss), 16:00 bar, Ridge, full predictor set
\providecommand{\rhrbPRidgeFullClose}{}\renewcommand{\rhrbPRidgeFullClose}{\ensuremath{p = 0.002}}  % ... its p statement
\providecommand{\rhrbDSRRidgeFullMid}{}\renewcommand{\rhrbDSRRidgeFullMid}{0.07}  % Sharpe ratio change, W = 3,500 minus 3,000, mid-quotes, Ridge, full predictor set
\providecommand{\rhrbPSRRidgeFullMid}{}\renewcommand{\rhrbPSRRidgeFullMid}{\ensuremath{p = 0.81}}  % ... its p statement (Ledoit-Wolf, D19; was the D6 bootstrap)
\providecommand{\rhrbDSRRidgeFullBa}{}\renewcommand{\rhrbDSRRidgeFullBa}{0.07}  % Sharpe ratio change, W = 3,500 minus 3,000, bid/ask prices, Ridge, full predictor set
\providecommand{\rhrbPSRRidgeFullBa}{}\renewcommand{\rhrbPSRRidgeFullBa}{\ensuremath{p = 0.81}}  % ... its p statement (Ledoit-Wolf, D19; was the D6 bootstrap)
\providecommand{\rhrbQLassoHarBarMain}{}\renewcommand{\rhrbQLassoHarBarMain}{0.1893}  % QLIKE, every bar, Lasso, HAR benchmark, W = 3,000
\providecommand{\rhrbQLassoHarCloseMain}{}\renewcommand{\rhrbQLassoHarCloseMain}{0.1269}  % QLIKE, 16:00 bar, Lasso, HAR benchmark, W = 3,000
\providecommand{\rhrbSRLassoHarMidMain}{}\renewcommand{\rhrbSRLassoHarMidMain}{0.95}  % Sharpe ratio (D19 pooled), sign strategy at mid-quotes, Lasso, HAR benchmark, W = 3,000
\providecommand{\rhrbSRLassoHarBaMain}{}\renewcommand{\rhrbSRLassoHarBaMain}{0.60}  % Sharpe ratio (D19 pooled), sign strategy at bid/ask prices, Lasso, HAR benchmark, W = 3,000
\providecommand{\rhrbQLassoHarBarLong}{}\renewcommand{\rhrbQLassoHarBarLong}{0.1897}  % QLIKE, every bar, Lasso, HAR benchmark, W = 3,500
\providecommand{\rhrbQLassoHarCloseLong}{}\renewcommand{\rhrbQLassoHarCloseLong}{0.1272}  % QLIKE, 16:00 bar, Lasso, HAR benchmark, W = 3,500
\providecommand{\rhrbSRLassoHarMidLong}{}\renewcommand{\rhrbSRLassoHarMidLong}{1.01}  % Sharpe ratio (D19 pooled), sign strategy at mid-quotes, Lasso, HAR benchmark, W = 3,500
\providecommand{\rhrbSRLassoHarBaLong}{}\renewcommand{\rhrbSRLassoHarBaLong}{0.65}  % Sharpe ratio (D19 pooled), sign strategy at bid/ask prices, Lasso, HAR benchmark, W = 3,500
\providecommand{\rhrbPctLassoHarBar}{}\renewcommand{\rhrbPctLassoHarBar}{0.2}  % QLIKE change, %, W = 3,500 vs 3,000, every bar, Lasso, HAR benchmark
\providecommand{\rhrbDMLassoHarBar}{}\renewcommand{\rhrbDMLassoHarBar}{2.48}  % DM, W = 3,500 vs 3,000 (positive: 3,500 has the higher loss), every bar, Lasso, HAR benchmark
\providecommand{\rhrbPLassoHarBar}{}\renewcommand{\rhrbPLassoHarBar}{\ensuremath{p = 0.013}}  % ... its p statement
\providecommand{\rhrbPctLassoHarClose}{}\renewcommand{\rhrbPctLassoHarClose}{0.3}  % QLIKE change, %, W = 3,500 vs 3,000, 16:00 bar, Lasso, HAR benchmark
\providecommand{\rhrbDMLassoHarClose}{}\renewcommand{\rhrbDMLassoHarClose}{1.31}  % DM, W = 3,500 vs 3,000 (positive: 3,500 has the higher loss), 16:00 bar, Lasso, HAR benchmark
\providecommand{\rhrbPLassoHarClose}{}\renewcommand{\rhrbPLassoHarClose}{\ensuremath{p = 0.19}}  % ... its p statement
\providecommand{\rhrbDSRLassoHarMid}{}\renewcommand{\rhrbDSRLassoHarMid}{0.05}  % Sharpe ratio change, W = 3,500 minus 3,000, mid-quotes, Lasso, HAR benchmark
\providecommand{\rhrbPSRLassoHarMid}{}\renewcommand{\rhrbPSRLassoHarMid}{\ensuremath{p = 0.58}}  % ... its p statement (Ledoit-Wolf, D19; was the D6 bootstrap)
\providecommand{\rhrbDSRLassoHarBa}{}\renewcommand{\rhrbDSRLassoHarBa}{0.05}  % Sharpe ratio change, W = 3,500 minus 3,000, bid/ask prices, Lasso, HAR benchmark
\providecommand{\rhrbPSRLassoHarBa}{}\renewcommand{\rhrbPSRLassoHarBa}{\ensuremath{p = 0.58}}  % ... its p statement (Ledoit-Wolf, D19; was the D6 bootstrap)
\providecommand{\rhrbQLassoFullBarMain}{}\renewcommand{\rhrbQLassoFullBarMain}{0.1775}  % QLIKE, every bar, Lasso, full predictor set, W = 3,000
\providecommand{\rhrbQLassoFullCloseMain}{}\renewcommand{\rhrbQLassoFullCloseMain}{0.1310}  % QLIKE, 16:00 bar, Lasso, full predictor set, W = 3,000
\providecommand{\rhrbSRLassoFullMidMain}{}\renewcommand{\rhrbSRLassoFullMidMain}{1.72}  % Sharpe ratio (D19 pooled), sign strategy at mid-quotes, Lasso, full predictor set, W = 3,000
\providecommand{\rhrbSRLassoFullBaMain}{}\renewcommand{\rhrbSRLassoFullBaMain}{1.37}  % Sharpe ratio (D19 pooled), sign strategy at bid/ask prices, Lasso, full predictor set, W = 3,000
\providecommand{\rhrbQLassoFullBarLong}{}\renewcommand{\rhrbQLassoFullBarLong}{0.1752}  % QLIKE, every bar, Lasso, full predictor set, W = 3,500
\providecommand{\rhrbQLassoFullCloseLong}{}\renewcommand{\rhrbQLassoFullCloseLong}{0.1253}  % QLIKE, 16:00 bar, Lasso, full predictor set, W = 3,500
\providecommand{\rhrbSRLassoFullMidLong}{}\renewcommand{\rhrbSRLassoFullMidLong}{1.80}  % Sharpe ratio (D19 pooled), sign strategy at mid-quotes, Lasso, full predictor set, W = 3,500
\providecommand{\rhrbSRLassoFullBaLong}{}\renewcommand{\rhrbSRLassoFullBaLong}{1.44}  % Sharpe ratio (D19 pooled), sign strategy at bid/ask prices, Lasso, full predictor set, W = 3,500
\providecommand{\rhrbPctLassoFullBar}{}\renewcommand{\rhrbPctLassoFullBar}{\ensuremath{-}1.3}  % QLIKE change, %, W = 3,500 vs 3,000, every bar, Lasso, full predictor set
\providecommand{\rhrbDMLassoFullBar}{}\renewcommand{\rhrbDMLassoFullBar}{\ensuremath{-}2.04}  % DM, W = 3,500 vs 3,000 (positive: 3,500 has the higher loss), every bar, Lasso, full predictor set
\providecommand{\rhrbPLassoFullBar}{}\renewcommand{\rhrbPLassoFullBar}{\ensuremath{p = 0.042}}  % ... its p statement
\providecommand{\rhrbPctLassoFullClose}{}\renewcommand{\rhrbPctLassoFullClose}{\ensuremath{-}4.4}  % QLIKE change, %, W = 3,500 vs 3,000, 16:00 bar, Lasso, full predictor set
\providecommand{\rhrbDMLassoFullClose}{}\renewcommand{\rhrbDMLassoFullClose}{\ensuremath{-}2.89}  % DM, W = 3,500 vs 3,000 (positive: 3,500 has the higher loss), 16:00 bar, Lasso, full predictor set
\providecommand{\rhrbPLassoFullClose}{}\renewcommand{\rhrbPLassoFullClose}{\ensuremath{p = 0.004}}  % ... its p statement
\providecommand{\rhrbDSRLassoFullMid}{}\renewcommand{\rhrbDSRLassoFullMid}{0.07}  % Sharpe ratio change, W = 3,500 minus 3,000, mid-quotes, Lasso, full predictor set
\providecommand{\rhrbPSRLassoFullMid}{}\renewcommand{\rhrbPSRLassoFullMid}{\ensuremath{p = 0.83}}  % ... its p statement (Ledoit-Wolf, D19; was the D6 bootstrap)
\providecommand{\rhrbDSRLassoFullBa}{}\renewcommand{\rhrbDSRLassoFullBa}{0.07}  % Sharpe ratio change, W = 3,500 minus 3,000, bid/ask prices, Lasso, full predictor set
\providecommand{\rhrbPSRLassoFullBa}{}\renewcommand{\rhrbPSRLassoFullBa}{\ensuremath{p = 0.83}}  % ... its p statement (Ledoit-Wolf, D19; was the D6 bootstrap)
\providecommand{\rhrbQEnetHarBarMain}{}\renewcommand{\rhrbQEnetHarBarMain}{0.1893}  % QLIKE, every bar, Elastic net, HAR benchmark, W = 3,000
\providecommand{\rhrbQEnetHarCloseMain}{}\renewcommand{\rhrbQEnetHarCloseMain}{0.1269}  % QLIKE, 16:00 bar, Elastic net, HAR benchmark, W = 3,000
\providecommand{\rhrbSREnetHarMidMain}{}\renewcommand{\rhrbSREnetHarMidMain}{0.95}  % Sharpe ratio (D19 pooled), sign strategy at mid-quotes, Elastic net, HAR benchmark, W = 3,000
\providecommand{\rhrbSREnetHarBaMain}{}\renewcommand{\rhrbSREnetHarBaMain}{0.60}  % Sharpe ratio (D19 pooled), sign strategy at bid/ask prices, Elastic net, HAR benchmark, W = 3,000
\providecommand{\rhrbQEnetHarBarLong}{}\renewcommand{\rhrbQEnetHarBarLong}{0.1897}  % QLIKE, every bar, Elastic net, HAR benchmark, W = 3,500
\providecommand{\rhrbQEnetHarCloseLong}{}\renewcommand{\rhrbQEnetHarCloseLong}{0.1272}  % QLIKE, 16:00 bar, Elastic net, HAR benchmark, W = 3,500
\providecommand{\rhrbSREnetHarMidLong}{}\renewcommand{\rhrbSREnetHarMidLong}{1.01}  % Sharpe ratio (D19 pooled), sign strategy at mid-quotes, Elastic net, HAR benchmark, W = 3,500
\providecommand{\rhrbSREnetHarBaLong}{}\renewcommand{\rhrbSREnetHarBaLong}{0.65}  % Sharpe ratio (D19 pooled), sign strategy at bid/ask prices, Elastic net, HAR benchmark, W = 3,500
\providecommand{\rhrbPctEnetHarBar}{}\renewcommand{\rhrbPctEnetHarBar}{0.2}  % QLIKE change, %, W = 3,500 vs 3,000, every bar, Elastic net, HAR benchmark
\providecommand{\rhrbDMEnetHarBar}{}\renewcommand{\rhrbDMEnetHarBar}{2.48}  % DM, W = 3,500 vs 3,000 (positive: 3,500 has the higher loss), every bar, Elastic net, HAR benchmark
\providecommand{\rhrbPEnetHarBar}{}\renewcommand{\rhrbPEnetHarBar}{\ensuremath{p = 0.013}}  % ... its p statement
\providecommand{\rhrbPctEnetHarClose}{}\renewcommand{\rhrbPctEnetHarClose}{0.3}  % QLIKE change, %, W = 3,500 vs 3,000, 16:00 bar, Elastic net, HAR benchmark
\providecommand{\rhrbDMEnetHarClose}{}\renewcommand{\rhrbDMEnetHarClose}{1.31}  % DM, W = 3,500 vs 3,000 (positive: 3,500 has the higher loss), 16:00 bar, Elastic net, HAR benchmark
\providecommand{\rhrbPEnetHarClose}{}\renewcommand{\rhrbPEnetHarClose}{\ensuremath{p = 0.19}}  % ... its p statement
\providecommand{\rhrbDSREnetHarMid}{}\renewcommand{\rhrbDSREnetHarMid}{0.05}  % Sharpe ratio change, W = 3,500 minus 3,000, mid-quotes, Elastic net, HAR benchmark
\providecommand{\rhrbPSREnetHarMid}{}\renewcommand{\rhrbPSREnetHarMid}{\ensuremath{p = 0.58}}  % ... its p statement (Ledoit-Wolf, D19; was the D6 bootstrap)
\providecommand{\rhrbDSREnetHarBa}{}\renewcommand{\rhrbDSREnetHarBa}{0.05}  % Sharpe ratio change, W = 3,500 minus 3,000, bid/ask prices, Elastic net, HAR benchmark
\providecommand{\rhrbPSREnetHarBa}{}\renewcommand{\rhrbPSREnetHarBa}{\ensuremath{p = 0.58}}  % ... its p statement (Ledoit-Wolf, D19; was the D6 bootstrap)
\providecommand{\rhrbQEnetFullBarMain}{}\renewcommand{\rhrbQEnetFullBarMain}{0.1778}  % QLIKE, every bar, Elastic net, full predictor set, W = 3,000
\providecommand{\rhrbQEnetFullCloseMain}{}\renewcommand{\rhrbQEnetFullCloseMain}{0.1262}  % QLIKE, 16:00 bar, Elastic net, full predictor set, W = 3,000
\providecommand{\rhrbSREnetFullMidMain}{}\renewcommand{\rhrbSREnetFullMidMain}{1.86}  % Sharpe ratio (D19 pooled), sign strategy at mid-quotes, Elastic net, full predictor set, W = 3,000
\providecommand{\rhrbSREnetFullBaMain}{}\renewcommand{\rhrbSREnetFullBaMain}{1.50}  % Sharpe ratio (D19 pooled), sign strategy at bid/ask prices, Elastic net, full predictor set, W = 3,000
\providecommand{\rhrbQEnetFullBarLong}{}\renewcommand{\rhrbQEnetFullBarLong}{0.1770}  % QLIKE, every bar, Elastic net, full predictor set, W = 3,500
\providecommand{\rhrbQEnetFullCloseLong}{}\renewcommand{\rhrbQEnetFullCloseLong}{0.1223}  % QLIKE, 16:00 bar, Elastic net, full predictor set, W = 3,500
\providecommand{\rhrbSREnetFullMidLong}{}\renewcommand{\rhrbSREnetFullMidLong}{2.00}  % Sharpe ratio (D19 pooled), sign strategy at mid-quotes, Elastic net, full predictor set, W = 3,500
\providecommand{\rhrbSREnetFullBaLong}{}\renewcommand{\rhrbSREnetFullBaLong}{1.65}  % Sharpe ratio (D19 pooled), sign strategy at bid/ask prices, Elastic net, full predictor set, W = 3,500
\providecommand{\rhrbPctEnetFullBar}{}\renewcommand{\rhrbPctEnetFullBar}{\ensuremath{-}0.5}  % QLIKE change, %, W = 3,500 vs 3,000, every bar, Elastic net, full predictor set
\providecommand{\rhrbDMEnetFullBar}{}\renewcommand{\rhrbDMEnetFullBar}{\ensuremath{-}1.47}  % DM, W = 3,500 vs 3,000 (positive: 3,500 has the higher loss), every bar, Elastic net, full predictor set
\providecommand{\rhrbPEnetFullBar}{}\renewcommand{\rhrbPEnetFullBar}{\ensuremath{p = 0.14}}  % ... its p statement
\providecommand{\rhrbPctEnetFullClose}{}\renewcommand{\rhrbPctEnetFullClose}{\ensuremath{-}3.1}  % QLIKE change, %, W = 3,500 vs 3,000, 16:00 bar, Elastic net, full predictor set
\providecommand{\rhrbDMEnetFullClose}{}\renewcommand{\rhrbDMEnetFullClose}{\ensuremath{-}2.93}  % DM, W = 3,500 vs 3,000 (positive: 3,500 has the higher loss), 16:00 bar, Elastic net, full predictor set
\providecommand{\rhrbPEnetFullClose}{}\renewcommand{\rhrbPEnetFullClose}{\ensuremath{p = 0.003}}  % ... its p statement
\providecommand{\rhrbDSREnetFullMid}{}\renewcommand{\rhrbDSREnetFullMid}{0.15}  % Sharpe ratio change, W = 3,500 minus 3,000, mid-quotes, Elastic net, full predictor set
\providecommand{\rhrbPSREnetFullMid}{}\renewcommand{\rhrbPSREnetFullMid}{\ensuremath{p = 0.56}}  % ... its p statement (Ledoit-Wolf, D19; was the D6 bootstrap)
\providecommand{\rhrbDSREnetFullBa}{}\renewcommand{\rhrbDSREnetFullBa}{0.15}  % Sharpe ratio change, W = 3,500 minus 3,000, bid/ask prices, Elastic net, full predictor set
\providecommand{\rhrbPSREnetFullBa}{}\renewcommand{\rhrbPSREnetFullBa}{\ensuremath{p = 0.55}}  % ... its p statement (Ledoit-Wolf, D19; was the D6 bootstrap)
\providecommand{\rhrbQLgbmHarBarMain}{}\renewcommand{\rhrbQLgbmHarBarMain}{0.1875}  % QLIKE, every bar, LightGBM, HAR benchmark, W = 3,000
\providecommand{\rhrbQLgbmHarCloseMain}{}\renewcommand{\rhrbQLgbmHarCloseMain}{0.1235}  % QLIKE, 16:00 bar, LightGBM, HAR benchmark, W = 3,000
\providecommand{\rhrbSRLgbmHarMidMain}{}\renewcommand{\rhrbSRLgbmHarMidMain}{0.65}  % Sharpe ratio (D19 pooled), sign strategy at mid-quotes, LightGBM, HAR benchmark, W = 3,000
\providecommand{\rhrbSRLgbmHarBaMain}{}\renewcommand{\rhrbSRLgbmHarBaMain}{0.31}  % Sharpe ratio (D19 pooled), sign strategy at bid/ask prices, LightGBM, HAR benchmark, W = 3,000
\providecommand{\rhrbQLgbmHarBarLong}{}\renewcommand{\rhrbQLgbmHarBarLong}{0.1876}  % QLIKE, every bar, LightGBM, HAR benchmark, W = 3,500
\providecommand{\rhrbQLgbmHarCloseLong}{}\renewcommand{\rhrbQLgbmHarCloseLong}{0.1233}  % QLIKE, 16:00 bar, LightGBM, HAR benchmark, W = 3,500
\providecommand{\rhrbSRLgbmHarMidLong}{}\renewcommand{\rhrbSRLgbmHarMidLong}{0.59}  % Sharpe ratio (D19 pooled), sign strategy at mid-quotes, LightGBM, HAR benchmark, W = 3,500
\providecommand{\rhrbSRLgbmHarBaLong}{}\renewcommand{\rhrbSRLgbmHarBaLong}{0.24}  % Sharpe ratio (D19 pooled), sign strategy at bid/ask prices, LightGBM, HAR benchmark, W = 3,500
\providecommand{\rhrbPctLgbmHarBar}{}\renewcommand{\rhrbPctLgbmHarBar}{0.0}  % QLIKE change, %, W = 3,500 vs 3,000, every bar, LightGBM, HAR benchmark
\providecommand{\rhrbDMLgbmHarBar}{}\renewcommand{\rhrbDMLgbmHarBar}{0.20}  % DM, W = 3,500 vs 3,000 (positive: 3,500 has the higher loss), every bar, LightGBM, HAR benchmark
\providecommand{\rhrbPLgbmHarBar}{}\renewcommand{\rhrbPLgbmHarBar}{\ensuremath{p = 0.84}}  % ... its p statement
\providecommand{\rhrbPctLgbmHarClose}{}\renewcommand{\rhrbPctLgbmHarClose}{\ensuremath{-}0.2}  % QLIKE change, %, W = 3,500 vs 3,000, 16:00 bar, LightGBM, HAR benchmark
\providecommand{\rhrbDMLgbmHarClose}{}\renewcommand{\rhrbDMLgbmHarClose}{\ensuremath{-}0.45}  % DM, W = 3,500 vs 3,000 (positive: 3,500 has the higher loss), 16:00 bar, LightGBM, HAR benchmark
\providecommand{\rhrbPLgbmHarClose}{}\renewcommand{\rhrbPLgbmHarClose}{\ensuremath{p = 0.65}}  % ... its p statement
\providecommand{\rhrbDSRLgbmHarMid}{}\renewcommand{\rhrbDSRLgbmHarMid}{\ensuremath{-}0.06}  % Sharpe ratio change, W = 3,500 minus 3,000, mid-quotes, LightGBM, HAR benchmark
\providecommand{\rhrbPSRLgbmHarMid}{}\renewcommand{\rhrbPSRLgbmHarMid}{\ensuremath{p = 0.57}}  % ... its p statement (Ledoit-Wolf, D19; was the D6 bootstrap)
\providecommand{\rhrbDSRLgbmHarBa}{}\renewcommand{\rhrbDSRLgbmHarBa}{\ensuremath{-}0.06}  % Sharpe ratio change, W = 3,500 minus 3,000, bid/ask prices, LightGBM, HAR benchmark
\providecommand{\rhrbPSRLgbmHarBa}{}\renewcommand{\rhrbPSRLgbmHarBa}{\ensuremath{p = 0.57}}  % ... its p statement (Ledoit-Wolf, D19; was the D6 bootstrap)
\providecommand{\rhrbQLgbmFullBarMain}{}\renewcommand{\rhrbQLgbmFullBarMain}{0.1613}  % QLIKE, every bar, LightGBM, full predictor set, W = 3,000
\providecommand{\rhrbQLgbmFullCloseMain}{}\renewcommand{\rhrbQLgbmFullCloseMain}{0.1201}  % QLIKE, 16:00 bar, LightGBM, full predictor set, W = 3,000
\providecommand{\rhrbSRLgbmFullMidMain}{}\renewcommand{\rhrbSRLgbmFullMidMain}{1.59}  % Sharpe ratio (D19 pooled), sign strategy at mid-quotes, LightGBM, full predictor set, W = 3,000
\providecommand{\rhrbSRLgbmFullBaMain}{}\renewcommand{\rhrbSRLgbmFullBaMain}{1.23}  % Sharpe ratio (D19 pooled), sign strategy at bid/ask prices, LightGBM, full predictor set, W = 3,000
\providecommand{\rhrbQLgbmFullBarLong}{}\renewcommand{\rhrbQLgbmFullBarLong}{0.1613}  % QLIKE, every bar, LightGBM, full predictor set, W = 3,500
\providecommand{\rhrbQLgbmFullCloseLong}{}\renewcommand{\rhrbQLgbmFullCloseLong}{0.1204}  % QLIKE, 16:00 bar, LightGBM, full predictor set, W = 3,500
\providecommand{\rhrbSRLgbmFullMidLong}{}\renewcommand{\rhrbSRLgbmFullMidLong}{1.54}  % Sharpe ratio (D19 pooled), sign strategy at mid-quotes, LightGBM, full predictor set, W = 3,500
\providecommand{\rhrbSRLgbmFullBaLong}{}\renewcommand{\rhrbSRLgbmFullBaLong}{1.19}  % Sharpe ratio (D19 pooled), sign strategy at bid/ask prices, LightGBM, full predictor set, W = 3,500
\providecommand{\rhrbPctLgbmFullBar}{}\renewcommand{\rhrbPctLgbmFullBar}{0.0}  % QLIKE change, %, W = 3,500 vs 3,000, every bar, LightGBM, full predictor set
\providecommand{\rhrbDMLgbmFullBar}{}\renewcommand{\rhrbDMLgbmFullBar}{0.20}  % DM, W = 3,500 vs 3,000 (positive: 3,500 has the higher loss), every bar, LightGBM, full predictor set
\providecommand{\rhrbPLgbmFullBar}{}\renewcommand{\rhrbPLgbmFullBar}{\ensuremath{p = 0.84}}  % ... its p statement
\providecommand{\rhrbPctLgbmFullClose}{}\renewcommand{\rhrbPctLgbmFullClose}{0.2}  % QLIKE change, %, W = 3,500 vs 3,000, 16:00 bar, LightGBM, full predictor set
\providecommand{\rhrbDMLgbmFullClose}{}\renewcommand{\rhrbDMLgbmFullClose}{0.55}  % DM, W = 3,500 vs 3,000 (positive: 3,500 has the higher loss), 16:00 bar, LightGBM, full predictor set
\providecommand{\rhrbPLgbmFullClose}{}\renewcommand{\rhrbPLgbmFullClose}{\ensuremath{p = 0.58}}  % ... its p statement
\providecommand{\rhrbDSRLgbmFullMid}{}\renewcommand{\rhrbDSRLgbmFullMid}{\ensuremath{-}0.05}  % Sharpe ratio change, W = 3,500 minus 3,000, mid-quotes, LightGBM, full predictor set
\providecommand{\rhrbPSRLgbmFullMid}{}\renewcommand{\rhrbPSRLgbmFullMid}{\ensuremath{p = 0.76}}  % ... its p statement (Ledoit-Wolf, D19; was the D6 bootstrap)
\providecommand{\rhrbDSRLgbmFullBa}{}\renewcommand{\rhrbDSRLgbmFullBa}{\ensuremath{-}0.05}  % Sharpe ratio change, W = 3,500 minus 3,000, bid/ask prices, LightGBM, full predictor set
\providecommand{\rhrbPSRLgbmFullBa}{}\renewcommand{\rhrbPSRLgbmFullBa}{\ensuremath{p = 0.76}}  % ... its p statement (Ledoit-Wolf, D19; was the D6 bootstrap)
\providecommand{\rhrbQXgbHarBarMain}{}\renewcommand{\rhrbQXgbHarBarMain}{0.1888}  % QLIKE, every bar, XGBoost, HAR benchmark, W = 3,000
\providecommand{\rhrbQXgbHarCloseMain}{}\renewcommand{\rhrbQXgbHarCloseMain}{0.1234}  % QLIKE, 16:00 bar, XGBoost, HAR benchmark, W = 3,000
\providecommand{\rhrbSRXgbHarMidMain}{}\renewcommand{\rhrbSRXgbHarMidMain}{0.98}  % Sharpe ratio (D19 pooled), sign strategy at mid-quotes, XGBoost, HAR benchmark, W = 3,000
\providecommand{\rhrbSRXgbHarBaMain}{}\renewcommand{\rhrbSRXgbHarBaMain}{0.63}  % Sharpe ratio (D19 pooled), sign strategy at bid/ask prices, XGBoost, HAR benchmark, W = 3,000
\providecommand{\rhrbQXgbHarBarLong}{}\renewcommand{\rhrbQXgbHarBarLong}{0.1890}  % QLIKE, every bar, XGBoost, HAR benchmark, W = 3,500
\providecommand{\rhrbQXgbHarCloseLong}{}\renewcommand{\rhrbQXgbHarCloseLong}{0.1232}  % QLIKE, 16:00 bar, XGBoost, HAR benchmark, W = 3,500
\providecommand{\rhrbSRXgbHarMidLong}{}\renewcommand{\rhrbSRXgbHarMidLong}{0.78}  % Sharpe ratio (D19 pooled), sign strategy at mid-quotes, XGBoost, HAR benchmark, W = 3,500
\providecommand{\rhrbSRXgbHarBaLong}{}\renewcommand{\rhrbSRXgbHarBaLong}{0.43}  % Sharpe ratio (D19 pooled), sign strategy at bid/ask prices, XGBoost, HAR benchmark, W = 3,500
\providecommand{\rhrbPctXgbHarBar}{}\renewcommand{\rhrbPctXgbHarBar}{0.1}  % QLIKE change, %, W = 3,500 vs 3,000, every bar, XGBoost, HAR benchmark
\providecommand{\rhrbDMXgbHarBar}{}\renewcommand{\rhrbDMXgbHarBar}{1.13}  % DM, W = 3,500 vs 3,000 (positive: 3,500 has the higher loss), every bar, XGBoost, HAR benchmark
\providecommand{\rhrbPXgbHarBar}{}\renewcommand{\rhrbPXgbHarBar}{\ensuremath{p = 0.26}}  % ... its p statement
\providecommand{\rhrbPctXgbHarClose}{}\renewcommand{\rhrbPctXgbHarClose}{\ensuremath{-}0.2}  % QLIKE change, %, W = 3,500 vs 3,000, 16:00 bar, XGBoost, HAR benchmark
\providecommand{\rhrbDMXgbHarClose}{}\renewcommand{\rhrbDMXgbHarClose}{\ensuremath{-}0.55}  % DM, W = 3,500 vs 3,000 (positive: 3,500 has the higher loss), 16:00 bar, XGBoost, HAR benchmark
\providecommand{\rhrbPXgbHarClose}{}\renewcommand{\rhrbPXgbHarClose}{\ensuremath{p = 0.58}}  % ... its p statement
\providecommand{\rhrbDSRXgbHarMid}{}\renewcommand{\rhrbDSRXgbHarMid}{\ensuremath{-}0.20}  % Sharpe ratio change, W = 3,500 minus 3,000, mid-quotes, XGBoost, HAR benchmark
\providecommand{\rhrbPSRXgbHarMid}{}\renewcommand{\rhrbPSRXgbHarMid}{\ensuremath{p = 0.40}}  % ... its p statement (Ledoit-Wolf, D19; was the D6 bootstrap)
\providecommand{\rhrbDSRXgbHarBa}{}\renewcommand{\rhrbDSRXgbHarBa}{\ensuremath{-}0.20}  % Sharpe ratio change, W = 3,500 minus 3,000, bid/ask prices, XGBoost, HAR benchmark
\providecommand{\rhrbPSRXgbHarBa}{}\renewcommand{\rhrbPSRXgbHarBa}{\ensuremath{p = 0.40}}  % ... its p statement (Ledoit-Wolf, D19; was the D6 bootstrap)
\providecommand{\rhrbQXgbFullBarMain}{}\renewcommand{\rhrbQXgbFullBarMain}{0.1627}  % QLIKE, every bar, XGBoost, full predictor set, W = 3,000
\providecommand{\rhrbQXgbFullCloseMain}{}\renewcommand{\rhrbQXgbFullCloseMain}{0.1192}  % QLIKE, 16:00 bar, XGBoost, full predictor set, W = 3,000
\providecommand{\rhrbSRXgbFullMidMain}{}\renewcommand{\rhrbSRXgbFullMidMain}{1.47}  % Sharpe ratio (D19 pooled), sign strategy at mid-quotes, XGBoost, full predictor set, W = 3,000
\providecommand{\rhrbSRXgbFullBaMain}{}\renewcommand{\rhrbSRXgbFullBaMain}{1.12}  % Sharpe ratio (D19 pooled), sign strategy at bid/ask prices, XGBoost, full predictor set, W = 3,000
\providecommand{\rhrbQXgbFullBarLong}{}\renewcommand{\rhrbQXgbFullBarLong}{0.1636}  % QLIKE, every bar, XGBoost, full predictor set, W = 3,500
\providecommand{\rhrbQXgbFullCloseLong}{}\renewcommand{\rhrbQXgbFullCloseLong}{0.1192}  % QLIKE, 16:00 bar, XGBoost, full predictor set, W = 3,500
\providecommand{\rhrbSRXgbFullMidLong}{}\renewcommand{\rhrbSRXgbFullMidLong}{1.35}  % Sharpe ratio (D19 pooled), sign strategy at mid-quotes, XGBoost, full predictor set, W = 3,500
\providecommand{\rhrbSRXgbFullBaLong}{}\renewcommand{\rhrbSRXgbFullBaLong}{0.99}  % Sharpe ratio (D19 pooled), sign strategy at bid/ask prices, XGBoost, full predictor set, W = 3,500
\providecommand{\rhrbPctXgbFullBar}{}\renewcommand{\rhrbPctXgbFullBar}{0.5}  % QLIKE change, %, W = 3,500 vs 3,000, every bar, XGBoost, full predictor set
\providecommand{\rhrbDMXgbFullBar}{}\renewcommand{\rhrbDMXgbFullBar}{1.87}  % DM, W = 3,500 vs 3,000 (positive: 3,500 has the higher loss), every bar, XGBoost, full predictor set
\providecommand{\rhrbPXgbFullBar}{}\renewcommand{\rhrbPXgbFullBar}{\ensuremath{p = 0.062}}  % ... its p statement
\providecommand{\rhrbPctXgbFullClose}{}\renewcommand{\rhrbPctXgbFullClose}{\ensuremath{-}0.1}  % QLIKE change, %, W = 3,500 vs 3,000, 16:00 bar, XGBoost, full predictor set
\providecommand{\rhrbDMXgbFullClose}{}\renewcommand{\rhrbDMXgbFullClose}{\ensuremath{-}0.15}  % DM, W = 3,500 vs 3,000 (positive: 3,500 has the higher loss), 16:00 bar, XGBoost, full predictor set
\providecommand{\rhrbPXgbFullClose}{}\renewcommand{\rhrbPXgbFullClose}{\ensuremath{p = 0.88}}  % ... its p statement
\providecommand{\rhrbDSRXgbFullMid}{}\renewcommand{\rhrbDSRXgbFullMid}{\ensuremath{-}0.12}  % Sharpe ratio change, W = 3,500 minus 3,000, mid-quotes, XGBoost, full predictor set
\providecommand{\rhrbPSRXgbFullMid}{}\renewcommand{\rhrbPSRXgbFullMid}{\ensuremath{p = 0.62}}  % ... its p statement (Ledoit-Wolf, D19; was the D6 bootstrap)
\providecommand{\rhrbDSRXgbFullBa}{}\renewcommand{\rhrbDSRXgbFullBa}{\ensuremath{-}0.12}  % Sharpe ratio change, W = 3,500 minus 3,000, bid/ask prices, XGBoost, full predictor set
\providecommand{\rhrbPSRXgbFullBa}{}\renewcommand{\rhrbPSRXgbFullBa}{\ensuremath{p = 0.62}}  % ... its p statement (Ledoit-Wolf, D19; was the D6 bootstrap)
\providecommand{\rhrbQRfHarBarMain}{}\renewcommand{\rhrbQRfHarBarMain}{0.1895}  % QLIKE, every bar, Random forest, HAR benchmark, W = 3,000
\providecommand{\rhrbQRfHarCloseMain}{}\renewcommand{\rhrbQRfHarCloseMain}{0.1266}  % QLIKE, 16:00 bar, Random forest, HAR benchmark, W = 3,000
\providecommand{\rhrbSRRfHarMidMain}{}\renewcommand{\rhrbSRRfHarMidMain}{0.32}  % Sharpe ratio (D19 pooled), sign strategy at mid-quotes, Random forest, HAR benchmark, W = 3,000
\providecommand{\rhrbSRRfHarBaMain}{}\renewcommand{\rhrbSRRfHarBaMain}{\ensuremath{-}0.03}  % Sharpe ratio (D19 pooled), sign strategy at bid/ask prices, Random forest, HAR benchmark, W = 3,000
\providecommand{\rhrbQRfHarBarLong}{}\renewcommand{\rhrbQRfHarBarLong}{0.1898}  % QLIKE, every bar, Random forest, HAR benchmark, W = 3,500
\providecommand{\rhrbQRfHarCloseLong}{}\renewcommand{\rhrbQRfHarCloseLong}{0.1264}  % QLIKE, 16:00 bar, Random forest, HAR benchmark, W = 3,500
\providecommand{\rhrbSRRfHarMidLong}{}\renewcommand{\rhrbSRRfHarMidLong}{0.80}  % Sharpe ratio (D19 pooled), sign strategy at mid-quotes, Random forest, HAR benchmark, W = 3,500
\providecommand{\rhrbSRRfHarBaLong}{}\renewcommand{\rhrbSRRfHarBaLong}{0.45}  % Sharpe ratio (D19 pooled), sign strategy at bid/ask prices, Random forest, HAR benchmark, W = 3,500
\providecommand{\rhrbPctRfHarBar}{}\renewcommand{\rhrbPctRfHarBar}{0.1}  % QLIKE change, %, W = 3,500 vs 3,000, every bar, Random forest, HAR benchmark
\providecommand{\rhrbDMRfHarBar}{}\renewcommand{\rhrbDMRfHarBar}{1.15}  % DM, W = 3,500 vs 3,000 (positive: 3,500 has the higher loss), every bar, Random forest, HAR benchmark
\providecommand{\rhrbPRfHarBar}{}\renewcommand{\rhrbPRfHarBar}{\ensuremath{p = 0.25}}  % ... its p statement
\providecommand{\rhrbPctRfHarClose}{}\renewcommand{\rhrbPctRfHarClose}{\ensuremath{-}0.1}  % QLIKE change, %, W = 3,500 vs 3,000, 16:00 bar, Random forest, HAR benchmark
\providecommand{\rhrbDMRfHarClose}{}\renewcommand{\rhrbDMRfHarClose}{\ensuremath{-}0.36}  % DM, W = 3,500 vs 3,000 (positive: 3,500 has the higher loss), 16:00 bar, Random forest, HAR benchmark
\providecommand{\rhrbPRfHarClose}{}\renewcommand{\rhrbPRfHarClose}{\ensuremath{p = 0.72}}  % ... its p statement
\providecommand{\rhrbDSRRfHarMid}{}\renewcommand{\rhrbDSRRfHarMid}{0.48}  % Sharpe ratio change, W = 3,500 minus 3,000, mid-quotes, Random forest, HAR benchmark
\providecommand{\rhrbPSRRfHarMid}{}\renewcommand{\rhrbPSRRfHarMid}{\ensuremath{p = 0.42}}  % ... its p statement (Ledoit-Wolf, D19; was the D6 bootstrap)
\providecommand{\rhrbDSRRfHarBa}{}\renewcommand{\rhrbDSRRfHarBa}{0.48}  % Sharpe ratio change, W = 3,500 minus 3,000, bid/ask prices, Random forest, HAR benchmark
\providecommand{\rhrbPSRRfHarBa}{}\renewcommand{\rhrbPSRRfHarBa}{\ensuremath{p = 0.42}}  % ... its p statement (Ledoit-Wolf, D19; was the D6 bootstrap)
\providecommand{\rhrbQRfFullBarMain}{}\renewcommand{\rhrbQRfFullBarMain}{0.1685}  % QLIKE, every bar, Random forest, full predictor set, W = 3,000
\providecommand{\rhrbQRfFullCloseMain}{}\renewcommand{\rhrbQRfFullCloseMain}{0.1248}  % QLIKE, 16:00 bar, Random forest, full predictor set, W = 3,000
\providecommand{\rhrbSRRfFullMidMain}{}\renewcommand{\rhrbSRRfFullMidMain}{1.88}  % Sharpe ratio (D19 pooled), sign strategy at mid-quotes, Random forest, full predictor set, W = 3,000
\providecommand{\rhrbSRRfFullBaMain}{}\renewcommand{\rhrbSRRfFullBaMain}{1.52}  % Sharpe ratio (D19 pooled), sign strategy at bid/ask prices, Random forest, full predictor set, W = 3,000
\providecommand{\rhrbQRfFullBarLong}{}\renewcommand{\rhrbQRfFullBarLong}{0.1685}  % QLIKE, every bar, Random forest, full predictor set, W = 3,500
\providecommand{\rhrbQRfFullCloseLong}{}\renewcommand{\rhrbQRfFullCloseLong}{0.1241}  % QLIKE, 16:00 bar, Random forest, full predictor set, W = 3,500
\providecommand{\rhrbSRRfFullMidLong}{}\renewcommand{\rhrbSRRfFullMidLong}{1.83}  % Sharpe ratio (D19 pooled), sign strategy at mid-quotes, Random forest, full predictor set, W = 3,500
\providecommand{\rhrbSRRfFullBaLong}{}\renewcommand{\rhrbSRRfFullBaLong}{1.47}  % Sharpe ratio (D19 pooled), sign strategy at bid/ask prices, Random forest, full predictor set, W = 3,500
\providecommand{\rhrbPctRfFullBar}{}\renewcommand{\rhrbPctRfFullBar}{0.0}  % QLIKE change, %, W = 3,500 vs 3,000, every bar, Random forest, full predictor set
\providecommand{\rhrbDMRfFullBar}{}\renewcommand{\rhrbDMRfFullBar}{0.02}  % DM, W = 3,500 vs 3,000 (positive: 3,500 has the higher loss), every bar, Random forest, full predictor set
\providecommand{\rhrbPRfFullBar}{}\renewcommand{\rhrbPRfFullBar}{\ensuremath{p = 0.99}}  % ... its p statement
\providecommand{\rhrbPctRfFullClose}{}\renewcommand{\rhrbPctRfFullClose}{\ensuremath{-}0.5}  % QLIKE change, %, W = 3,500 vs 3,000, 16:00 bar, Random forest, full predictor set
\providecommand{\rhrbDMRfFullClose}{}\renewcommand{\rhrbDMRfFullClose}{\ensuremath{-}1.43}  % DM, W = 3,500 vs 3,000 (positive: 3,500 has the higher loss), 16:00 bar, Random forest, full predictor set
\providecommand{\rhrbPRfFullClose}{}\renewcommand{\rhrbPRfFullClose}{\ensuremath{p = 0.15}}  % ... its p statement
\providecommand{\rhrbDSRRfFullMid}{}\renewcommand{\rhrbDSRRfFullMid}{\ensuremath{-}0.05}  % Sharpe ratio change, W = 3,500 minus 3,000, mid-quotes, Random forest, full predictor set
\providecommand{\rhrbPSRRfFullMid}{}\renewcommand{\rhrbPSRRfFullMid}{\ensuremath{p = 0.58}}  % ... its p statement (Ledoit-Wolf, D19; was the D6 bootstrap)
\providecommand{\rhrbDSRRfFullBa}{}\renewcommand{\rhrbDSRRfFullBa}{\ensuremath{-}0.05}  % Sharpe ratio change, W = 3,500 minus 3,000, bid/ask prices, Random forest, full predictor set
\providecommand{\rhrbPSRRfFullBa}{}\renewcommand{\rhrbPSRRfFullBa}{\ensuremath{p = 0.59}}  % ... its p statement (Ledoit-Wolf, D19; was the D6 bootstrap)
\providecommand{\rhrbPctFRidgeBarMain}{}\renewcommand{\rhrbPctFRidgeBarMain}{\ensuremath{-}5.4}  % QLIKE, full set vs HAR benchmark, %, every bar, Ridge, W = 3,000
\providecommand{\rhrbDMFRidgeBarMain}{}\renewcommand{\rhrbDMFRidgeBarMain}{\ensuremath{-}5.63}  % DM, full set vs HAR benchmark (negative: full more accurate), every bar, Ridge, W = 3,000
\providecommand{\rhrbPFRidgeBarMain}{}\renewcommand{\rhrbPFRidgeBarMain}{\ensuremath{p < 0.001}}  % ... its p statement
\providecommand{\rhrbPctFRidgeCloseMain}{}\renewcommand{\rhrbPctFRidgeCloseMain}{0.8}  % QLIKE, full set vs HAR benchmark, %, 16:00 bar, Ridge, W = 3,000
\providecommand{\rhrbDMFRidgeCloseMain}{}\renewcommand{\rhrbDMFRidgeCloseMain}{0.31}  % DM, full set vs HAR benchmark (negative: full more accurate), 16:00 bar, Ridge, W = 3,000
\providecommand{\rhrbPFRidgeCloseMain}{}\renewcommand{\rhrbPFRidgeCloseMain}{\ensuremath{p = 0.76}}  % ... its p statement
\providecommand{\rhrbDSRFRidgeMidMain}{}\renewcommand{\rhrbDSRFRidgeMidMain}{0.82}  % Sharpe ratio, full set minus HAR benchmark, mid-quotes, Ridge, W = 3,000
\providecommand{\rhrbPSRFRidgeMidMain}{}\renewcommand{\rhrbPSRFRidgeMidMain}{\ensuremath{p = 0.11}}  % ... its p statement (Ledoit-Wolf, D19; was the D6 bootstrap)
\providecommand{\rhrbDSRFRidgeBaMain}{}\renewcommand{\rhrbDSRFRidgeBaMain}{0.81}  % Sharpe ratio, full set minus HAR benchmark, bid/ask prices, Ridge, W = 3,000
\providecommand{\rhrbPSRFRidgeBaMain}{}\renewcommand{\rhrbPSRFRidgeBaMain}{\ensuremath{p = 0.11}}  % ... its p statement (Ledoit-Wolf, D19; was the D6 bootstrap)
\providecommand{\rhrbPctFRidgeBarLong}{}\renewcommand{\rhrbPctFRidgeBarLong}{\ensuremath{-}6.3}  % QLIKE, full set vs HAR benchmark, %, every bar, Ridge, W = 3,500
\providecommand{\rhrbDMFRidgeBarLong}{}\renewcommand{\rhrbDMFRidgeBarLong}{\ensuremath{-}7.06}  % DM, full set vs HAR benchmark (negative: full more accurate), every bar, Ridge, W = 3,500
\providecommand{\rhrbPFRidgeBarLong}{}\renewcommand{\rhrbPFRidgeBarLong}{\ensuremath{p < 0.001}}  % ... its p statement
\providecommand{\rhrbPctFRidgeCloseLong}{}\renewcommand{\rhrbPctFRidgeCloseLong}{\ensuremath{-}3.1}  % QLIKE, full set vs HAR benchmark, %, 16:00 bar, Ridge, W = 3,500
\providecommand{\rhrbDMFRidgeCloseLong}{}\renewcommand{\rhrbDMFRidgeCloseLong}{\ensuremath{-}1.55}  % DM, full set vs HAR benchmark (negative: full more accurate), 16:00 bar, Ridge, W = 3,500
\providecommand{\rhrbPFRidgeCloseLong}{}\renewcommand{\rhrbPFRidgeCloseLong}{\ensuremath{p = 0.12}}  % ... its p statement
\providecommand{\rhrbDSRFRidgeMidLong}{}\renewcommand{\rhrbDSRFRidgeMidLong}{0.83}  % Sharpe ratio, full set minus HAR benchmark, mid-quotes, Ridge, W = 3,500
\providecommand{\rhrbPSRFRidgeMidLong}{}\renewcommand{\rhrbPSRFRidgeMidLong}{\ensuremath{p = 0.069}}  % ... its p statement (Ledoit-Wolf, D19; was the D6 bootstrap)
\providecommand{\rhrbDSRFRidgeBaLong}{}\renewcommand{\rhrbDSRFRidgeBaLong}{0.83}  % Sharpe ratio, full set minus HAR benchmark, bid/ask prices, Ridge, W = 3,500
\providecommand{\rhrbPSRFRidgeBaLong}{}\renewcommand{\rhrbPSRFRidgeBaLong}{\ensuremath{p = 0.068}}  % ... its p statement (Ledoit-Wolf, D19; was the D6 bootstrap)
\providecommand{\rhrbPctFLassoBarMain}{}\renewcommand{\rhrbPctFLassoBarMain}{\ensuremath{-}6.2}  % QLIKE, full set vs HAR benchmark, %, every bar, Lasso, W = 3,000
\providecommand{\rhrbDMFLassoBarMain}{}\renewcommand{\rhrbDMFLassoBarMain}{\ensuremath{-}4.72}  % DM, full set vs HAR benchmark (negative: full more accurate), every bar, Lasso, W = 3,000
\providecommand{\rhrbPFLassoBarMain}{}\renewcommand{\rhrbPFLassoBarMain}{\ensuremath{p < 0.001}}  % ... its p statement
\providecommand{\rhrbPctFLassoCloseMain}{}\renewcommand{\rhrbPctFLassoCloseMain}{3.3}  % QLIKE, full set vs HAR benchmark, %, 16:00 bar, Lasso, W = 3,000
\providecommand{\rhrbDMFLassoCloseMain}{}\renewcommand{\rhrbDMFLassoCloseMain}{1.04}  % DM, full set vs HAR benchmark (negative: full more accurate), 16:00 bar, Lasso, W = 3,000
\providecommand{\rhrbPFLassoCloseMain}{}\renewcommand{\rhrbPFLassoCloseMain}{\ensuremath{p = 0.30}}  % ... its p statement
\providecommand{\rhrbDSRFLassoMidMain}{}\renewcommand{\rhrbDSRFLassoMidMain}{0.77}  % Sharpe ratio, full set minus HAR benchmark, mid-quotes, Lasso, W = 3,000
\providecommand{\rhrbPSRFLassoMidMain}{}\renewcommand{\rhrbPSRFLassoMidMain}{\ensuremath{p = 0.12}}  % ... its p statement (Ledoit-Wolf, D19; was the D6 bootstrap)
\providecommand{\rhrbDSRFLassoBaMain}{}\renewcommand{\rhrbDSRFLassoBaMain}{0.77}  % Sharpe ratio, full set minus HAR benchmark, bid/ask prices, Lasso, W = 3,000
\providecommand{\rhrbPSRFLassoBaMain}{}\renewcommand{\rhrbPSRFLassoBaMain}{\ensuremath{p = 0.12}}  % ... its p statement (Ledoit-Wolf, D19; was the D6 bootstrap)
\providecommand{\rhrbPctFLassoBarLong}{}\renewcommand{\rhrbPctFLassoBarLong}{\ensuremath{-}7.6}  % QLIKE, full set vs HAR benchmark, %, every bar, Lasso, W = 3,500
\providecommand{\rhrbDMFLassoBarLong}{}\renewcommand{\rhrbDMFLassoBarLong}{\ensuremath{-}6.39}  % DM, full set vs HAR benchmark (negative: full more accurate), every bar, Lasso, W = 3,500
\providecommand{\rhrbPFLassoBarLong}{}\renewcommand{\rhrbPFLassoBarLong}{\ensuremath{p < 0.001}}  % ... its p statement
\providecommand{\rhrbPctFLassoCloseLong}{}\renewcommand{\rhrbPctFLassoCloseLong}{\ensuremath{-}1.6}  % QLIKE, full set vs HAR benchmark, %, 16:00 bar, Lasso, W = 3,500
\providecommand{\rhrbDMFLassoCloseLong}{}\renewcommand{\rhrbDMFLassoCloseLong}{\ensuremath{-}0.69}  % DM, full set vs HAR benchmark (negative: full more accurate), 16:00 bar, Lasso, W = 3,500
\providecommand{\rhrbPFLassoCloseLong}{}\renewcommand{\rhrbPFLassoCloseLong}{\ensuremath{p = 0.49}}  % ... its p statement
\providecommand{\rhrbDSRFLassoMidLong}{}\renewcommand{\rhrbDSRFLassoMidLong}{0.79}  % Sharpe ratio, full set minus HAR benchmark, mid-quotes, Lasso, W = 3,500
\providecommand{\rhrbPSRFLassoMidLong}{}\renewcommand{\rhrbPSRFLassoMidLong}{\ensuremath{p = 0.1006}}  % ... its p statement (Ledoit-Wolf, D19; was the D6 bootstrap)
\providecommand{\rhrbDSRFLassoBaLong}{}\renewcommand{\rhrbDSRFLassoBaLong}{0.79}  % Sharpe ratio, full set minus HAR benchmark, bid/ask prices, Lasso, W = 3,500
\providecommand{\rhrbPSRFLassoBaLong}{}\renewcommand{\rhrbPSRFLassoBaLong}{\ensuremath{p = 0.1008}}  % ... its p statement (Ledoit-Wolf, D19; was the D6 bootstrap)
\providecommand{\rhrbPctFEnetBarMain}{}\renewcommand{\rhrbPctFEnetBarMain}{\ensuremath{-}6.0}  % QLIKE, full set vs HAR benchmark, %, every bar, Elastic net, W = 3,000
\providecommand{\rhrbDMFEnetBarMain}{}\renewcommand{\rhrbDMFEnetBarMain}{\ensuremath{-}6.98}  % DM, full set vs HAR benchmark (negative: full more accurate), every bar, Elastic net, W = 3,000
\providecommand{\rhrbPFEnetBarMain}{}\renewcommand{\rhrbPFEnetBarMain}{\ensuremath{p < 0.001}}  % ... its p statement
\providecommand{\rhrbPctFEnetCloseMain}{}\renewcommand{\rhrbPctFEnetCloseMain}{\ensuremath{-}0.5}  % QLIKE, full set vs HAR benchmark, %, 16:00 bar, Elastic net, W = 3,000
\providecommand{\rhrbDMFEnetCloseMain}{}\renewcommand{\rhrbDMFEnetCloseMain}{\ensuremath{-}0.21}  % DM, full set vs HAR benchmark (negative: full more accurate), 16:00 bar, Elastic net, W = 3,000
\providecommand{\rhrbPFEnetCloseMain}{}\renewcommand{\rhrbPFEnetCloseMain}{\ensuremath{p = 0.83}}  % ... its p statement
\providecommand{\rhrbDSRFEnetMidMain}{}\renewcommand{\rhrbDSRFEnetMidMain}{0.90}  % Sharpe ratio, full set minus HAR benchmark, mid-quotes, Elastic net, W = 3,000
\providecommand{\rhrbPSRFEnetMidMain}{}\renewcommand{\rhrbPSRFEnetMidMain}{\ensuremath{p = 0.077}}  % ... its p statement (Ledoit-Wolf, D19; was the D6 bootstrap)
\providecommand{\rhrbDSRFEnetBaMain}{}\renewcommand{\rhrbDSRFEnetBaMain}{0.90}  % Sharpe ratio, full set minus HAR benchmark, bid/ask prices, Elastic net, W = 3,000
\providecommand{\rhrbPSRFEnetBaMain}{}\renewcommand{\rhrbPSRFEnetBaMain}{\ensuremath{p = 0.076}}  % ... its p statement (Ledoit-Wolf, D19; was the D6 bootstrap)
\providecommand{\rhrbPctFEnetBarLong}{}\renewcommand{\rhrbPctFEnetBarLong}{\ensuremath{-}6.7}  % QLIKE, full set vs HAR benchmark, %, every bar, Elastic net, W = 3,500
\providecommand{\rhrbDMFEnetBarLong}{}\renewcommand{\rhrbDMFEnetBarLong}{\ensuremath{-}8.45}  % DM, full set vs HAR benchmark (negative: full more accurate), every bar, Elastic net, W = 3,500
\providecommand{\rhrbPFEnetBarLong}{}\renewcommand{\rhrbPFEnetBarLong}{\ensuremath{p < 0.001}}  % ... its p statement
\providecommand{\rhrbPctFEnetCloseLong}{}\renewcommand{\rhrbPctFEnetCloseLong}{\ensuremath{-}3.9}  % QLIKE, full set vs HAR benchmark, %, 16:00 bar, Elastic net, W = 3,500
\providecommand{\rhrbDMFEnetCloseLong}{}\renewcommand{\rhrbDMFEnetCloseLong}{\ensuremath{-}2.15}  % DM, full set vs HAR benchmark (negative: full more accurate), 16:00 bar, Elastic net, W = 3,500
\providecommand{\rhrbPFEnetCloseLong}{}\renewcommand{\rhrbPFEnetCloseLong}{\ensuremath{p = 0.032}}  % ... its p statement
\providecommand{\rhrbDSRFEnetMidLong}{}\renewcommand{\rhrbDSRFEnetMidLong}{1.00}  % Sharpe ratio, full set minus HAR benchmark, mid-quotes, Elastic net, W = 3,500
\providecommand{\rhrbPSRFEnetMidLong}{}\renewcommand{\rhrbPSRFEnetMidLong}{\ensuremath{p = 0.027}}  % ... its p statement (Ledoit-Wolf, D19; was the D6 bootstrap)
\providecommand{\rhrbDSRFEnetBaLong}{}\renewcommand{\rhrbDSRFEnetBaLong}{0.99}  % Sharpe ratio, full set minus HAR benchmark, bid/ask prices, Elastic net, W = 3,500
\providecommand{\rhrbPSRFEnetBaLong}{}\renewcommand{\rhrbPSRFEnetBaLong}{\ensuremath{p = 0.027}}  % ... its p statement (Ledoit-Wolf, D19; was the D6 bootstrap)
\providecommand{\rhrbPctFLgbmBarMain}{}\renewcommand{\rhrbPctFLgbmBarMain}{\ensuremath{-}14.0}  % QLIKE, full set vs HAR benchmark, %, every bar, LightGBM, W = 3,000
\providecommand{\rhrbDMFLgbmBarMain}{}\renewcommand{\rhrbDMFLgbmBarMain}{\ensuremath{-}10.32}  % DM, full set vs HAR benchmark (negative: full more accurate), every bar, LightGBM, W = 3,000
\providecommand{\rhrbPFLgbmBarMain}{}\renewcommand{\rhrbPFLgbmBarMain}{\ensuremath{p < 0.001}}  % ... its p statement
\providecommand{\rhrbPctFLgbmCloseMain}{}\renewcommand{\rhrbPctFLgbmCloseMain}{\ensuremath{-}2.7}  % QLIKE, full set vs HAR benchmark, %, 16:00 bar, LightGBM, W = 3,000
\providecommand{\rhrbDMFLgbmCloseMain}{}\renewcommand{\rhrbDMFLgbmCloseMain}{\ensuremath{-}1.81}  % DM, full set vs HAR benchmark (negative: full more accurate), 16:00 bar, LightGBM, W = 3,000
\providecommand{\rhrbPFLgbmCloseMain}{}\renewcommand{\rhrbPFLgbmCloseMain}{\ensuremath{p = 0.071}}  % ... its p statement
\providecommand{\rhrbDSRFLgbmMidMain}{}\renewcommand{\rhrbDSRFLgbmMidMain}{0.93}  % Sharpe ratio, full set minus HAR benchmark, mid-quotes, LightGBM, W = 3,000
\providecommand{\rhrbPSRFLgbmMidMain}{}\renewcommand{\rhrbPSRFLgbmMidMain}{\ensuremath{p = 0.040}}  % ... its p statement (Ledoit-Wolf, D19; was the D6 bootstrap)
\providecommand{\rhrbDSRFLgbmBaMain}{}\renewcommand{\rhrbDSRFLgbmBaMain}{0.93}  % Sharpe ratio, full set minus HAR benchmark, bid/ask prices, LightGBM, W = 3,000
\providecommand{\rhrbPSRFLgbmBaMain}{}\renewcommand{\rhrbPSRFLgbmBaMain}{\ensuremath{p = 0.042}}  % ... its p statement (Ledoit-Wolf, D19; was the D6 bootstrap)
\providecommand{\rhrbPctFLgbmBarLong}{}\renewcommand{\rhrbPctFLgbmBarLong}{\ensuremath{-}14.0}  % QLIKE, full set vs HAR benchmark, %, every bar, LightGBM, W = 3,500
\providecommand{\rhrbDMFLgbmBarLong}{}\renewcommand{\rhrbDMFLgbmBarLong}{\ensuremath{-}10.02}  % DM, full set vs HAR benchmark (negative: full more accurate), every bar, LightGBM, W = 3,500
\providecommand{\rhrbPFLgbmBarLong}{}\renewcommand{\rhrbPFLgbmBarLong}{\ensuremath{p < 0.001}}  % ... its p statement
\providecommand{\rhrbPctFLgbmCloseLong}{}\renewcommand{\rhrbPctFLgbmCloseLong}{\ensuremath{-}2.4}  % QLIKE, full set vs HAR benchmark, %, 16:00 bar, LightGBM, W = 3,500
\providecommand{\rhrbDMFLgbmCloseLong}{}\renewcommand{\rhrbDMFLgbmCloseLong}{\ensuremath{-}1.62}  % DM, full set vs HAR benchmark (negative: full more accurate), 16:00 bar, LightGBM, W = 3,500
\providecommand{\rhrbPFLgbmCloseLong}{}\renewcommand{\rhrbPFLgbmCloseLong}{\ensuremath{p = 0.11}}  % ... its p statement
\providecommand{\rhrbDSRFLgbmMidLong}{}\renewcommand{\rhrbDSRFLgbmMidLong}{0.95}  % Sharpe ratio, full set minus HAR benchmark, mid-quotes, LightGBM, W = 3,500
\providecommand{\rhrbPSRFLgbmMidLong}{}\renewcommand{\rhrbPSRFLgbmMidLong}{\ensuremath{p = 0.035}}  % ... its p statement (Ledoit-Wolf, D19; was the D6 bootstrap)
\providecommand{\rhrbDSRFLgbmBaLong}{}\renewcommand{\rhrbDSRFLgbmBaLong}{0.94}  % Sharpe ratio, full set minus HAR benchmark, bid/ask prices, LightGBM, W = 3,500
\providecommand{\rhrbPSRFLgbmBaLong}{}\renewcommand{\rhrbPSRFLgbmBaLong}{\ensuremath{p = 0.036}}  % ... its p statement (Ledoit-Wolf, D19; was the D6 bootstrap)
\providecommand{\rhrbPctFXgbBarMain}{}\renewcommand{\rhrbPctFXgbBarMain}{\ensuremath{-}13.8}  % QLIKE, full set vs HAR benchmark, %, every bar, XGBoost, W = 3,000
\providecommand{\rhrbDMFXgbBarMain}{}\renewcommand{\rhrbDMFXgbBarMain}{\ensuremath{-}9.94}  % DM, full set vs HAR benchmark (negative: full more accurate), every bar, XGBoost, W = 3,000
\providecommand{\rhrbPFXgbBarMain}{}\renewcommand{\rhrbPFXgbBarMain}{\ensuremath{p < 0.001}}  % ... its p statement
\providecommand{\rhrbPctFXgbCloseMain}{}\renewcommand{\rhrbPctFXgbCloseMain}{\ensuremath{-}3.4}  % QLIKE, full set vs HAR benchmark, %, 16:00 bar, XGBoost, W = 3,000
\providecommand{\rhrbDMFXgbCloseMain}{}\renewcommand{\rhrbDMFXgbCloseMain}{\ensuremath{-}2.18}  % DM, full set vs HAR benchmark (negative: full more accurate), 16:00 bar, XGBoost, W = 3,000
\providecommand{\rhrbPFXgbCloseMain}{}\renewcommand{\rhrbPFXgbCloseMain}{\ensuremath{p = 0.030}}  % ... its p statement
\providecommand{\rhrbDSRFXgbMidMain}{}\renewcommand{\rhrbDSRFXgbMidMain}{0.49}  % Sharpe ratio, full set minus HAR benchmark, mid-quotes, XGBoost, W = 3,000
\providecommand{\rhrbPSRFXgbMidMain}{}\renewcommand{\rhrbPSRFXgbMidMain}{\ensuremath{p = 0.27}}  % ... its p statement (Ledoit-Wolf, D19; was the D6 bootstrap)
\providecommand{\rhrbDSRFXgbBaMain}{}\renewcommand{\rhrbDSRFXgbBaMain}{0.49}  % Sharpe ratio, full set minus HAR benchmark, bid/ask prices, XGBoost, W = 3,000
\providecommand{\rhrbPSRFXgbBaMain}{}\renewcommand{\rhrbPSRFXgbBaMain}{\ensuremath{p = 0.27}}  % ... its p statement (Ledoit-Wolf, D19; was the D6 bootstrap)
\providecommand{\rhrbPctFXgbBarLong}{}\renewcommand{\rhrbPctFXgbBarLong}{\ensuremath{-}13.5}  % QLIKE, full set vs HAR benchmark, %, every bar, XGBoost, W = 3,500
\providecommand{\rhrbDMFXgbBarLong}{}\renewcommand{\rhrbDMFXgbBarLong}{\ensuremath{-}9.97}  % DM, full set vs HAR benchmark (negative: full more accurate), every bar, XGBoost, W = 3,500
\providecommand{\rhrbPFXgbBarLong}{}\renewcommand{\rhrbPFXgbBarLong}{\ensuremath{p < 0.001}}  % ... its p statement
\providecommand{\rhrbPctFXgbCloseLong}{}\renewcommand{\rhrbPctFXgbCloseLong}{\ensuremath{-}3.2}  % QLIKE, full set vs HAR benchmark, %, 16:00 bar, XGBoost, W = 3,500
\providecommand{\rhrbDMFXgbCloseLong}{}\renewcommand{\rhrbDMFXgbCloseLong}{\ensuremath{-}2.19}  % DM, full set vs HAR benchmark (negative: full more accurate), 16:00 bar, XGBoost, W = 3,500
\providecommand{\rhrbPFXgbCloseLong}{}\renewcommand{\rhrbPFXgbCloseLong}{\ensuremath{p = 0.029}}  % ... its p statement
\providecommand{\rhrbDSRFXgbMidLong}{}\renewcommand{\rhrbDSRFXgbMidLong}{0.57}  % Sharpe ratio, full set minus HAR benchmark, mid-quotes, XGBoost, W = 3,500
\providecommand{\rhrbPSRFXgbMidLong}{}\renewcommand{\rhrbPSRFXgbMidLong}{\ensuremath{p = 0.15}}  % ... its p statement (Ledoit-Wolf, D19; was the D6 bootstrap)
\providecommand{\rhrbDSRFXgbBaLong}{}\renewcommand{\rhrbDSRFXgbBaLong}{0.56}  % Sharpe ratio, full set minus HAR benchmark, bid/ask prices, XGBoost, W = 3,500
\providecommand{\rhrbPSRFXgbBaLong}{}\renewcommand{\rhrbPSRFXgbBaLong}{\ensuremath{p = 0.15}}  % ... its p statement (Ledoit-Wolf, D19; was the D6 bootstrap)
\providecommand{\rhrbPctFRfBarMain}{}\renewcommand{\rhrbPctFRfBarMain}{\ensuremath{-}11.1}  % QLIKE, full set vs HAR benchmark, %, every bar, Random forest, W = 3,000
\providecommand{\rhrbDMFRfBarMain}{}\renewcommand{\rhrbDMFRfBarMain}{\ensuremath{-}9.85}  % DM, full set vs HAR benchmark (negative: full more accurate), every bar, Random forest, W = 3,000
\providecommand{\rhrbPFRfBarMain}{}\renewcommand{\rhrbPFRfBarMain}{\ensuremath{p < 0.001}}  % ... its p statement
\providecommand{\rhrbPctFRfCloseMain}{}\renewcommand{\rhrbPctFRfCloseMain}{\ensuremath{-}1.4}  % QLIKE, full set vs HAR benchmark, %, 16:00 bar, Random forest, W = 3,000
\providecommand{\rhrbDMFRfCloseMain}{}\renewcommand{\rhrbDMFRfCloseMain}{\ensuremath{-}0.92}  % DM, full set vs HAR benchmark (negative: full more accurate), 16:00 bar, Random forest, W = 3,000
\providecommand{\rhrbPFRfCloseMain}{}\renewcommand{\rhrbPFRfCloseMain}{\ensuremath{p = 0.36}}  % ... its p statement
\providecommand{\rhrbDSRFRfMidMain}{}\renewcommand{\rhrbDSRFRfMidMain}{1.56}  % Sharpe ratio, full set minus HAR benchmark, mid-quotes, Random forest, W = 3,000
\providecommand{\rhrbPSRFRfMidMain}{}\renewcommand{\rhrbPSRFRfMidMain}{\ensuremath{p = 0.020}}  % ... its p statement (Ledoit-Wolf, D19; was the D6 bootstrap)
\providecommand{\rhrbDSRFRfBaMain}{}\renewcommand{\rhrbDSRFRfBaMain}{1.55}  % Sharpe ratio, full set minus HAR benchmark, bid/ask prices, Random forest, W = 3,000
\providecommand{\rhrbPSRFRfBaMain}{}\renewcommand{\rhrbPSRFRfBaMain}{\ensuremath{p = 0.023}}  % ... its p statement (Ledoit-Wolf, D19; was the D6 bootstrap)
\providecommand{\rhrbPctFRfBarLong}{}\renewcommand{\rhrbPctFRfBarLong}{\ensuremath{-}11.2}  % QLIKE, full set vs HAR benchmark, %, every bar, Random forest, W = 3,500
\providecommand{\rhrbDMFRfBarLong}{}\renewcommand{\rhrbDMFRfBarLong}{\ensuremath{-}9.40}  % DM, full set vs HAR benchmark (negative: full more accurate), every bar, Random forest, W = 3,500
\providecommand{\rhrbPFRfBarLong}{}\renewcommand{\rhrbPFRfBarLong}{\ensuremath{p < 0.001}}  % ... its p statement
\providecommand{\rhrbPctFRfCloseLong}{}\renewcommand{\rhrbPctFRfCloseLong}{\ensuremath{-}1.8}  % QLIKE, full set vs HAR benchmark, %, 16:00 bar, Random forest, W = 3,500
\providecommand{\rhrbDMFRfCloseLong}{}\renewcommand{\rhrbDMFRfCloseLong}{\ensuremath{-}1.15}  % DM, full set vs HAR benchmark (negative: full more accurate), 16:00 bar, Random forest, W = 3,500
\providecommand{\rhrbPFRfCloseLong}{}\renewcommand{\rhrbPFRfCloseLong}{\ensuremath{p = 0.25}}  % ... its p statement
\providecommand{\rhrbDSRFRfMidLong}{}\renewcommand{\rhrbDSRFRfMidLong}{1.03}  % Sharpe ratio, full set minus HAR benchmark, mid-quotes, Random forest, W = 3,500
\providecommand{\rhrbPSRFRfMidLong}{}\renewcommand{\rhrbPSRFRfMidLong}{\ensuremath{p = 0.0098}}  % ... its p statement (Ledoit-Wolf, D19; was the D6 bootstrap)
\providecommand{\rhrbDSRFRfBaLong}{}\renewcommand{\rhrbDSRFRfBaLong}{1.02}  % Sharpe ratio, full set minus HAR benchmark, bid/ask prices, Random forest, W = 3,500
\providecommand{\rhrbPSRFRfBaLong}{}\renewcommand{\rhrbPSRFRfBaLong}{\ensuremath{p = 0.0100}}  % ... its p statement (Ledoit-Wolf, D19; was the D6 bootstrap)
\providecommand{\rhrbNArms}{}\renewcommand{\rhrbNArms}{12}  % pooled forecasts compared at both windows (finished; OLS reference row out)
\providecommand{\rhrbNArmsAll}{}\renewcommand{\rhrbNArmsAll}{12}  % pooled forecasts in all (finished or pending)
\providecommand{\rhrbNArmsPending}{}\renewcommand{\rhrbNArmsPending}{0}  % pooled forecasts pending at a window
\providecommand{\rhrbNSigBar}{}\renewcommand{\rhrbNSigBar}{5}  % of rhrbNArms: QLIKE differs between the windows at 5 % (DM), every bar
\providecommand{\rhrbNLowerBar}{}\renewcommand{\rhrbNLowerBar}{3}  % of rhrbNArms: lower QLIKE at 3,500 (DM < 0), every bar
\providecommand{\rhrbNHigherSigBar}{}\renewcommand{\rhrbNHigherSigBar}{3}  % of rhrbNArms: significantly higher QLIKE at 3,500 (DM > 0, 5 %), every bar
\providecommand{\rhrbNLowerSigBar}{}\renewcommand{\rhrbNLowerSigBar}{2}  % of rhrbNArms: significantly lower QLIKE at 3,500 (DM < 0, 5 %), every bar
\providecommand{\rhrbMaxAbsPctBar}{}\renewcommand{\rhrbMaxAbsPctBar}{1.3}  % largest absolute QLIKE change between the windows, %, every bar (finished arms)
\providecommand{\rhrbNSigClose}{}\renewcommand{\rhrbNSigClose}{3}  % of rhrbNArms: QLIKE differs between the windows at 5 % (DM), 16:00 bar
\providecommand{\rhrbNLowerClose}{}\renewcommand{\rhrbNLowerClose}{8}  % of rhrbNArms: lower QLIKE at 3,500 (DM < 0), 16:00 bar
\providecommand{\rhrbNHigherSigClose}{}\renewcommand{\rhrbNHigherSigClose}{0}  % of rhrbNArms: significantly higher QLIKE at 3,500 (DM > 0, 5 %), 16:00 bar
\providecommand{\rhrbNLowerSigClose}{}\renewcommand{\rhrbNLowerSigClose}{3}  % of rhrbNArms: significantly lower QLIKE at 3,500 (DM < 0, 5 %), 16:00 bar
\providecommand{\rhrbMaxAbsPctClose}{}\renewcommand{\rhrbMaxAbsPctClose}{4.4}  % largest absolute QLIKE change between the windows, %, 16:00 bar (finished arms)
\providecommand{\rhrbNModelsMain}{}\renewcommand{\rhrbNModelsMain}{6}  % model classes with both predictor sets finished, W = 3,000
\providecommand{\rhrbNFullLowerSigBarMain}{}\renewcommand{\rhrbNFullLowerSigBarMain}{6}  % of rhrbNModelsMain: full set significantly more accurate than HAR (DM < 0, 5 %), every bar
\providecommand{\rhrbNFullLowerSigCloseMain}{}\renewcommand{\rhrbNFullLowerSigCloseMain}{1}  % of rhrbNModelsMain: full set significantly more accurate than HAR (DM < 0, 5 %), 16:00 bar
\providecommand{\rhrbNDSRFPosMidMain}{}\renewcommand{\rhrbNDSRFPosMidMain}{6}  % of rhrbNModelsMain: full-set Sharpe ratio above HAR, mid-quotes
\providecommand{\rhrbNDSRFSigMidMain}{}\renewcommand{\rhrbNDSRFSigMidMain}{2}  % of rhrbNModelsMain: ... and Ledoit-Wolf p < 0.05 (D19), mid-quotes
\providecommand{\rhrbNDSRFPosBaMain}{}\renewcommand{\rhrbNDSRFPosBaMain}{6}  % of rhrbNModelsMain: full-set Sharpe ratio above HAR, bid/ask prices
\providecommand{\rhrbNDSRFSigBaMain}{}\renewcommand{\rhrbNDSRFSigBaMain}{2}  % of rhrbNModelsMain: ... and Ledoit-Wolf p < 0.05 (D19), bid/ask prices
\providecommand{\rhrbNModelsLong}{}\renewcommand{\rhrbNModelsLong}{6}  % model classes with both predictor sets finished, W = 3,500
\providecommand{\rhrbNFullLowerSigBarLong}{}\renewcommand{\rhrbNFullLowerSigBarLong}{6}  % of rhrbNModelsLong: full set significantly more accurate than HAR (DM < 0, 5 %), every bar
\providecommand{\rhrbNFullLowerSigCloseLong}{}\renewcommand{\rhrbNFullLowerSigCloseLong}{2}  % of rhrbNModelsLong: full set significantly more accurate than HAR (DM < 0, 5 %), 16:00 bar
\providecommand{\rhrbNDSRFPosMidLong}{}\renewcommand{\rhrbNDSRFPosMidLong}{6}  % of rhrbNModelsLong: full-set Sharpe ratio above HAR, mid-quotes
\providecommand{\rhrbNDSRFSigMidLong}{}\renewcommand{\rhrbNDSRFSigMidLong}{3}  % of rhrbNModelsLong: ... and Ledoit-Wolf p < 0.05 (D19), mid-quotes
\providecommand{\rhrbNDSRFPosBaLong}{}\renewcommand{\rhrbNDSRFPosBaLong}{6}  % of rhrbNModelsLong: full-set Sharpe ratio above HAR, bid/ask prices
\providecommand{\rhrbNDSRFSigBaLong}{}\renewcommand{\rhrbNDSRFSigBaLong}{3}  % of rhrbNModelsLong: ... and Ledoit-Wolf p < 0.05 (D19), bid/ask prices
\providecommand{\rhrbNModelsAll}{}\renewcommand{\rhrbNModelsAll}{6}  % model classes in all
\providecommand{\rhrbNSRSigMid}{}\renewcommand{\rhrbNSRSigMid}{0}  % of rhrbNArms: Sharpe ratio differs between the windows (Ledoit-Wolf p < 0.05, D19), mid-quotes
\providecommand{\rhrbMaxAbsDSRMid}{}\renewcommand{\rhrbMaxAbsDSRMid}{0.48}  % largest absolute Sharpe-ratio change between the windows, mid-quotes (finished arms)
\providecommand{\rhrbMinPSRMid}{}\renewcommand{\rhrbMinPSRMid}{\ensuremath{p = 0.40}}  % smallest Ledoit-Wolf p (D19) of a Sharpe-ratio change between the windows, mid-quotes
\providecommand{\rhrbNSRSigBa}{}\renewcommand{\rhrbNSRSigBa}{0}  % of rhrbNArms: Sharpe ratio differs between the windows (Ledoit-Wolf p < 0.05, D19), bid/ask prices
\providecommand{\rhrbMaxAbsDSRBa}{}\renewcommand{\rhrbMaxAbsDSRBa}{0.48}  % largest absolute Sharpe-ratio change between the windows, bid/ask prices (finished arms)
\providecommand{\rhrbMinPSRBa}{}\renewcommand{\rhrbMinPSRBa}{\ensuremath{p = 0.40}}  % smallest Ledoit-Wolf p (D19) of a Sharpe-ratio change between the windows, bid/ask prices
\fi
\begin{table}[tbp]
\centering
\small
\setlength{\tabcolsep}{4pt}
\begin{threeparttable}
\caption{Robustness: a rolling window of \rhrbWinLong{} trading days}
\label{tab:rh_robust_w3500}
\begin{tabular*}{\linewidth}{@{\extracolsep{\fill}}llrrrrrr@{}}
\toprule
 & & \multicolumn{3}{c}{Every bar} & \multicolumn{3}{c}{16:00 bar} \\
\cmidrule(lr){3-5}\cmidrule(l){6-8}
Model & Window & HAR & Full & $\DM$ & HAR & Full & $\DM$ \\
\midrule
\multicolumn{8}{@{}l}{\textit{Panel A: QLIKE}} \\
Ridge & \rhrbWinMain & $0.1893^{\phantom{***}}$ & $0.1791^{\phantom{***}}$ & $-5.63^{***}$ & $0.1269^{\phantom{***}}$ & $0.1279^{\phantom{***}}$ & $0.31^{\phantom{***}}$ \\
 & \rhrbWinLong & $0.1897^{\phantom{***}}$ & $0.1777^{\phantom{***}}$ & $-7.06^{***}$ & $0.1272^{\phantom{***}}$ & $0.1233^{\phantom{***}}$ & $-1.55^{\phantom{***}}$ \\
 & Change & $(2.48)^{**\phantom{*}}$ & $(-2.26)^{**\phantom{*}}$ &  & $(1.31)^{\phantom{***}}$ & $(-3.10)^{***}$ &  \\
Lasso & \rhrbWinMain & $0.1893^{\phantom{***}}$ & $0.1775^{\phantom{***}}$ & $-4.72^{***}$ & $0.1269^{\phantom{***}}$ & $0.1310^{\phantom{***}}$ & $1.04^{\phantom{***}}$ \\
 & \rhrbWinLong & $0.1897^{\phantom{***}}$ & $0.1752^{\phantom{***}}$ & $-6.39^{***}$ & $0.1272^{\phantom{***}}$ & $0.1253^{\phantom{***}}$ & $-0.69^{\phantom{***}}$ \\
 & Change & $(2.48)^{**\phantom{*}}$ & $(-2.04)^{**\phantom{*}}$ &  & $(1.31)^{\phantom{***}}$ & $(-2.89)^{***}$ &  \\
Elastic net & \rhrbWinMain & $0.1893^{\phantom{***}}$ & $0.1778^{\phantom{***}}$ & $-6.98^{***}$ & $0.1269^{\phantom{***}}$ & $0.1262^{\phantom{***}}$ & $-0.21^{\phantom{***}}$ \\
 & \rhrbWinLong & $0.1897^{\phantom{***}}$ & $0.1770^{\phantom{***}}$ & $-8.45^{***}$ & $0.1272^{\phantom{***}}$ & $0.1223^{\phantom{***}}$ & $-2.15^{**\phantom{*}}$ \\
 & Change & $(2.48)^{**\phantom{*}}$ & $(-1.47)^{\phantom{***}}$ &  & $(1.31)^{\phantom{***}}$ & $(-2.93)^{***}$ &  \\
LightGBM & \rhrbWinMain & $0.1875^{\phantom{***}}$ & $0.1613^{\phantom{***}}$ & $-10.32^{***}$ & $0.1235^{\phantom{***}}$ & $0.1201^{\phantom{***}}$ & $-1.81^{*\phantom{**}}$ \\
 & \rhrbWinLong & $0.1876^{\phantom{***}}$ & $0.1613^{\phantom{***}}$ & $-10.02^{***}$ & $0.1233^{\phantom{***}}$ & $0.1204^{\phantom{***}}$ & $-1.62^{\phantom{***}}$ \\
 & Change & $(0.20)^{\phantom{***}}$ & $(0.20)^{\phantom{***}}$ &  & $(-0.45)^{\phantom{***}}$ & $(0.55)^{\phantom{***}}$ &  \\
XGBoost & \rhrbWinMain & $0.1888^{\phantom{***}}$ & $0.1627^{\phantom{***}}$ & $-9.94^{***}$ & $0.1234^{\phantom{***}}$ & $0.1192^{\phantom{***}}$ & $-2.18^{**\phantom{*}}$ \\
 & \rhrbWinLong & $0.1890^{\phantom{***}}$ & $0.1636^{\phantom{***}}$ & $-9.97^{***}$ & $0.1232^{\phantom{***}}$ & $0.1192^{\phantom{***}}$ & $-2.19^{**\phantom{*}}$ \\
 & Change & $(1.13)^{\phantom{***}}$ & $(1.87)^{*\phantom{**}}$ &  & $(-0.55)^{\phantom{***}}$ & $(-0.15)^{\phantom{***}}$ &  \\
Random forest & \rhrbWinMain & $0.1895^{\phantom{***}}$ & $0.1685^{\phantom{***}}$ & $-9.85^{***}$ & $0.1266^{\phantom{***}}$ & $0.1248^{\phantom{***}}$ & $-0.92^{\phantom{***}}$ \\
 & \rhrbWinLong & $0.1898^{\phantom{***}}$ & $0.1685^{\phantom{***}}$ & $-9.40^{***}$ & $0.1264^{\phantom{***}}$ & $0.1241^{\phantom{***}}$ & $-1.15^{\phantom{***}}$ \\
 & Change & $(1.15)^{\phantom{***}}$ & $(0.02)^{\phantom{***}}$ &  & $(-0.36)^{\phantom{***}}$ & $(-1.43)^{\phantom{***}}$ &  \\
\addlinespace[4pt]
\multicolumn{8}{@{}l}{\textit{Panel B: Sharpe ratio of the sign strategy}} \\
 & & \multicolumn{3}{c}{Mid-quotes} & \multicolumn{3}{c}{Bid/ask prices} \\
\cmidrule(lr){3-5}\cmidrule(l){6-8}
 & & HAR & Full & $\Delta\SR$ & HAR & Full & $\Delta\SR$ \\
\addlinespace[2pt]
Ridge & \rhrbWinMain & $0.98^{\phantom{***}}$ & $1.85^{\phantom{***}}$ & $0.86^{\phantom{***}}$ & $0.62^{\phantom{***}}$ & $1.49^{\phantom{***}}$ & $0.87^{\phantom{***}}$ \\
 & \rhrbWinLong & $1.00^{\phantom{***}}$ & $1.74^{\phantom{***}}$ & $0.74^{\phantom{***}}$ & $0.64^{\phantom{***}}$ & $1.38^{\phantom{***}}$ & $0.74^{\phantom{***}}$ \\
 & Change & $0.02^{\phantom{***}}$ & $-0.10^{\phantom{***}}$ &  & $0.02^{\phantom{***}}$ & $-0.11^{\phantom{***}}$ &  \\
Lasso & \rhrbWinMain & $0.98^{\phantom{***}}$ & $1.72^{\phantom{***}}$ & $0.74^{\phantom{***}}$ & $0.62^{\phantom{***}}$ & $1.36^{\phantom{***}}$ & $0.74^{\phantom{***}}$ \\
 & \rhrbWinLong & $1.00^{\phantom{***}}$ & $1.73^{\phantom{***}}$ & $0.74^{\phantom{***}}$ & $0.64^{\phantom{***}}$ & $1.37^{\phantom{***}}$ & $0.73^{\phantom{***}}$ \\
 & Change & $0.02^{\phantom{***}}$ & $0.01^{\phantom{***}}$ &  & $0.02^{\phantom{***}}$ & $0.01^{\phantom{***}}$ &  \\
Elastic net & \rhrbWinMain & $0.98^{\phantom{***}}$ & $1.94^{\phantom{***}}$ & $0.96^{\phantom{***}}$ & $0.62^{\phantom{***}}$ & $1.58^{\phantom{***}}$ & $0.96^{\phantom{***}}$ \\
 & \rhrbWinLong & $1.00^{\phantom{***}}$ & $1.97^{\phantom{***}}$ & $0.98^{\phantom{***}}$ & $0.64^{\phantom{***}}$ & $1.61^{\phantom{***}}$ & $0.98^{*\phantom{**}}$ \\
 & Change & $0.02^{\phantom{***}}$ & $0.03^{\phantom{***}}$ &  & $0.02^{\phantom{***}}$ & $0.03^{\phantom{***}}$ &  \\
LightGBM & \rhrbWinMain & $0.61^{\phantom{***}}$ & $1.75^{\phantom{***}}$ & $1.14^{*\phantom{**}}$ & $0.25^{\phantom{***}}$ & $1.40^{\phantom{***}}$ & $1.15^{**\phantom{*}}$ \\
 & \rhrbWinLong & $0.47^{\phantom{***}}$ & $1.62^{\phantom{***}}$ & $1.16^{**\phantom{*}}$ & $0.11^{\phantom{***}}$ & $1.27^{\phantom{***}}$ & $1.16^{**\phantom{*}}$ \\
 & Change & $-0.14^{\phantom{***}}$ & $-0.12^{\phantom{***}}$ &  & $-0.14^{\phantom{***}}$ & $-0.13^{\phantom{***}}$ &  \\
XGBoost & \rhrbWinMain & $0.91^{\phantom{***}}$ & $1.67^{\phantom{***}}$ & $0.76^{\phantom{***}}$ & $0.55^{\phantom{***}}$ & $1.31^{\phantom{***}}$ & $0.76^{\phantom{***}}$ \\
 & \rhrbWinLong & $0.66^{\phantom{***}}$ & $1.54^{\phantom{***}}$ & $0.88^{*\phantom{**}}$ & $0.31^{\phantom{***}}$ & $1.19^{\phantom{***}}$ & $0.88^{*\phantom{**}}$ \\
 & Change & $-0.25^{\phantom{***}}$ & $-0.13^{\phantom{***}}$ &  & $-0.24^{\phantom{***}}$ & $-0.13^{\phantom{***}}$ &  \\
Random forest & \rhrbWinMain & $0.01^{\phantom{***}}$ & $1.96^{\phantom{***}}$ & $1.95^{***}$ & $-0.32^{\phantom{***}}$ & $1.60^{\phantom{***}}$ & $1.93^{***}$ \\
 & \rhrbWinLong & $0.74^{\phantom{***}}$ & $1.92^{\phantom{***}}$ & $1.18^{**\phantom{*}}$ & $0.39^{\phantom{***}}$ & $1.57^{\phantom{***}}$ & $1.18^{**\phantom{*}}$ \\
 & Change & $0.74^{\phantom{***}}$ & $-0.03^{\phantom{***}}$ &  & $0.71^{\phantom{***}}$ & $-0.03^{\phantom{***}}$ &  \\
\bottomrule
\end{tabular*}
\begin{tablenotes}[flushleft]
\footnotesize
\item \textit{Notes:} The models of Table~\ref{tab:rh_oos} on a rolling window of \rhrbWinLong{} instead of
\rhrbWinMain{} trading days, all else equal; HAR: HAR benchmark; Full: full predictor set. Panel A: mean QLIKE on
every bar (\rhrbNbars{} bars) and on the 16:00 bar of the $T_{\mathrm{oos}} = \rhrbNdays$ trading days,
\rhrbSpanStart{} to \rhrbSpanEnd{}. $\DM$: Diebold--Mariano statistic (Harvey--Leybourne--Newbold,
$q_{\mathrm{NW}} = \rhrbDMLag$) of the daily loss, full minus HAR (negative: full more accurate); Change, in
parentheses: \rhrbWinLong-day minus \rhrbWinMain-day loss (positive: the longer window is less accurate). Panel B:
Sharpe ratio of the sign strategy (Section~\ref{sec:rh_econ_measures}) on the \rhrbNpw{} post-warmup option days;
$\Delta\SR$: full minus HAR; Change: \rhrbWinLong-day minus \rhrbWinMain-day ratio; stars from the paired
block-bootstrap test of equal Sharpe ratios (Section~\ref{sec:rh_methods_tests}). $^{*}$, $^{**}$, $^{***}$:
significance at the 10\%, 5\% and 1\% levels. \rhrbPendingNote
\end{tablenotes}
\end{threeparttable}
\end{table}

\item \textbf{Models trained on the closing half hour.} Panel~F of Table~\ref{tab:rh_train_arms_cont} re-estimates
  them on windows of \rhtaWsMin{} to \rhtaWsMax{} closes (Appendix~\ref{sec:rh_econ_last30}).
\end{itemize}

% Moved from drafts/2026-10-26_train_arms.tex (subsection "Training sample", label app:rh_robust_train), lightly edited:
% the CLOSE results (Panels A, C, D and E for CLOSE; Panel F) are in rehaul/A_closing.tex (D3), and the pre-patch
% design is marked (D13). Every class name, direction and "all / every / none" below is asserted by
% writeup/make_rehaul_train_arms.py (claims()).
\subsection{Training sample}
\label{app:rh_robust_train}
\begin{itemize}
\item \textbf{Main design.} Every pooled model of Table~\ref{tab:rh_oos} (ridge regression, the lasso, the elastic
  net, LightGBM, XGBoost and the random forest, each on the HAR benchmark and on the full predictor set) is estimated
  on the \rhtaNBarsDay{} regular-hours bars of each day, $\mathcal I = \mathcal I_{\mathrm{rh}}$ (bar ends 10:00 to
  16:00~ET), over a rolling window of \rhtaWinDays{} trading days (\rhtaRowsRth{} bars), with its predictors
  computed on the full grid of \rhtaNIntervals{} half-hour intervals, overnight included
  (Appendix~\ref{sec:rh_exp_schemes}). Estimation on the regular session follows the intraday volatility
  forecasting literature \citep{EngleSokalska2012,ZhangZhangCucuringuQian2024}. Below, this sample is RTH13.
\item \textbf{A pre-registered comparison.} Before any model was estimated on the other samples, a written rule (one
  for the linear models, one for the trees) fixed two other training samples, the models, their tuning and the tests:
  \begin{itemize}
  \item ALL48: all \rhtaNIntervals{} half-hour intervals of the futures' day, overnight, weekend and holiday
    intervals and early-close sessions included, over the same \rhtaWinDays-day window (\rhtaRowsAll{} rows);
  \item CLOSE: the closing half hour alone (the 16:00 rows of the main design), over \rhtaRowsCl{} closes; its
    results are in Appendix~\ref{sec:rh_econ_last30}.
  \end{itemize}
  Everything else is the main design's: the predictors, the frozen tree configurations, a refit at every training
  row, and the rule that chooses the linear penalties, applied to each sample's own pre-sample rows (validation year
  \rhtaValYear). Every sample is scored on the same \rhtaNdays{} trading days from \rhtaStart{} to \rhtaEnd, on the
  16:00 bar and on every regular-hours bar. The linear rule also fixed secondary samples for the penalized linear
  models (a natively built CLOSE, and penalties chosen on the 16:00 validation rows, ALL48@16 and RTH13@16). Samples
  with equal numbers of training rows were added before their models were estimated (for the linear models by a
  dated addendum to the rule).
\item \textbf{Design and table.} The comparison ran on the design before the FOMC fixes of
  Appendix~\ref{app:rh_dp_fomc}; its RTH13 rows are the forecasts of Table~\ref{tab:rh_oos} on that design.
  Table~\ref{tab:rh_train_arms} reports every comparison.
\item \textbf{The closing half hour.} Trained on all \rhtaNIntervals{} intervals, the full predictor set gains more
  over the HAR benchmark at 16:00 than in the main design (Table~\ref{tab:rh_train_arms}, Panel~A).
  \begin{itemize}
  \item ALL48: the change in 16:00 QLIKE is $\Delta\% = \rhtaTreePctAllCloseMin$ to $\rhtaTreePctAllCloseMax$ for the
    three tree ensembles, each with $p < 0.001$, and $\Delta\% = \rhtaLinPctAllCloseMin$ to $\rhtaLinPctAllCloseMax$
    for the penalized linear models, none significant at the 10\% level (the smallest, \rhtaLinAllClosePMin).
  \item RTH13: $\Delta\% = \rhtaPctRthCloseMin$ to $\rhtaPctRthCloseMax$, significant at the 5\% level for XGBoost
    only ($\Delta\% = \rhtaPctXgbRthClose$, \rhtaPXgbRthClose). Choosing the linear penalties on the 16:00 validation
    rows leaves every linear change insignificant at the 10\% level, in ALL48@16 and in RTH13@16, although in RTH13
    it lowers the 16:00 loss of the lasso's full-set forecast ($\Delta\% = \rhtaXPctLassoRthCtVsRthClose$,
    \rhtaXPLassoRthCtVsRthClose; Panel~C), whose HAR benchmark changes with it.
  \item A full-versus-HAR change compares two forecasts trained on the same sample, so it also moves with the HAR
    benchmark. Compared directly (Panel~C), the full-set forecasts trained on ALL48 have a lower 16:00 loss than
    those trained on RTH13 in all six classes ($\Delta\% = \rhtaAllVsRthClosePctMin$ to
    $\rhtaAllVsRthClosePctMax$), significantly at the 5\% level in \rhtaNAllVsRthCloseSig.
  \end{itemize}
\item \textbf{The average regular-hours bar.} On every bar, RTH13 is the most accurate sample: no other sample
  gives the full predictor set a lower every-bar loss in any class (Panels~C and~D test the pre-registered pairs).
  \begin{itemize}
  \item The full-set forecasts trained on ALL48 have a higher every-bar QLIKE than those trained on RTH13 in all six
    classes ($\Delta\% = \rhtaAllVsRthBarPctMin$ to $\rhtaAllVsRthBarPctMax$), each significant at the 1\% level
    (the largest, \rhtaAllVsRthBarPMax).
  \item The full predictor set lowers the every-bar loss significantly in every class under both samples, and by more
    under RTH13 in all six: $\Delta\% = \rhtaLinPctRthBarMin$ to $\rhtaLinPctRthBarMax$ for the penalized linear models
    and $\rhtaTreePctRthBarMin$ to $\rhtaTreePctRthBarMax$ for the tree ensembles, against $\rhtaLinPctAllBarMin$ to
    $\rhtaLinPctAllBarMax$ and $\rhtaTreePctAllBarMin$ to $\rhtaTreePctAllBarMax$ under ALL48.
  \end{itemize}
\item \textbf{Data volume.} At equal calendar the samples differ in their numbers of training rows (\rhtaRowsAll{}
  for ALL48, \rhtaRowsRth{} for RTH13). The equal-row samples (Panel~D) hold the number of rows fixed and change the
  span and recency of the window instead.
  \begin{itemize}
  \item With \rhtaRowsLarge{} rows, as many as RTH13, ALL48-39k has a significantly lower 16:00 loss than RTH13 in
    \rhtaNAllThirtyNineKVsRthCloseSig{} classes, the lasso and the random forest
    ($\Delta\% = \rhtaAllThirtyNineKVsRthClosePctMin$ to $\rhtaAllThirtyNineKVsRthClosePctMax$ across the six),
    against \rhtaNAllVsRthCloseSig{} for ALL48 with all its rows, and in every class its 16:00 loss relative to
    RTH13's is higher than ALL48's. Cutting ALL48 from \rhtaRowsAll{} to \rhtaRowsLarge{} rows raises its 16:00 loss
    in all six classes ($\Delta\% = \rhtaAllThirtyNineKVsAllClosePctMin$ to $\rhtaAllThirtyNineKVsAllClosePctMax$),
    significantly at the 5\% level in \rhtaNAllThirtyNineKVsAllCloseSig. Much of ALL48's advantage at 16:00 thus
    goes with its larger number of rows. On every bar, ALL48-39k remains less accurate than RTH13 in all six classes
    ($\Delta\% = \rhtaAllThirtyNineKVsRthBarPctMin$ to $\rhtaAllThirtyNineKVsRthBarPctMax$, each with $p < 0.001$).
  \item With \rhtaRowsSmall{} rows, ALL48-3k has a lower 16:00 loss than RTH13-3k in all six classes
    ($\Delta\% = \rhtaAllThreeKVsRthThreeKClosePctMin$ to $\rhtaAllThreeKVsRthThreeKClosePctMax$, significantly at the
    5\% level in \rhtaNAllThreeKVsRthThreeKCloseSig) and a higher every-bar loss in all six
    ($\Delta\% = \rhtaAllThreeKVsRthThreeKBarPctMin$ to $\rhtaAllThreeKVsRthThreeKBarPctMax$, each with
    $p < 0.001$), the same pattern as at equal calendar; the two windows span about \rhtaSpanAllSmall{} and
    \rhtaSpanRthSmall{} trading days.
  \item Fewer rows also cost accuracy within the main sample: RTH13-3k has a higher loss than RTH13 at 16:00
    ($\Delta\% = \rhtaRthThreeKVsRthClosePctMin$ to $\rhtaRthThreeKVsRthClosePctMax$) and on every bar
    ($\Delta\% = \rhtaRthThreeKVsRthBarPctMin$ to $\rhtaRthThreeKVsRthBarPctMax$), each significant at the 1\% level.
  \end{itemize}
\item \textbf{The trade.} The Sharpe ratio $\SR$ of the full-set sign strategy (Section~\ref{sec:rh_econ_measures};
  Panel~E) lies between \rhtaShFullRthMin{} and \rhtaShFullRthMax{} under RTH13 and between \rhtaShFullAllMin{} and
  \rhtaShFullAllMax{} under ALL48. None of the \rhtaNDShTests{} differences in $\SR$ between the full-set strategies
  of ALL48 or CLOSE and those of RTH13 is significant at the 5\% level (the smallest, \rhtaDShPMin). For the tree
  ensembles, the full-set strategy's gain over the HAR benchmark's is significant at the 5\% level in
  \rhtaNTreeShSigRth{} of the three classes under RTH13 and \rhtaNTreeShSigAll{} under ALL48.
\item \textbf{Caveats.} The tree configurations were chosen on an earlier design (Appendix~\ref{app:rh_tuning}) and,
  by the rule, are not re-tuned; with \rhtaRowsSmall{} rows (ALL48-3k, RTH13-3k) or \rhtaRowsAll{} (ALL48) their
  sample-size parameters may not suit the window. The linear penalties are re-chosen in every sample. ALL48 differs
  from RTH13 in more than its rows: it adds the overnight indicator, keeps the early-close sessions, and computes its
  rolling scaling on its own rows.
\item \textbf{The main design is unchanged.} RTH13 was chosen before the comparison, following the literature cited
  above, and it is the sample under which the full predictor set is most accurate on the average regular-hours bar,
  the score of Table~\ref{tab:rh_oos}. At the 16:00 bar, training on all \rhtaNIntervals{} intervals or on the close
  alone gives the full predictor set a larger gain over the HAR benchmark in every class
  (Appendix~\ref{sec:rh_econ_last30} for the close), a gain significant for the tree ensembles under ALL48 and for
  every class under CLOSE. The closing-bar results of Sections~\ref{sec:rh_results} and~\ref{sec:rh_econ} are those
  of models estimated for the regular session as a whole, not for its last half hour.
\end{itemize}
% AUTO-GENERATED by writeup/make_rehaul_train_arms.py -- do not edit; re-run the script.
% tab:rh_train_arms (appendix): the pre-registered three-way comparison of training samples (RTH13 = the main
% design, ALL48, CLOSE; the linear secondary arms; the equal-row arms) and the close-only window sweep, scored on
% the 2,579 days 2014-01-02 .. 2024-04-30; DM (HLN, Newey-West lag 8) on the daily loss; Sharpe ratio D6 of the
% sign strategy on the 614 post-warmup option days.  Two floats: Panels A-C, and D-F (\ContinuedFloat).
% Sources (git show, pinned): insurance-selling-2026-09-19-e34-train3way 0106ccb; insurance-selling-2026-09-19-e34-train3way-trees 3501ab7; insurance-selling-2026-09-19-e34-closeonly-wsweep f4abb69.
\ifdefined\rhtaNdays\else% AUTO-GENERATED by writeup/make_rehaul_train_arms.py -- do not edit; re-run the script.
% \rhta... macros: the pre-registered comparison of training samples (tab:rh_train_arms; draft
% writeup/drafts/2026-10-26_train_arms.tex).  Names: <Stat><Model><Arm><Bar>; Arm Rth (RTH13, the main design),
% All (ALL48), Cl (CLOSE), ClNat (CLOSE_NAT), AllCt / RthCt (penalties chosen at 16:00), AllThreeK / RthThreeK /
% AllThirtyNineK (the last 3,000 / 39,000 rows); Bar Close (16:00) or Bar (every bar); X... = first arm vs second
% (full set); Sh... = Sharpe ratio (D19: the paper's pooled ratio; p by Ledoit-Wolf); Ws... = the close-only
% window sweep.
% Sources (git show, pinned): insurance-selling-2026-09-19-e34-train3way 0106ccb; insurance-selling-2026-09-19-e34-train3way-trees 3501ab7; insurance-selling-2026-09-19-e34-closeonly-wsweep f4abb69.
\providecommand{\rhtaNdays}{}\renewcommand{\rhtaNdays}{2{,}579}  % trading days scored, 2 January 2014 .. 30 April 2024 (early closes dropped)
\providecommand{\rhtaNbars}{}\renewcommand{\rhtaNbars}{33{,}527}  % regular-hours bars scored (every bar; RV > 0)
\providecommand{\rhtaStart}{}\renewcommand{\rhtaStart}{2 January 2014}  % first evaluation day
\providecommand{\rhtaEnd}{}\renewcommand{\rhtaEnd}{30 April 2024}  % last evaluation day
\providecommand{\rhtaHacLag}{}\renewcommand{\rhtaHacLag}{8}  % Newey-West lag of the daily Diebold-Mariano tests
\providecommand{\rhtaNIntervals}{}\renewcommand{\rhtaNIntervals}{48}  % half-hour intervals of the futures' day (ALL48)
\providecommand{\rhtaNBarsDay}{}\renewcommand{\rhtaNBarsDay}{13}  % regular-hours bars of a day, bar ends 10:00 .. 16:00 ET (RTH13)
\providecommand{\rhtaWinDays}{}\renewcommand{\rhtaWinDays}{3{,}000}  % rolling window of the equal-calendar arms, trading days
\providecommand{\rhtaRowsAll}{}\renewcommand{\rhtaRowsAll}{144{,}000}  % training rows of ALL48 (3,000 days x 48 intervals)
\providecommand{\rhtaRowsRth}{}\renewcommand{\rhtaRowsRth}{39{,}000}  % training rows of RTH13 (3,000 days x 13 bars)
\providecommand{\rhtaRowsCl}{}\renewcommand{\rhtaRowsCl}{3{,}000}  % training rows of CLOSE (3,000 closes)
\providecommand{\rhtaRowsSmall}{}\renewcommand{\rhtaRowsSmall}{3{,}000}  % training rows of the small equal-row arms (ALL48_N3k, RTH13_N3k; CLOSE)
\providecommand{\rhtaRowsLarge}{}\renewcommand{\rhtaRowsLarge}{39{,}000}  % training rows of the large equal-row arm ALL48_N39k (= RTH13's)
\providecommand{\rhtaSpanAllSmall}{}\renewcommand{\rhtaSpanAllSmall}{62.5}  % days of 48 intervals spanned by 3,000 rows (about)
\providecommand{\rhtaSpanRthSmall}{}\renewcommand{\rhtaSpanRthSmall}{231}  % sessions of 13 bars spanned by 3,000 rows (about)
\providecommand{\rhtaSpanClSmall}{}\renewcommand{\rhtaSpanClSmall}{12}  % years of closes spanned by 3,000 closes (about)
\providecommand{\rhtaSpanAllLarge}{}\renewcommand{\rhtaSpanAllLarge}{812.5}  % days of 48 intervals spanned by 39,000 rows (about)
\providecommand{\rhtaValYear}{}\renewcommand{\rhtaValYear}{2013}  % validation year of the pre-sample fold on which every arm's linear penalties are chosen
\providecommand{\rhtaOptDays}{}\renewcommand{\rhtaOptDays}{614}  % post-warmup option days of the Sharpe ratio (Section 7; D19 pooled ratio)
\providecommand{\rhtaNClasses}{}\renewcommand{\rhtaNClasses}{6}  % model classes (ridge, lasso, elastic net, LightGBM, XGBoost, random forest)
\providecommand{\rhtaWsMin}{}\renewcommand{\rhtaWsMin}{1{,}000}  % smallest window of the close-only sweep, closes
\providecommand{\rhtaWsMax}{}\renewcommand{\rhtaWsMax}{3{,}500}  % largest window of the close-only sweep, closes
\providecommand{\rhtaPctRidgeRthClose}{}\renewcommand{\rhtaPctRidgeRthClose}{0.8}  % QLIKE % change, full set vs HAR benchmark: Ridge, RTH13 (main design: the 13 regular-hours bars, 3,000 days), 16:00 bar
\providecommand{\rhtaDmRidgeRthClose}{}\renewcommand{\rhtaDmRidgeRthClose}{0.31}  % ... its Diebold-Mariano statistic (negative: full more accurate)
\providecommand{\rhtaPRidgeRthClose}{}\renewcommand{\rhtaPRidgeRthClose}{\ensuremath{p = 0.76}}  % ... its p statement
\providecommand{\rhtaQHarRidgeRthClose}{}\renewcommand{\rhtaQHarRidgeRthClose}{0.1269}  % ... mean QLIKE, HAR benchmark
\providecommand{\rhtaQFullRidgeRthClose}{}\renewcommand{\rhtaQFullRidgeRthClose}{0.1279}  % ... mean QLIKE, full set
\providecommand{\rhtaPctLassoRthClose}{}\renewcommand{\rhtaPctLassoRthClose}{3.3}  % QLIKE % change, full set vs HAR benchmark: Lasso, RTH13 (main design: the 13 regular-hours bars, 3,000 days), 16:00 bar
\providecommand{\rhtaDmLassoRthClose}{}\renewcommand{\rhtaDmLassoRthClose}{1.04}  % ... its Diebold-Mariano statistic (negative: full more accurate)
\providecommand{\rhtaPLassoRthClose}{}\renewcommand{\rhtaPLassoRthClose}{\ensuremath{p = 0.30}}  % ... its p statement
\providecommand{\rhtaQHarLassoRthClose}{}\renewcommand{\rhtaQHarLassoRthClose}{0.1269}  % ... mean QLIKE, HAR benchmark
\providecommand{\rhtaQFullLassoRthClose}{}\renewcommand{\rhtaQFullLassoRthClose}{0.1310}  % ... mean QLIKE, full set
\providecommand{\rhtaPctEnetRthClose}{}\renewcommand{\rhtaPctEnetRthClose}{\ensuremath{-}0.5}  % QLIKE % change, full set vs HAR benchmark: Elastic net, RTH13 (main design: the 13 regular-hours bars, 3,000 days), 16:00 bar
\providecommand{\rhtaDmEnetRthClose}{}\renewcommand{\rhtaDmEnetRthClose}{\ensuremath{-}0.21}  % ... its Diebold-Mariano statistic (negative: full more accurate)
\providecommand{\rhtaPEnetRthClose}{}\renewcommand{\rhtaPEnetRthClose}{\ensuremath{p = 0.83}}  % ... its p statement
\providecommand{\rhtaQHarEnetRthClose}{}\renewcommand{\rhtaQHarEnetRthClose}{0.1269}  % ... mean QLIKE, HAR benchmark
\providecommand{\rhtaQFullEnetRthClose}{}\renewcommand{\rhtaQFullEnetRthClose}{0.1262}  % ... mean QLIKE, full set
\providecommand{\rhtaPctLgbmRthClose}{}\renewcommand{\rhtaPctLgbmRthClose}{\ensuremath{-}2.7}  % QLIKE % change, full set vs HAR benchmark: LightGBM, RTH13 (main design: the 13 regular-hours bars, 3,000 days), 16:00 bar
\providecommand{\rhtaDmLgbmRthClose}{}\renewcommand{\rhtaDmLgbmRthClose}{\ensuremath{-}1.81}  % ... its Diebold-Mariano statistic (negative: full more accurate)
\providecommand{\rhtaPLgbmRthClose}{}\renewcommand{\rhtaPLgbmRthClose}{\ensuremath{p = 0.071}}  % ... its p statement
\providecommand{\rhtaQHarLgbmRthClose}{}\renewcommand{\rhtaQHarLgbmRthClose}{0.1235}  % ... mean QLIKE, HAR benchmark
\providecommand{\rhtaQFullLgbmRthClose}{}\renewcommand{\rhtaQFullLgbmRthClose}{0.1201}  % ... mean QLIKE, full set
\providecommand{\rhtaPctXgbRthClose}{}\renewcommand{\rhtaPctXgbRthClose}{\ensuremath{-}3.4}  % QLIKE % change, full set vs HAR benchmark: XGBoost, RTH13 (main design: the 13 regular-hours bars, 3,000 days), 16:00 bar
\providecommand{\rhtaDmXgbRthClose}{}\renewcommand{\rhtaDmXgbRthClose}{\ensuremath{-}2.18}  % ... its Diebold-Mariano statistic (negative: full more accurate)
\providecommand{\rhtaPXgbRthClose}{}\renewcommand{\rhtaPXgbRthClose}{\ensuremath{p = 0.030}}  % ... its p statement
\providecommand{\rhtaQHarXgbRthClose}{}\renewcommand{\rhtaQHarXgbRthClose}{0.1234}  % ... mean QLIKE, HAR benchmark
\providecommand{\rhtaQFullXgbRthClose}{}\renewcommand{\rhtaQFullXgbRthClose}{0.1192}  % ... mean QLIKE, full set
\providecommand{\rhtaPctRfRthClose}{}\renewcommand{\rhtaPctRfRthClose}{\ensuremath{-}1.4}  % QLIKE % change, full set vs HAR benchmark: Random forest, RTH13 (main design: the 13 regular-hours bars, 3,000 days), 16:00 bar
\providecommand{\rhtaDmRfRthClose}{}\renewcommand{\rhtaDmRfRthClose}{\ensuremath{-}0.92}  % ... its Diebold-Mariano statistic (negative: full more accurate)
\providecommand{\rhtaPRfRthClose}{}\renewcommand{\rhtaPRfRthClose}{\ensuremath{p = 0.36}}  % ... its p statement
\providecommand{\rhtaQHarRfRthClose}{}\renewcommand{\rhtaQHarRfRthClose}{0.1266}  % ... mean QLIKE, HAR benchmark
\providecommand{\rhtaQFullRfRthClose}{}\renewcommand{\rhtaQFullRfRthClose}{0.1248}  % ... mean QLIKE, full set
\providecommand{\rhtaPctRidgeRthBar}{}\renewcommand{\rhtaPctRidgeRthBar}{\ensuremath{-}5.4}  % QLIKE % change, full set vs HAR benchmark: Ridge, RTH13 (main design: the 13 regular-hours bars, 3,000 days), every bar
\providecommand{\rhtaDmRidgeRthBar}{}\renewcommand{\rhtaDmRidgeRthBar}{\ensuremath{-}5.63}  % ... its Diebold-Mariano statistic (negative: full more accurate)
\providecommand{\rhtaPRidgeRthBar}{}\renewcommand{\rhtaPRidgeRthBar}{\ensuremath{p < 0.001}}  % ... its p statement
\providecommand{\rhtaQHarRidgeRthBar}{}\renewcommand{\rhtaQHarRidgeRthBar}{0.1893}  % ... mean QLIKE, HAR benchmark
\providecommand{\rhtaQFullRidgeRthBar}{}\renewcommand{\rhtaQFullRidgeRthBar}{0.1791}  % ... mean QLIKE, full set
\providecommand{\rhtaPctLassoRthBar}{}\renewcommand{\rhtaPctLassoRthBar}{\ensuremath{-}6.2}  % QLIKE % change, full set vs HAR benchmark: Lasso, RTH13 (main design: the 13 regular-hours bars, 3,000 days), every bar
\providecommand{\rhtaDmLassoRthBar}{}\renewcommand{\rhtaDmLassoRthBar}{\ensuremath{-}4.72}  % ... its Diebold-Mariano statistic (negative: full more accurate)
\providecommand{\rhtaPLassoRthBar}{}\renewcommand{\rhtaPLassoRthBar}{\ensuremath{p < 0.001}}  % ... its p statement
\providecommand{\rhtaQHarLassoRthBar}{}\renewcommand{\rhtaQHarLassoRthBar}{0.1893}  % ... mean QLIKE, HAR benchmark
\providecommand{\rhtaQFullLassoRthBar}{}\renewcommand{\rhtaQFullLassoRthBar}{0.1775}  % ... mean QLIKE, full set
\providecommand{\rhtaPctEnetRthBar}{}\renewcommand{\rhtaPctEnetRthBar}{\ensuremath{-}6.0}  % QLIKE % change, full set vs HAR benchmark: Elastic net, RTH13 (main design: the 13 regular-hours bars, 3,000 days), every bar
\providecommand{\rhtaDmEnetRthBar}{}\renewcommand{\rhtaDmEnetRthBar}{\ensuremath{-}6.98}  % ... its Diebold-Mariano statistic (negative: full more accurate)
\providecommand{\rhtaPEnetRthBar}{}\renewcommand{\rhtaPEnetRthBar}{\ensuremath{p < 0.001}}  % ... its p statement
\providecommand{\rhtaQHarEnetRthBar}{}\renewcommand{\rhtaQHarEnetRthBar}{0.1893}  % ... mean QLIKE, HAR benchmark
\providecommand{\rhtaQFullEnetRthBar}{}\renewcommand{\rhtaQFullEnetRthBar}{0.1778}  % ... mean QLIKE, full set
\providecommand{\rhtaPctLgbmRthBar}{}\renewcommand{\rhtaPctLgbmRthBar}{\ensuremath{-}14.0}  % QLIKE % change, full set vs HAR benchmark: LightGBM, RTH13 (main design: the 13 regular-hours bars, 3,000 days), every bar
\providecommand{\rhtaDmLgbmRthBar}{}\renewcommand{\rhtaDmLgbmRthBar}{\ensuremath{-}10.32}  % ... its Diebold-Mariano statistic (negative: full more accurate)
\providecommand{\rhtaPLgbmRthBar}{}\renewcommand{\rhtaPLgbmRthBar}{\ensuremath{p < 0.001}}  % ... its p statement
\providecommand{\rhtaQHarLgbmRthBar}{}\renewcommand{\rhtaQHarLgbmRthBar}{0.1875}  % ... mean QLIKE, HAR benchmark
\providecommand{\rhtaQFullLgbmRthBar}{}\renewcommand{\rhtaQFullLgbmRthBar}{0.1613}  % ... mean QLIKE, full set
\providecommand{\rhtaPctXgbRthBar}{}\renewcommand{\rhtaPctXgbRthBar}{\ensuremath{-}13.8}  % QLIKE % change, full set vs HAR benchmark: XGBoost, RTH13 (main design: the 13 regular-hours bars, 3,000 days), every bar
\providecommand{\rhtaDmXgbRthBar}{}\renewcommand{\rhtaDmXgbRthBar}{\ensuremath{-}9.94}  % ... its Diebold-Mariano statistic (negative: full more accurate)
\providecommand{\rhtaPXgbRthBar}{}\renewcommand{\rhtaPXgbRthBar}{\ensuremath{p < 0.001}}  % ... its p statement
\providecommand{\rhtaQHarXgbRthBar}{}\renewcommand{\rhtaQHarXgbRthBar}{0.1888}  % ... mean QLIKE, HAR benchmark
\providecommand{\rhtaQFullXgbRthBar}{}\renewcommand{\rhtaQFullXgbRthBar}{0.1627}  % ... mean QLIKE, full set
\providecommand{\rhtaPctRfRthBar}{}\renewcommand{\rhtaPctRfRthBar}{\ensuremath{-}11.1}  % QLIKE % change, full set vs HAR benchmark: Random forest, RTH13 (main design: the 13 regular-hours bars, 3,000 days), every bar
\providecommand{\rhtaDmRfRthBar}{}\renewcommand{\rhtaDmRfRthBar}{\ensuremath{-}9.85}  % ... its Diebold-Mariano statistic (negative: full more accurate)
\providecommand{\rhtaPRfRthBar}{}\renewcommand{\rhtaPRfRthBar}{\ensuremath{p < 0.001}}  % ... its p statement
\providecommand{\rhtaQHarRfRthBar}{}\renewcommand{\rhtaQHarRfRthBar}{0.1895}  % ... mean QLIKE, HAR benchmark
\providecommand{\rhtaQFullRfRthBar}{}\renewcommand{\rhtaQFullRfRthBar}{0.1685}  % ... mean QLIKE, full set
\providecommand{\rhtaPctRidgeAllClose}{}\renewcommand{\rhtaPctRidgeAllClose}{\ensuremath{-}2.2}  % QLIKE % change, full set vs HAR benchmark: Ridge, ALL48 (all 48 intervals, 3,000 days), 16:00 bar
\providecommand{\rhtaDmRidgeAllClose}{}\renewcommand{\rhtaDmRidgeAllClose}{\ensuremath{-}1.18}  % ... its Diebold-Mariano statistic (negative: full more accurate)
\providecommand{\rhtaPRidgeAllClose}{}\renewcommand{\rhtaPRidgeAllClose}{\ensuremath{p = 0.24}}  % ... its p statement
\providecommand{\rhtaQHarRidgeAllClose}{}\renewcommand{\rhtaQHarRidgeAllClose}{0.1262}  % ... mean QLIKE, HAR benchmark
\providecommand{\rhtaQFullRidgeAllClose}{}\renewcommand{\rhtaQFullRidgeAllClose}{0.1234}  % ... mean QLIKE, full set
\providecommand{\rhtaPctLassoAllClose}{}\renewcommand{\rhtaPctLassoAllClose}{\ensuremath{-}2.6}  % QLIKE % change, full set vs HAR benchmark: Lasso, ALL48 (all 48 intervals, 3,000 days), 16:00 bar
\providecommand{\rhtaDmLassoAllClose}{}\renewcommand{\rhtaDmLassoAllClose}{\ensuremath{-}1.45}  % ... its Diebold-Mariano statistic (negative: full more accurate)
\providecommand{\rhtaPLassoAllClose}{}\renewcommand{\rhtaPLassoAllClose}{\ensuremath{p = 0.15}}  % ... its p statement
\providecommand{\rhtaQHarLassoAllClose}{}\renewcommand{\rhtaQHarLassoAllClose}{0.1264}  % ... mean QLIKE, HAR benchmark
\providecommand{\rhtaQFullLassoAllClose}{}\renewcommand{\rhtaQFullLassoAllClose}{0.1232}  % ... mean QLIKE, full set
\providecommand{\rhtaPctEnetAllClose}{}\renewcommand{\rhtaPctEnetAllClose}{\ensuremath{-}1.8}  % QLIKE % change, full set vs HAR benchmark: Elastic net, ALL48 (all 48 intervals, 3,000 days), 16:00 bar
\providecommand{\rhtaDmEnetAllClose}{}\renewcommand{\rhtaDmEnetAllClose}{\ensuremath{-}0.96}  % ... its Diebold-Mariano statistic (negative: full more accurate)
\providecommand{\rhtaPEnetAllClose}{}\renewcommand{\rhtaPEnetAllClose}{\ensuremath{p = 0.34}}  % ... its p statement
\providecommand{\rhtaQHarEnetAllClose}{}\renewcommand{\rhtaQHarEnetAllClose}{0.1259}  % ... mean QLIKE, HAR benchmark
\providecommand{\rhtaQFullEnetAllClose}{}\renewcommand{\rhtaQFullEnetAllClose}{0.1236}  % ... mean QLIKE, full set
\providecommand{\rhtaPctLgbmAllClose}{}\renewcommand{\rhtaPctLgbmAllClose}{\ensuremath{-}6.1}  % QLIKE % change, full set vs HAR benchmark: LightGBM, ALL48 (all 48 intervals, 3,000 days), 16:00 bar
\providecommand{\rhtaDmLgbmAllClose}{}\renewcommand{\rhtaDmLgbmAllClose}{\ensuremath{-}4.45}  % ... its Diebold-Mariano statistic (negative: full more accurate)
\providecommand{\rhtaPLgbmAllClose}{}\renewcommand{\rhtaPLgbmAllClose}{\ensuremath{p < 0.001}}  % ... its p statement
\providecommand{\rhtaQHarLgbmAllClose}{}\renewcommand{\rhtaQHarLgbmAllClose}{0.1250}  % ... mean QLIKE, HAR benchmark
\providecommand{\rhtaQFullLgbmAllClose}{}\renewcommand{\rhtaQFullLgbmAllClose}{0.1174}  % ... mean QLIKE, full set
\providecommand{\rhtaPctXgbAllClose}{}\renewcommand{\rhtaPctXgbAllClose}{\ensuremath{-}6.5}  % QLIKE % change, full set vs HAR benchmark: XGBoost, ALL48 (all 48 intervals, 3,000 days), 16:00 bar
\providecommand{\rhtaDmXgbAllClose}{}\renewcommand{\rhtaDmXgbAllClose}{\ensuremath{-}5.03}  % ... its Diebold-Mariano statistic (negative: full more accurate)
\providecommand{\rhtaPXgbAllClose}{}\renewcommand{\rhtaPXgbAllClose}{\ensuremath{p < 0.001}}  % ... its p statement
\providecommand{\rhtaQHarXgbAllClose}{}\renewcommand{\rhtaQHarXgbAllClose}{0.1251}  % ... mean QLIKE, HAR benchmark
\providecommand{\rhtaQFullXgbAllClose}{}\renewcommand{\rhtaQFullXgbAllClose}{0.1170}  % ... mean QLIKE, full set
\providecommand{\rhtaPctRfAllClose}{}\renewcommand{\rhtaPctRfAllClose}{\ensuremath{-}5.6}  % QLIKE % change, full set vs HAR benchmark: Random forest, ALL48 (all 48 intervals, 3,000 days), 16:00 bar
\providecommand{\rhtaDmRfAllClose}{}\renewcommand{\rhtaDmRfAllClose}{\ensuremath{-}4.11}  % ... its Diebold-Mariano statistic (negative: full more accurate)
\providecommand{\rhtaPRfAllClose}{}\renewcommand{\rhtaPRfAllClose}{\ensuremath{p < 0.001}}  % ... its p statement
\providecommand{\rhtaQHarRfAllClose}{}\renewcommand{\rhtaQHarRfAllClose}{0.1265}  % ... mean QLIKE, HAR benchmark
\providecommand{\rhtaQFullRfAllClose}{}\renewcommand{\rhtaQFullRfAllClose}{0.1194}  % ... mean QLIKE, full set
\providecommand{\rhtaPctRidgeAllBar}{}\renewcommand{\rhtaPctRidgeAllBar}{\ensuremath{-}4.9}  % QLIKE % change, full set vs HAR benchmark: Ridge, ALL48 (all 48 intervals, 3,000 days), every bar
\providecommand{\rhtaDmRidgeAllBar}{}\renewcommand{\rhtaDmRidgeAllBar}{\ensuremath{-}6.71}  % ... its Diebold-Mariano statistic (negative: full more accurate)
\providecommand{\rhtaPRidgeAllBar}{}\renewcommand{\rhtaPRidgeAllBar}{\ensuremath{p < 0.001}}  % ... its p statement
\providecommand{\rhtaQHarRidgeAllBar}{}\renewcommand{\rhtaQHarRidgeAllBar}{0.1919}  % ... mean QLIKE, HAR benchmark
\providecommand{\rhtaQFullRidgeAllBar}{}\renewcommand{\rhtaQFullRidgeAllBar}{0.1825}  % ... mean QLIKE, full set
\providecommand{\rhtaPctLassoAllBar}{}\renewcommand{\rhtaPctLassoAllBar}{\ensuremath{-}5.1}  % QLIKE % change, full set vs HAR benchmark: Lasso, ALL48 (all 48 intervals, 3,000 days), every bar
\providecommand{\rhtaDmLassoAllBar}{}\renewcommand{\rhtaDmLassoAllBar}{\ensuremath{-}7.14}  % ... its Diebold-Mariano statistic (negative: full more accurate)
\providecommand{\rhtaPLassoAllBar}{}\renewcommand{\rhtaPLassoAllBar}{\ensuremath{p < 0.001}}  % ... its p statement
\providecommand{\rhtaQHarLassoAllBar}{}\renewcommand{\rhtaQHarLassoAllBar}{0.1917}  % ... mean QLIKE, HAR benchmark
\providecommand{\rhtaQFullLassoAllBar}{}\renewcommand{\rhtaQFullLassoAllBar}{0.1820}  % ... mean QLIKE, full set
\providecommand{\rhtaPctEnetAllBar}{}\renewcommand{\rhtaPctEnetAllBar}{\ensuremath{-}5.2}  % QLIKE % change, full set vs HAR benchmark: Elastic net, ALL48 (all 48 intervals, 3,000 days), every bar
\providecommand{\rhtaDmEnetAllBar}{}\renewcommand{\rhtaDmEnetAllBar}{\ensuremath{-}6.68}  % ... its Diebold-Mariano statistic (negative: full more accurate)
\providecommand{\rhtaPEnetAllBar}{}\renewcommand{\rhtaPEnetAllBar}{\ensuremath{p < 0.001}}  % ... its p statement
\providecommand{\rhtaQHarEnetAllBar}{}\renewcommand{\rhtaQHarEnetAllBar}{0.1916}  % ... mean QLIKE, HAR benchmark
\providecommand{\rhtaQFullEnetAllBar}{}\renewcommand{\rhtaQFullEnetAllBar}{0.1817}  % ... mean QLIKE, full set
\providecommand{\rhtaPctLgbmAllBar}{}\renewcommand{\rhtaPctLgbmAllBar}{\ensuremath{-}11.2}  % QLIKE % change, full set vs HAR benchmark: LightGBM, ALL48 (all 48 intervals, 3,000 days), every bar
\providecommand{\rhtaDmLgbmAllBar}{}\renewcommand{\rhtaDmLgbmAllBar}{\ensuremath{-}9.84}  % ... its Diebold-Mariano statistic (negative: full more accurate)
\providecommand{\rhtaPLgbmAllBar}{}\renewcommand{\rhtaPLgbmAllBar}{\ensuremath{p < 0.001}}  % ... its p statement
\providecommand{\rhtaQHarLgbmAllBar}{}\renewcommand{\rhtaQHarLgbmAllBar}{0.1898}  % ... mean QLIKE, HAR benchmark
\providecommand{\rhtaQFullLgbmAllBar}{}\renewcommand{\rhtaQFullLgbmAllBar}{0.1685}  % ... mean QLIKE, full set
\providecommand{\rhtaPctXgbAllBar}{}\renewcommand{\rhtaPctXgbAllBar}{\ensuremath{-}12.0}  % QLIKE % change, full set vs HAR benchmark: XGBoost, ALL48 (all 48 intervals, 3,000 days), every bar
\providecommand{\rhtaDmXgbAllBar}{}\renewcommand{\rhtaDmXgbAllBar}{\ensuremath{-}9.70}  % ... its Diebold-Mariano statistic (negative: full more accurate)
\providecommand{\rhtaPXgbAllBar}{}\renewcommand{\rhtaPXgbAllBar}{\ensuremath{p < 0.001}}  % ... its p statement
\providecommand{\rhtaQHarXgbAllBar}{}\renewcommand{\rhtaQHarXgbAllBar}{0.1908}  % ... mean QLIKE, HAR benchmark
\providecommand{\rhtaQFullXgbAllBar}{}\renewcommand{\rhtaQFullXgbAllBar}{0.1679}  % ... mean QLIKE, full set
\providecommand{\rhtaPctRfAllBar}{}\renewcommand{\rhtaPctRfAllBar}{\ensuremath{-}7.0}  % QLIKE % change, full set vs HAR benchmark: Random forest, ALL48 (all 48 intervals, 3,000 days), every bar
\providecommand{\rhtaDmRfAllBar}{}\renewcommand{\rhtaDmRfAllBar}{\ensuremath{-}9.54}  % ... its Diebold-Mariano statistic (negative: full more accurate)
\providecommand{\rhtaPRfAllBar}{}\renewcommand{\rhtaPRfAllBar}{\ensuremath{p < 0.001}}  % ... its p statement
\providecommand{\rhtaQHarRfAllBar}{}\renewcommand{\rhtaQHarRfAllBar}{0.1905}  % ... mean QLIKE, HAR benchmark
\providecommand{\rhtaQFullRfAllBar}{}\renewcommand{\rhtaQFullRfAllBar}{0.1773}  % ... mean QLIKE, full set
\providecommand{\rhtaPctRidgeClClose}{}\renewcommand{\rhtaPctRidgeClClose}{\ensuremath{-}6.1}  % QLIKE % change, full set vs HAR benchmark: Ridge, CLOSE (the closing half hour only, 3,000 closes), 16:00 bar
\providecommand{\rhtaDmRidgeClClose}{}\renewcommand{\rhtaDmRidgeClClose}{\ensuremath{-}3.62}  % ... its Diebold-Mariano statistic (negative: full more accurate)
\providecommand{\rhtaPRidgeClClose}{}\renewcommand{\rhtaPRidgeClClose}{\ensuremath{p < 0.001}}  % ... its p statement
\providecommand{\rhtaQHarRidgeClClose}{}\renewcommand{\rhtaQHarRidgeClClose}{0.1223}  % ... mean QLIKE, HAR benchmark
\providecommand{\rhtaQFullRidgeClClose}{}\renewcommand{\rhtaQFullRidgeClClose}{0.1148}  % ... mean QLIKE, full set
\providecommand{\rhtaPctLassoClClose}{}\renewcommand{\rhtaPctLassoClClose}{\ensuremath{-}5.4}  % QLIKE % change, full set vs HAR benchmark: Lasso, CLOSE (the closing half hour only, 3,000 closes), 16:00 bar
\providecommand{\rhtaDmLassoClClose}{}\renewcommand{\rhtaDmLassoClClose}{\ensuremath{-}5.19}  % ... its Diebold-Mariano statistic (negative: full more accurate)
\providecommand{\rhtaPLassoClClose}{}\renewcommand{\rhtaPLassoClClose}{\ensuremath{p < 0.001}}  % ... its p statement
\providecommand{\rhtaQHarLassoClClose}{}\renewcommand{\rhtaQHarLassoClClose}{0.1223}  % ... mean QLIKE, HAR benchmark
\providecommand{\rhtaQFullLassoClClose}{}\renewcommand{\rhtaQFullLassoClClose}{0.1156}  % ... mean QLIKE, full set
\providecommand{\rhtaPctEnetClClose}{}\renewcommand{\rhtaPctEnetClClose}{\ensuremath{-}6.5}  % QLIKE % change, full set vs HAR benchmark: Elastic net, CLOSE (the closing half hour only, 3,000 closes), 16:00 bar
\providecommand{\rhtaDmEnetClClose}{}\renewcommand{\rhtaDmEnetClClose}{\ensuremath{-}4.59}  % ... its Diebold-Mariano statistic (negative: full more accurate)
\providecommand{\rhtaPEnetClClose}{}\renewcommand{\rhtaPEnetClClose}{\ensuremath{p < 0.001}}  % ... its p statement
\providecommand{\rhtaQHarEnetClClose}{}\renewcommand{\rhtaQHarEnetClClose}{0.1223}  % ... mean QLIKE, HAR benchmark
\providecommand{\rhtaQFullEnetClClose}{}\renewcommand{\rhtaQFullEnetClClose}{0.1144}  % ... mean QLIKE, full set
\providecommand{\rhtaPctLgbmClClose}{}\renewcommand{\rhtaPctLgbmClClose}{\ensuremath{-}7.0}  % QLIKE % change, full set vs HAR benchmark: LightGBM, CLOSE (the closing half hour only, 3,000 closes), 16:00 bar
\providecommand{\rhtaDmLgbmClClose}{}\renewcommand{\rhtaDmLgbmClClose}{\ensuremath{-}4.50}  % ... its Diebold-Mariano statistic (negative: full more accurate)
\providecommand{\rhtaPLgbmClClose}{}\renewcommand{\rhtaPLgbmClClose}{\ensuremath{p < 0.001}}  % ... its p statement
\providecommand{\rhtaQHarLgbmClClose}{}\renewcommand{\rhtaQHarLgbmClClose}{0.1313}  % ... mean QLIKE, HAR benchmark
\providecommand{\rhtaQFullLgbmClClose}{}\renewcommand{\rhtaQFullLgbmClClose}{0.1221}  % ... mean QLIKE, full set
\providecommand{\rhtaPctXgbClClose}{}\renewcommand{\rhtaPctXgbClClose}{\ensuremath{-}8.5}  % QLIKE % change, full set vs HAR benchmark: XGBoost, CLOSE (the closing half hour only, 3,000 closes), 16:00 bar
\providecommand{\rhtaDmXgbClClose}{}\renewcommand{\rhtaDmXgbClClose}{\ensuremath{-}5.95}  % ... its Diebold-Mariano statistic (negative: full more accurate)
\providecommand{\rhtaPXgbClClose}{}\renewcommand{\rhtaPXgbClClose}{\ensuremath{p < 0.001}}  % ... its p statement
\providecommand{\rhtaQHarXgbClClose}{}\renewcommand{\rhtaQHarXgbClClose}{0.1290}  % ... mean QLIKE, HAR benchmark
\providecommand{\rhtaQFullXgbClClose}{}\renewcommand{\rhtaQFullXgbClClose}{0.1180}  % ... mean QLIKE, full set
\providecommand{\rhtaPctRfClClose}{}\renewcommand{\rhtaPctRfClClose}{\ensuremath{-}3.7}  % QLIKE % change, full set vs HAR benchmark: Random forest, CLOSE (the closing half hour only, 3,000 closes), 16:00 bar
\providecommand{\rhtaDmRfClClose}{}\renewcommand{\rhtaDmRfClClose}{\ensuremath{-}3.66}  % ... its Diebold-Mariano statistic (negative: full more accurate)
\providecommand{\rhtaPRfClClose}{}\renewcommand{\rhtaPRfClClose}{\ensuremath{p < 0.001}}  % ... its p statement
\providecommand{\rhtaQHarRfClClose}{}\renewcommand{\rhtaQHarRfClClose}{0.1267}  % ... mean QLIKE, HAR benchmark
\providecommand{\rhtaQFullRfClClose}{}\renewcommand{\rhtaQFullRfClClose}{0.1220}  % ... mean QLIKE, full set
\providecommand{\rhtaPctRidgeClNatClose}{}\renewcommand{\rhtaPctRidgeClNatClose}{\ensuremath{-}6.5}  % QLIKE % change, full set vs HAR benchmark: Ridge, CLOSE_NAT (closing rows, design built natively on them; linear only), 16:00 bar
\providecommand{\rhtaDmRidgeClNatClose}{}\renewcommand{\rhtaDmRidgeClNatClose}{\ensuremath{-}3.76}  % ... its Diebold-Mariano statistic (negative: full more accurate)
\providecommand{\rhtaPRidgeClNatClose}{}\renewcommand{\rhtaPRidgeClNatClose}{\ensuremath{p < 0.001}}  % ... its p statement
\providecommand{\rhtaQHarRidgeClNatClose}{}\renewcommand{\rhtaQHarRidgeClNatClose}{0.1226}  % ... mean QLIKE, HAR benchmark
\providecommand{\rhtaQFullRidgeClNatClose}{}\renewcommand{\rhtaQFullRidgeClNatClose}{0.1146}  % ... mean QLIKE, full set
\providecommand{\rhtaPctLassoClNatClose}{}\renewcommand{\rhtaPctLassoClNatClose}{\ensuremath{-}5.4}  % QLIKE % change, full set vs HAR benchmark: Lasso, CLOSE_NAT (closing rows, design built natively on them; linear only), 16:00 bar
\providecommand{\rhtaDmLassoClNatClose}{}\renewcommand{\rhtaDmLassoClNatClose}{\ensuremath{-}4.95}  % ... its Diebold-Mariano statistic (negative: full more accurate)
\providecommand{\rhtaPLassoClNatClose}{}\renewcommand{\rhtaPLassoClNatClose}{\ensuremath{p < 0.001}}  % ... its p statement
\providecommand{\rhtaQHarLassoClNatClose}{}\renewcommand{\rhtaQHarLassoClNatClose}{0.1226}  % ... mean QLIKE, HAR benchmark
\providecommand{\rhtaQFullLassoClNatClose}{}\renewcommand{\rhtaQFullLassoClNatClose}{0.1159}  % ... mean QLIKE, full set
\providecommand{\rhtaPctEnetClNatClose}{}\renewcommand{\rhtaPctEnetClNatClose}{\ensuremath{-}6.6}  % QLIKE % change, full set vs HAR benchmark: Elastic net, CLOSE_NAT (closing rows, design built natively on them; linear only), 16:00 bar
\providecommand{\rhtaDmEnetClNatClose}{}\renewcommand{\rhtaDmEnetClNatClose}{\ensuremath{-}4.48}  % ... its Diebold-Mariano statistic (negative: full more accurate)
\providecommand{\rhtaPEnetClNatClose}{}\renewcommand{\rhtaPEnetClNatClose}{\ensuremath{p < 0.001}}  % ... its p statement
\providecommand{\rhtaQHarEnetClNatClose}{}\renewcommand{\rhtaQHarEnetClNatClose}{0.1226}  % ... mean QLIKE, HAR benchmark
\providecommand{\rhtaQFullEnetClNatClose}{}\renewcommand{\rhtaQFullEnetClNatClose}{0.1145}  % ... mean QLIKE, full set
\providecommand{\rhtaPctRidgeAllCtClose}{}\renewcommand{\rhtaPctRidgeAllCtClose}{\ensuremath{-}1.9}  % QLIKE % change, full set vs HAR benchmark: Ridge, ALL48@16 (ALL48, penalties chosen on the 16:00 validation rows; linear only), 16:00 bar
\providecommand{\rhtaDmRidgeAllCtClose}{}\renewcommand{\rhtaDmRidgeAllCtClose}{\ensuremath{-}0.97}  % ... its Diebold-Mariano statistic (negative: full more accurate)
\providecommand{\rhtaPRidgeAllCtClose}{}\renewcommand{\rhtaPRidgeAllCtClose}{\ensuremath{p = 0.33}}  % ... its p statement
\providecommand{\rhtaQHarRidgeAllCtClose}{}\renewcommand{\rhtaQHarRidgeAllCtClose}{0.1262}  % ... mean QLIKE, HAR benchmark
\providecommand{\rhtaQFullRidgeAllCtClose}{}\renewcommand{\rhtaQFullRidgeAllCtClose}{0.1238}  % ... mean QLIKE, full set
\providecommand{\rhtaPctLassoAllCtClose}{}\renewcommand{\rhtaPctLassoAllCtClose}{\ensuremath{-}1.1}  % QLIKE % change, full set vs HAR benchmark: Lasso, ALL48@16 (ALL48, penalties chosen on the 16:00 validation rows; linear only), 16:00 bar
\providecommand{\rhtaDmLassoAllCtClose}{}\renewcommand{\rhtaDmLassoAllCtClose}{\ensuremath{-}0.53}  % ... its Diebold-Mariano statistic (negative: full more accurate)
\providecommand{\rhtaPLassoAllCtClose}{}\renewcommand{\rhtaPLassoAllCtClose}{\ensuremath{p = 0.59}}  % ... its p statement
\providecommand{\rhtaQHarLassoAllCtClose}{}\renewcommand{\rhtaQHarLassoAllCtClose}{0.1253}  % ... mean QLIKE, HAR benchmark
\providecommand{\rhtaQFullLassoAllCtClose}{}\renewcommand{\rhtaQFullLassoAllCtClose}{0.1239}  % ... mean QLIKE, full set
\providecommand{\rhtaPctEnetAllCtClose}{}\renewcommand{\rhtaPctEnetAllCtClose}{\ensuremath{-}0.8}  % QLIKE % change, full set vs HAR benchmark: Elastic net, ALL48@16 (ALL48, penalties chosen on the 16:00 validation rows; linear only), 16:00 bar
\providecommand{\rhtaDmEnetAllCtClose}{}\renewcommand{\rhtaDmEnetAllCtClose}{\ensuremath{-}0.40}  % ... its Diebold-Mariano statistic (negative: full more accurate)
\providecommand{\rhtaPEnetAllCtClose}{}\renewcommand{\rhtaPEnetAllCtClose}{\ensuremath{p = 0.69}}  % ... its p statement
\providecommand{\rhtaQHarEnetAllCtClose}{}\renewcommand{\rhtaQHarEnetAllCtClose}{0.1246}  % ... mean QLIKE, HAR benchmark
\providecommand{\rhtaQFullEnetAllCtClose}{}\renewcommand{\rhtaQFullEnetAllCtClose}{0.1236}  % ... mean QLIKE, full set
\providecommand{\rhtaPctRidgeAllCtBar}{}\renewcommand{\rhtaPctRidgeAllCtBar}{\ensuremath{-}5.2}  % QLIKE % change, full set vs HAR benchmark: Ridge, ALL48@16 (ALL48, penalties chosen on the 16:00 validation rows; linear only), every bar
\providecommand{\rhtaDmRidgeAllCtBar}{}\renewcommand{\rhtaDmRidgeAllCtBar}{\ensuremath{-}6.22}  % ... its Diebold-Mariano statistic (negative: full more accurate)
\providecommand{\rhtaPRidgeAllCtBar}{}\renewcommand{\rhtaPRidgeAllCtBar}{\ensuremath{p < 0.001}}  % ... its p statement
\providecommand{\rhtaQHarRidgeAllCtBar}{}\renewcommand{\rhtaQHarRidgeAllCtBar}{0.1919}  % ... mean QLIKE, HAR benchmark
\providecommand{\rhtaQFullRidgeAllCtBar}{}\renewcommand{\rhtaQFullRidgeAllCtBar}{0.1820}  % ... mean QLIKE, full set
\providecommand{\rhtaPctLassoAllCtBar}{}\renewcommand{\rhtaPctLassoAllCtBar}{\ensuremath{-}5.3}  % QLIKE % change, full set vs HAR benchmark: Lasso, ALL48@16 (ALL48, penalties chosen on the 16:00 validation rows; linear only), every bar
\providecommand{\rhtaDmLassoAllCtBar}{}\renewcommand{\rhtaDmLassoAllCtBar}{\ensuremath{-}6.18}  % ... its Diebold-Mariano statistic (negative: full more accurate)
\providecommand{\rhtaPLassoAllCtBar}{}\renewcommand{\rhtaPLassoAllCtBar}{\ensuremath{p < 0.001}}  % ... its p statement
\providecommand{\rhtaQHarLassoAllCtBar}{}\renewcommand{\rhtaQHarLassoAllCtBar}{0.1915}  % ... mean QLIKE, HAR benchmark
\providecommand{\rhtaQFullLassoAllCtBar}{}\renewcommand{\rhtaQFullLassoAllCtBar}{0.1814}  % ... mean QLIKE, full set
\providecommand{\rhtaPctEnetAllCtBar}{}\renewcommand{\rhtaPctEnetAllCtBar}{\ensuremath{-}5.5}  % QLIKE % change, full set vs HAR benchmark: Elastic net, ALL48@16 (ALL48, penalties chosen on the 16:00 validation rows; linear only), every bar
\providecommand{\rhtaDmEnetAllCtBar}{}\renewcommand{\rhtaDmEnetAllCtBar}{\ensuremath{-}7.11}  % ... its Diebold-Mariano statistic (negative: full more accurate)
\providecommand{\rhtaPEnetAllCtBar}{}\renewcommand{\rhtaPEnetAllCtBar}{\ensuremath{p < 0.001}}  % ... its p statement
\providecommand{\rhtaQHarEnetAllCtBar}{}\renewcommand{\rhtaQHarEnetAllCtBar}{0.1922}  % ... mean QLIKE, HAR benchmark
\providecommand{\rhtaQFullEnetAllCtBar}{}\renewcommand{\rhtaQFullEnetAllCtBar}{0.1817}  % ... mean QLIKE, full set
\providecommand{\rhtaPctRidgeRthCtClose}{}\renewcommand{\rhtaPctRidgeRthCtClose}{0.8}  % QLIKE % change, full set vs HAR benchmark: Ridge, RTH13@16 (RTH13, penalties chosen on the 16:00 validation rows; linear only), 16:00 bar
\providecommand{\rhtaDmRidgeRthCtClose}{}\renewcommand{\rhtaDmRidgeRthCtClose}{0.31}  % ... its Diebold-Mariano statistic (negative: full more accurate)
\providecommand{\rhtaPRidgeRthCtClose}{}\renewcommand{\rhtaPRidgeRthCtClose}{\ensuremath{p = 0.76}}  % ... its p statement
\providecommand{\rhtaQHarRidgeRthCtClose}{}\renewcommand{\rhtaQHarRidgeRthCtClose}{0.1269}  % ... mean QLIKE, HAR benchmark
\providecommand{\rhtaQFullRidgeRthCtClose}{}\renewcommand{\rhtaQFullRidgeRthCtClose}{0.1279}  % ... mean QLIKE, full set
\providecommand{\rhtaPctLassoRthCtClose}{}\renewcommand{\rhtaPctLassoRthCtClose}{\ensuremath{-}1.8}  % QLIKE % change, full set vs HAR benchmark: Lasso, RTH13@16 (RTH13, penalties chosen on the 16:00 validation rows; linear only), 16:00 bar
\providecommand{\rhtaDmLassoRthCtClose}{}\renewcommand{\rhtaDmLassoRthCtClose}{\ensuremath{-}0.87}  % ... its Diebold-Mariano statistic (negative: full more accurate)
\providecommand{\rhtaPLassoRthCtClose}{}\renewcommand{\rhtaPLassoRthCtClose}{\ensuremath{p = 0.38}}  % ... its p statement
\providecommand{\rhtaQHarLassoRthCtClose}{}\renewcommand{\rhtaQHarLassoRthCtClose}{0.1262}  % ... mean QLIKE, HAR benchmark
\providecommand{\rhtaQFullLassoRthCtClose}{}\renewcommand{\rhtaQFullLassoRthCtClose}{0.1239}  % ... mean QLIKE, full set
\providecommand{\rhtaPctEnetRthCtClose}{}\renewcommand{\rhtaPctEnetRthCtClose}{1.5}  % QLIKE % change, full set vs HAR benchmark: Elastic net, RTH13@16 (RTH13, penalties chosen on the 16:00 validation rows; linear only), 16:00 bar
\providecommand{\rhtaDmEnetRthCtClose}{}\renewcommand{\rhtaDmEnetRthCtClose}{0.59}  % ... its Diebold-Mariano statistic (negative: full more accurate)
\providecommand{\rhtaPEnetRthCtClose}{}\renewcommand{\rhtaPEnetRthCtClose}{\ensuremath{p = 0.56}}  % ... its p statement
\providecommand{\rhtaQHarEnetRthCtClose}{}\renewcommand{\rhtaQHarEnetRthCtClose}{0.1243}  % ... mean QLIKE, HAR benchmark
\providecommand{\rhtaQFullEnetRthCtClose}{}\renewcommand{\rhtaQFullEnetRthCtClose}{0.1262}  % ... mean QLIKE, full set
\providecommand{\rhtaPctRidgeRthCtBar}{}\renewcommand{\rhtaPctRidgeRthCtBar}{\ensuremath{-}5.4}  % QLIKE % change, full set vs HAR benchmark: Ridge, RTH13@16 (RTH13, penalties chosen on the 16:00 validation rows; linear only), every bar
\providecommand{\rhtaDmRidgeRthCtBar}{}\renewcommand{\rhtaDmRidgeRthCtBar}{\ensuremath{-}5.63}  % ... its Diebold-Mariano statistic (negative: full more accurate)
\providecommand{\rhtaPRidgeRthCtBar}{}\renewcommand{\rhtaPRidgeRthCtBar}{\ensuremath{p < 0.001}}  % ... its p statement
\providecommand{\rhtaQHarRidgeRthCtBar}{}\renewcommand{\rhtaQHarRidgeRthCtBar}{0.1893}  % ... mean QLIKE, HAR benchmark
\providecommand{\rhtaQFullRidgeRthCtBar}{}\renewcommand{\rhtaQFullRidgeRthCtBar}{0.1791}  % ... mean QLIKE, full set
\providecommand{\rhtaPctLassoRthCtBar}{}\renewcommand{\rhtaPctLassoRthCtBar}{\ensuremath{-}6.2}  % QLIKE % change, full set vs HAR benchmark: Lasso, RTH13@16 (RTH13, penalties chosen on the 16:00 validation rows; linear only), every bar
\providecommand{\rhtaDmLassoRthCtBar}{}\renewcommand{\rhtaDmLassoRthCtBar}{\ensuremath{-}8.49}  % ... its Diebold-Mariano statistic (negative: full more accurate)
\providecommand{\rhtaPLassoRthCtBar}{}\renewcommand{\rhtaPLassoRthCtBar}{\ensuremath{p < 0.001}}  % ... its p statement
\providecommand{\rhtaQHarLassoRthCtBar}{}\renewcommand{\rhtaQHarLassoRthCtBar}{0.1896}  % ... mean QLIKE, HAR benchmark
\providecommand{\rhtaQFullLassoRthCtBar}{}\renewcommand{\rhtaQFullLassoRthCtBar}{0.1777}  % ... mean QLIKE, full set
\providecommand{\rhtaPctEnetRthCtBar}{}\renewcommand{\rhtaPctEnetRthCtBar}{\ensuremath{-}6.9}  % QLIKE % change, full set vs HAR benchmark: Elastic net, RTH13@16 (RTH13, penalties chosen on the 16:00 validation rows; linear only), every bar
\providecommand{\rhtaDmEnetRthCtBar}{}\renewcommand{\rhtaDmEnetRthCtBar}{\ensuremath{-}7.94}  % ... its Diebold-Mariano statistic (negative: full more accurate)
\providecommand{\rhtaPEnetRthCtBar}{}\renewcommand{\rhtaPEnetRthCtBar}{\ensuremath{p < 0.001}}  % ... its p statement
\providecommand{\rhtaQHarEnetRthCtBar}{}\renewcommand{\rhtaQHarEnetRthCtBar}{0.1911}  % ... mean QLIKE, HAR benchmark
\providecommand{\rhtaQFullEnetRthCtBar}{}\renewcommand{\rhtaQFullEnetRthCtBar}{0.1778}  % ... mean QLIKE, full set
\providecommand{\rhtaPctRidgeAllThreeKClose}{}\renewcommand{\rhtaPctRidgeAllThreeKClose}{\ensuremath{-}4.7}  % QLIKE % change, full set vs HAR benchmark: Ridge, ALL48_N3k (ALL48, the last 3,000 rows), 16:00 bar
\providecommand{\rhtaDmRidgeAllThreeKClose}{}\renewcommand{\rhtaDmRidgeAllThreeKClose}{\ensuremath{-}4.44}  % ... its Diebold-Mariano statistic (negative: full more accurate)
\providecommand{\rhtaPRidgeAllThreeKClose}{}\renewcommand{\rhtaPRidgeAllThreeKClose}{\ensuremath{p < 0.001}}  % ... its p statement
\providecommand{\rhtaQHarRidgeAllThreeKClose}{}\renewcommand{\rhtaQHarRidgeAllThreeKClose}{0.1341}  % ... mean QLIKE, HAR benchmark
\providecommand{\rhtaQFullRidgeAllThreeKClose}{}\renewcommand{\rhtaQFullRidgeAllThreeKClose}{0.1278}  % ... mean QLIKE, full set
\providecommand{\rhtaPctLassoAllThreeKClose}{}\renewcommand{\rhtaPctLassoAllThreeKClose}{\ensuremath{-}1.6}  % QLIKE % change, full set vs HAR benchmark: Lasso, ALL48_N3k (ALL48, the last 3,000 rows), 16:00 bar
\providecommand{\rhtaDmLassoAllThreeKClose}{}\renewcommand{\rhtaDmLassoAllThreeKClose}{\ensuremath{-}1.25}  % ... its Diebold-Mariano statistic (negative: full more accurate)
\providecommand{\rhtaPLassoAllThreeKClose}{}\renewcommand{\rhtaPLassoAllThreeKClose}{\ensuremath{p = 0.21}}  % ... its p statement
\providecommand{\rhtaQHarLassoAllThreeKClose}{}\renewcommand{\rhtaQHarLassoAllThreeKClose}{0.1322}  % ... mean QLIKE, HAR benchmark
\providecommand{\rhtaQFullLassoAllThreeKClose}{}\renewcommand{\rhtaQFullLassoAllThreeKClose}{0.1301}  % ... mean QLIKE, full set
\providecommand{\rhtaPctEnetAllThreeKClose}{}\renewcommand{\rhtaPctEnetAllThreeKClose}{0.5}  % QLIKE % change, full set vs HAR benchmark: Elastic net, ALL48_N3k (ALL48, the last 3,000 rows), 16:00 bar
\providecommand{\rhtaDmEnetAllThreeKClose}{}\renewcommand{\rhtaDmEnetAllThreeKClose}{0.27}  % ... its Diebold-Mariano statistic (negative: full more accurate)
\providecommand{\rhtaPEnetAllThreeKClose}{}\renewcommand{\rhtaPEnetAllThreeKClose}{\ensuremath{p = 0.78}}  % ... its p statement
\providecommand{\rhtaQHarEnetAllThreeKClose}{}\renewcommand{\rhtaQHarEnetAllThreeKClose}{0.1331}  % ... mean QLIKE, HAR benchmark
\providecommand{\rhtaQFullEnetAllThreeKClose}{}\renewcommand{\rhtaQFullEnetAllThreeKClose}{0.1338}  % ... mean QLIKE, full set
\providecommand{\rhtaPctLgbmAllThreeKClose}{}\renewcommand{\rhtaPctLgbmAllThreeKClose}{\ensuremath{-}8.0}  % QLIKE % change, full set vs HAR benchmark: LightGBM, ALL48_N3k (ALL48, the last 3,000 rows), 16:00 bar
\providecommand{\rhtaDmLgbmAllThreeKClose}{}\renewcommand{\rhtaDmLgbmAllThreeKClose}{\ensuremath{-}3.83}  % ... its Diebold-Mariano statistic (negative: full more accurate)
\providecommand{\rhtaPLgbmAllThreeKClose}{}\renewcommand{\rhtaPLgbmAllThreeKClose}{\ensuremath{p < 0.001}}  % ... its p statement
\providecommand{\rhtaQHarLgbmAllThreeKClose}{}\renewcommand{\rhtaQHarLgbmAllThreeKClose}{0.1409}  % ... mean QLIKE, HAR benchmark
\providecommand{\rhtaQFullLgbmAllThreeKClose}{}\renewcommand{\rhtaQFullLgbmAllThreeKClose}{0.1295}  % ... mean QLIKE, full set
\providecommand{\rhtaPctXgbAllThreeKClose}{}\renewcommand{\rhtaPctXgbAllThreeKClose}{\ensuremath{-}6.6}  % QLIKE % change, full set vs HAR benchmark: XGBoost, ALL48_N3k (ALL48, the last 3,000 rows), 16:00 bar
\providecommand{\rhtaDmXgbAllThreeKClose}{}\renewcommand{\rhtaDmXgbAllThreeKClose}{\ensuremath{-}4.14}  % ... its Diebold-Mariano statistic (negative: full more accurate)
\providecommand{\rhtaPXgbAllThreeKClose}{}\renewcommand{\rhtaPXgbAllThreeKClose}{\ensuremath{p < 0.001}}  % ... its p statement
\providecommand{\rhtaQHarXgbAllThreeKClose}{}\renewcommand{\rhtaQHarXgbAllThreeKClose}{0.1348}  % ... mean QLIKE, HAR benchmark
\providecommand{\rhtaQFullXgbAllThreeKClose}{}\renewcommand{\rhtaQFullXgbAllThreeKClose}{0.1260}  % ... mean QLIKE, full set
\providecommand{\rhtaPctRfAllThreeKClose}{}\renewcommand{\rhtaPctRfAllThreeKClose}{\ensuremath{-}6.9}  % QLIKE % change, full set vs HAR benchmark: Random forest, ALL48_N3k (ALL48, the last 3,000 rows), 16:00 bar
\providecommand{\rhtaDmRfAllThreeKClose}{}\renewcommand{\rhtaDmRfAllThreeKClose}{\ensuremath{-}6.46}  % ... its Diebold-Mariano statistic (negative: full more accurate)
\providecommand{\rhtaPRfAllThreeKClose}{}\renewcommand{\rhtaPRfAllThreeKClose}{\ensuremath{p < 0.001}}  % ... its p statement
\providecommand{\rhtaQHarRfAllThreeKClose}{}\renewcommand{\rhtaQHarRfAllThreeKClose}{0.1355}  % ... mean QLIKE, HAR benchmark
\providecommand{\rhtaQFullRfAllThreeKClose}{}\renewcommand{\rhtaQFullRfAllThreeKClose}{0.1260}  % ... mean QLIKE, full set
\providecommand{\rhtaPctRidgeAllThreeKBar}{}\renewcommand{\rhtaPctRidgeAllThreeKBar}{\ensuremath{-}2.7}  % QLIKE % change, full set vs HAR benchmark: Ridge, ALL48_N3k (ALL48, the last 3,000 rows), every bar
\providecommand{\rhtaDmRidgeAllThreeKBar}{}\renewcommand{\rhtaDmRidgeAllThreeKBar}{\ensuremath{-}7.15}  % ... its Diebold-Mariano statistic (negative: full more accurate)
\providecommand{\rhtaPRidgeAllThreeKBar}{}\renewcommand{\rhtaPRidgeAllThreeKBar}{\ensuremath{p < 0.001}}  % ... its p statement
\providecommand{\rhtaQHarRidgeAllThreeKBar}{}\renewcommand{\rhtaQHarRidgeAllThreeKBar}{0.1976}  % ... mean QLIKE, HAR benchmark
\providecommand{\rhtaQFullRidgeAllThreeKBar}{}\renewcommand{\rhtaQFullRidgeAllThreeKBar}{0.1924}  % ... mean QLIKE, full set
\providecommand{\rhtaPctLassoAllThreeKBar}{}\renewcommand{\rhtaPctLassoAllThreeKBar}{\ensuremath{-}2.1}  % QLIKE % change, full set vs HAR benchmark: Lasso, ALL48_N3k (ALL48, the last 3,000 rows), every bar
\providecommand{\rhtaDmLassoAllThreeKBar}{}\renewcommand{\rhtaDmLassoAllThreeKBar}{\ensuremath{-}4.17}  % ... its Diebold-Mariano statistic (negative: full more accurate)
\providecommand{\rhtaPLassoAllThreeKBar}{}\renewcommand{\rhtaPLassoAllThreeKBar}{\ensuremath{p < 0.001}}  % ... its p statement
\providecommand{\rhtaQHarLassoAllThreeKBar}{}\renewcommand{\rhtaQHarLassoAllThreeKBar}{0.1967}  % ... mean QLIKE, HAR benchmark
\providecommand{\rhtaQFullLassoAllThreeKBar}{}\renewcommand{\rhtaQFullLassoAllThreeKBar}{0.1927}  % ... mean QLIKE, full set
\providecommand{\rhtaPctEnetAllThreeKBar}{}\renewcommand{\rhtaPctEnetAllThreeKBar}{\ensuremath{-}1.8}  % QLIKE % change, full set vs HAR benchmark: Elastic net, ALL48_N3k (ALL48, the last 3,000 rows), every bar
\providecommand{\rhtaDmEnetAllThreeKBar}{}\renewcommand{\rhtaDmEnetAllThreeKBar}{\ensuremath{-}2.57}  % ... its Diebold-Mariano statistic (negative: full more accurate)
\providecommand{\rhtaPEnetAllThreeKBar}{}\renewcommand{\rhtaPEnetAllThreeKBar}{\ensuremath{p = 0.0101}}  % ... its p statement
\providecommand{\rhtaQHarEnetAllThreeKBar}{}\renewcommand{\rhtaQHarEnetAllThreeKBar}{0.1972}  % ... mean QLIKE, HAR benchmark
\providecommand{\rhtaQFullEnetAllThreeKBar}{}\renewcommand{\rhtaQFullEnetAllThreeKBar}{0.1936}  % ... mean QLIKE, full set
\providecommand{\rhtaPctLgbmAllThreeKBar}{}\renewcommand{\rhtaPctLgbmAllThreeKBar}{\ensuremath{-}5.9}  % QLIKE % change, full set vs HAR benchmark: LightGBM, ALL48_N3k (ALL48, the last 3,000 rows), every bar
\providecommand{\rhtaDmLgbmAllThreeKBar}{}\renewcommand{\rhtaDmLgbmAllThreeKBar}{\ensuremath{-}8.86}  % ... its Diebold-Mariano statistic (negative: full more accurate)
\providecommand{\rhtaPLgbmAllThreeKBar}{}\renewcommand{\rhtaPLgbmAllThreeKBar}{\ensuremath{p < 0.001}}  % ... its p statement
\providecommand{\rhtaQHarLgbmAllThreeKBar}{}\renewcommand{\rhtaQHarLgbmAllThreeKBar}{0.2062}  % ... mean QLIKE, HAR benchmark
\providecommand{\rhtaQFullLgbmAllThreeKBar}{}\renewcommand{\rhtaQFullLgbmAllThreeKBar}{0.1941}  % ... mean QLIKE, full set
\providecommand{\rhtaPctXgbAllThreeKBar}{}\renewcommand{\rhtaPctXgbAllThreeKBar}{\ensuremath{-}5.9}  % QLIKE % change, full set vs HAR benchmark: XGBoost, ALL48_N3k (ALL48, the last 3,000 rows), every bar
\providecommand{\rhtaDmXgbAllThreeKBar}{}\renewcommand{\rhtaDmXgbAllThreeKBar}{\ensuremath{-}8.32}  % ... its Diebold-Mariano statistic (negative: full more accurate)
\providecommand{\rhtaPXgbAllThreeKBar}{}\renewcommand{\rhtaPXgbAllThreeKBar}{\ensuremath{p < 0.001}}  % ... its p statement
\providecommand{\rhtaQHarXgbAllThreeKBar}{}\renewcommand{\rhtaQHarXgbAllThreeKBar}{0.2011}  % ... mean QLIKE, HAR benchmark
\providecommand{\rhtaQFullXgbAllThreeKBar}{}\renewcommand{\rhtaQFullXgbAllThreeKBar}{0.1893}  % ... mean QLIKE, full set
\providecommand{\rhtaPctRfAllThreeKBar}{}\renewcommand{\rhtaPctRfAllThreeKBar}{\ensuremath{-}5.0}  % QLIKE % change, full set vs HAR benchmark: Random forest, ALL48_N3k (ALL48, the last 3,000 rows), every bar
\providecommand{\rhtaDmRfAllThreeKBar}{}\renewcommand{\rhtaDmRfAllThreeKBar}{\ensuremath{-}8.83}  % ... its Diebold-Mariano statistic (negative: full more accurate)
\providecommand{\rhtaPRfAllThreeKBar}{}\renewcommand{\rhtaPRfAllThreeKBar}{\ensuremath{p < 0.001}}  % ... its p statement
\providecommand{\rhtaQHarRfAllThreeKBar}{}\renewcommand{\rhtaQHarRfAllThreeKBar}{0.1994}  % ... mean QLIKE, HAR benchmark
\providecommand{\rhtaQFullRfAllThreeKBar}{}\renewcommand{\rhtaQFullRfAllThreeKBar}{0.1894}  % ... mean QLIKE, full set
\providecommand{\rhtaPctRidgeRthThreeKClose}{}\renewcommand{\rhtaPctRidgeRthThreeKClose}{\ensuremath{-}1.9}  % QLIKE % change, full set vs HAR benchmark: Ridge, RTH13_N3k (RTH13, the last 3,000 rows), 16:00 bar
\providecommand{\rhtaDmRidgeRthThreeKClose}{}\renewcommand{\rhtaDmRidgeRthThreeKClose}{\ensuremath{-}1.55}  % ... its Diebold-Mariano statistic (negative: full more accurate)
\providecommand{\rhtaPRidgeRthThreeKClose}{}\renewcommand{\rhtaPRidgeRthThreeKClose}{\ensuremath{p = 0.12}}  % ... its p statement
\providecommand{\rhtaQHarRidgeRthThreeKClose}{}\renewcommand{\rhtaQHarRidgeRthThreeKClose}{0.1418}  % ... mean QLIKE, HAR benchmark
\providecommand{\rhtaQFullRidgeRthThreeKClose}{}\renewcommand{\rhtaQFullRidgeRthThreeKClose}{0.1391}  % ... mean QLIKE, full set
\providecommand{\rhtaPctLassoRthThreeKClose}{}\renewcommand{\rhtaPctLassoRthThreeKClose}{\ensuremath{-}1.9}  % QLIKE % change, full set vs HAR benchmark: Lasso, RTH13_N3k (RTH13, the last 3,000 rows), 16:00 bar
\providecommand{\rhtaDmLassoRthThreeKClose}{}\renewcommand{\rhtaDmLassoRthThreeKClose}{\ensuremath{-}1.29}  % ... its Diebold-Mariano statistic (negative: full more accurate)
\providecommand{\rhtaPLassoRthThreeKClose}{}\renewcommand{\rhtaPLassoRthThreeKClose}{\ensuremath{p = 0.20}}  % ... its p statement
\providecommand{\rhtaQHarLassoRthThreeKClose}{}\renewcommand{\rhtaQHarLassoRthThreeKClose}{0.1418}  % ... mean QLIKE, HAR benchmark
\providecommand{\rhtaQFullLassoRthThreeKClose}{}\renewcommand{\rhtaQFullLassoRthThreeKClose}{0.1392}  % ... mean QLIKE, full set
\providecommand{\rhtaPctEnetRthThreeKClose}{}\renewcommand{\rhtaPctEnetRthThreeKClose}{\ensuremath{-}1.6}  % QLIKE % change, full set vs HAR benchmark: Elastic net, RTH13_N3k (RTH13, the last 3,000 rows), 16:00 bar
\providecommand{\rhtaDmEnetRthThreeKClose}{}\renewcommand{\rhtaDmEnetRthThreeKClose}{\ensuremath{-}0.85}  % ... its Diebold-Mariano statistic (negative: full more accurate)
\providecommand{\rhtaPEnetRthThreeKClose}{}\renewcommand{\rhtaPEnetRthThreeKClose}{\ensuremath{p = 0.39}}  % ... its p statement
\providecommand{\rhtaQHarEnetRthThreeKClose}{}\renewcommand{\rhtaQHarEnetRthThreeKClose}{0.1418}  % ... mean QLIKE, HAR benchmark
\providecommand{\rhtaQFullEnetRthThreeKClose}{}\renewcommand{\rhtaQFullEnetRthThreeKClose}{0.1396}  % ... mean QLIKE, full set
\providecommand{\rhtaPctLgbmRthThreeKClose}{}\renewcommand{\rhtaPctLgbmRthThreeKClose}{\ensuremath{-}3.7}  % QLIKE % change, full set vs HAR benchmark: LightGBM, RTH13_N3k (RTH13, the last 3,000 rows), 16:00 bar
\providecommand{\rhtaDmLgbmRthThreeKClose}{}\renewcommand{\rhtaDmLgbmRthThreeKClose}{\ensuremath{-}2.15}  % ... its Diebold-Mariano statistic (negative: full more accurate)
\providecommand{\rhtaPLgbmRthThreeKClose}{}\renewcommand{\rhtaPLgbmRthThreeKClose}{\ensuremath{p = 0.031}}  % ... its p statement
\providecommand{\rhtaQHarLgbmRthThreeKClose}{}\renewcommand{\rhtaQHarLgbmRthThreeKClose}{0.1401}  % ... mean QLIKE, HAR benchmark
\providecommand{\rhtaQFullLgbmRthThreeKClose}{}\renewcommand{\rhtaQFullLgbmRthThreeKClose}{0.1349}  % ... mean QLIKE, full set
\providecommand{\rhtaPctXgbRthThreeKClose}{}\renewcommand{\rhtaPctXgbRthThreeKClose}{\ensuremath{-}2.3}  % QLIKE % change, full set vs HAR benchmark: XGBoost, RTH13_N3k (RTH13, the last 3,000 rows), 16:00 bar
\providecommand{\rhtaDmXgbRthThreeKClose}{}\renewcommand{\rhtaDmXgbRthThreeKClose}{\ensuremath{-}1.26}  % ... its Diebold-Mariano statistic (negative: full more accurate)
\providecommand{\rhtaPXgbRthThreeKClose}{}\renewcommand{\rhtaPXgbRthThreeKClose}{\ensuremath{p = 0.21}}  % ... its p statement
\providecommand{\rhtaQHarXgbRthThreeKClose}{}\renewcommand{\rhtaQHarXgbRthThreeKClose}{0.1334}  % ... mean QLIKE, HAR benchmark
\providecommand{\rhtaQFullXgbRthThreeKClose}{}\renewcommand{\rhtaQFullXgbRthThreeKClose}{0.1303}  % ... mean QLIKE, full set
\providecommand{\rhtaPctRfRthThreeKClose}{}\renewcommand{\rhtaPctRfRthThreeKClose}{\ensuremath{-}4.5}  % QLIKE % change, full set vs HAR benchmark: Random forest, RTH13_N3k (RTH13, the last 3,000 rows), 16:00 bar
\providecommand{\rhtaDmRfRthThreeKClose}{}\renewcommand{\rhtaDmRfRthThreeKClose}{\ensuremath{-}3.54}  % ... its Diebold-Mariano statistic (negative: full more accurate)
\providecommand{\rhtaPRfRthThreeKClose}{}\renewcommand{\rhtaPRfRthThreeKClose}{\ensuremath{p < 0.001}}  % ... its p statement
\providecommand{\rhtaQHarRfRthThreeKClose}{}\renewcommand{\rhtaQHarRfRthThreeKClose}{0.1400}  % ... mean QLIKE, HAR benchmark
\providecommand{\rhtaQFullRfRthThreeKClose}{}\renewcommand{\rhtaQFullRfRthThreeKClose}{0.1336}  % ... mean QLIKE, full set
\providecommand{\rhtaPctRidgeRthThreeKBar}{}\renewcommand{\rhtaPctRidgeRthThreeKBar}{\ensuremath{-}3.6}  % QLIKE % change, full set vs HAR benchmark: Ridge, RTH13_N3k (RTH13, the last 3,000 rows), every bar
\providecommand{\rhtaDmRidgeRthThreeKBar}{}\renewcommand{\rhtaDmRidgeRthThreeKBar}{\ensuremath{-}9.44}  % ... its Diebold-Mariano statistic (negative: full more accurate)
\providecommand{\rhtaPRidgeRthThreeKBar}{}\renewcommand{\rhtaPRidgeRthThreeKBar}{\ensuremath{p < 0.001}}  % ... its p statement
\providecommand{\rhtaQHarRidgeRthThreeKBar}{}\renewcommand{\rhtaQHarRidgeRthThreeKBar}{0.1910}  % ... mean QLIKE, HAR benchmark
\providecommand{\rhtaQFullRidgeRthThreeKBar}{}\renewcommand{\rhtaQFullRidgeRthThreeKBar}{0.1840}  % ... mean QLIKE, full set
\providecommand{\rhtaPctLassoRthThreeKBar}{}\renewcommand{\rhtaPctLassoRthThreeKBar}{\ensuremath{-}3.8}  % QLIKE % change, full set vs HAR benchmark: Lasso, RTH13_N3k (RTH13, the last 3,000 rows), every bar
\providecommand{\rhtaDmLassoRthThreeKBar}{}\renewcommand{\rhtaDmLassoRthThreeKBar}{\ensuremath{-}7.70}  % ... its Diebold-Mariano statistic (negative: full more accurate)
\providecommand{\rhtaPLassoRthThreeKBar}{}\renewcommand{\rhtaPLassoRthThreeKBar}{\ensuremath{p < 0.001}}  % ... its p statement
\providecommand{\rhtaQHarLassoRthThreeKBar}{}\renewcommand{\rhtaQHarLassoRthThreeKBar}{0.1910}  % ... mean QLIKE, HAR benchmark
\providecommand{\rhtaQFullLassoRthThreeKBar}{}\renewcommand{\rhtaQFullLassoRthThreeKBar}{0.1837}  % ... mean QLIKE, full set
\providecommand{\rhtaPctEnetRthThreeKBar}{}\renewcommand{\rhtaPctEnetRthThreeKBar}{\ensuremath{-}4.2}  % QLIKE % change, full set vs HAR benchmark: Elastic net, RTH13_N3k (RTH13, the last 3,000 rows), every bar
\providecommand{\rhtaDmEnetRthThreeKBar}{}\renewcommand{\rhtaDmEnetRthThreeKBar}{\ensuremath{-}6.30}  % ... its Diebold-Mariano statistic (negative: full more accurate)
\providecommand{\rhtaPEnetRthThreeKBar}{}\renewcommand{\rhtaPEnetRthThreeKBar}{\ensuremath{p < 0.001}}  % ... its p statement
\providecommand{\rhtaQHarEnetRthThreeKBar}{}\renewcommand{\rhtaQHarEnetRthThreeKBar}{0.1910}  % ... mean QLIKE, HAR benchmark
\providecommand{\rhtaQFullEnetRthThreeKBar}{}\renewcommand{\rhtaQFullEnetRthThreeKBar}{0.1830}  % ... mean QLIKE, full set
\providecommand{\rhtaPctLgbmRthThreeKBar}{}\renewcommand{\rhtaPctLgbmRthThreeKBar}{\ensuremath{-}6.0}  % QLIKE % change, full set vs HAR benchmark: LightGBM, RTH13_N3k (RTH13, the last 3,000 rows), every bar
\providecommand{\rhtaDmLgbmRthThreeKBar}{}\renewcommand{\rhtaDmLgbmRthThreeKBar}{\ensuremath{-}7.85}  % ... its Diebold-Mariano statistic (negative: full more accurate)
\providecommand{\rhtaPLgbmRthThreeKBar}{}\renewcommand{\rhtaPLgbmRthThreeKBar}{\ensuremath{p < 0.001}}  % ... its p statement
\providecommand{\rhtaQHarLgbmRthThreeKBar}{}\renewcommand{\rhtaQHarLgbmRthThreeKBar}{0.1972}  % ... mean QLIKE, HAR benchmark
\providecommand{\rhtaQFullLgbmRthThreeKBar}{}\renewcommand{\rhtaQFullLgbmRthThreeKBar}{0.1854}  % ... mean QLIKE, full set
\providecommand{\rhtaPctXgbRthThreeKBar}{}\renewcommand{\rhtaPctXgbRthThreeKBar}{\ensuremath{-}6.8}  % QLIKE % change, full set vs HAR benchmark: XGBoost, RTH13_N3k (RTH13, the last 3,000 rows), every bar
\providecommand{\rhtaDmXgbRthThreeKBar}{}\renewcommand{\rhtaDmXgbRthThreeKBar}{\ensuremath{-}9.13}  % ... its Diebold-Mariano statistic (negative: full more accurate)
\providecommand{\rhtaPXgbRthThreeKBar}{}\renewcommand{\rhtaPXgbRthThreeKBar}{\ensuremath{p < 0.001}}  % ... its p statement
\providecommand{\rhtaQHarXgbRthThreeKBar}{}\renewcommand{\rhtaQHarXgbRthThreeKBar}{0.1931}  % ... mean QLIKE, HAR benchmark
\providecommand{\rhtaQFullXgbRthThreeKBar}{}\renewcommand{\rhtaQFullXgbRthThreeKBar}{0.1801}  % ... mean QLIKE, full set
\providecommand{\rhtaPctRfRthThreeKBar}{}\renewcommand{\rhtaPctRfRthThreeKBar}{\ensuremath{-}4.7}  % QLIKE % change, full set vs HAR benchmark: Random forest, RTH13_N3k (RTH13, the last 3,000 rows), every bar
\providecommand{\rhtaDmRfRthThreeKBar}{}\renewcommand{\rhtaDmRfRthThreeKBar}{\ensuremath{-}8.98}  % ... its Diebold-Mariano statistic (negative: full more accurate)
\providecommand{\rhtaPRfRthThreeKBar}{}\renewcommand{\rhtaPRfRthThreeKBar}{\ensuremath{p < 0.001}}  % ... its p statement
\providecommand{\rhtaQHarRfRthThreeKBar}{}\renewcommand{\rhtaQHarRfRthThreeKBar}{0.1935}  % ... mean QLIKE, HAR benchmark
\providecommand{\rhtaQFullRfRthThreeKBar}{}\renewcommand{\rhtaQFullRfRthThreeKBar}{0.1844}  % ... mean QLIKE, full set
\providecommand{\rhtaPctRidgeAllThirtyNineKClose}{}\renewcommand{\rhtaPctRidgeAllThirtyNineKClose}{\ensuremath{-}0.9}  % QLIKE % change, full set vs HAR benchmark: Ridge, ALL48_N39k (ALL48, the last 39,000 rows), 16:00 bar
\providecommand{\rhtaDmRidgeAllThirtyNineKClose}{}\renewcommand{\rhtaDmRidgeAllThirtyNineKClose}{\ensuremath{-}0.41}  % ... its Diebold-Mariano statistic (negative: full more accurate)
\providecommand{\rhtaPRidgeAllThirtyNineKClose}{}\renewcommand{\rhtaPRidgeAllThirtyNineKClose}{\ensuremath{p = 0.68}}  % ... its p statement
\providecommand{\rhtaQHarRidgeAllThirtyNineKClose}{}\renewcommand{\rhtaQHarRidgeAllThirtyNineKClose}{0.1296}  % ... mean QLIKE, HAR benchmark
\providecommand{\rhtaQFullRidgeAllThirtyNineKClose}{}\renewcommand{\rhtaQFullRidgeAllThirtyNineKClose}{0.1285}  % ... mean QLIKE, full set
\providecommand{\rhtaPctLassoAllThirtyNineKClose}{}\renewcommand{\rhtaPctLassoAllThirtyNineKClose}{\ensuremath{-}3.6}  % QLIKE % change, full set vs HAR benchmark: Lasso, ALL48_N39k (ALL48, the last 39,000 rows), 16:00 bar
\providecommand{\rhtaDmLassoAllThirtyNineKClose}{}\renewcommand{\rhtaDmLassoAllThirtyNineKClose}{\ensuremath{-}2.19}  % ... its Diebold-Mariano statistic (negative: full more accurate)
\providecommand{\rhtaPLassoAllThirtyNineKClose}{}\renewcommand{\rhtaPLassoAllThirtyNineKClose}{\ensuremath{p = 0.028}}  % ... its p statement
\providecommand{\rhtaQHarLassoAllThirtyNineKClose}{}\renewcommand{\rhtaQHarLassoAllThirtyNineKClose}{0.1311}  % ... mean QLIKE, HAR benchmark
\providecommand{\rhtaQFullLassoAllThirtyNineKClose}{}\renewcommand{\rhtaQFullLassoAllThirtyNineKClose}{0.1264}  % ... mean QLIKE, full set
\providecommand{\rhtaPctEnetAllThirtyNineKClose}{}\renewcommand{\rhtaPctEnetAllThirtyNineKClose}{\ensuremath{-}3.7}  % QLIKE % change, full set vs HAR benchmark: Elastic net, ALL48_N39k (ALL48, the last 39,000 rows), 16:00 bar
\providecommand{\rhtaDmEnetAllThirtyNineKClose}{}\renewcommand{\rhtaDmEnetAllThirtyNineKClose}{\ensuremath{-}3.22}  % ... its Diebold-Mariano statistic (negative: full more accurate)
\providecommand{\rhtaPEnetAllThirtyNineKClose}{}\renewcommand{\rhtaPEnetAllThirtyNineKClose}{\ensuremath{p = 0.001}}  % ... its p statement
\providecommand{\rhtaQHarEnetAllThirtyNineKClose}{}\renewcommand{\rhtaQHarEnetAllThirtyNineKClose}{0.1312}  % ... mean QLIKE, HAR benchmark
\providecommand{\rhtaQFullEnetAllThirtyNineKClose}{}\renewcommand{\rhtaQFullEnetAllThirtyNineKClose}{0.1263}  % ... mean QLIKE, full set
\providecommand{\rhtaPctLgbmAllThirtyNineKClose}{}\renewcommand{\rhtaPctLgbmAllThirtyNineKClose}{\ensuremath{-}5.5}  % QLIKE % change, full set vs HAR benchmark: LightGBM, ALL48_N39k (ALL48, the last 39,000 rows), 16:00 bar
\providecommand{\rhtaDmLgbmAllThirtyNineKClose}{}\renewcommand{\rhtaDmLgbmAllThirtyNineKClose}{\ensuremath{-}4.12}  % ... its Diebold-Mariano statistic (negative: full more accurate)
\providecommand{\rhtaPLgbmAllThirtyNineKClose}{}\renewcommand{\rhtaPLgbmAllThirtyNineKClose}{\ensuremath{p < 0.001}}  % ... its p statement
\providecommand{\rhtaQHarLgbmAllThirtyNineKClose}{}\renewcommand{\rhtaQHarLgbmAllThirtyNineKClose}{0.1276}  % ... mean QLIKE, HAR benchmark
\providecommand{\rhtaQFullLgbmAllThirtyNineKClose}{}\renewcommand{\rhtaQFullLgbmAllThirtyNineKClose}{0.1206}  % ... mean QLIKE, full set
\providecommand{\rhtaPctXgbAllThirtyNineKClose}{}\renewcommand{\rhtaPctXgbAllThirtyNineKClose}{\ensuremath{-}4.7}  % QLIKE % change, full set vs HAR benchmark: XGBoost, ALL48_N39k (ALL48, the last 39,000 rows), 16:00 bar
\providecommand{\rhtaDmXgbAllThirtyNineKClose}{}\renewcommand{\rhtaDmXgbAllThirtyNineKClose}{\ensuremath{-}3.66}  % ... its Diebold-Mariano statistic (negative: full more accurate)
\providecommand{\rhtaPXgbAllThirtyNineKClose}{}\renewcommand{\rhtaPXgbAllThirtyNineKClose}{\ensuremath{p < 0.001}}  % ... its p statement
\providecommand{\rhtaQHarXgbAllThirtyNineKClose}{}\renewcommand{\rhtaQHarXgbAllThirtyNineKClose}{0.1267}  % ... mean QLIKE, HAR benchmark
\providecommand{\rhtaQFullXgbAllThirtyNineKClose}{}\renewcommand{\rhtaQFullXgbAllThirtyNineKClose}{0.1208}  % ... mean QLIKE, full set
\providecommand{\rhtaPctRfAllThirtyNineKClose}{}\renewcommand{\rhtaPctRfAllThirtyNineKClose}{\ensuremath{-}4.9}  % QLIKE % change, full set vs HAR benchmark: Random forest, ALL48_N39k (ALL48, the last 39,000 rows), 16:00 bar
\providecommand{\rhtaDmRfAllThirtyNineKClose}{}\renewcommand{\rhtaDmRfAllThirtyNineKClose}{\ensuremath{-}3.89}  % ... its Diebold-Mariano statistic (negative: full more accurate)
\providecommand{\rhtaPRfAllThirtyNineKClose}{}\renewcommand{\rhtaPRfAllThirtyNineKClose}{\ensuremath{p < 0.001}}  % ... its p statement
\providecommand{\rhtaQHarRfAllThirtyNineKClose}{}\renewcommand{\rhtaQHarRfAllThirtyNineKClose}{0.1277}  % ... mean QLIKE, HAR benchmark
\providecommand{\rhtaQFullRfAllThirtyNineKClose}{}\renewcommand{\rhtaQFullRfAllThirtyNineKClose}{0.1214}  % ... mean QLIKE, full set
\providecommand{\rhtaPctRidgeAllThirtyNineKBar}{}\renewcommand{\rhtaPctRidgeAllThirtyNineKBar}{\ensuremath{-}4.4}  % QLIKE % change, full set vs HAR benchmark: Ridge, ALL48_N39k (ALL48, the last 39,000 rows), every bar
\providecommand{\rhtaDmRidgeAllThirtyNineKBar}{}\renewcommand{\rhtaDmRidgeAllThirtyNineKBar}{\ensuremath{-}4.51}  % ... its Diebold-Mariano statistic (negative: full more accurate)
\providecommand{\rhtaPRidgeAllThirtyNineKBar}{}\renewcommand{\rhtaPRidgeAllThirtyNineKBar}{\ensuremath{p < 0.001}}  % ... its p statement
\providecommand{\rhtaQHarRidgeAllThirtyNineKBar}{}\renewcommand{\rhtaQHarRidgeAllThirtyNineKBar}{0.1925}  % ... mean QLIKE, HAR benchmark
\providecommand{\rhtaQFullRidgeAllThirtyNineKBar}{}\renewcommand{\rhtaQFullRidgeAllThirtyNineKBar}{0.1840}  % ... mean QLIKE, full set
\providecommand{\rhtaPctLassoAllThirtyNineKBar}{}\renewcommand{\rhtaPctLassoAllThirtyNineKBar}{\ensuremath{-}3.8}  % QLIKE % change, full set vs HAR benchmark: Lasso, ALL48_N39k (ALL48, the last 39,000 rows), every bar
\providecommand{\rhtaDmLassoAllThirtyNineKBar}{}\renewcommand{\rhtaDmLassoAllThirtyNineKBar}{\ensuremath{-}7.03}  % ... its Diebold-Mariano statistic (negative: full more accurate)
\providecommand{\rhtaPLassoAllThirtyNineKBar}{}\renewcommand{\rhtaPLassoAllThirtyNineKBar}{\ensuremath{p < 0.001}}  % ... its p statement
\providecommand{\rhtaQHarLassoAllThirtyNineKBar}{}\renewcommand{\rhtaQHarLassoAllThirtyNineKBar}{0.1922}  % ... mean QLIKE, HAR benchmark
\providecommand{\rhtaQFullLassoAllThirtyNineKBar}{}\renewcommand{\rhtaQFullLassoAllThirtyNineKBar}{0.1849}  % ... mean QLIKE, full set
\providecommand{\rhtaPctEnetAllThirtyNineKBar}{}\renewcommand{\rhtaPctEnetAllThirtyNineKBar}{\ensuremath{-}3.0}  % QLIKE % change, full set vs HAR benchmark: Elastic net, ALL48_N39k (ALL48, the last 39,000 rows), every bar
\providecommand{\rhtaDmEnetAllThirtyNineKBar}{}\renewcommand{\rhtaDmEnetAllThirtyNineKBar}{\ensuremath{-}8.37}  % ... its Diebold-Mariano statistic (negative: full more accurate)
\providecommand{\rhtaPEnetAllThirtyNineKBar}{}\renewcommand{\rhtaPEnetAllThirtyNineKBar}{\ensuremath{p < 0.001}}  % ... its p statement
\providecommand{\rhtaQHarEnetAllThirtyNineKBar}{}\renewcommand{\rhtaQHarEnetAllThirtyNineKBar}{0.1922}  % ... mean QLIKE, HAR benchmark
\providecommand{\rhtaQFullEnetAllThirtyNineKBar}{}\renewcommand{\rhtaQFullEnetAllThirtyNineKBar}{0.1865}  % ... mean QLIKE, full set
\providecommand{\rhtaPctLgbmAllThirtyNineKBar}{}\renewcommand{\rhtaPctLgbmAllThirtyNineKBar}{\ensuremath{-}7.2}  % QLIKE % change, full set vs HAR benchmark: LightGBM, ALL48_N39k (ALL48, the last 39,000 rows), every bar
\providecommand{\rhtaDmLgbmAllThirtyNineKBar}{}\renewcommand{\rhtaDmLgbmAllThirtyNineKBar}{\ensuremath{-}10.73}  % ... its Diebold-Mariano statistic (negative: full more accurate)
\providecommand{\rhtaPLgbmAllThirtyNineKBar}{}\renewcommand{\rhtaPLgbmAllThirtyNineKBar}{\ensuremath{p < 0.001}}  % ... its p statement
\providecommand{\rhtaQHarLgbmAllThirtyNineKBar}{}\renewcommand{\rhtaQHarLgbmAllThirtyNineKBar}{0.1920}  % ... mean QLIKE, HAR benchmark
\providecommand{\rhtaQFullLgbmAllThirtyNineKBar}{}\renewcommand{\rhtaQFullLgbmAllThirtyNineKBar}{0.1781}  % ... mean QLIKE, full set
\providecommand{\rhtaPctXgbAllThirtyNineKBar}{}\renewcommand{\rhtaPctXgbAllThirtyNineKBar}{\ensuremath{-}10.1}  % QLIKE % change, full set vs HAR benchmark: XGBoost, ALL48_N39k (ALL48, the last 39,000 rows), every bar
\providecommand{\rhtaDmXgbAllThirtyNineKBar}{}\renewcommand{\rhtaDmXgbAllThirtyNineKBar}{\ensuremath{-}10.02}  % ... its Diebold-Mariano statistic (negative: full more accurate)
\providecommand{\rhtaPXgbAllThirtyNineKBar}{}\renewcommand{\rhtaPXgbAllThirtyNineKBar}{\ensuremath{p < 0.001}}  % ... its p statement
\providecommand{\rhtaQHarXgbAllThirtyNineKBar}{}\renewcommand{\rhtaQHarXgbAllThirtyNineKBar}{0.1927}  % ... mean QLIKE, HAR benchmark
\providecommand{\rhtaQFullXgbAllThirtyNineKBar}{}\renewcommand{\rhtaQFullXgbAllThirtyNineKBar}{0.1732}  % ... mean QLIKE, full set
\providecommand{\rhtaPctRfAllThirtyNineKBar}{}\renewcommand{\rhtaPctRfAllThirtyNineKBar}{\ensuremath{-}5.6}  % QLIKE % change, full set vs HAR benchmark: Random forest, ALL48_N39k (ALL48, the last 39,000 rows), every bar
\providecommand{\rhtaDmRfAllThirtyNineKBar}{}\renewcommand{\rhtaDmRfAllThirtyNineKBar}{\ensuremath{-}12.63}  % ... its Diebold-Mariano statistic (negative: full more accurate)
\providecommand{\rhtaPRfAllThirtyNineKBar}{}\renewcommand{\rhtaPRfAllThirtyNineKBar}{\ensuremath{p < 0.001}}  % ... its p statement
\providecommand{\rhtaQHarRfAllThirtyNineKBar}{}\renewcommand{\rhtaQHarRfAllThirtyNineKBar}{0.1926}  % ... mean QLIKE, HAR benchmark
\providecommand{\rhtaQFullRfAllThirtyNineKBar}{}\renewcommand{\rhtaQFullRfAllThirtyNineKBar}{0.1818}  % ... mean QLIKE, full set
\providecommand{\rhtaXPctRidgeAllVsRthClose}{}\renewcommand{\rhtaXPctRidgeAllVsRthClose}{\ensuremath{-}3.5}  % QLIKE % change, ALL48 vs RTH13 (full set; negative: ALL48 more accurate), Ridge, 16:00 bar
\providecommand{\rhtaXDmRidgeAllVsRthClose}{}\renewcommand{\rhtaXDmRidgeAllVsRthClose}{\ensuremath{-}2.72}  % ... its Diebold-Mariano statistic
\providecommand{\rhtaXPRidgeAllVsRthClose}{}\renewcommand{\rhtaXPRidgeAllVsRthClose}{\ensuremath{p = 0.006}}  % ... its p statement
\providecommand{\rhtaXPctLassoAllVsRthClose}{}\renewcommand{\rhtaXPctLassoAllVsRthClose}{\ensuremath{-}6.0}  % QLIKE % change, ALL48 vs RTH13 (full set; negative: ALL48 more accurate), Lasso, 16:00 bar
\providecommand{\rhtaXDmLassoAllVsRthClose}{}\renewcommand{\rhtaXDmLassoAllVsRthClose}{\ensuremath{-}3.31}  % ... its Diebold-Mariano statistic
\providecommand{\rhtaXPLassoAllVsRthClose}{}\renewcommand{\rhtaXPLassoAllVsRthClose}{\ensuremath{p < 0.001}}  % ... its p statement
\providecommand{\rhtaXPctEnetAllVsRthClose}{}\renewcommand{\rhtaXPctEnetAllVsRthClose}{\ensuremath{-}2.1}  % QLIKE % change, ALL48 vs RTH13 (full set; negative: ALL48 more accurate), Elastic net, 16:00 bar
\providecommand{\rhtaXDmEnetAllVsRthClose}{}\renewcommand{\rhtaXDmEnetAllVsRthClose}{\ensuremath{-}1.64}  % ... its Diebold-Mariano statistic
\providecommand{\rhtaXPEnetAllVsRthClose}{}\renewcommand{\rhtaXPEnetAllVsRthClose}{\ensuremath{p = 0.1006}}  % ... its p statement
\providecommand{\rhtaXPctLgbmAllVsRthClose}{}\renewcommand{\rhtaXPctLgbmAllVsRthClose}{\ensuremath{-}2.3}  % QLIKE % change, ALL48 vs RTH13 (full set; negative: ALL48 more accurate), LightGBM, 16:00 bar
\providecommand{\rhtaXDmLgbmAllVsRthClose}{}\renewcommand{\rhtaXDmLgbmAllVsRthClose}{\ensuremath{-}2.16}  % ... its Diebold-Mariano statistic
\providecommand{\rhtaXPLgbmAllVsRthClose}{}\renewcommand{\rhtaXPLgbmAllVsRthClose}{\ensuremath{p = 0.030}}  % ... its p statement
\providecommand{\rhtaXPctXgbAllVsRthClose}{}\renewcommand{\rhtaXPctXgbAllVsRthClose}{\ensuremath{-}1.9}  % QLIKE % change, ALL48 vs RTH13 (full set; negative: ALL48 more accurate), XGBoost, 16:00 bar
\providecommand{\rhtaXDmXgbAllVsRthClose}{}\renewcommand{\rhtaXDmXgbAllVsRthClose}{\ensuremath{-}2.21}  % ... its Diebold-Mariano statistic
\providecommand{\rhtaXPXgbAllVsRthClose}{}\renewcommand{\rhtaXPXgbAllVsRthClose}{\ensuremath{p = 0.027}}  % ... its p statement
\providecommand{\rhtaXPctRfAllVsRthClose}{}\renewcommand{\rhtaXPctRfAllVsRthClose}{\ensuremath{-}4.3}  % QLIKE % change, ALL48 vs RTH13 (full set; negative: ALL48 more accurate), Random forest, 16:00 bar
\providecommand{\rhtaXDmRfAllVsRthClose}{}\renewcommand{\rhtaXDmRfAllVsRthClose}{\ensuremath{-}4.26}  % ... its Diebold-Mariano statistic
\providecommand{\rhtaXPRfAllVsRthClose}{}\renewcommand{\rhtaXPRfAllVsRthClose}{\ensuremath{p < 0.001}}  % ... its p statement
\providecommand{\rhtaXPctRidgeAllVsRthBar}{}\renewcommand{\rhtaXPctRidgeAllVsRthBar}{1.9}  % QLIKE % change, ALL48 vs RTH13 (full set; negative: ALL48 more accurate), Ridge, every bar
\providecommand{\rhtaXDmRidgeAllVsRthBar}{}\renewcommand{\rhtaXDmRidgeAllVsRthBar}{3.53}  % ... its Diebold-Mariano statistic
\providecommand{\rhtaXPRidgeAllVsRthBar}{}\renewcommand{\rhtaXPRidgeAllVsRthBar}{\ensuremath{p < 0.001}}  % ... its p statement
\providecommand{\rhtaXPctLassoAllVsRthBar}{}\renewcommand{\rhtaXPctLassoAllVsRthBar}{2.5}  % QLIKE % change, ALL48 vs RTH13 (full set; negative: ALL48 more accurate), Lasso, every bar
\providecommand{\rhtaXDmLassoAllVsRthBar}{}\renewcommand{\rhtaXDmLassoAllVsRthBar}{2.64}  % ... its Diebold-Mariano statistic
\providecommand{\rhtaXPLassoAllVsRthBar}{}\renewcommand{\rhtaXPLassoAllVsRthBar}{\ensuremath{p = 0.008}}  % ... its p statement
\providecommand{\rhtaXPctEnetAllVsRthBar}{}\renewcommand{\rhtaXPctEnetAllVsRthBar}{2.2}  % QLIKE % change, ALL48 vs RTH13 (full set; negative: ALL48 more accurate), Elastic net, every bar
\providecommand{\rhtaXDmEnetAllVsRthBar}{}\renewcommand{\rhtaXDmEnetAllVsRthBar}{4.68}  % ... its Diebold-Mariano statistic
\providecommand{\rhtaXPEnetAllVsRthBar}{}\renewcommand{\rhtaXPEnetAllVsRthBar}{\ensuremath{p < 0.001}}  % ... its p statement
\providecommand{\rhtaXPctLgbmAllVsRthBar}{}\renewcommand{\rhtaXPctLgbmAllVsRthBar}{4.5}  % QLIKE % change, ALL48 vs RTH13 (full set; negative: ALL48 more accurate), LightGBM, every bar
\providecommand{\rhtaXDmLgbmAllVsRthBar}{}\renewcommand{\rhtaXDmLgbmAllVsRthBar}{7.39}  % ... its Diebold-Mariano statistic
\providecommand{\rhtaXPLgbmAllVsRthBar}{}\renewcommand{\rhtaXPLgbmAllVsRthBar}{\ensuremath{p < 0.001}}  % ... its p statement
\providecommand{\rhtaXPctXgbAllVsRthBar}{}\renewcommand{\rhtaXPctXgbAllVsRthBar}{3.2}  % QLIKE % change, ALL48 vs RTH13 (full set; negative: ALL48 more accurate), XGBoost, every bar
\providecommand{\rhtaXDmXgbAllVsRthBar}{}\renewcommand{\rhtaXDmXgbAllVsRthBar}{5.64}  % ... its Diebold-Mariano statistic
\providecommand{\rhtaXPXgbAllVsRthBar}{}\renewcommand{\rhtaXPXgbAllVsRthBar}{\ensuremath{p < 0.001}}  % ... its p statement
\providecommand{\rhtaXPctRfAllVsRthBar}{}\renewcommand{\rhtaXPctRfAllVsRthBar}{5.2}  % QLIKE % change, ALL48 vs RTH13 (full set; negative: ALL48 more accurate), Random forest, every bar
\providecommand{\rhtaXDmRfAllVsRthBar}{}\renewcommand{\rhtaXDmRfAllVsRthBar}{6.61}  % ... its Diebold-Mariano statistic
\providecommand{\rhtaXPRfAllVsRthBar}{}\renewcommand{\rhtaXPRfAllVsRthBar}{\ensuremath{p < 0.001}}  % ... its p statement
\providecommand{\rhtaXPctRidgeClVsRthClose}{}\renewcommand{\rhtaXPctRidgeClVsRthClose}{\ensuremath{-}10.2}  % QLIKE % change, CLOSE vs RTH13 (full set; negative: CLOSE more accurate), Ridge, 16:00 bar
\providecommand{\rhtaXDmRidgeClVsRthClose}{}\renewcommand{\rhtaXDmRidgeClVsRthClose}{\ensuremath{-}3.98}  % ... its Diebold-Mariano statistic
\providecommand{\rhtaXPRidgeClVsRthClose}{}\renewcommand{\rhtaXPRidgeClVsRthClose}{\ensuremath{p < 0.001}}  % ... its p statement
\providecommand{\rhtaXPctLassoClVsRthClose}{}\renewcommand{\rhtaXPctLassoClVsRthClose}{\ensuremath{-}11.8}  % QLIKE % change, CLOSE vs RTH13 (full set; negative: CLOSE more accurate), Lasso, 16:00 bar
\providecommand{\rhtaXDmLassoClVsRthClose}{}\renewcommand{\rhtaXDmLassoClVsRthClose}{\ensuremath{-}3.79}  % ... its Diebold-Mariano statistic
\providecommand{\rhtaXPLassoClVsRthClose}{}\renewcommand{\rhtaXPLassoClVsRthClose}{\ensuremath{p < 0.001}}  % ... its p statement
\providecommand{\rhtaXPctEnetClVsRthClose}{}\renewcommand{\rhtaXPctEnetClVsRthClose}{\ensuremath{-}9.4}  % QLIKE % change, CLOSE vs RTH13 (full set; negative: CLOSE more accurate), Elastic net, 16:00 bar
\providecommand{\rhtaXDmEnetClVsRthClose}{}\renewcommand{\rhtaXDmEnetClVsRthClose}{\ensuremath{-}3.86}  % ... its Diebold-Mariano statistic
\providecommand{\rhtaXPEnetClVsRthClose}{}\renewcommand{\rhtaXPEnetClVsRthClose}{\ensuremath{p < 0.001}}  % ... its p statement
\providecommand{\rhtaXPctLgbmClVsRthClose}{}\renewcommand{\rhtaXPctLgbmClVsRthClose}{1.6}  % QLIKE % change, CLOSE vs RTH13 (full set; negative: CLOSE more accurate), LightGBM, 16:00 bar
\providecommand{\rhtaXDmLgbmClVsRthClose}{}\renewcommand{\rhtaXDmLgbmClVsRthClose}{0.45}  % ... its Diebold-Mariano statistic
\providecommand{\rhtaXPLgbmClVsRthClose}{}\renewcommand{\rhtaXPLgbmClVsRthClose}{\ensuremath{p = 0.65}}  % ... its p statement
\providecommand{\rhtaXPctXgbClVsRthClose}{}\renewcommand{\rhtaXPctXgbClVsRthClose}{\ensuremath{-}1.1}  % QLIKE % change, CLOSE vs RTH13 (full set; negative: CLOSE more accurate), XGBoost, 16:00 bar
\providecommand{\rhtaXDmXgbClVsRthClose}{}\renewcommand{\rhtaXDmXgbClVsRthClose}{\ensuremath{-}0.40}  % ... its Diebold-Mariano statistic
\providecommand{\rhtaXPXgbClVsRthClose}{}\renewcommand{\rhtaXPXgbClVsRthClose}{\ensuremath{p = 0.69}}  % ... its p statement
\providecommand{\rhtaXPctRfClVsRthClose}{}\renewcommand{\rhtaXPctRfClVsRthClose}{\ensuremath{-}2.2}  % QLIKE % change, CLOSE vs RTH13 (full set; negative: CLOSE more accurate), Random forest, 16:00 bar
\providecommand{\rhtaXDmRfClVsRthClose}{}\renewcommand{\rhtaXDmRfClVsRthClose}{\ensuremath{-}0.81}  % ... its Diebold-Mariano statistic
\providecommand{\rhtaXPRfClVsRthClose}{}\renewcommand{\rhtaXPRfClVsRthClose}{\ensuremath{p = 0.42}}  % ... its p statement
\providecommand{\rhtaXPctRidgeClVsAllClose}{}\renewcommand{\rhtaXPctRidgeClVsAllClose}{\ensuremath{-}7.0}  % QLIKE % change, CLOSE vs ALL48 (full set; negative: CLOSE more accurate), Ridge, 16:00 bar
\providecommand{\rhtaXDmRidgeClVsAllClose}{}\renewcommand{\rhtaXDmRidgeClVsAllClose}{\ensuremath{-}3.31}  % ... its Diebold-Mariano statistic
\providecommand{\rhtaXPRidgeClVsAllClose}{}\renewcommand{\rhtaXPRidgeClVsAllClose}{\ensuremath{p < 0.001}}  % ... its p statement
\providecommand{\rhtaXPctLassoClVsAllClose}{}\renewcommand{\rhtaXPctLassoClVsAllClose}{\ensuremath{-}6.2}  % QLIKE % change, CLOSE vs ALL48 (full set; negative: CLOSE more accurate), Lasso, 16:00 bar
\providecommand{\rhtaXDmLassoClVsAllClose}{}\renewcommand{\rhtaXDmLassoClVsAllClose}{\ensuremath{-}3.16}  % ... its Diebold-Mariano statistic
\providecommand{\rhtaXPLassoClVsAllClose}{}\renewcommand{\rhtaXPLassoClVsAllClose}{\ensuremath{p = 0.002}}  % ... its p statement
\providecommand{\rhtaXPctEnetClVsAllClose}{}\renewcommand{\rhtaXPctEnetClVsAllClose}{\ensuremath{-}7.5}  % QLIKE % change, CLOSE vs ALL48 (full set; negative: CLOSE more accurate), Elastic net, 16:00 bar
\providecommand{\rhtaXDmEnetClVsAllClose}{}\renewcommand{\rhtaXDmEnetClVsAllClose}{\ensuremath{-}3.70}  % ... its Diebold-Mariano statistic
\providecommand{\rhtaXPEnetClVsAllClose}{}\renewcommand{\rhtaXPEnetClVsAllClose}{\ensuremath{p < 0.001}}  % ... its p statement
\providecommand{\rhtaXPctLgbmClVsAllClose}{}\renewcommand{\rhtaXPctLgbmClVsAllClose}{4.0}  % QLIKE % change, CLOSE vs ALL48 (full set; negative: CLOSE more accurate), LightGBM, 16:00 bar
\providecommand{\rhtaXDmLgbmClVsAllClose}{}\renewcommand{\rhtaXDmLgbmClVsAllClose}{0.94}  % ... its Diebold-Mariano statistic
\providecommand{\rhtaXPLgbmClVsAllClose}{}\renewcommand{\rhtaXPLgbmClVsAllClose}{\ensuremath{p = 0.35}}  % ... its p statement
\providecommand{\rhtaXPctXgbClVsAllClose}{}\renewcommand{\rhtaXPctXgbClVsAllClose}{0.8}  % QLIKE % change, CLOSE vs ALL48 (full set; negative: CLOSE more accurate), XGBoost, 16:00 bar
\providecommand{\rhtaXDmXgbClVsAllClose}{}\renewcommand{\rhtaXDmXgbClVsAllClose}{0.28}  % ... its Diebold-Mariano statistic
\providecommand{\rhtaXPXgbClVsAllClose}{}\renewcommand{\rhtaXPXgbClVsAllClose}{\ensuremath{p = 0.78}}  % ... its p statement
\providecommand{\rhtaXPctRfClVsAllClose}{}\renewcommand{\rhtaXPctRfClVsAllClose}{2.2}  % QLIKE % change, CLOSE vs ALL48 (full set; negative: CLOSE more accurate), Random forest, 16:00 bar
\providecommand{\rhtaXDmRfClVsAllClose}{}\renewcommand{\rhtaXDmRfClVsAllClose}{0.74}  % ... its Diebold-Mariano statistic
\providecommand{\rhtaXPRfClVsAllClose}{}\renewcommand{\rhtaXPRfClVsAllClose}{\ensuremath{p = 0.46}}  % ... its p statement
\providecommand{\rhtaXPctRidgeClNatVsRthClose}{}\renewcommand{\rhtaXPctRidgeClNatVsRthClose}{\ensuremath{-}10.4}  % QLIKE % change, CLOSE_NAT vs RTH13 (full set; negative: CLOSE_NAT more accurate), Ridge, 16:00 bar
\providecommand{\rhtaXDmRidgeClNatVsRthClose}{}\renewcommand{\rhtaXDmRidgeClNatVsRthClose}{\ensuremath{-}3.91}  % ... its Diebold-Mariano statistic
\providecommand{\rhtaXPRidgeClNatVsRthClose}{}\renewcommand{\rhtaXPRidgeClNatVsRthClose}{\ensuremath{p < 0.001}}  % ... its p statement
\providecommand{\rhtaXPctLassoClNatVsRthClose}{}\renewcommand{\rhtaXPctLassoClNatVsRthClose}{\ensuremath{-}11.5}  % QLIKE % change, CLOSE_NAT vs RTH13 (full set; negative: CLOSE_NAT more accurate), Lasso, 16:00 bar
\providecommand{\rhtaXDmLassoClNatVsRthClose}{}\renewcommand{\rhtaXDmLassoClNatVsRthClose}{\ensuremath{-}3.57}  % ... its Diebold-Mariano statistic
\providecommand{\rhtaXPLassoClNatVsRthClose}{}\renewcommand{\rhtaXPLassoClNatVsRthClose}{\ensuremath{p < 0.001}}  % ... its p statement
\providecommand{\rhtaXPctEnetClNatVsRthClose}{}\renewcommand{\rhtaXPctEnetClNatVsRthClose}{\ensuremath{-}9.3}  % QLIKE % change, CLOSE_NAT vs RTH13 (full set; negative: CLOSE_NAT more accurate), Elastic net, 16:00 bar
\providecommand{\rhtaXDmEnetClNatVsRthClose}{}\renewcommand{\rhtaXDmEnetClNatVsRthClose}{\ensuremath{-}3.58}  % ... its Diebold-Mariano statistic
\providecommand{\rhtaXPEnetClNatVsRthClose}{}\renewcommand{\rhtaXPEnetClNatVsRthClose}{\ensuremath{p < 0.001}}  % ... its p statement
\providecommand{\rhtaXPctRidgeAllCtVsRthClose}{}\renewcommand{\rhtaXPctRidgeAllCtVsRthClose}{\ensuremath{-}3.2}  % QLIKE % change, ALL48@16 vs RTH13 (full set; negative: ALL48@16 more accurate), Ridge, 16:00 bar
\providecommand{\rhtaXDmRidgeAllCtVsRthClose}{}\renewcommand{\rhtaXDmRidgeAllCtVsRthClose}{\ensuremath{-}2.35}  % ... its Diebold-Mariano statistic
\providecommand{\rhtaXPRidgeAllCtVsRthClose}{}\renewcommand{\rhtaXPRidgeAllCtVsRthClose}{\ensuremath{p = 0.019}}  % ... its p statement
\providecommand{\rhtaXPctLassoAllCtVsRthClose}{}\renewcommand{\rhtaXPctLassoAllCtVsRthClose}{\ensuremath{-}5.4}  % QLIKE % change, ALL48@16 vs RTH13 (full set; negative: ALL48@16 more accurate), Lasso, 16:00 bar
\providecommand{\rhtaXDmLassoAllCtVsRthClose}{}\renewcommand{\rhtaXDmLassoAllCtVsRthClose}{\ensuremath{-}3.04}  % ... its Diebold-Mariano statistic
\providecommand{\rhtaXPLassoAllCtVsRthClose}{}\renewcommand{\rhtaXPLassoAllCtVsRthClose}{\ensuremath{p = 0.002}}  % ... its p statement
\providecommand{\rhtaXPctEnetAllCtVsRthClose}{}\renewcommand{\rhtaXPctEnetAllCtVsRthClose}{\ensuremath{-}2.1}  % QLIKE % change, ALL48@16 vs RTH13 (full set; negative: ALL48@16 more accurate), Elastic net, 16:00 bar
\providecommand{\rhtaXDmEnetAllCtVsRthClose}{}\renewcommand{\rhtaXDmEnetAllCtVsRthClose}{\ensuremath{-}1.64}  % ... its Diebold-Mariano statistic
\providecommand{\rhtaXPEnetAllCtVsRthClose}{}\renewcommand{\rhtaXPEnetAllCtVsRthClose}{\ensuremath{p = 0.1006}}  % ... its p statement
\providecommand{\rhtaXPctRidgeRthCtVsRthClose}{}\renewcommand{\rhtaXPctRidgeRthCtVsRthClose}{0.0}  % QLIKE % change, RTH13@16 vs RTH13 (full set; negative: RTH13@16 more accurate), Ridge, 16:00 bar
\providecommand{\rhtaXDmRidgeRthCtVsRthClose}{}\renewcommand{\rhtaXDmRidgeRthCtVsRthClose}{0.00}  % ... its Diebold-Mariano statistic
\providecommand{\rhtaXPRidgeRthCtVsRthClose}{}\renewcommand{\rhtaXPRidgeRthCtVsRthClose}{\ensuremath{p = 1.00}}  % ... its p statement
\providecommand{\rhtaXPctLassoRthCtVsRthClose}{}\renewcommand{\rhtaXPctLassoRthCtVsRthClose}{\ensuremath{-}5.4}  % QLIKE % change, RTH13@16 vs RTH13 (full set; negative: RTH13@16 more accurate), Lasso, 16:00 bar
\providecommand{\rhtaXDmLassoRthCtVsRthClose}{}\renewcommand{\rhtaXDmLassoRthCtVsRthClose}{\ensuremath{-}3.40}  % ... its Diebold-Mariano statistic
\providecommand{\rhtaXPLassoRthCtVsRthClose}{}\renewcommand{\rhtaXPLassoRthCtVsRthClose}{\ensuremath{p < 0.001}}  % ... its p statement
\providecommand{\rhtaXPctEnetRthCtVsRthClose}{}\renewcommand{\rhtaXPctEnetRthCtVsRthClose}{0.0}  % QLIKE % change, RTH13@16 vs RTH13 (full set; negative: RTH13@16 more accurate), Elastic net, 16:00 bar
\providecommand{\rhtaXDmEnetRthCtVsRthClose}{}\renewcommand{\rhtaXDmEnetRthCtVsRthClose}{0.00}  % ... its Diebold-Mariano statistic
\providecommand{\rhtaXPEnetRthCtVsRthClose}{}\renewcommand{\rhtaXPEnetRthCtVsRthClose}{\ensuremath{p = 1.00}}  % ... its p statement
\providecommand{\rhtaXPctRidgeAllCtVsRthBar}{}\renewcommand{\rhtaXPctRidgeAllCtVsRthBar}{1.6}  % QLIKE % change, ALL48@16 vs RTH13 (full set; negative: ALL48@16 more accurate), Ridge, every bar
\providecommand{\rhtaXDmRidgeAllCtVsRthBar}{}\renewcommand{\rhtaXDmRidgeAllCtVsRthBar}{3.24}  % ... its Diebold-Mariano statistic
\providecommand{\rhtaXPRidgeAllCtVsRthBar}{}\renewcommand{\rhtaXPRidgeAllCtVsRthBar}{\ensuremath{p = 0.001}}  % ... its p statement
\providecommand{\rhtaXPctLassoAllCtVsRthBar}{}\renewcommand{\rhtaXPctLassoAllCtVsRthBar}{2.2}  % QLIKE % change, ALL48@16 vs RTH13 (full set; negative: ALL48@16 more accurate), Lasso, every bar
\providecommand{\rhtaXDmLassoAllCtVsRthBar}{}\renewcommand{\rhtaXDmLassoAllCtVsRthBar}{2.65}  % ... its Diebold-Mariano statistic
\providecommand{\rhtaXPLassoAllCtVsRthBar}{}\renewcommand{\rhtaXPLassoAllCtVsRthBar}{\ensuremath{p = 0.008}}  % ... its p statement
\providecommand{\rhtaXPctEnetAllCtVsRthBar}{}\renewcommand{\rhtaXPctEnetAllCtVsRthBar}{2.2}  % QLIKE % change, ALL48@16 vs RTH13 (full set; negative: ALL48@16 more accurate), Elastic net, every bar
\providecommand{\rhtaXDmEnetAllCtVsRthBar}{}\renewcommand{\rhtaXDmEnetAllCtVsRthBar}{4.68}  % ... its Diebold-Mariano statistic
\providecommand{\rhtaXPEnetAllCtVsRthBar}{}\renewcommand{\rhtaXPEnetAllCtVsRthBar}{\ensuremath{p < 0.001}}  % ... its p statement
\providecommand{\rhtaXPctRidgeRthCtVsRthBar}{}\renewcommand{\rhtaXPctRidgeRthCtVsRthBar}{0.0}  % QLIKE % change, RTH13@16 vs RTH13 (full set; negative: RTH13@16 more accurate), Ridge, every bar
\providecommand{\rhtaXDmRidgeRthCtVsRthBar}{}\renewcommand{\rhtaXDmRidgeRthCtVsRthBar}{0.00}  % ... its Diebold-Mariano statistic
\providecommand{\rhtaXPRidgeRthCtVsRthBar}{}\renewcommand{\rhtaXPRidgeRthCtVsRthBar}{\ensuremath{p = 1.00}}  % ... its p statement
\providecommand{\rhtaXPctLassoRthCtVsRthBar}{}\renewcommand{\rhtaXPctLassoRthCtVsRthBar}{0.2}  % QLIKE % change, RTH13@16 vs RTH13 (full set; negative: RTH13@16 more accurate), Lasso, every bar
\providecommand{\rhtaXDmLassoRthCtVsRthBar}{}\renewcommand{\rhtaXDmLassoRthCtVsRthBar}{0.19}  % ... its Diebold-Mariano statistic
\providecommand{\rhtaXPLassoRthCtVsRthBar}{}\renewcommand{\rhtaXPLassoRthCtVsRthBar}{\ensuremath{p = 0.85}}  % ... its p statement
\providecommand{\rhtaXPctEnetRthCtVsRthBar}{}\renewcommand{\rhtaXPctEnetRthCtVsRthBar}{0.0}  % QLIKE % change, RTH13@16 vs RTH13 (full set; negative: RTH13@16 more accurate), Elastic net, every bar
\providecommand{\rhtaXDmEnetRthCtVsRthBar}{}\renewcommand{\rhtaXDmEnetRthCtVsRthBar}{0.00}  % ... its Diebold-Mariano statistic
\providecommand{\rhtaXPEnetRthCtVsRthBar}{}\renewcommand{\rhtaXPEnetRthCtVsRthBar}{\ensuremath{p = 1.00}}  % ... its p statement
\providecommand{\rhtaXPctRidgeAllThreeKVsRthThreeKClose}{}\renewcommand{\rhtaXPctRidgeAllThreeKVsRthThreeKClose}{\ensuremath{-}8.1}  % QLIKE % change, ALL48_N3k vs RTH13_N3k (full set; negative: ALL48_N3k more accurate), Ridge, 16:00 bar
\providecommand{\rhtaXDmRidgeAllThreeKVsRthThreeKClose}{}\renewcommand{\rhtaXDmRidgeAllThreeKVsRthThreeKClose}{\ensuremath{-}5.43}  % ... its Diebold-Mariano statistic
\providecommand{\rhtaXPRidgeAllThreeKVsRthThreeKClose}{}\renewcommand{\rhtaXPRidgeAllThreeKVsRthThreeKClose}{\ensuremath{p < 0.001}}  % ... its p statement
\providecommand{\rhtaXPctLassoAllThreeKVsRthThreeKClose}{}\renewcommand{\rhtaXPctLassoAllThreeKVsRthThreeKClose}{\ensuremath{-}6.5}  % QLIKE % change, ALL48_N3k vs RTH13_N3k (full set; negative: ALL48_N3k more accurate), Lasso, 16:00 bar
\providecommand{\rhtaXDmLassoAllThreeKVsRthThreeKClose}{}\renewcommand{\rhtaXDmLassoAllThreeKVsRthThreeKClose}{\ensuremath{-}3.99}  % ... its Diebold-Mariano statistic
\providecommand{\rhtaXPLassoAllThreeKVsRthThreeKClose}{}\renewcommand{\rhtaXPLassoAllThreeKVsRthThreeKClose}{\ensuremath{p < 0.001}}  % ... its p statement
\providecommand{\rhtaXPctEnetAllThreeKVsRthThreeKClose}{}\renewcommand{\rhtaXPctEnetAllThreeKVsRthThreeKClose}{\ensuremath{-}4.2}  % QLIKE % change, ALL48_N3k vs RTH13_N3k (full set; negative: ALL48_N3k more accurate), Elastic net, 16:00 bar
\providecommand{\rhtaXDmEnetAllThreeKVsRthThreeKClose}{}\renewcommand{\rhtaXDmEnetAllThreeKVsRthThreeKClose}{\ensuremath{-}2.13}  % ... its Diebold-Mariano statistic
\providecommand{\rhtaXPEnetAllThreeKVsRthThreeKClose}{}\renewcommand{\rhtaXPEnetAllThreeKVsRthThreeKClose}{\ensuremath{p = 0.033}}  % ... its p statement
\providecommand{\rhtaXPctLgbmAllThreeKVsRthThreeKClose}{}\renewcommand{\rhtaXPctLgbmAllThreeKVsRthThreeKClose}{\ensuremath{-}3.9}  % QLIKE % change, ALL48_N3k vs RTH13_N3k (full set; negative: ALL48_N3k more accurate), LightGBM, 16:00 bar
\providecommand{\rhtaXDmLgbmAllThreeKVsRthThreeKClose}{}\renewcommand{\rhtaXDmLgbmAllThreeKVsRthThreeKClose}{\ensuremath{-}2.07}  % ... its Diebold-Mariano statistic
\providecommand{\rhtaXPLgbmAllThreeKVsRthThreeKClose}{}\renewcommand{\rhtaXPLgbmAllThreeKVsRthThreeKClose}{\ensuremath{p = 0.038}}  % ... its p statement
\providecommand{\rhtaXPctXgbAllThreeKVsRthThreeKClose}{}\renewcommand{\rhtaXPctXgbAllThreeKVsRthThreeKClose}{\ensuremath{-}3.3}  % QLIKE % change, ALL48_N3k vs RTH13_N3k (full set; negative: ALL48_N3k more accurate), XGBoost, 16:00 bar
\providecommand{\rhtaXDmXgbAllThreeKVsRthThreeKClose}{}\renewcommand{\rhtaXDmXgbAllThreeKVsRthThreeKClose}{\ensuremath{-}1.80}  % ... its Diebold-Mariano statistic
\providecommand{\rhtaXPXgbAllThreeKVsRthThreeKClose}{}\renewcommand{\rhtaXPXgbAllThreeKVsRthThreeKClose}{\ensuremath{p = 0.072}}  % ... its p statement
\providecommand{\rhtaXPctRfAllThreeKVsRthThreeKClose}{}\renewcommand{\rhtaXPctRfAllThreeKVsRthThreeKClose}{\ensuremath{-}5.7}  % QLIKE % change, ALL48_N3k vs RTH13_N3k (full set; negative: ALL48_N3k more accurate), Random forest, 16:00 bar
\providecommand{\rhtaXDmRfAllThreeKVsRthThreeKClose}{}\renewcommand{\rhtaXDmRfAllThreeKVsRthThreeKClose}{\ensuremath{-}3.60}  % ... its Diebold-Mariano statistic
\providecommand{\rhtaXPRfAllThreeKVsRthThreeKClose}{}\renewcommand{\rhtaXPRfAllThreeKVsRthThreeKClose}{\ensuremath{p < 0.001}}  % ... its p statement
\providecommand{\rhtaXPctRidgeAllThreeKVsRthThreeKBar}{}\renewcommand{\rhtaXPctRidgeAllThreeKVsRthThreeKBar}{4.5}  % QLIKE % change, ALL48_N3k vs RTH13_N3k (full set; negative: ALL48_N3k more accurate), Ridge, every bar
\providecommand{\rhtaXDmRidgeAllThreeKVsRthThreeKBar}{}\renewcommand{\rhtaXDmRidgeAllThreeKVsRthThreeKBar}{7.72}  % ... its Diebold-Mariano statistic
\providecommand{\rhtaXPRidgeAllThreeKVsRthThreeKBar}{}\renewcommand{\rhtaXPRidgeAllThreeKVsRthThreeKBar}{\ensuremath{p < 0.001}}  % ... its p statement
\providecommand{\rhtaXPctLassoAllThreeKVsRthThreeKBar}{}\renewcommand{\rhtaXPctLassoAllThreeKVsRthThreeKBar}{4.9}  % QLIKE % change, ALL48_N3k vs RTH13_N3k (full set; negative: ALL48_N3k more accurate), Lasso, every bar
\providecommand{\rhtaXDmLassoAllThreeKVsRthThreeKBar}{}\renewcommand{\rhtaXDmLassoAllThreeKVsRthThreeKBar}{7.96}  % ... its Diebold-Mariano statistic
\providecommand{\rhtaXPLassoAllThreeKVsRthThreeKBar}{}\renewcommand{\rhtaXPLassoAllThreeKVsRthThreeKBar}{\ensuremath{p < 0.001}}  % ... its p statement
\providecommand{\rhtaXPctEnetAllThreeKVsRthThreeKBar}{}\renewcommand{\rhtaXPctEnetAllThreeKVsRthThreeKBar}{5.8}  % QLIKE % change, ALL48_N3k vs RTH13_N3k (full set; negative: ALL48_N3k more accurate), Elastic net, every bar
\providecommand{\rhtaXDmEnetAllThreeKVsRthThreeKBar}{}\renewcommand{\rhtaXDmEnetAllThreeKVsRthThreeKBar}{7.58}  % ... its Diebold-Mariano statistic
\providecommand{\rhtaXPEnetAllThreeKVsRthThreeKBar}{}\renewcommand{\rhtaXPEnetAllThreeKVsRthThreeKBar}{\ensuremath{p < 0.001}}  % ... its p statement
\providecommand{\rhtaXPctLgbmAllThreeKVsRthThreeKBar}{}\renewcommand{\rhtaXPctLgbmAllThreeKVsRthThreeKBar}{4.7}  % QLIKE % change, ALL48_N3k vs RTH13_N3k (full set; negative: ALL48_N3k more accurate), LightGBM, every bar
\providecommand{\rhtaXDmLgbmAllThreeKVsRthThreeKBar}{}\renewcommand{\rhtaXDmLgbmAllThreeKVsRthThreeKBar}{5.00}  % ... its Diebold-Mariano statistic
\providecommand{\rhtaXPLgbmAllThreeKVsRthThreeKBar}{}\renewcommand{\rhtaXPLgbmAllThreeKVsRthThreeKBar}{\ensuremath{p < 0.001}}  % ... its p statement
\providecommand{\rhtaXPctXgbAllThreeKVsRthThreeKBar}{}\renewcommand{\rhtaXPctXgbAllThreeKVsRthThreeKBar}{5.1}  % QLIKE % change, ALL48_N3k vs RTH13_N3k (full set; negative: ALL48_N3k more accurate), XGBoost, every bar
\providecommand{\rhtaXDmXgbAllThreeKVsRthThreeKBar}{}\renewcommand{\rhtaXDmXgbAllThreeKVsRthThreeKBar}{5.92}  % ... its Diebold-Mariano statistic
\providecommand{\rhtaXPXgbAllThreeKVsRthThreeKBar}{}\renewcommand{\rhtaXPXgbAllThreeKVsRthThreeKBar}{\ensuremath{p < 0.001}}  % ... its p statement
\providecommand{\rhtaXPctRfAllThreeKVsRthThreeKBar}{}\renewcommand{\rhtaXPctRfAllThreeKVsRthThreeKBar}{2.7}  % QLIKE % change, ALL48_N3k vs RTH13_N3k (full set; negative: ALL48_N3k more accurate), Random forest, every bar
\providecommand{\rhtaXDmRfAllThreeKVsRthThreeKBar}{}\renewcommand{\rhtaXDmRfAllThreeKVsRthThreeKBar}{3.94}  % ... its Diebold-Mariano statistic
\providecommand{\rhtaXPRfAllThreeKVsRthThreeKBar}{}\renewcommand{\rhtaXPRfAllThreeKVsRthThreeKBar}{\ensuremath{p < 0.001}}  % ... its p statement
\providecommand{\rhtaXPctRidgeClVsRthThreeKClose}{}\renewcommand{\rhtaXPctRidgeClVsRthThreeKClose}{\ensuremath{-}17.4}  % QLIKE % change, CLOSE vs RTH13_N3k (full set; negative: CLOSE more accurate), Ridge, 16:00 bar
\providecommand{\rhtaXDmRidgeClVsRthThreeKClose}{}\renewcommand{\rhtaXDmRidgeClVsRthThreeKClose}{\ensuremath{-}6.84}  % ... its Diebold-Mariano statistic
\providecommand{\rhtaXPRidgeClVsRthThreeKClose}{}\renewcommand{\rhtaXPRidgeClVsRthThreeKClose}{\ensuremath{p < 0.001}}  % ... its p statement
\providecommand{\rhtaXPctLassoClVsRthThreeKClose}{}\renewcommand{\rhtaXPctLassoClVsRthThreeKClose}{\ensuremath{-}16.9}  % QLIKE % change, CLOSE vs RTH13_N3k (full set; negative: CLOSE more accurate), Lasso, 16:00 bar
\providecommand{\rhtaXDmLassoClVsRthThreeKClose}{}\renewcommand{\rhtaXDmLassoClVsRthThreeKClose}{\ensuremath{-}6.57}  % ... its Diebold-Mariano statistic
\providecommand{\rhtaXPLassoClVsRthThreeKClose}{}\renewcommand{\rhtaXPLassoClVsRthThreeKClose}{\ensuremath{p < 0.001}}  % ... its p statement
\providecommand{\rhtaXPctEnetClVsRthThreeKClose}{}\renewcommand{\rhtaXPctEnetClVsRthThreeKClose}{\ensuremath{-}18.0}  % QLIKE % change, CLOSE vs RTH13_N3k (full set; negative: CLOSE more accurate), Elastic net, 16:00 bar
\providecommand{\rhtaXDmEnetClVsRthThreeKClose}{}\renewcommand{\rhtaXDmEnetClVsRthThreeKClose}{\ensuremath{-}6.66}  % ... its Diebold-Mariano statistic
\providecommand{\rhtaXPEnetClVsRthThreeKClose}{}\renewcommand{\rhtaXPEnetClVsRthThreeKClose}{\ensuremath{p < 0.001}}  % ... its p statement
\providecommand{\rhtaXPctLgbmClVsRthThreeKClose}{}\renewcommand{\rhtaXPctLgbmClVsRthThreeKClose}{\ensuremath{-}9.5}  % QLIKE % change, CLOSE vs RTH13_N3k (full set; negative: CLOSE more accurate), LightGBM, 16:00 bar
\providecommand{\rhtaXDmLgbmClVsRthThreeKClose}{}\renewcommand{\rhtaXDmLgbmClVsRthThreeKClose}{\ensuremath{-}2.88}  % ... its Diebold-Mariano statistic
\providecommand{\rhtaXPLgbmClVsRthThreeKClose}{}\renewcommand{\rhtaXPLgbmClVsRthThreeKClose}{\ensuremath{p = 0.004}}  % ... its p statement
\providecommand{\rhtaXPctXgbClVsRthThreeKClose}{}\renewcommand{\rhtaXPctXgbClVsRthThreeKClose}{\ensuremath{-}9.5}  % QLIKE % change, CLOSE vs RTH13_N3k (full set; negative: CLOSE more accurate), XGBoost, 16:00 bar
\providecommand{\rhtaXDmXgbClVsRthThreeKClose}{}\renewcommand{\rhtaXDmXgbClVsRthThreeKClose}{\ensuremath{-}3.40}  % ... its Diebold-Mariano statistic
\providecommand{\rhtaXPXgbClVsRthThreeKClose}{}\renewcommand{\rhtaXPXgbClVsRthThreeKClose}{\ensuremath{p < 0.001}}  % ... its p statement
\providecommand{\rhtaXPctRfClVsRthThreeKClose}{}\renewcommand{\rhtaXPctRfClVsRthThreeKClose}{\ensuremath{-}8.7}  % QLIKE % change, CLOSE vs RTH13_N3k (full set; negative: CLOSE more accurate), Random forest, 16:00 bar
\providecommand{\rhtaXDmRfClVsRthThreeKClose}{}\renewcommand{\rhtaXDmRfClVsRthThreeKClose}{\ensuremath{-}3.08}  % ... its Diebold-Mariano statistic
\providecommand{\rhtaXPRfClVsRthThreeKClose}{}\renewcommand{\rhtaXPRfClVsRthThreeKClose}{\ensuremath{p = 0.002}}  % ... its p statement
\providecommand{\rhtaXPctRidgeClVsAllThreeKClose}{}\renewcommand{\rhtaXPctRidgeClVsAllThreeKClose}{\ensuremath{-}10.2}  % QLIKE % change, CLOSE vs ALL48_N3k (full set; negative: CLOSE more accurate), Ridge, 16:00 bar
\providecommand{\rhtaXDmRidgeClVsAllThreeKClose}{}\renewcommand{\rhtaXDmRidgeClVsAllThreeKClose}{\ensuremath{-}4.63}  % ... its Diebold-Mariano statistic
\providecommand{\rhtaXPRidgeClVsAllThreeKClose}{}\renewcommand{\rhtaXPRidgeClVsAllThreeKClose}{\ensuremath{p < 0.001}}  % ... its p statement
\providecommand{\rhtaXPctLassoClVsAllThreeKClose}{}\renewcommand{\rhtaXPctLassoClVsAllThreeKClose}{\ensuremath{-}11.1}  % QLIKE % change, CLOSE vs ALL48_N3k (full set; negative: CLOSE more accurate), Lasso, 16:00 bar
\providecommand{\rhtaXDmLassoClVsAllThreeKClose}{}\renewcommand{\rhtaXDmLassoClVsAllThreeKClose}{\ensuremath{-}4.62}  % ... its Diebold-Mariano statistic
\providecommand{\rhtaXPLassoClVsAllThreeKClose}{}\renewcommand{\rhtaXPLassoClVsAllThreeKClose}{\ensuremath{p < 0.001}}  % ... its p statement
\providecommand{\rhtaXPctEnetClVsAllThreeKClose}{}\renewcommand{\rhtaXPctEnetClVsAllThreeKClose}{\ensuremath{-}14.5}  % QLIKE % change, CLOSE vs ALL48_N3k (full set; negative: CLOSE more accurate), Elastic net, 16:00 bar
\providecommand{\rhtaXDmEnetClVsAllThreeKClose}{}\renewcommand{\rhtaXDmEnetClVsAllThreeKClose}{\ensuremath{-}5.41}  % ... its Diebold-Mariano statistic
\providecommand{\rhtaXPEnetClVsAllThreeKClose}{}\renewcommand{\rhtaXPEnetClVsAllThreeKClose}{\ensuremath{p < 0.001}}  % ... its p statement
\providecommand{\rhtaXPctLgbmClVsAllThreeKClose}{}\renewcommand{\rhtaXPctLgbmClVsAllThreeKClose}{\ensuremath{-}5.8}  % QLIKE % change, CLOSE vs ALL48_N3k (full set; negative: CLOSE more accurate), LightGBM, 16:00 bar
\providecommand{\rhtaXDmLgbmClVsAllThreeKClose}{}\renewcommand{\rhtaXDmLgbmClVsAllThreeKClose}{\ensuremath{-}1.68}  % ... its Diebold-Mariano statistic
\providecommand{\rhtaXPLgbmClVsAllThreeKClose}{}\renewcommand{\rhtaXPLgbmClVsAllThreeKClose}{\ensuremath{p = 0.093}}  % ... its p statement
\providecommand{\rhtaXPctXgbClVsAllThreeKClose}{}\renewcommand{\rhtaXPctXgbClVsAllThreeKClose}{\ensuremath{-}6.4}  % QLIKE % change, CLOSE vs ALL48_N3k (full set; negative: CLOSE more accurate), XGBoost, 16:00 bar
\providecommand{\rhtaXDmXgbClVsAllThreeKClose}{}\renewcommand{\rhtaXDmXgbClVsAllThreeKClose}{\ensuremath{-}2.23}  % ... its Diebold-Mariano statistic
\providecommand{\rhtaXPXgbClVsAllThreeKClose}{}\renewcommand{\rhtaXPXgbClVsAllThreeKClose}{\ensuremath{p = 0.026}}  % ... its p statement
\providecommand{\rhtaXPctRfClVsAllThreeKClose}{}\renewcommand{\rhtaXPctRfClVsAllThreeKClose}{\ensuremath{-}3.2}  % QLIKE % change, CLOSE vs ALL48_N3k (full set; negative: CLOSE more accurate), Random forest, 16:00 bar
\providecommand{\rhtaXDmRfClVsAllThreeKClose}{}\renewcommand{\rhtaXDmRfClVsAllThreeKClose}{\ensuremath{-}1.10}  % ... its Diebold-Mariano statistic
\providecommand{\rhtaXPRfClVsAllThreeKClose}{}\renewcommand{\rhtaXPRfClVsAllThreeKClose}{\ensuremath{p = 0.27}}  % ... its p statement
\providecommand{\rhtaXPctRidgeAllThirtyNineKVsRthClose}{}\renewcommand{\rhtaXPctRidgeAllThirtyNineKVsRthClose}{0.5}  % QLIKE % change, ALL48_N39k vs RTH13 (full set; negative: ALL48_N39k more accurate), Ridge, 16:00 bar
\providecommand{\rhtaXDmRidgeAllThirtyNineKVsRthClose}{}\renewcommand{\rhtaXDmRidgeAllThirtyNineKVsRthClose}{0.34}  % ... its Diebold-Mariano statistic
\providecommand{\rhtaXPRidgeAllThirtyNineKVsRthClose}{}\renewcommand{\rhtaXPRidgeAllThirtyNineKVsRthClose}{\ensuremath{p = 0.73}}  % ... its p statement
\providecommand{\rhtaXPctLassoAllThirtyNineKVsRthClose}{}\renewcommand{\rhtaXPctLassoAllThirtyNineKVsRthClose}{\ensuremath{-}3.6}  % QLIKE % change, ALL48_N39k vs RTH13 (full set; negative: ALL48_N39k more accurate), Lasso, 16:00 bar
\providecommand{\rhtaXDmLassoAllThirtyNineKVsRthClose}{}\renewcommand{\rhtaXDmLassoAllThirtyNineKVsRthClose}{\ensuremath{-}2.08}  % ... its Diebold-Mariano statistic
\providecommand{\rhtaXPLassoAllThirtyNineKVsRthClose}{}\renewcommand{\rhtaXPLassoAllThirtyNineKVsRthClose}{\ensuremath{p = 0.038}}  % ... its p statement
\providecommand{\rhtaXPctEnetAllThirtyNineKVsRthClose}{}\renewcommand{\rhtaXPctEnetAllThirtyNineKVsRthClose}{0.1}  % QLIKE % change, ALL48_N39k vs RTH13 (full set; negative: ALL48_N39k more accurate), Elastic net, 16:00 bar
\providecommand{\rhtaXDmEnetAllThirtyNineKVsRthClose}{}\renewcommand{\rhtaXDmEnetAllThirtyNineKVsRthClose}{0.07}  % ... its Diebold-Mariano statistic
\providecommand{\rhtaXPEnetAllThirtyNineKVsRthClose}{}\renewcommand{\rhtaXPEnetAllThirtyNineKVsRthClose}{\ensuremath{p = 0.95}}  % ... its p statement
\providecommand{\rhtaXPctLgbmAllThirtyNineKVsRthClose}{}\renewcommand{\rhtaXPctLgbmAllThirtyNineKVsRthClose}{0.4}  % QLIKE % change, ALL48_N39k vs RTH13 (full set; negative: ALL48_N39k more accurate), LightGBM, 16:00 bar
\providecommand{\rhtaXDmLgbmAllThirtyNineKVsRthClose}{}\renewcommand{\rhtaXDmLgbmAllThirtyNineKVsRthClose}{0.36}  % ... its Diebold-Mariano statistic
\providecommand{\rhtaXPLgbmAllThirtyNineKVsRthClose}{}\renewcommand{\rhtaXPLgbmAllThirtyNineKVsRthClose}{\ensuremath{p = 0.72}}  % ... its p statement
\providecommand{\rhtaXPctXgbAllThirtyNineKVsRthClose}{}\renewcommand{\rhtaXPctXgbAllThirtyNineKVsRthClose}{1.3}  % QLIKE % change, ALL48_N39k vs RTH13 (full set; negative: ALL48_N39k more accurate), XGBoost, 16:00 bar
\providecommand{\rhtaXDmXgbAllThirtyNineKVsRthClose}{}\renewcommand{\rhtaXDmXgbAllThirtyNineKVsRthClose}{1.45}  % ... its Diebold-Mariano statistic
\providecommand{\rhtaXPXgbAllThirtyNineKVsRthClose}{}\renewcommand{\rhtaXPXgbAllThirtyNineKVsRthClose}{\ensuremath{p = 0.15}}  % ... its p statement
\providecommand{\rhtaXPctRfAllThirtyNineKVsRthClose}{}\renewcommand{\rhtaXPctRfAllThirtyNineKVsRthClose}{\ensuremath{-}2.7}  % QLIKE % change, ALL48_N39k vs RTH13 (full set; negative: ALL48_N39k more accurate), Random forest, 16:00 bar
\providecommand{\rhtaXDmRfAllThirtyNineKVsRthClose}{}\renewcommand{\rhtaXDmRfAllThirtyNineKVsRthClose}{\ensuremath{-}2.72}  % ... its Diebold-Mariano statistic
\providecommand{\rhtaXPRfAllThirtyNineKVsRthClose}{}\renewcommand{\rhtaXPRfAllThirtyNineKVsRthClose}{\ensuremath{p = 0.007}}  % ... its p statement
\providecommand{\rhtaXPctRidgeAllThirtyNineKVsRthBar}{}\renewcommand{\rhtaXPctRidgeAllThirtyNineKVsRthBar}{2.7}  % QLIKE % change, ALL48_N39k vs RTH13 (full set; negative: ALL48_N39k more accurate), Ridge, every bar
\providecommand{\rhtaXDmRidgeAllThirtyNineKVsRthBar}{}\renewcommand{\rhtaXDmRidgeAllThirtyNineKVsRthBar}{4.54}  % ... its Diebold-Mariano statistic
\providecommand{\rhtaXPRidgeAllThirtyNineKVsRthBar}{}\renewcommand{\rhtaXPRidgeAllThirtyNineKVsRthBar}{\ensuremath{p < 0.001}}  % ... its p statement
\providecommand{\rhtaXPctLassoAllThirtyNineKVsRthBar}{}\renewcommand{\rhtaXPctLassoAllThirtyNineKVsRthBar}{4.2}  % QLIKE % change, ALL48_N39k vs RTH13 (full set; negative: ALL48_N39k more accurate), Lasso, every bar
\providecommand{\rhtaXDmLassoAllThirtyNineKVsRthBar}{}\renewcommand{\rhtaXDmLassoAllThirtyNineKVsRthBar}{3.44}  % ... its Diebold-Mariano statistic
\providecommand{\rhtaXPLassoAllThirtyNineKVsRthBar}{}\renewcommand{\rhtaXPLassoAllThirtyNineKVsRthBar}{\ensuremath{p < 0.001}}  % ... its p statement
\providecommand{\rhtaXPctEnetAllThirtyNineKVsRthBar}{}\renewcommand{\rhtaXPctEnetAllThirtyNineKVsRthBar}{4.9}  % QLIKE % change, ALL48_N39k vs RTH13 (full set; negative: ALL48_N39k more accurate), Elastic net, every bar
\providecommand{\rhtaXDmEnetAllThirtyNineKVsRthBar}{}\renewcommand{\rhtaXDmEnetAllThirtyNineKVsRthBar}{5.04}  % ... its Diebold-Mariano statistic
\providecommand{\rhtaXPEnetAllThirtyNineKVsRthBar}{}\renewcommand{\rhtaXPEnetAllThirtyNineKVsRthBar}{\ensuremath{p < 0.001}}  % ... its p statement
\providecommand{\rhtaXPctLgbmAllThirtyNineKVsRthBar}{}\renewcommand{\rhtaXPctLgbmAllThirtyNineKVsRthBar}{10.4}  % QLIKE % change, ALL48_N39k vs RTH13 (full set; negative: ALL48_N39k more accurate), LightGBM, every bar
\providecommand{\rhtaXDmLgbmAllThirtyNineKVsRthBar}{}\renewcommand{\rhtaXDmLgbmAllThirtyNineKVsRthBar}{8.35}  % ... its Diebold-Mariano statistic
\providecommand{\rhtaXPLgbmAllThirtyNineKVsRthBar}{}\renewcommand{\rhtaXPLgbmAllThirtyNineKVsRthBar}{\ensuremath{p < 0.001}}  % ... its p statement
\providecommand{\rhtaXPctXgbAllThirtyNineKVsRthBar}{}\renewcommand{\rhtaXPctXgbAllThirtyNineKVsRthBar}{6.4}  % QLIKE % change, ALL48_N39k vs RTH13 (full set; negative: ALL48_N39k more accurate), XGBoost, every bar
\providecommand{\rhtaXDmXgbAllThirtyNineKVsRthBar}{}\renewcommand{\rhtaXDmXgbAllThirtyNineKVsRthBar}{8.06}  % ... its Diebold-Mariano statistic
\providecommand{\rhtaXPXgbAllThirtyNineKVsRthBar}{}\renewcommand{\rhtaXPXgbAllThirtyNineKVsRthBar}{\ensuremath{p < 0.001}}  % ... its p statement
\providecommand{\rhtaXPctRfAllThirtyNineKVsRthBar}{}\renewcommand{\rhtaXPctRfAllThirtyNineKVsRthBar}{7.9}  % QLIKE % change, ALL48_N39k vs RTH13 (full set; negative: ALL48_N39k more accurate), Random forest, every bar
\providecommand{\rhtaXDmRfAllThirtyNineKVsRthBar}{}\renewcommand{\rhtaXDmRfAllThirtyNineKVsRthBar}{6.79}  % ... its Diebold-Mariano statistic
\providecommand{\rhtaXPRfAllThirtyNineKVsRthBar}{}\renewcommand{\rhtaXPRfAllThirtyNineKVsRthBar}{\ensuremath{p < 0.001}}  % ... its p statement
\providecommand{\rhtaXPctRidgeAllThreeKVsAllClose}{}\renewcommand{\rhtaXPctRidgeAllThreeKVsAllClose}{3.5}  % QLIKE % change, ALL48_N3k vs ALL48 (full set; negative: ALL48_N3k more accurate), Ridge, 16:00 bar
\providecommand{\rhtaXDmRidgeAllThreeKVsAllClose}{}\renewcommand{\rhtaXDmRidgeAllThreeKVsAllClose}{2.28}  % ... its Diebold-Mariano statistic
\providecommand{\rhtaXPRidgeAllThreeKVsAllClose}{}\renewcommand{\rhtaXPRidgeAllThreeKVsAllClose}{\ensuremath{p = 0.022}}  % ... its p statement
\providecommand{\rhtaXPctLassoAllThreeKVsAllClose}{}\renewcommand{\rhtaXPctLassoAllThreeKVsAllClose}{5.6}  % QLIKE % change, ALL48_N3k vs ALL48 (full set; negative: ALL48_N3k more accurate), Lasso, 16:00 bar
\providecommand{\rhtaXDmLassoAllThreeKVsAllClose}{}\renewcommand{\rhtaXDmLassoAllThreeKVsAllClose}{3.42}  % ... its Diebold-Mariano statistic
\providecommand{\rhtaXPLassoAllThreeKVsAllClose}{}\renewcommand{\rhtaXPLassoAllThreeKVsAllClose}{\ensuremath{p < 0.001}}  % ... its p statement
\providecommand{\rhtaXPctEnetAllThreeKVsAllClose}{}\renewcommand{\rhtaXPctEnetAllThreeKVsAllClose}{8.2}  % QLIKE % change, ALL48_N3k vs ALL48 (full set; negative: ALL48_N3k more accurate), Elastic net, 16:00 bar
\providecommand{\rhtaXDmEnetAllThreeKVsAllClose}{}\renewcommand{\rhtaXDmEnetAllThreeKVsAllClose}{4.01}  % ... its Diebold-Mariano statistic
\providecommand{\rhtaXPEnetAllThreeKVsAllClose}{}\renewcommand{\rhtaXPEnetAllThreeKVsAllClose}{\ensuremath{p < 0.001}}  % ... its p statement
\providecommand{\rhtaXPctLgbmAllThreeKVsAllClose}{}\renewcommand{\rhtaXPctLgbmAllThreeKVsAllClose}{10.4}  % QLIKE % change, ALL48_N3k vs ALL48 (full set; negative: ALL48_N3k more accurate), LightGBM, 16:00 bar
\providecommand{\rhtaXDmLgbmAllThreeKVsAllClose}{}\renewcommand{\rhtaXDmLgbmAllThreeKVsAllClose}{4.29}  % ... its Diebold-Mariano statistic
\providecommand{\rhtaXPLgbmAllThreeKVsAllClose}{}\renewcommand{\rhtaXPLgbmAllThreeKVsAllClose}{\ensuremath{p < 0.001}}  % ... its p statement
\providecommand{\rhtaXPctXgbAllThreeKVsAllClose}{}\renewcommand{\rhtaXPctXgbAllThreeKVsAllClose}{7.7}  % QLIKE % change, ALL48_N3k vs ALL48 (full set; negative: ALL48_N3k more accurate), XGBoost, 16:00 bar
\providecommand{\rhtaXDmXgbAllThreeKVsAllClose}{}\renewcommand{\rhtaXDmXgbAllThreeKVsAllClose}{4.20}  % ... its Diebold-Mariano statistic
\providecommand{\rhtaXPXgbAllThreeKVsAllClose}{}\renewcommand{\rhtaXPXgbAllThreeKVsAllClose}{\ensuremath{p < 0.001}}  % ... its p statement
\providecommand{\rhtaXPctRfAllThreeKVsAllClose}{}\renewcommand{\rhtaXPctRfAllThreeKVsAllClose}{5.6}  % QLIKE % change, ALL48_N3k vs ALL48 (full set; negative: ALL48_N3k more accurate), Random forest, 16:00 bar
\providecommand{\rhtaXDmRfAllThreeKVsAllClose}{}\renewcommand{\rhtaXDmRfAllThreeKVsAllClose}{3.47}  % ... its Diebold-Mariano statistic
\providecommand{\rhtaXPRfAllThreeKVsAllClose}{}\renewcommand{\rhtaXPRfAllThreeKVsAllClose}{\ensuremath{p < 0.001}}  % ... its p statement
\providecommand{\rhtaXPctRidgeAllThreeKVsAllBar}{}\renewcommand{\rhtaXPctRidgeAllThreeKVsAllBar}{5.4}  % QLIKE % change, ALL48_N3k vs ALL48 (full set; negative: ALL48_N3k more accurate), Ridge, every bar
\providecommand{\rhtaXDmRidgeAllThreeKVsAllBar}{}\renewcommand{\rhtaXDmRidgeAllThreeKVsAllBar}{6.67}  % ... its Diebold-Mariano statistic
\providecommand{\rhtaXPRidgeAllThreeKVsAllBar}{}\renewcommand{\rhtaXPRidgeAllThreeKVsAllBar}{\ensuremath{p < 0.001}}  % ... its p statement
\providecommand{\rhtaXPctLassoAllThreeKVsAllBar}{}\renewcommand{\rhtaXPctLassoAllThreeKVsAllBar}{5.9}  % QLIKE % change, ALL48_N3k vs ALL48 (full set; negative: ALL48_N3k more accurate), Lasso, every bar
\providecommand{\rhtaXDmLassoAllThreeKVsAllBar}{}\renewcommand{\rhtaXDmLassoAllThreeKVsAllBar}{7.47}  % ... its Diebold-Mariano statistic
\providecommand{\rhtaXPLassoAllThreeKVsAllBar}{}\renewcommand{\rhtaXPLassoAllThreeKVsAllBar}{\ensuremath{p < 0.001}}  % ... its p statement
\providecommand{\rhtaXPctEnetAllThreeKVsAllBar}{}\renewcommand{\rhtaXPctEnetAllThreeKVsAllBar}{6.6}  % QLIKE % change, ALL48_N3k vs ALL48 (full set; negative: ALL48_N3k more accurate), Elastic net, every bar
\providecommand{\rhtaXDmEnetAllThreeKVsAllBar}{}\renewcommand{\rhtaXDmEnetAllThreeKVsAllBar}{7.48}  % ... its Diebold-Mariano statistic
\providecommand{\rhtaXPEnetAllThreeKVsAllBar}{}\renewcommand{\rhtaXPEnetAllThreeKVsAllBar}{\ensuremath{p < 0.001}}  % ... its p statement
\providecommand{\rhtaXPctLgbmAllThreeKVsAllBar}{}\renewcommand{\rhtaXPctLgbmAllThreeKVsAllBar}{15.2}  % QLIKE % change, ALL48_N3k vs ALL48 (full set; negative: ALL48_N3k more accurate), LightGBM, every bar
\providecommand{\rhtaXDmLgbmAllThreeKVsAllBar}{}\renewcommand{\rhtaXDmLgbmAllThreeKVsAllBar}{8.62}  % ... its Diebold-Mariano statistic
\providecommand{\rhtaXPLgbmAllThreeKVsAllBar}{}\renewcommand{\rhtaXPLgbmAllThreeKVsAllBar}{\ensuremath{p < 0.001}}  % ... its p statement
\providecommand{\rhtaXPctXgbAllThreeKVsAllBar}{}\renewcommand{\rhtaXPctXgbAllThreeKVsAllBar}{12.7}  % QLIKE % change, ALL48_N3k vs ALL48 (full set; negative: ALL48_N3k more accurate), XGBoost, every bar
\providecommand{\rhtaXDmXgbAllThreeKVsAllBar}{}\renewcommand{\rhtaXDmXgbAllThreeKVsAllBar}{8.10}  % ... its Diebold-Mariano statistic
\providecommand{\rhtaXPXgbAllThreeKVsAllBar}{}\renewcommand{\rhtaXPXgbAllThreeKVsAllBar}{\ensuremath{p < 0.001}}  % ... its p statement
\providecommand{\rhtaXPctRfAllThreeKVsAllBar}{}\renewcommand{\rhtaXPctRfAllThreeKVsAllBar}{6.8}  % QLIKE % change, ALL48_N3k vs ALL48 (full set; negative: ALL48_N3k more accurate), Random forest, every bar
\providecommand{\rhtaXDmRfAllThreeKVsAllBar}{}\renewcommand{\rhtaXDmRfAllThreeKVsAllBar}{7.80}  % ... its Diebold-Mariano statistic
\providecommand{\rhtaXPRfAllThreeKVsAllBar}{}\renewcommand{\rhtaXPRfAllThreeKVsAllBar}{\ensuremath{p < 0.001}}  % ... its p statement
\providecommand{\rhtaXPctRidgeRthThreeKVsRthClose}{}\renewcommand{\rhtaXPctRidgeRthThreeKVsRthClose}{8.7}  % QLIKE % change, RTH13_N3k vs RTH13 (full set; negative: RTH13_N3k more accurate), Ridge, 16:00 bar
\providecommand{\rhtaXDmRidgeRthThreeKVsRthClose}{}\renewcommand{\rhtaXDmRidgeRthThreeKVsRthClose}{4.61}  % ... its Diebold-Mariano statistic
\providecommand{\rhtaXPRidgeRthThreeKVsRthClose}{}\renewcommand{\rhtaXPRidgeRthThreeKVsRthClose}{\ensuremath{p < 0.001}}  % ... its p statement
\providecommand{\rhtaXPctLassoRthThreeKVsRthClose}{}\renewcommand{\rhtaXPctLassoRthThreeKVsRthClose}{6.2}  % QLIKE % change, RTH13_N3k vs RTH13 (full set; negative: RTH13_N3k more accurate), Lasso, 16:00 bar
\providecommand{\rhtaXDmLassoRthThreeKVsRthClose}{}\renewcommand{\rhtaXDmLassoRthThreeKVsRthClose}{2.63}  % ... its Diebold-Mariano statistic
\providecommand{\rhtaXPLassoRthThreeKVsRthClose}{}\renewcommand{\rhtaXPLassoRthThreeKVsRthClose}{\ensuremath{p = 0.009}}  % ... its p statement
\providecommand{\rhtaXPctEnetRthThreeKVsRthClose}{}\renewcommand{\rhtaXPctEnetRthThreeKVsRthClose}{10.6}  % QLIKE % change, RTH13_N3k vs RTH13 (full set; negative: RTH13_N3k more accurate), Elastic net, 16:00 bar
\providecommand{\rhtaXDmEnetRthThreeKVsRthClose}{}\renewcommand{\rhtaXDmEnetRthThreeKVsRthClose}{5.13}  % ... its Diebold-Mariano statistic
\providecommand{\rhtaXPEnetRthThreeKVsRthClose}{}\renewcommand{\rhtaXPEnetRthThreeKVsRthClose}{\ensuremath{p < 0.001}}  % ... its p statement
\providecommand{\rhtaXPctLgbmRthThreeKVsRthClose}{}\renewcommand{\rhtaXPctLgbmRthThreeKVsRthClose}{12.3}  % QLIKE % change, RTH13_N3k vs RTH13 (full set; negative: RTH13_N3k more accurate), LightGBM, 16:00 bar
\providecommand{\rhtaXDmLgbmRthThreeKVsRthClose}{}\renewcommand{\rhtaXDmLgbmRthThreeKVsRthClose}{5.35}  % ... its Diebold-Mariano statistic
\providecommand{\rhtaXPLgbmRthThreeKVsRthClose}{}\renewcommand{\rhtaXPLgbmRthThreeKVsRthClose}{\ensuremath{p < 0.001}}  % ... its p statement
\providecommand{\rhtaXPctXgbRthThreeKVsRthClose}{}\renewcommand{\rhtaXPctXgbRthThreeKVsRthClose}{9.3}  % QLIKE % change, RTH13_N3k vs RTH13 (full set; negative: RTH13_N3k more accurate), XGBoost, 16:00 bar
\providecommand{\rhtaXDmXgbRthThreeKVsRthClose}{}\renewcommand{\rhtaXDmXgbRthThreeKVsRthClose}{5.11}  % ... its Diebold-Mariano statistic
\providecommand{\rhtaXPXgbRthThreeKVsRthClose}{}\renewcommand{\rhtaXPXgbRthThreeKVsRthClose}{\ensuremath{p < 0.001}}  % ... its p statement
\providecommand{\rhtaXPctRfRthThreeKVsRthClose}{}\renewcommand{\rhtaXPctRfRthThreeKVsRthClose}{7.1}  % QLIKE % change, RTH13_N3k vs RTH13 (full set; negative: RTH13_N3k more accurate), Random forest, 16:00 bar
\providecommand{\rhtaXDmRfRthThreeKVsRthClose}{}\renewcommand{\rhtaXDmRfRthThreeKVsRthClose}{5.16}  % ... its Diebold-Mariano statistic
\providecommand{\rhtaXPRfRthThreeKVsRthClose}{}\renewcommand{\rhtaXPRfRthThreeKVsRthClose}{\ensuremath{p < 0.001}}  % ... its p statement
\providecommand{\rhtaXPctRidgeRthThreeKVsRthBar}{}\renewcommand{\rhtaXPctRidgeRthThreeKVsRthBar}{2.7}  % QLIKE % change, RTH13_N3k vs RTH13 (full set; negative: RTH13_N3k more accurate), Ridge, every bar
\providecommand{\rhtaXDmRidgeRthThreeKVsRthBar}{}\renewcommand{\rhtaXDmRidgeRthThreeKVsRthBar}{3.01}  % ... its Diebold-Mariano statistic
\providecommand{\rhtaXPRidgeRthThreeKVsRthBar}{}\renewcommand{\rhtaXPRidgeRthThreeKVsRthBar}{\ensuremath{p = 0.003}}  % ... its p statement
\providecommand{\rhtaXPctLassoRthThreeKVsRthBar}{}\renewcommand{\rhtaXPctLassoRthThreeKVsRthBar}{3.5}  % QLIKE % change, RTH13_N3k vs RTH13 (full set; negative: RTH13_N3k more accurate), Lasso, every bar
\providecommand{\rhtaXDmLassoRthThreeKVsRthBar}{}\renewcommand{\rhtaXDmLassoRthThreeKVsRthBar}{2.71}  % ... its Diebold-Mariano statistic
\providecommand{\rhtaXPLassoRthThreeKVsRthBar}{}\renewcommand{\rhtaXPLassoRthThreeKVsRthBar}{\ensuremath{p = 0.007}}  % ... its p statement
\providecommand{\rhtaXPctEnetRthThreeKVsRthBar}{}\renewcommand{\rhtaXPctEnetRthThreeKVsRthBar}{2.9}  % QLIKE % change, RTH13_N3k vs RTH13 (full set; negative: RTH13_N3k more accurate), Elastic net, every bar
\providecommand{\rhtaXDmEnetRthThreeKVsRthBar}{}\renewcommand{\rhtaXDmEnetRthThreeKVsRthBar}{3.54}  % ... its Diebold-Mariano statistic
\providecommand{\rhtaXPEnetRthThreeKVsRthBar}{}\renewcommand{\rhtaXPEnetRthThreeKVsRthBar}{\ensuremath{p < 0.001}}  % ... its p statement
\providecommand{\rhtaXPctLgbmRthThreeKVsRthBar}{}\renewcommand{\rhtaXPctLgbmRthThreeKVsRthBar}{14.9}  % QLIKE % change, RTH13_N3k vs RTH13 (full set; negative: RTH13_N3k more accurate), LightGBM, every bar
\providecommand{\rhtaXDmLgbmRthThreeKVsRthBar}{}\renewcommand{\rhtaXDmLgbmRthThreeKVsRthBar}{8.05}  % ... its Diebold-Mariano statistic
\providecommand{\rhtaXPLgbmRthThreeKVsRthBar}{}\renewcommand{\rhtaXPLgbmRthThreeKVsRthBar}{\ensuremath{p < 0.001}}  % ... its p statement
\providecommand{\rhtaXPctXgbRthThreeKVsRthBar}{}\renewcommand{\rhtaXPctXgbRthThreeKVsRthBar}{10.7}  % QLIKE % change, RTH13_N3k vs RTH13 (full set; negative: RTH13_N3k more accurate), XGBoost, every bar
\providecommand{\rhtaXDmXgbRthThreeKVsRthBar}{}\renewcommand{\rhtaXDmXgbRthThreeKVsRthBar}{7.31}  % ... its Diebold-Mariano statistic
\providecommand{\rhtaXPXgbRthThreeKVsRthBar}{}\renewcommand{\rhtaXPXgbRthThreeKVsRthBar}{\ensuremath{p < 0.001}}  % ... its p statement
\providecommand{\rhtaXPctRfRthThreeKVsRthBar}{}\renewcommand{\rhtaXPctRfRthThreeKVsRthBar}{9.4}  % QLIKE % change, RTH13_N3k vs RTH13 (full set; negative: RTH13_N3k more accurate), Random forest, every bar
\providecommand{\rhtaXDmRfRthThreeKVsRthBar}{}\renewcommand{\rhtaXDmRfRthThreeKVsRthBar}{7.22}  % ... its Diebold-Mariano statistic
\providecommand{\rhtaXPRfRthThreeKVsRthBar}{}\renewcommand{\rhtaXPRfRthThreeKVsRthBar}{\ensuremath{p < 0.001}}  % ... its p statement
\providecommand{\rhtaXPctRidgeAllThirtyNineKVsAllClose}{}\renewcommand{\rhtaXPctRidgeAllThirtyNineKVsAllClose}{4.1}  % QLIKE % change, ALL48_N39k vs ALL48 (full set; negative: ALL48_N39k more accurate), Ridge, 16:00 bar
\providecommand{\rhtaXDmRidgeAllThirtyNineKVsAllClose}{}\renewcommand{\rhtaXDmRidgeAllThirtyNineKVsAllClose}{4.03}  % ... its Diebold-Mariano statistic
\providecommand{\rhtaXPRidgeAllThirtyNineKVsAllClose}{}\renewcommand{\rhtaXPRidgeAllThirtyNineKVsAllClose}{\ensuremath{p < 0.001}}  % ... its p statement
\providecommand{\rhtaXPctLassoAllThirtyNineKVsAllClose}{}\renewcommand{\rhtaXPctLassoAllThirtyNineKVsAllClose}{2.6}  % QLIKE % change, ALL48_N39k vs ALL48 (full set; negative: ALL48_N39k more accurate), Lasso, 16:00 bar
\providecommand{\rhtaXDmLassoAllThirtyNineKVsAllClose}{}\renewcommand{\rhtaXDmLassoAllThirtyNineKVsAllClose}{2.67}  % ... its Diebold-Mariano statistic
\providecommand{\rhtaXPLassoAllThirtyNineKVsAllClose}{}\renewcommand{\rhtaXPLassoAllThirtyNineKVsAllClose}{\ensuremath{p = 0.008}}  % ... its p statement
\providecommand{\rhtaXPctEnetAllThirtyNineKVsAllClose}{}\renewcommand{\rhtaXPctEnetAllThirtyNineKVsAllClose}{2.2}  % QLIKE % change, ALL48_N39k vs ALL48 (full set; negative: ALL48_N39k more accurate), Elastic net, 16:00 bar
\providecommand{\rhtaXDmEnetAllThirtyNineKVsAllClose}{}\renewcommand{\rhtaXDmEnetAllThirtyNineKVsAllClose}{1.89}  % ... its Diebold-Mariano statistic
\providecommand{\rhtaXPEnetAllThirtyNineKVsAllClose}{}\renewcommand{\rhtaXPEnetAllThirtyNineKVsAllClose}{\ensuremath{p = 0.059}}  % ... its p statement
\providecommand{\rhtaXPctLgbmAllThirtyNineKVsAllClose}{}\renewcommand{\rhtaXPctLgbmAllThirtyNineKVsAllClose}{2.7}  % QLIKE % change, ALL48_N39k vs ALL48 (full set; negative: ALL48_N39k more accurate), LightGBM, 16:00 bar
\providecommand{\rhtaXDmLgbmAllThirtyNineKVsAllClose}{}\renewcommand{\rhtaXDmLgbmAllThirtyNineKVsAllClose}{2.87}  % ... its Diebold-Mariano statistic
\providecommand{\rhtaXPLgbmAllThirtyNineKVsAllClose}{}\renewcommand{\rhtaXPLgbmAllThirtyNineKVsAllClose}{\ensuremath{p = 0.004}}  % ... its p statement
\providecommand{\rhtaXPctXgbAllThirtyNineKVsAllClose}{}\renewcommand{\rhtaXPctXgbAllThirtyNineKVsAllClose}{3.3}  % QLIKE % change, ALL48_N39k vs ALL48 (full set; negative: ALL48_N39k more accurate), XGBoost, 16:00 bar
\providecommand{\rhtaXDmXgbAllThirtyNineKVsAllClose}{}\renewcommand{\rhtaXDmXgbAllThirtyNineKVsAllClose}{4.05}  % ... its Diebold-Mariano statistic
\providecommand{\rhtaXPXgbAllThirtyNineKVsAllClose}{}\renewcommand{\rhtaXPXgbAllThirtyNineKVsAllClose}{\ensuremath{p < 0.001}}  % ... its p statement
\providecommand{\rhtaXPctRfAllThirtyNineKVsAllClose}{}\renewcommand{\rhtaXPctRfAllThirtyNineKVsAllClose}{1.7}  % QLIKE % change, ALL48_N39k vs ALL48 (full set; negative: ALL48_N39k more accurate), Random forest, 16:00 bar
\providecommand{\rhtaXDmRfAllThirtyNineKVsAllClose}{}\renewcommand{\rhtaXDmRfAllThirtyNineKVsAllClose}{2.38}  % ... its Diebold-Mariano statistic
\providecommand{\rhtaXPRfAllThirtyNineKVsAllClose}{}\renewcommand{\rhtaXPRfAllThirtyNineKVsAllClose}{\ensuremath{p = 0.017}}  % ... its p statement
\providecommand{\rhtaXPctRidgeAllThirtyNineKVsAllBar}{}\renewcommand{\rhtaXPctRidgeAllThirtyNineKVsAllBar}{0.8}  % QLIKE % change, ALL48_N39k vs ALL48 (full set; negative: ALL48_N39k more accurate), Ridge, every bar
\providecommand{\rhtaXDmRidgeAllThirtyNineKVsAllBar}{}\renewcommand{\rhtaXDmRidgeAllThirtyNineKVsAllBar}{1.56}  % ... its Diebold-Mariano statistic
\providecommand{\rhtaXPRidgeAllThirtyNineKVsAllBar}{}\renewcommand{\rhtaXPRidgeAllThirtyNineKVsAllBar}{\ensuremath{p = 0.12}}  % ... its p statement
\providecommand{\rhtaXPctLassoAllThirtyNineKVsAllBar}{}\renewcommand{\rhtaXPctLassoAllThirtyNineKVsAllBar}{1.6}  % QLIKE % change, ALL48_N39k vs ALL48 (full set; negative: ALL48_N39k more accurate), Lasso, every bar
\providecommand{\rhtaXDmLassoAllThirtyNineKVsAllBar}{}\renewcommand{\rhtaXDmLassoAllThirtyNineKVsAllBar}{3.26}  % ... its Diebold-Mariano statistic
\providecommand{\rhtaXPLassoAllThirtyNineKVsAllBar}{}\renewcommand{\rhtaXPLassoAllThirtyNineKVsAllBar}{\ensuremath{p = 0.001}}  % ... its p statement
\providecommand{\rhtaXPctEnetAllThirtyNineKVsAllBar}{}\renewcommand{\rhtaXPctEnetAllThirtyNineKVsAllBar}{2.7}  % QLIKE % change, ALL48_N39k vs ALL48 (full set; negative: ALL48_N39k more accurate), Elastic net, every bar
\providecommand{\rhtaXDmEnetAllThirtyNineKVsAllBar}{}\renewcommand{\rhtaXDmEnetAllThirtyNineKVsAllBar}{3.42}  % ... its Diebold-Mariano statistic
\providecommand{\rhtaXPEnetAllThirtyNineKVsAllBar}{}\renewcommand{\rhtaXPEnetAllThirtyNineKVsAllBar}{\ensuremath{p < 0.001}}  % ... its p statement
\providecommand{\rhtaXPctLgbmAllThirtyNineKVsAllBar}{}\renewcommand{\rhtaXPctLgbmAllThirtyNineKVsAllBar}{5.7}  % QLIKE % change, ALL48_N39k vs ALL48 (full set; negative: ALL48_N39k more accurate), LightGBM, every bar
\providecommand{\rhtaXDmLgbmAllThirtyNineKVsAllBar}{}\renewcommand{\rhtaXDmLgbmAllThirtyNineKVsAllBar}{6.76}  % ... its Diebold-Mariano statistic
\providecommand{\rhtaXPLgbmAllThirtyNineKVsAllBar}{}\renewcommand{\rhtaXPLgbmAllThirtyNineKVsAllBar}{\ensuremath{p < 0.001}}  % ... its p statement
\providecommand{\rhtaXPctXgbAllThirtyNineKVsAllBar}{}\renewcommand{\rhtaXPctXgbAllThirtyNineKVsAllBar}{3.2}  % QLIKE % change, ALL48_N39k vs ALL48 (full set; negative: ALL48_N39k more accurate), XGBoost, every bar
\providecommand{\rhtaXDmXgbAllThirtyNineKVsAllBar}{}\renewcommand{\rhtaXDmXgbAllThirtyNineKVsAllBar}{6.86}  % ... its Diebold-Mariano statistic
\providecommand{\rhtaXPXgbAllThirtyNineKVsAllBar}{}\renewcommand{\rhtaXPXgbAllThirtyNineKVsAllBar}{\ensuremath{p < 0.001}}  % ... its p statement
\providecommand{\rhtaXPctRfAllThirtyNineKVsAllBar}{}\renewcommand{\rhtaXPctRfAllThirtyNineKVsAllBar}{2.6}  % QLIKE % change, ALL48_N39k vs ALL48 (full set; negative: ALL48_N39k more accurate), Random forest, every bar
\providecommand{\rhtaXDmRfAllThirtyNineKVsAllBar}{}\renewcommand{\rhtaXDmRfAllThirtyNineKVsAllBar}{4.91}  % ... its Diebold-Mariano statistic
\providecommand{\rhtaXPRfAllThirtyNineKVsAllBar}{}\renewcommand{\rhtaXPRfAllThirtyNineKVsAllBar}{\ensuremath{p < 0.001}}  % ... its p statement
\providecommand{\rhtaShRidgeHarRth}{}\renewcommand{\rhtaShRidgeHarRth}{0.95}  % Sharpe ratio (D19), mid-quotes: Ridge, baseline, RTH13 (main design: the 13 regular-hours bars, 3,000 days)
\providecommand{\rhtaShXRidgeHarRth}{}\renewcommand{\rhtaShXRidgeHarRth}{0.60}  % ... at quoted bid/ask prices
\providecommand{\rhtaShRidgeFullRth}{}\renewcommand{\rhtaShRidgeFullRth}{1.77}  % Sharpe ratio (D19), mid-quotes: Ridge, all_features, RTH13 (main design: the 13 regular-hours bars, 3,000 days)
\providecommand{\rhtaShXRidgeFullRth}{}\renewcommand{\rhtaShXRidgeFullRth}{1.41}  % ... at quoted bid/ask prices
\providecommand{\rhtaPShRidgeRth}{}\renewcommand{\rhtaPShRidgeRth}{\ensuremath{p = 0.11}}  % Ledoit-Wolf p, Sharpe ratio full minus HAR (mid): Ridge, RTH13
\providecommand{\rhtaShLassoHarRth}{}\renewcommand{\rhtaShLassoHarRth}{0.95}  % Sharpe ratio (D19), mid-quotes: Lasso, baseline, RTH13 (main design: the 13 regular-hours bars, 3,000 days)
\providecommand{\rhtaShXLassoHarRth}{}\renewcommand{\rhtaShXLassoHarRth}{0.60}  % ... at quoted bid/ask prices
\providecommand{\rhtaShLassoFullRth}{}\renewcommand{\rhtaShLassoFullRth}{1.72}  % Sharpe ratio (D19), mid-quotes: Lasso, all_features, RTH13 (main design: the 13 regular-hours bars, 3,000 days)
\providecommand{\rhtaShXLassoFullRth}{}\renewcommand{\rhtaShXLassoFullRth}{1.37}  % ... at quoted bid/ask prices
\providecommand{\rhtaPShLassoRth}{}\renewcommand{\rhtaPShLassoRth}{\ensuremath{p = 0.12}}  % Ledoit-Wolf p, Sharpe ratio full minus HAR (mid): Lasso, RTH13
\providecommand{\rhtaShEnetHarRth}{}\renewcommand{\rhtaShEnetHarRth}{0.95}  % Sharpe ratio (D19), mid-quotes: Elastic net, baseline, RTH13 (main design: the 13 regular-hours bars, 3,000 days)
\providecommand{\rhtaShXEnetHarRth}{}\renewcommand{\rhtaShXEnetHarRth}{0.60}  % ... at quoted bid/ask prices
\providecommand{\rhtaShEnetFullRth}{}\renewcommand{\rhtaShEnetFullRth}{1.86}  % Sharpe ratio (D19), mid-quotes: Elastic net, all_features, RTH13 (main design: the 13 regular-hours bars, 3,000 days)
\providecommand{\rhtaShXEnetFullRth}{}\renewcommand{\rhtaShXEnetFullRth}{1.50}  % ... at quoted bid/ask prices
\providecommand{\rhtaPShEnetRth}{}\renewcommand{\rhtaPShEnetRth}{\ensuremath{p = 0.077}}  % Ledoit-Wolf p, Sharpe ratio full minus HAR (mid): Elastic net, RTH13
\providecommand{\rhtaShLgbmHarRth}{}\renewcommand{\rhtaShLgbmHarRth}{0.65}  % Sharpe ratio (D19), mid-quotes: LightGBM, baseline, RTH13 (main design: the 13 regular-hours bars, 3,000 days)
\providecommand{\rhtaShXLgbmHarRth}{}\renewcommand{\rhtaShXLgbmHarRth}{0.31}  % ... at quoted bid/ask prices
\providecommand{\rhtaShLgbmFullRth}{}\renewcommand{\rhtaShLgbmFullRth}{1.59}  % Sharpe ratio (D19), mid-quotes: LightGBM, all_features, RTH13 (main design: the 13 regular-hours bars, 3,000 days)
\providecommand{\rhtaShXLgbmFullRth}{}\renewcommand{\rhtaShXLgbmFullRth}{1.23}  % ... at quoted bid/ask prices
\providecommand{\rhtaPShLgbmRth}{}\renewcommand{\rhtaPShLgbmRth}{\ensuremath{p = 0.040}}  % Ledoit-Wolf p, Sharpe ratio full minus HAR (mid): LightGBM, RTH13
\providecommand{\rhtaShXgbHarRth}{}\renewcommand{\rhtaShXgbHarRth}{0.98}  % Sharpe ratio (D19), mid-quotes: XGBoost, baseline, RTH13 (main design: the 13 regular-hours bars, 3,000 days)
\providecommand{\rhtaShXXgbHarRth}{}\renewcommand{\rhtaShXXgbHarRth}{0.63}  % ... at quoted bid/ask prices
\providecommand{\rhtaShXgbFullRth}{}\renewcommand{\rhtaShXgbFullRth}{1.47}  % Sharpe ratio (D19), mid-quotes: XGBoost, all_features, RTH13 (main design: the 13 regular-hours bars, 3,000 days)
\providecommand{\rhtaShXXgbFullRth}{}\renewcommand{\rhtaShXXgbFullRth}{1.12}  % ... at quoted bid/ask prices
\providecommand{\rhtaPShXgbRth}{}\renewcommand{\rhtaPShXgbRth}{\ensuremath{p = 0.27}}  % Ledoit-Wolf p, Sharpe ratio full minus HAR (mid): XGBoost, RTH13
\providecommand{\rhtaShRfHarRth}{}\renewcommand{\rhtaShRfHarRth}{0.32}  % Sharpe ratio (D19), mid-quotes: Random forest, baseline, RTH13 (main design: the 13 regular-hours bars, 3,000 days)
\providecommand{\rhtaShXRfHarRth}{}\renewcommand{\rhtaShXRfHarRth}{\ensuremath{-}0.03}  % ... at quoted bid/ask prices
\providecommand{\rhtaShRfFullRth}{}\renewcommand{\rhtaShRfFullRth}{1.88}  % Sharpe ratio (D19), mid-quotes: Random forest, all_features, RTH13 (main design: the 13 regular-hours bars, 3,000 days)
\providecommand{\rhtaShXRfFullRth}{}\renewcommand{\rhtaShXRfFullRth}{1.52}  % ... at quoted bid/ask prices
\providecommand{\rhtaPShRfRth}{}\renewcommand{\rhtaPShRfRth}{\ensuremath{p = 0.020}}  % Ledoit-Wolf p, Sharpe ratio full minus HAR (mid): Random forest, RTH13
\providecommand{\rhtaShRidgeHarAll}{}\renewcommand{\rhtaShRidgeHarAll}{0.95}  % Sharpe ratio (D19), mid-quotes: Ridge, baseline, ALL48 (all 48 intervals, 3,000 days)
\providecommand{\rhtaShXRidgeHarAll}{}\renewcommand{\rhtaShXRidgeHarAll}{0.59}  % ... at quoted bid/ask prices
\providecommand{\rhtaShRidgeFullAll}{}\renewcommand{\rhtaShRidgeFullAll}{1.61}  % Sharpe ratio (D19), mid-quotes: Ridge, all_features, ALL48 (all 48 intervals, 3,000 days)
\providecommand{\rhtaShXRidgeFullAll}{}\renewcommand{\rhtaShXRidgeFullAll}{1.26}  % ... at quoted bid/ask prices
\providecommand{\rhtaShLassoHarAll}{}\renewcommand{\rhtaShLassoHarAll}{0.93}  % Sharpe ratio (D19), mid-quotes: Lasso, baseline, ALL48 (all 48 intervals, 3,000 days)
\providecommand{\rhtaShXLassoHarAll}{}\renewcommand{\rhtaShXLassoHarAll}{0.58}  % ... at quoted bid/ask prices
\providecommand{\rhtaShLassoFullAll}{}\renewcommand{\rhtaShLassoFullAll}{1.50}  % Sharpe ratio (D19), mid-quotes: Lasso, all_features, ALL48 (all 48 intervals, 3,000 days)
\providecommand{\rhtaShXLassoFullAll}{}\renewcommand{\rhtaShXLassoFullAll}{1.14}  % ... at quoted bid/ask prices
\providecommand{\rhtaShEnetHarAll}{}\renewcommand{\rhtaShEnetHarAll}{0.89}  % Sharpe ratio (D19), mid-quotes: Elastic net, baseline, ALL48 (all 48 intervals, 3,000 days)
\providecommand{\rhtaShXEnetHarAll}{}\renewcommand{\rhtaShXEnetHarAll}{0.54}  % ... at quoted bid/ask prices
\providecommand{\rhtaShEnetFullAll}{}\renewcommand{\rhtaShEnetFullAll}{1.52}  % Sharpe ratio (D19), mid-quotes: Elastic net, all_features, ALL48 (all 48 intervals, 3,000 days)
\providecommand{\rhtaShXEnetFullAll}{}\renewcommand{\rhtaShXEnetFullAll}{1.17}  % ... at quoted bid/ask prices
\providecommand{\rhtaShLgbmHarAll}{}\renewcommand{\rhtaShLgbmHarAll}{0.21}  % Sharpe ratio (D19), mid-quotes: LightGBM, baseline, ALL48 (all 48 intervals, 3,000 days)
\providecommand{\rhtaShXLgbmHarAll}{}\renewcommand{\rhtaShXLgbmHarAll}{\ensuremath{-}0.14}  % ... at quoted bid/ask prices
\providecommand{\rhtaShLgbmFullAll}{}\renewcommand{\rhtaShLgbmFullAll}{1.25}  % Sharpe ratio (D19), mid-quotes: LightGBM, all_features, ALL48 (all 48 intervals, 3,000 days)
\providecommand{\rhtaShXLgbmFullAll}{}\renewcommand{\rhtaShXLgbmFullAll}{0.90}  % ... at quoted bid/ask prices
\providecommand{\rhtaPShLgbmAll}{}\renewcommand{\rhtaPShLgbmAll}{\ensuremath{p = 0.11}}  % Ledoit-Wolf p, Sharpe ratio full minus HAR (mid): LightGBM, ALL48
\providecommand{\rhtaShXgbHarAll}{}\renewcommand{\rhtaShXgbHarAll}{0.73}  % Sharpe ratio (D19), mid-quotes: XGBoost, baseline, ALL48 (all 48 intervals, 3,000 days)
\providecommand{\rhtaShXXgbHarAll}{}\renewcommand{\rhtaShXXgbHarAll}{0.38}  % ... at quoted bid/ask prices
\providecommand{\rhtaShXgbFullAll}{}\renewcommand{\rhtaShXgbFullAll}{1.50}  % Sharpe ratio (D19), mid-quotes: XGBoost, all_features, ALL48 (all 48 intervals, 3,000 days)
\providecommand{\rhtaShXXgbFullAll}{}\renewcommand{\rhtaShXXgbFullAll}{1.15}  % ... at quoted bid/ask prices
\providecommand{\rhtaPShXgbAll}{}\renewcommand{\rhtaPShXgbAll}{\ensuremath{p = 0.11}}  % Ledoit-Wolf p, Sharpe ratio full minus HAR (mid): XGBoost, ALL48
\providecommand{\rhtaShRfHarAll}{}\renewcommand{\rhtaShRfHarAll}{0.65}  % Sharpe ratio (D19), mid-quotes: Random forest, baseline, ALL48 (all 48 intervals, 3,000 days)
\providecommand{\rhtaShXRfHarAll}{}\renewcommand{\rhtaShXRfHarAll}{0.29}  % ... at quoted bid/ask prices
\providecommand{\rhtaShRfFullAll}{}\renewcommand{\rhtaShRfFullAll}{1.82}  % Sharpe ratio (D19), mid-quotes: Random forest, all_features, ALL48 (all 48 intervals, 3,000 days)
\providecommand{\rhtaShXRfFullAll}{}\renewcommand{\rhtaShXRfFullAll}{1.46}  % ... at quoted bid/ask prices
\providecommand{\rhtaPShRfAll}{}\renewcommand{\rhtaPShRfAll}{\ensuremath{p = 0.066}}  % Ledoit-Wolf p, Sharpe ratio full minus HAR (mid): Random forest, ALL48
\providecommand{\rhtaShRidgeHarCl}{}\renewcommand{\rhtaShRidgeHarCl}{1.03}  % Sharpe ratio (D19), mid-quotes: Ridge, baseline, CLOSE (the closing half hour only, 3,000 closes)
\providecommand{\rhtaShXRidgeHarCl}{}\renewcommand{\rhtaShXRidgeHarCl}{0.68}  % ... at quoted bid/ask prices
\providecommand{\rhtaShRidgeFullCl}{}\renewcommand{\rhtaShRidgeFullCl}{1.26}  % Sharpe ratio (D19), mid-quotes: Ridge, all_features, CLOSE (the closing half hour only, 3,000 closes)
\providecommand{\rhtaShXRidgeFullCl}{}\renewcommand{\rhtaShXRidgeFullCl}{0.91}  % ... at quoted bid/ask prices
\providecommand{\rhtaShLassoHarCl}{}\renewcommand{\rhtaShLassoHarCl}{1.03}  % Sharpe ratio (D19), mid-quotes: Lasso, baseline, CLOSE (the closing half hour only, 3,000 closes)
\providecommand{\rhtaShXLassoHarCl}{}\renewcommand{\rhtaShXLassoHarCl}{0.68}  % ... at quoted bid/ask prices
\providecommand{\rhtaShLassoFullCl}{}\renewcommand{\rhtaShLassoFullCl}{1.58}  % Sharpe ratio (D19), mid-quotes: Lasso, all_features, CLOSE (the closing half hour only, 3,000 closes)
\providecommand{\rhtaShXLassoFullCl}{}\renewcommand{\rhtaShXLassoFullCl}{1.23}  % ... at quoted bid/ask prices
\providecommand{\rhtaShEnetHarCl}{}\renewcommand{\rhtaShEnetHarCl}{1.03}  % Sharpe ratio (D19), mid-quotes: Elastic net, baseline, CLOSE (the closing half hour only, 3,000 closes)
\providecommand{\rhtaShXEnetHarCl}{}\renewcommand{\rhtaShXEnetHarCl}{0.68}  % ... at quoted bid/ask prices
\providecommand{\rhtaShEnetFullCl}{}\renewcommand{\rhtaShEnetFullCl}{1.71}  % Sharpe ratio (D19), mid-quotes: Elastic net, all_features, CLOSE (the closing half hour only, 3,000 closes)
\providecommand{\rhtaShXEnetFullCl}{}\renewcommand{\rhtaShXEnetFullCl}{1.36}  % ... at quoted bid/ask prices
\providecommand{\rhtaShLgbmHarCl}{}\renewcommand{\rhtaShLgbmHarCl}{\ensuremath{-}0.03}  % Sharpe ratio (D19), mid-quotes: LightGBM, baseline, CLOSE (the closing half hour only, 3,000 closes)
\providecommand{\rhtaShXLgbmHarCl}{}\renewcommand{\rhtaShXLgbmHarCl}{\ensuremath{-}0.37}  % ... at quoted bid/ask prices
\providecommand{\rhtaShLgbmFullCl}{}\renewcommand{\rhtaShLgbmFullCl}{1.94}  % Sharpe ratio (D19), mid-quotes: LightGBM, all_features, CLOSE (the closing half hour only, 3,000 closes)
\providecommand{\rhtaShXLgbmFullCl}{}\renewcommand{\rhtaShXLgbmFullCl}{1.60}  % ... at quoted bid/ask prices
\providecommand{\rhtaPShLgbmCl}{}\renewcommand{\rhtaPShLgbmCl}{\ensuremath{p = 0.003}}  % Ledoit-Wolf p, Sharpe ratio full minus HAR (mid): LightGBM, CLOSE
\providecommand{\rhtaShXgbHarCl}{}\renewcommand{\rhtaShXgbHarCl}{0.40}  % Sharpe ratio (D19), mid-quotes: XGBoost, baseline, CLOSE (the closing half hour only, 3,000 closes)
\providecommand{\rhtaShXXgbHarCl}{}\renewcommand{\rhtaShXXgbHarCl}{0.05}  % ... at quoted bid/ask prices
\providecommand{\rhtaShXgbFullCl}{}\renewcommand{\rhtaShXgbFullCl}{1.62}  % Sharpe ratio (D19), mid-quotes: XGBoost, all_features, CLOSE (the closing half hour only, 3,000 closes)
\providecommand{\rhtaShXXgbFullCl}{}\renewcommand{\rhtaShXXgbFullCl}{1.27}  % ... at quoted bid/ask prices
\providecommand{\rhtaPShXgbCl}{}\renewcommand{\rhtaPShXgbCl}{\ensuremath{p = 0.055}}  % Ledoit-Wolf p, Sharpe ratio full minus HAR (mid): XGBoost, CLOSE
\providecommand{\rhtaShRfHarCl}{}\renewcommand{\rhtaShRfHarCl}{0.29}  % Sharpe ratio (D19), mid-quotes: Random forest, baseline, CLOSE (the closing half hour only, 3,000 closes)
\providecommand{\rhtaShXRfHarCl}{}\renewcommand{\rhtaShXRfHarCl}{\ensuremath{-}0.05}  % ... at quoted bid/ask prices
\providecommand{\rhtaShRfFullCl}{}\renewcommand{\rhtaShRfFullCl}{1.29}  % Sharpe ratio (D19), mid-quotes: Random forest, all_features, CLOSE (the closing half hour only, 3,000 closes)
\providecommand{\rhtaShXRfFullCl}{}\renewcommand{\rhtaShXRfFullCl}{0.94}  % ... at quoted bid/ask prices
\providecommand{\rhtaPShRfCl}{}\renewcommand{\rhtaPShRfCl}{\ensuremath{p = 0.14}}  % Ledoit-Wolf p, Sharpe ratio full minus HAR (mid): Random forest, CLOSE
\providecommand{\rhtaShRidgeHarClNat}{}\renewcommand{\rhtaShRidgeHarClNat}{0.94}  % Sharpe ratio (D19), mid-quotes: Ridge, baseline, CLOSE_NAT (closing rows, design built natively on them; linear only)
\providecommand{\rhtaShXRidgeHarClNat}{}\renewcommand{\rhtaShXRidgeHarClNat}{0.59}  % ... at quoted bid/ask prices
\providecommand{\rhtaShRidgeFullClNat}{}\renewcommand{\rhtaShRidgeFullClNat}{1.24}  % Sharpe ratio (D19), mid-quotes: Ridge, all_features, CLOSE_NAT (closing rows, design built natively on them; linear only)
\providecommand{\rhtaShXRidgeFullClNat}{}\renewcommand{\rhtaShXRidgeFullClNat}{0.89}  % ... at quoted bid/ask prices
\providecommand{\rhtaShLassoHarClNat}{}\renewcommand{\rhtaShLassoHarClNat}{0.94}  % Sharpe ratio (D19), mid-quotes: Lasso, baseline, CLOSE_NAT (closing rows, design built natively on them; linear only)
\providecommand{\rhtaShXLassoHarClNat}{}\renewcommand{\rhtaShXLassoHarClNat}{0.59}  % ... at quoted bid/ask prices
\providecommand{\rhtaShLassoFullClNat}{}\renewcommand{\rhtaShLassoFullClNat}{1.23}  % Sharpe ratio (D19), mid-quotes: Lasso, all_features, CLOSE_NAT (closing rows, design built natively on them; linear only)
\providecommand{\rhtaShXLassoFullClNat}{}\renewcommand{\rhtaShXLassoFullClNat}{0.89}  % ... at quoted bid/ask prices
\providecommand{\rhtaShEnetHarClNat}{}\renewcommand{\rhtaShEnetHarClNat}{0.94}  % Sharpe ratio (D19), mid-quotes: Elastic net, baseline, CLOSE_NAT (closing rows, design built natively on them; linear only)
\providecommand{\rhtaShXEnetHarClNat}{}\renewcommand{\rhtaShXEnetHarClNat}{0.59}  % ... at quoted bid/ask prices
\providecommand{\rhtaShEnetFullClNat}{}\renewcommand{\rhtaShEnetFullClNat}{1.71}  % Sharpe ratio (D19), mid-quotes: Elastic net, all_features, CLOSE_NAT (closing rows, design built natively on them; linear only)
\providecommand{\rhtaShXEnetFullClNat}{}\renewcommand{\rhtaShXEnetFullClNat}{1.37}  % ... at quoted bid/ask prices
\providecommand{\rhtaShRidgeHarAllCt}{}\renewcommand{\rhtaShRidgeHarAllCt}{0.95}  % Sharpe ratio (D19), mid-quotes: Ridge, baseline, ALL48@16 (ALL48, penalties chosen on the 16:00 validation rows; linear only)
\providecommand{\rhtaShXRidgeHarAllCt}{}\renewcommand{\rhtaShXRidgeHarAllCt}{0.59}  % ... at quoted bid/ask prices
\providecommand{\rhtaShRidgeFullAllCt}{}\renewcommand{\rhtaShRidgeFullAllCt}{1.40}  % Sharpe ratio (D19), mid-quotes: Ridge, all_features, ALL48@16 (ALL48, penalties chosen on the 16:00 validation rows; linear only)
\providecommand{\rhtaShXRidgeFullAllCt}{}\renewcommand{\rhtaShXRidgeFullAllCt}{1.05}  % ... at quoted bid/ask prices
\providecommand{\rhtaShLassoHarAllCt}{}\renewcommand{\rhtaShLassoHarAllCt}{0.96}  % Sharpe ratio (D19), mid-quotes: Lasso, baseline, ALL48@16 (ALL48, penalties chosen on the 16:00 validation rows; linear only)
\providecommand{\rhtaShXLassoHarAllCt}{}\renewcommand{\rhtaShXLassoHarAllCt}{0.60}  % ... at quoted bid/ask prices
\providecommand{\rhtaShLassoFullAllCt}{}\renewcommand{\rhtaShLassoFullAllCt}{1.35}  % Sharpe ratio (D19), mid-quotes: Lasso, all_features, ALL48@16 (ALL48, penalties chosen on the 16:00 validation rows; linear only)
\providecommand{\rhtaShXLassoFullAllCt}{}\renewcommand{\rhtaShXLassoFullAllCt}{1.00}  % ... at quoted bid/ask prices
\providecommand{\rhtaShEnetHarAllCt}{}\renewcommand{\rhtaShEnetHarAllCt}{0.80}  % Sharpe ratio (D19), mid-quotes: Elastic net, baseline, ALL48@16 (ALL48, penalties chosen on the 16:00 validation rows; linear only)
\providecommand{\rhtaShXEnetHarAllCt}{}\renewcommand{\rhtaShXEnetHarAllCt}{0.45}  % ... at quoted bid/ask prices
\providecommand{\rhtaShEnetFullAllCt}{}\renewcommand{\rhtaShEnetFullAllCt}{1.52}  % Sharpe ratio (D19), mid-quotes: Elastic net, all_features, ALL48@16 (ALL48, penalties chosen on the 16:00 validation rows; linear only)
\providecommand{\rhtaShXEnetFullAllCt}{}\renewcommand{\rhtaShXEnetFullAllCt}{1.17}  % ... at quoted bid/ask prices
\providecommand{\rhtaShRidgeHarRthCt}{}\renewcommand{\rhtaShRidgeHarRthCt}{0.95}  % Sharpe ratio (D19), mid-quotes: Ridge, baseline, RTH13@16 (RTH13, penalties chosen on the 16:00 validation rows; linear only)
\providecommand{\rhtaShXRidgeHarRthCt}{}\renewcommand{\rhtaShXRidgeHarRthCt}{0.60}  % ... at quoted bid/ask prices
\providecommand{\rhtaShRidgeFullRthCt}{}\renewcommand{\rhtaShRidgeFullRthCt}{1.77}  % Sharpe ratio (D19), mid-quotes: Ridge, all_features, RTH13@16 (RTH13, penalties chosen on the 16:00 validation rows; linear only)
\providecommand{\rhtaShXRidgeFullRthCt}{}\renewcommand{\rhtaShXRidgeFullRthCt}{1.41}  % ... at quoted bid/ask prices
\providecommand{\rhtaShLassoHarRthCt}{}\renewcommand{\rhtaShLassoHarRthCt}{1.01}  % Sharpe ratio (D19), mid-quotes: Lasso, baseline, RTH13@16 (RTH13, penalties chosen on the 16:00 validation rows; linear only)
\providecommand{\rhtaShXLassoHarRthCt}{}\renewcommand{\rhtaShXLassoHarRthCt}{0.65}  % ... at quoted bid/ask prices
\providecommand{\rhtaShLassoFullRthCt}{}\renewcommand{\rhtaShLassoFullRthCt}{1.88}  % Sharpe ratio (D19), mid-quotes: Lasso, all_features, RTH13@16 (RTH13, penalties chosen on the 16:00 validation rows; linear only)
\providecommand{\rhtaShXLassoFullRthCt}{}\renewcommand{\rhtaShXLassoFullRthCt}{1.52}  % ... at quoted bid/ask prices
\providecommand{\rhtaShEnetHarRthCt}{}\renewcommand{\rhtaShEnetHarRthCt}{1.01}  % Sharpe ratio (D19), mid-quotes: Elastic net, baseline, RTH13@16 (RTH13, penalties chosen on the 16:00 validation rows; linear only)
\providecommand{\rhtaShXEnetHarRthCt}{}\renewcommand{\rhtaShXEnetHarRthCt}{0.66}  % ... at quoted bid/ask prices
\providecommand{\rhtaShEnetFullRthCt}{}\renewcommand{\rhtaShEnetFullRthCt}{1.86}  % Sharpe ratio (D19), mid-quotes: Elastic net, all_features, RTH13@16 (RTH13, penalties chosen on the 16:00 validation rows; linear only)
\providecommand{\rhtaShXEnetFullRthCt}{}\renewcommand{\rhtaShXEnetFullRthCt}{1.50}  % ... at quoted bid/ask prices
\providecommand{\rhtaShRidgeHarAllThreeK}{}\renewcommand{\rhtaShRidgeHarAllThreeK}{1.56}  % Sharpe ratio (D19), mid-quotes: Ridge, baseline, ALL48_N3k (ALL48, the last 3,000 rows)
\providecommand{\rhtaShXRidgeHarAllThreeK}{}\renewcommand{\rhtaShXRidgeHarAllThreeK}{1.21}  % ... at quoted bid/ask prices
\providecommand{\rhtaShRidgeFullAllThreeK}{}\renewcommand{\rhtaShRidgeFullAllThreeK}{1.45}  % Sharpe ratio (D19), mid-quotes: Ridge, all_features, ALL48_N3k (ALL48, the last 3,000 rows)
\providecommand{\rhtaShXRidgeFullAllThreeK}{}\renewcommand{\rhtaShXRidgeFullAllThreeK}{1.10}  % ... at quoted bid/ask prices
\providecommand{\rhtaShLassoHarAllThreeK}{}\renewcommand{\rhtaShLassoHarAllThreeK}{1.57}  % Sharpe ratio (D19), mid-quotes: Lasso, baseline, ALL48_N3k (ALL48, the last 3,000 rows)
\providecommand{\rhtaShXLassoHarAllThreeK}{}\renewcommand{\rhtaShXLassoHarAllThreeK}{1.21}  % ... at quoted bid/ask prices
\providecommand{\rhtaShLassoFullAllThreeK}{}\renewcommand{\rhtaShLassoFullAllThreeK}{1.41}  % Sharpe ratio (D19), mid-quotes: Lasso, all_features, ALL48_N3k (ALL48, the last 3,000 rows)
\providecommand{\rhtaShXLassoFullAllThreeK}{}\renewcommand{\rhtaShXLassoFullAllThreeK}{1.06}  % ... at quoted bid/ask prices
\providecommand{\rhtaShEnetHarAllThreeK}{}\renewcommand{\rhtaShEnetHarAllThreeK}{1.52}  % Sharpe ratio (D19), mid-quotes: Elastic net, baseline, ALL48_N3k (ALL48, the last 3,000 rows)
\providecommand{\rhtaShXEnetHarAllThreeK}{}\renewcommand{\rhtaShXEnetHarAllThreeK}{1.16}  % ... at quoted bid/ask prices
\providecommand{\rhtaShEnetFullAllThreeK}{}\renewcommand{\rhtaShEnetFullAllThreeK}{1.63}  % Sharpe ratio (D19), mid-quotes: Elastic net, all_features, ALL48_N3k (ALL48, the last 3,000 rows)
\providecommand{\rhtaShXEnetFullAllThreeK}{}\renewcommand{\rhtaShXEnetFullAllThreeK}{1.27}  % ... at quoted bid/ask prices
\providecommand{\rhtaShLgbmHarAllThreeK}{}\renewcommand{\rhtaShLgbmHarAllThreeK}{0.92}  % Sharpe ratio (D19), mid-quotes: LightGBM, baseline, ALL48_N3k (ALL48, the last 3,000 rows)
\providecommand{\rhtaShXLgbmHarAllThreeK}{}\renewcommand{\rhtaShXLgbmHarAllThreeK}{0.57}  % ... at quoted bid/ask prices
\providecommand{\rhtaShLgbmFullAllThreeK}{}\renewcommand{\rhtaShLgbmFullAllThreeK}{1.35}  % Sharpe ratio (D19), mid-quotes: LightGBM, all_features, ALL48_N3k (ALL48, the last 3,000 rows)
\providecommand{\rhtaShXLgbmFullAllThreeK}{}\renewcommand{\rhtaShXLgbmFullAllThreeK}{0.99}  % ... at quoted bid/ask prices
\providecommand{\rhtaPShLgbmAllThreeK}{}\renewcommand{\rhtaPShLgbmAllThreeK}{\ensuremath{p = 0.20}}  % Ledoit-Wolf p, Sharpe ratio full minus HAR (mid): LightGBM, ALL48_N3k
\providecommand{\rhtaShXgbHarAllThreeK}{}\renewcommand{\rhtaShXgbHarAllThreeK}{0.93}  % Sharpe ratio (D19), mid-quotes: XGBoost, baseline, ALL48_N3k (ALL48, the last 3,000 rows)
\providecommand{\rhtaShXXgbHarAllThreeK}{}\renewcommand{\rhtaShXXgbHarAllThreeK}{0.58}  % ... at quoted bid/ask prices
\providecommand{\rhtaShXgbFullAllThreeK}{}\renewcommand{\rhtaShXgbFullAllThreeK}{1.24}  % Sharpe ratio (D19), mid-quotes: XGBoost, all_features, ALL48_N3k (ALL48, the last 3,000 rows)
\providecommand{\rhtaShXXgbFullAllThreeK}{}\renewcommand{\rhtaShXXgbFullAllThreeK}{0.88}  % ... at quoted bid/ask prices
\providecommand{\rhtaPShXgbAllThreeK}{}\renewcommand{\rhtaPShXgbAllThreeK}{\ensuremath{p = 0.41}}  % Ledoit-Wolf p, Sharpe ratio full minus HAR (mid): XGBoost, ALL48_N3k
\providecommand{\rhtaShRfHarAllThreeK}{}\renewcommand{\rhtaShRfHarAllThreeK}{0.87}  % Sharpe ratio (D19), mid-quotes: Random forest, baseline, ALL48_N3k (ALL48, the last 3,000 rows)
\providecommand{\rhtaShXRfHarAllThreeK}{}\renewcommand{\rhtaShXRfHarAllThreeK}{0.52}  % ... at quoted bid/ask prices
\providecommand{\rhtaShRfFullAllThreeK}{}\renewcommand{\rhtaShRfFullAllThreeK}{1.30}  % Sharpe ratio (D19), mid-quotes: Random forest, all_features, ALL48_N3k (ALL48, the last 3,000 rows)
\providecommand{\rhtaShXRfFullAllThreeK}{}\renewcommand{\rhtaShXRfFullAllThreeK}{0.94}  % ... at quoted bid/ask prices
\providecommand{\rhtaPShRfAllThreeK}{}\renewcommand{\rhtaPShRfAllThreeK}{\ensuremath{p = 0.21}}  % Ledoit-Wolf p, Sharpe ratio full minus HAR (mid): Random forest, ALL48_N3k
\providecommand{\rhtaShRidgeHarRthThreeK}{}\renewcommand{\rhtaShRidgeHarRthThreeK}{1.16}  % Sharpe ratio (D19), mid-quotes: Ridge, baseline, RTH13_N3k (RTH13, the last 3,000 rows)
\providecommand{\rhtaShXRidgeHarRthThreeK}{}\renewcommand{\rhtaShXRidgeHarRthThreeK}{0.80}  % ... at quoted bid/ask prices
\providecommand{\rhtaShRidgeFullRthThreeK}{}\renewcommand{\rhtaShRidgeFullRthThreeK}{1.12}  % Sharpe ratio (D19), mid-quotes: Ridge, all_features, RTH13_N3k (RTH13, the last 3,000 rows)
\providecommand{\rhtaShXRidgeFullRthThreeK}{}\renewcommand{\rhtaShXRidgeFullRthThreeK}{0.77}  % ... at quoted bid/ask prices
\providecommand{\rhtaShLassoHarRthThreeK}{}\renewcommand{\rhtaShLassoHarRthThreeK}{1.16}  % Sharpe ratio (D19), mid-quotes: Lasso, baseline, RTH13_N3k (RTH13, the last 3,000 rows)
\providecommand{\rhtaShXLassoHarRthThreeK}{}\renewcommand{\rhtaShXLassoHarRthThreeK}{0.80}  % ... at quoted bid/ask prices
\providecommand{\rhtaShLassoFullRthThreeK}{}\renewcommand{\rhtaShLassoFullRthThreeK}{1.55}  % Sharpe ratio (D19), mid-quotes: Lasso, all_features, RTH13_N3k (RTH13, the last 3,000 rows)
\providecommand{\rhtaShXLassoFullRthThreeK}{}\renewcommand{\rhtaShXLassoFullRthThreeK}{1.19}  % ... at quoted bid/ask prices
\providecommand{\rhtaShEnetHarRthThreeK}{}\renewcommand{\rhtaShEnetHarRthThreeK}{1.16}  % Sharpe ratio (D19), mid-quotes: Elastic net, baseline, RTH13_N3k (RTH13, the last 3,000 rows)
\providecommand{\rhtaShXEnetHarRthThreeK}{}\renewcommand{\rhtaShXEnetHarRthThreeK}{0.80}  % ... at quoted bid/ask prices
\providecommand{\rhtaShEnetFullRthThreeK}{}\renewcommand{\rhtaShEnetFullRthThreeK}{1.25}  % Sharpe ratio (D19), mid-quotes: Elastic net, all_features, RTH13_N3k (RTH13, the last 3,000 rows)
\providecommand{\rhtaShXEnetFullRthThreeK}{}\renewcommand{\rhtaShXEnetFullRthThreeK}{0.89}  % ... at quoted bid/ask prices
\providecommand{\rhtaShLgbmHarRthThreeK}{}\renewcommand{\rhtaShLgbmHarRthThreeK}{0.25}  % Sharpe ratio (D19), mid-quotes: LightGBM, baseline, RTH13_N3k (RTH13, the last 3,000 rows)
\providecommand{\rhtaShXLgbmHarRthThreeK}{}\renewcommand{\rhtaShXLgbmHarRthThreeK}{\ensuremath{-}0.10}  % ... at quoted bid/ask prices
\providecommand{\rhtaShLgbmFullRthThreeK}{}\renewcommand{\rhtaShLgbmFullRthThreeK}{0.83}  % Sharpe ratio (D19), mid-quotes: LightGBM, all_features, RTH13_N3k (RTH13, the last 3,000 rows)
\providecommand{\rhtaShXLgbmFullRthThreeK}{}\renewcommand{\rhtaShXLgbmFullRthThreeK}{0.47}  % ... at quoted bid/ask prices
\providecommand{\rhtaPShLgbmRthThreeK}{}\renewcommand{\rhtaPShLgbmRthThreeK}{\ensuremath{p = 0.22}}  % Ledoit-Wolf p, Sharpe ratio full minus HAR (mid): LightGBM, RTH13_N3k
\providecommand{\rhtaShXgbHarRthThreeK}{}\renewcommand{\rhtaShXgbHarRthThreeK}{0.87}  % Sharpe ratio (D19), mid-quotes: XGBoost, baseline, RTH13_N3k (RTH13, the last 3,000 rows)
\providecommand{\rhtaShXXgbHarRthThreeK}{}\renewcommand{\rhtaShXXgbHarRthThreeK}{0.52}  % ... at quoted bid/ask prices
\providecommand{\rhtaShXgbFullRthThreeK}{}\renewcommand{\rhtaShXgbFullRthThreeK}{0.86}  % Sharpe ratio (D19), mid-quotes: XGBoost, all_features, RTH13_N3k (RTH13, the last 3,000 rows)
\providecommand{\rhtaShXXgbFullRthThreeK}{}\renewcommand{\rhtaShXXgbFullRthThreeK}{0.51}  % ... at quoted bid/ask prices
\providecommand{\rhtaPShXgbRthThreeK}{}\renewcommand{\rhtaPShXgbRthThreeK}{\ensuremath{p = 0.98}}  % Ledoit-Wolf p, Sharpe ratio full minus HAR (mid): XGBoost, RTH13_N3k
\providecommand{\rhtaShRfHarRthThreeK}{}\renewcommand{\rhtaShRfHarRthThreeK}{0.67}  % Sharpe ratio (D19), mid-quotes: Random forest, baseline, RTH13_N3k (RTH13, the last 3,000 rows)
\providecommand{\rhtaShXRfHarRthThreeK}{}\renewcommand{\rhtaShXRfHarRthThreeK}{0.31}  % ... at quoted bid/ask prices
\providecommand{\rhtaShRfFullRthThreeK}{}\renewcommand{\rhtaShRfFullRthThreeK}{0.91}  % Sharpe ratio (D19), mid-quotes: Random forest, all_features, RTH13_N3k (RTH13, the last 3,000 rows)
\providecommand{\rhtaShXRfFullRthThreeK}{}\renewcommand{\rhtaShXRfFullRthThreeK}{0.56}  % ... at quoted bid/ask prices
\providecommand{\rhtaPShRfRthThreeK}{}\renewcommand{\rhtaPShRfRthThreeK}{\ensuremath{p = 0.53}}  % Ledoit-Wolf p, Sharpe ratio full minus HAR (mid): Random forest, RTH13_N3k
\providecommand{\rhtaShRidgeHarAllThirtyNineK}{}\renewcommand{\rhtaShRidgeHarAllThirtyNineK}{1.23}  % Sharpe ratio (D19), mid-quotes: Ridge, baseline, ALL48_N39k (ALL48, the last 39,000 rows)
\providecommand{\rhtaShXRidgeHarAllThirtyNineK}{}\renewcommand{\rhtaShXRidgeHarAllThirtyNineK}{0.87}  % ... at quoted bid/ask prices
\providecommand{\rhtaShRidgeFullAllThirtyNineK}{}\renewcommand{\rhtaShRidgeFullAllThirtyNineK}{0.76}  % Sharpe ratio (D19), mid-quotes: Ridge, all_features, ALL48_N39k (ALL48, the last 39,000 rows)
\providecommand{\rhtaShXRidgeFullAllThirtyNineK}{}\renewcommand{\rhtaShXRidgeFullAllThirtyNineK}{0.40}  % ... at quoted bid/ask prices
\providecommand{\rhtaShLassoHarAllThirtyNineK}{}\renewcommand{\rhtaShLassoHarAllThirtyNineK}{1.20}  % Sharpe ratio (D19), mid-quotes: Lasso, baseline, ALL48_N39k (ALL48, the last 39,000 rows)
\providecommand{\rhtaShXLassoHarAllThirtyNineK}{}\renewcommand{\rhtaShXLassoHarAllThirtyNineK}{0.84}  % ... at quoted bid/ask prices
\providecommand{\rhtaShLassoFullAllThirtyNineK}{}\renewcommand{\rhtaShLassoFullAllThirtyNineK}{1.33}  % Sharpe ratio (D19), mid-quotes: Lasso, all_features, ALL48_N39k (ALL48, the last 39,000 rows)
\providecommand{\rhtaShXLassoFullAllThirtyNineK}{}\renewcommand{\rhtaShXLassoFullAllThirtyNineK}{0.98}  % ... at quoted bid/ask prices
\providecommand{\rhtaShEnetHarAllThirtyNineK}{}\renewcommand{\rhtaShEnetHarAllThirtyNineK}{1.22}  % Sharpe ratio (D19), mid-quotes: Elastic net, baseline, ALL48_N39k (ALL48, the last 39,000 rows)
\providecommand{\rhtaShXEnetHarAllThirtyNineK}{}\renewcommand{\rhtaShXEnetHarAllThirtyNineK}{0.87}  % ... at quoted bid/ask prices
\providecommand{\rhtaShEnetFullAllThirtyNineK}{}\renewcommand{\rhtaShEnetFullAllThirtyNineK}{1.63}  % Sharpe ratio (D19), mid-quotes: Elastic net, all_features, ALL48_N39k (ALL48, the last 39,000 rows)
\providecommand{\rhtaShXEnetFullAllThirtyNineK}{}\renewcommand{\rhtaShXEnetFullAllThirtyNineK}{1.27}  % ... at quoted bid/ask prices
\providecommand{\rhtaShLgbmHarAllThirtyNineK}{}\renewcommand{\rhtaShLgbmHarAllThirtyNineK}{0.48}  % Sharpe ratio (D19), mid-quotes: LightGBM, baseline, ALL48_N39k (ALL48, the last 39,000 rows)
\providecommand{\rhtaShXLgbmHarAllThirtyNineK}{}\renewcommand{\rhtaShXLgbmHarAllThirtyNineK}{0.13}  % ... at quoted bid/ask prices
\providecommand{\rhtaShLgbmFullAllThirtyNineK}{}\renewcommand{\rhtaShLgbmFullAllThirtyNineK}{1.73}  % Sharpe ratio (D19), mid-quotes: LightGBM, all_features, ALL48_N39k (ALL48, the last 39,000 rows)
\providecommand{\rhtaShXLgbmFullAllThirtyNineK}{}\renewcommand{\rhtaShXLgbmFullAllThirtyNineK}{1.37}  % ... at quoted bid/ask prices
\providecommand{\rhtaPShLgbmAllThirtyNineK}{}\renewcommand{\rhtaPShLgbmAllThirtyNineK}{\ensuremath{p = 0.060}}  % Ledoit-Wolf p, Sharpe ratio full minus HAR (mid): LightGBM, ALL48_N39k
\providecommand{\rhtaShXgbHarAllThirtyNineK}{}\renewcommand{\rhtaShXgbHarAllThirtyNineK}{0.62}  % Sharpe ratio (D19), mid-quotes: XGBoost, baseline, ALL48_N39k (ALL48, the last 39,000 rows)
\providecommand{\rhtaShXXgbHarAllThirtyNineK}{}\renewcommand{\rhtaShXXgbHarAllThirtyNineK}{0.27}  % ... at quoted bid/ask prices
\providecommand{\rhtaShXgbFullAllThirtyNineK}{}\renewcommand{\rhtaShXgbFullAllThirtyNineK}{1.64}  % Sharpe ratio (D19), mid-quotes: XGBoost, all_features, ALL48_N39k (ALL48, the last 39,000 rows)
\providecommand{\rhtaShXXgbFullAllThirtyNineK}{}\renewcommand{\rhtaShXXgbFullAllThirtyNineK}{1.28}  % ... at quoted bid/ask prices
\providecommand{\rhtaPShXgbAllThirtyNineK}{}\renewcommand{\rhtaPShXgbAllThirtyNineK}{\ensuremath{p = 0.12}}  % Ledoit-Wolf p, Sharpe ratio full minus HAR (mid): XGBoost, ALL48_N39k
\providecommand{\rhtaShRfHarAllThirtyNineK}{}\renewcommand{\rhtaShRfHarAllThirtyNineK}{0.87}  % Sharpe ratio (D19), mid-quotes: Random forest, baseline, ALL48_N39k (ALL48, the last 39,000 rows)
\providecommand{\rhtaShXRfHarAllThirtyNineK}{}\renewcommand{\rhtaShXRfHarAllThirtyNineK}{0.50}  % ... at quoted bid/ask prices
\providecommand{\rhtaShRfFullAllThirtyNineK}{}\renewcommand{\rhtaShRfFullAllThirtyNineK}{0.95}  % Sharpe ratio (D19), mid-quotes: Random forest, all_features, ALL48_N39k (ALL48, the last 39,000 rows)
\providecommand{\rhtaShXRfFullAllThirtyNineK}{}\renewcommand{\rhtaShXRfFullAllThirtyNineK}{0.58}  % ... at quoted bid/ask prices
\providecommand{\rhtaPShRfAllThirtyNineK}{}\renewcommand{\rhtaPShRfAllThirtyNineK}{\ensuremath{p = 0.83}}  % Ledoit-Wolf p, Sharpe ratio full minus HAR (mid): Random forest, ALL48_N39k
\providecommand{\rhtaDShRidgeAllVsRth}{}\renewcommand{\rhtaDShRidgeAllVsRth}{\ensuremath{-}0.16}  % Sharpe ratio (D19, mid), ALL48 minus RTH13, full set, Ridge
\providecommand{\rhtaPDShRidgeAllVsRth}{}\renewcommand{\rhtaPDShRidgeAllVsRth}{\ensuremath{p = 0.63}}  % ... its Ledoit-Wolf p
\providecommand{\rhtaDShLassoAllVsRth}{}\renewcommand{\rhtaDShLassoAllVsRth}{\ensuremath{-}0.23}  % Sharpe ratio (D19, mid), ALL48 minus RTH13, full set, Lasso
\providecommand{\rhtaPDShLassoAllVsRth}{}\renewcommand{\rhtaPDShLassoAllVsRth}{\ensuremath{p = 0.51}}  % ... its Ledoit-Wolf p
\providecommand{\rhtaDShEnetAllVsRth}{}\renewcommand{\rhtaDShEnetAllVsRth}{\ensuremath{-}0.33}  % Sharpe ratio (D19, mid), ALL48 minus RTH13, full set, Elastic net
\providecommand{\rhtaPDShEnetAllVsRth}{}\renewcommand{\rhtaPDShEnetAllVsRth}{\ensuremath{p = 0.32}}  % ... its Ledoit-Wolf p
\providecommand{\rhtaDShLgbmAllVsRth}{}\renewcommand{\rhtaDShLgbmAllVsRth}{\ensuremath{-}0.34}  % Sharpe ratio (D19, mid), ALL48 minus RTH13, full set, LightGBM
\providecommand{\rhtaPDShLgbmAllVsRth}{}\renewcommand{\rhtaPDShLgbmAllVsRth}{\ensuremath{p = 0.17}}  % ... its Ledoit-Wolf p
\providecommand{\rhtaDShXgbAllVsRth}{}\renewcommand{\rhtaDShXgbAllVsRth}{0.03}  % Sharpe ratio (D19, mid), ALL48 minus RTH13, full set, XGBoost
\providecommand{\rhtaPDShXgbAllVsRth}{}\renewcommand{\rhtaPDShXgbAllVsRth}{\ensuremath{p = 0.90}}  % ... its Ledoit-Wolf p
\providecommand{\rhtaDShRfAllVsRth}{}\renewcommand{\rhtaDShRfAllVsRth}{\ensuremath{-}0.06}  % Sharpe ratio (D19, mid), ALL48 minus RTH13, full set, Random forest
\providecommand{\rhtaPDShRfAllVsRth}{}\renewcommand{\rhtaPDShRfAllVsRth}{\ensuremath{p = 0.85}}  % ... its Ledoit-Wolf p
\providecommand{\rhtaDShRidgeClVsRth}{}\renewcommand{\rhtaDShRidgeClVsRth}{\ensuremath{-}0.51}  % Sharpe ratio (D19, mid), CLOSE minus RTH13, full set, Ridge
\providecommand{\rhtaPDShRidgeClVsRth}{}\renewcommand{\rhtaPDShRidgeClVsRth}{\ensuremath{p = 0.22}}  % ... its Ledoit-Wolf p
\providecommand{\rhtaDShLassoClVsRth}{}\renewcommand{\rhtaDShLassoClVsRth}{\ensuremath{-}0.15}  % Sharpe ratio (D19, mid), CLOSE minus RTH13, full set, Lasso
\providecommand{\rhtaPDShLassoClVsRth}{}\renewcommand{\rhtaPDShLassoClVsRth}{\ensuremath{p = 0.72}}  % ... its Ledoit-Wolf p
\providecommand{\rhtaDShEnetClVsRth}{}\renewcommand{\rhtaDShEnetClVsRth}{\ensuremath{-}0.15}  % Sharpe ratio (D19, mid), CLOSE minus RTH13, full set, Elastic net
\providecommand{\rhtaPDShEnetClVsRth}{}\renewcommand{\rhtaPDShEnetClVsRth}{\ensuremath{p = 0.70}}  % ... its Ledoit-Wolf p
\providecommand{\rhtaDShLgbmClVsRth}{}\renewcommand{\rhtaDShLgbmClVsRth}{0.36}  % Sharpe ratio (D19, mid), CLOSE minus RTH13, full set, LightGBM
\providecommand{\rhtaPDShLgbmClVsRth}{}\renewcommand{\rhtaPDShLgbmClVsRth}{\ensuremath{p = 0.41}}  % ... its Ledoit-Wolf p
\providecommand{\rhtaDShXgbClVsRth}{}\renewcommand{\rhtaDShXgbClVsRth}{0.15}  % Sharpe ratio (D19, mid), CLOSE minus RTH13, full set, XGBoost
\providecommand{\rhtaPDShXgbClVsRth}{}\renewcommand{\rhtaPDShXgbClVsRth}{\ensuremath{p = 0.60}}  % ... its Ledoit-Wolf p
\providecommand{\rhtaDShRfClVsRth}{}\renewcommand{\rhtaDShRfClVsRth}{\ensuremath{-}0.59}  % Sharpe ratio (D19, mid), CLOSE minus RTH13, full set, Random forest
\providecommand{\rhtaPDShRfClVsRth}{}\renewcommand{\rhtaPDShRfClVsRth}{\ensuremath{p = 0.22}}  % ... its Ledoit-Wolf p
\providecommand{\rhtaWsPctRidgeWOne}{}\renewcommand{\rhtaWsPctRidgeWOne}{\ensuremath{-}4.8}  % CLOSE at W = 1,000 closes, 16:00, QLIKE % change full vs HAR, Ridge
\providecommand{\rhtaWsPRidgeWOne}{}\renewcommand{\rhtaWsPRidgeWOne}{\ensuremath{p < 0.001}}  % ... its p statement
\providecommand{\rhtaWsPctLassoWOne}{}\renewcommand{\rhtaWsPctLassoWOne}{\ensuremath{-}2.6}  % CLOSE at W = 1,000 closes, 16:00, QLIKE % change full vs HAR, Lasso
\providecommand{\rhtaWsPLassoWOne}{}\renewcommand{\rhtaWsPLassoWOne}{\ensuremath{p = 0.12}}  % ... its p statement
\providecommand{\rhtaWsPctEnetWOne}{}\renewcommand{\rhtaWsPctEnetWOne}{\ensuremath{-}1.5}  % CLOSE at W = 1,000 closes, 16:00, QLIKE % change full vs HAR, Elastic net
\providecommand{\rhtaWsPEnetWOne}{}\renewcommand{\rhtaWsPEnetWOne}{\ensuremath{p = 0.51}}  % ... its p statement
\providecommand{\rhtaWsPctLgbmWOne}{}\renewcommand{\rhtaWsPctLgbmWOne}{\ensuremath{-}7.8}  % CLOSE at W = 1,000 closes, 16:00, QLIKE % change full vs HAR, LightGBM
\providecommand{\rhtaWsPLgbmWOne}{}\renewcommand{\rhtaWsPLgbmWOne}{\ensuremath{p < 0.001}}  % ... its p statement
\providecommand{\rhtaWsPctXgbWOne}{}\renewcommand{\rhtaWsPctXgbWOne}{\ensuremath{-}6.8}  % CLOSE at W = 1,000 closes, 16:00, QLIKE % change full vs HAR, XGBoost
\providecommand{\rhtaWsPXgbWOne}{}\renewcommand{\rhtaWsPXgbWOne}{\ensuremath{p < 0.001}}  % ... its p statement
\providecommand{\rhtaWsPctRfWOne}{}\renewcommand{\rhtaWsPctRfWOne}{\ensuremath{-}2.1}  % CLOSE at W = 1,000 closes, 16:00, QLIKE % change full vs HAR, Random forest
\providecommand{\rhtaWsPRfWOne}{}\renewcommand{\rhtaWsPRfWOne}{\ensuremath{p = 0.069}}  % ... its p statement
\providecommand{\rhtaWsPctRidgeWOneHalf}{}\renewcommand{\rhtaWsPctRidgeWOneHalf}{\ensuremath{-}4.1}  % CLOSE at W = 1,500 closes, 16:00, QLIKE % change full vs HAR, Ridge
\providecommand{\rhtaWsPRidgeWOneHalf}{}\renewcommand{\rhtaWsPRidgeWOneHalf}{\ensuremath{p = 0.036}}  % ... its p statement
\providecommand{\rhtaWsPctLassoWOneHalf}{}\renewcommand{\rhtaWsPctLassoWOneHalf}{\ensuremath{-}5.0}  % CLOSE at W = 1,500 closes, 16:00, QLIKE % change full vs HAR, Lasso
\providecommand{\rhtaWsPLassoWOneHalf}{}\renewcommand{\rhtaWsPLassoWOneHalf}{\ensuremath{p < 0.001}}  % ... its p statement
\providecommand{\rhtaWsPctEnetWOneHalf}{}\renewcommand{\rhtaWsPctEnetWOneHalf}{\ensuremath{-}4.3}  % CLOSE at W = 1,500 closes, 16:00, QLIKE % change full vs HAR, Elastic net
\providecommand{\rhtaWsPEnetWOneHalf}{}\renewcommand{\rhtaWsPEnetWOneHalf}{\ensuremath{p = 0.035}}  % ... its p statement
\providecommand{\rhtaWsPctLgbmWOneHalf}{}\renewcommand{\rhtaWsPctLgbmWOneHalf}{\ensuremath{-}5.8}  % CLOSE at W = 1,500 closes, 16:00, QLIKE % change full vs HAR, LightGBM
\providecommand{\rhtaWsPLgbmWOneHalf}{}\renewcommand{\rhtaWsPLgbmWOneHalf}{\ensuremath{p = 0.001}}  % ... its p statement
\providecommand{\rhtaWsPctXgbWOneHalf}{}\renewcommand{\rhtaWsPctXgbWOneHalf}{\ensuremath{-}7.2}  % CLOSE at W = 1,500 closes, 16:00, QLIKE % change full vs HAR, XGBoost
\providecommand{\rhtaWsPXgbWOneHalf}{}\renewcommand{\rhtaWsPXgbWOneHalf}{\ensuremath{p < 0.001}}  % ... its p statement
\providecommand{\rhtaWsPctRfWOneHalf}{}\renewcommand{\rhtaWsPctRfWOneHalf}{\ensuremath{-}2.2}  % CLOSE at W = 1,500 closes, 16:00, QLIKE % change full vs HAR, Random forest
\providecommand{\rhtaWsPRfWOneHalf}{}\renewcommand{\rhtaWsPRfWOneHalf}{\ensuremath{p = 0.052}}  % ... its p statement
\providecommand{\rhtaWsPctRidgeWTwo}{}\renewcommand{\rhtaWsPctRidgeWTwo}{\ensuremath{-}3.9}  % CLOSE at W = 2,000 closes, 16:00, QLIKE % change full vs HAR, Ridge
\providecommand{\rhtaWsPRidgeWTwo}{}\renewcommand{\rhtaWsPRidgeWTwo}{\ensuremath{p = 0.049}}  % ... its p statement
\providecommand{\rhtaWsPctLassoWTwo}{}\renewcommand{\rhtaWsPctLassoWTwo}{\ensuremath{-}4.6}  % CLOSE at W = 2,000 closes, 16:00, QLIKE % change full vs HAR, Lasso
\providecommand{\rhtaWsPLassoWTwo}{}\renewcommand{\rhtaWsPLassoWTwo}{\ensuremath{p < 0.001}}  % ... its p statement
\providecommand{\rhtaWsPctEnetWTwo}{}\renewcommand{\rhtaWsPctEnetWTwo}{\ensuremath{-}4.6}  % CLOSE at W = 2,000 closes, 16:00, QLIKE % change full vs HAR, Elastic net
\providecommand{\rhtaWsPEnetWTwo}{}\renewcommand{\rhtaWsPEnetWTwo}{\ensuremath{p = 0.007}}  % ... its p statement
\providecommand{\rhtaWsPctLgbmWTwo}{}\renewcommand{\rhtaWsPctLgbmWTwo}{\ensuremath{-}7.4}  % CLOSE at W = 2,000 closes, 16:00, QLIKE % change full vs HAR, LightGBM
\providecommand{\rhtaWsPLgbmWTwo}{}\renewcommand{\rhtaWsPLgbmWTwo}{\ensuremath{p < 0.001}}  % ... its p statement
\providecommand{\rhtaWsPctXgbWTwo}{}\renewcommand{\rhtaWsPctXgbWTwo}{\ensuremath{-}7.8}  % CLOSE at W = 2,000 closes, 16:00, QLIKE % change full vs HAR, XGBoost
\providecommand{\rhtaWsPXgbWTwo}{}\renewcommand{\rhtaWsPXgbWTwo}{\ensuremath{p < 0.001}}  % ... its p statement
\providecommand{\rhtaWsPctRfWTwo}{}\renewcommand{\rhtaWsPctRfWTwo}{\ensuremath{-}3.2}  % CLOSE at W = 2,000 closes, 16:00, QLIKE % change full vs HAR, Random forest
\providecommand{\rhtaWsPRfWTwo}{}\renewcommand{\rhtaWsPRfWTwo}{\ensuremath{p = 0.002}}  % ... its p statement
\providecommand{\rhtaWsPctRidgeWTwoHalf}{}\renewcommand{\rhtaWsPctRidgeWTwoHalf}{\ensuremath{-}4.8}  % CLOSE at W = 2,500 closes, 16:00, QLIKE % change full vs HAR, Ridge
\providecommand{\rhtaWsPRidgeWTwoHalf}{}\renewcommand{\rhtaWsPRidgeWTwoHalf}{\ensuremath{p = 0.008}}  % ... its p statement
\providecommand{\rhtaWsPctLassoWTwoHalf}{}\renewcommand{\rhtaWsPctLassoWTwoHalf}{\ensuremath{-}5.3}  % CLOSE at W = 2,500 closes, 16:00, QLIKE % change full vs HAR, Lasso
\providecommand{\rhtaWsPLassoWTwoHalf}{}\renewcommand{\rhtaWsPLassoWTwoHalf}{\ensuremath{p < 0.001}}  % ... its p statement
\providecommand{\rhtaWsPctEnetWTwoHalf}{}\renewcommand{\rhtaWsPctEnetWTwoHalf}{\ensuremath{-}5.8}  % CLOSE at W = 2,500 closes, 16:00, QLIKE % change full vs HAR, Elastic net
\providecommand{\rhtaWsPEnetWTwoHalf}{}\renewcommand{\rhtaWsPEnetWTwoHalf}{\ensuremath{p < 0.001}}  % ... its p statement
\providecommand{\rhtaWsPctLgbmWTwoHalf}{}\renewcommand{\rhtaWsPctLgbmWTwoHalf}{\ensuremath{-}7.0}  % CLOSE at W = 2,500 closes, 16:00, QLIKE % change full vs HAR, LightGBM
\providecommand{\rhtaWsPLgbmWTwoHalf}{}\renewcommand{\rhtaWsPLgbmWTwoHalf}{\ensuremath{p < 0.001}}  % ... its p statement
\providecommand{\rhtaWsPctXgbWTwoHalf}{}\renewcommand{\rhtaWsPctXgbWTwoHalf}{\ensuremath{-}7.6}  % CLOSE at W = 2,500 closes, 16:00, QLIKE % change full vs HAR, XGBoost
\providecommand{\rhtaWsPXgbWTwoHalf}{}\renewcommand{\rhtaWsPXgbWTwoHalf}{\ensuremath{p < 0.001}}  % ... its p statement
\providecommand{\rhtaWsPctRfWTwoHalf}{}\renewcommand{\rhtaWsPctRfWTwoHalf}{\ensuremath{-}3.7}  % CLOSE at W = 2,500 closes, 16:00, QLIKE % change full vs HAR, Random forest
\providecommand{\rhtaWsPRfWTwoHalf}{}\renewcommand{\rhtaWsPRfWTwoHalf}{\ensuremath{p < 0.001}}  % ... its p statement
\providecommand{\rhtaWsPctRidgeWThree}{}\renewcommand{\rhtaWsPctRidgeWThree}{\ensuremath{-}6.1}  % CLOSE at W = 3,000 closes, 16:00, QLIKE % change full vs HAR, Ridge
\providecommand{\rhtaWsPRidgeWThree}{}\renewcommand{\rhtaWsPRidgeWThree}{\ensuremath{p < 0.001}}  % ... its p statement
\providecommand{\rhtaWsPctLassoWThree}{}\renewcommand{\rhtaWsPctLassoWThree}{\ensuremath{-}5.4}  % CLOSE at W = 3,000 closes, 16:00, QLIKE % change full vs HAR, Lasso
\providecommand{\rhtaWsPLassoWThree}{}\renewcommand{\rhtaWsPLassoWThree}{\ensuremath{p < 0.001}}  % ... its p statement
\providecommand{\rhtaWsPctEnetWThree}{}\renewcommand{\rhtaWsPctEnetWThree}{\ensuremath{-}6.5}  % CLOSE at W = 3,000 closes, 16:00, QLIKE % change full vs HAR, Elastic net
\providecommand{\rhtaWsPEnetWThree}{}\renewcommand{\rhtaWsPEnetWThree}{\ensuremath{p < 0.001}}  % ... its p statement
\providecommand{\rhtaWsPctLgbmWThree}{}\renewcommand{\rhtaWsPctLgbmWThree}{\ensuremath{-}7.0}  % CLOSE at W = 3,000 closes, 16:00, QLIKE % change full vs HAR, LightGBM
\providecommand{\rhtaWsPLgbmWThree}{}\renewcommand{\rhtaWsPLgbmWThree}{\ensuremath{p < 0.001}}  % ... its p statement
\providecommand{\rhtaWsPctXgbWThree}{}\renewcommand{\rhtaWsPctXgbWThree}{\ensuremath{-}8.5}  % CLOSE at W = 3,000 closes, 16:00, QLIKE % change full vs HAR, XGBoost
\providecommand{\rhtaWsPXgbWThree}{}\renewcommand{\rhtaWsPXgbWThree}{\ensuremath{p < 0.001}}  % ... its p statement
\providecommand{\rhtaWsPctRfWThree}{}\renewcommand{\rhtaWsPctRfWThree}{\ensuremath{-}3.7}  % CLOSE at W = 3,000 closes, 16:00, QLIKE % change full vs HAR, Random forest
\providecommand{\rhtaWsPRfWThree}{}\renewcommand{\rhtaWsPRfWThree}{\ensuremath{p < 0.001}}  % ... its p statement
\providecommand{\rhtaWsPctRidgeWThreeHalf}{}\renewcommand{\rhtaWsPctRidgeWThreeHalf}{\ensuremath{-}6.5}  % CLOSE at W = 3,500 closes, 16:00, QLIKE % change full vs HAR, Ridge
\providecommand{\rhtaWsPRidgeWThreeHalf}{}\renewcommand{\rhtaWsPRidgeWThreeHalf}{\ensuremath{p < 0.001}}  % ... its p statement
\providecommand{\rhtaWsPctLassoWThreeHalf}{}\renewcommand{\rhtaWsPctLassoWThreeHalf}{\ensuremath{-}5.1}  % CLOSE at W = 3,500 closes, 16:00, QLIKE % change full vs HAR, Lasso
\providecommand{\rhtaWsPLassoWThreeHalf}{}\renewcommand{\rhtaWsPLassoWThreeHalf}{\ensuremath{p = 0.011}}  % ... its p statement
\providecommand{\rhtaWsPctEnetWThreeHalf}{}\renewcommand{\rhtaWsPctEnetWThreeHalf}{\ensuremath{-}6.7}  % CLOSE at W = 3,500 closes, 16:00, QLIKE % change full vs HAR, Elastic net
\providecommand{\rhtaWsPEnetWThreeHalf}{}\renewcommand{\rhtaWsPEnetWThreeHalf}{\ensuremath{p < 0.001}}  % ... its p statement
\providecommand{\rhtaWsPctLgbmWThreeHalf}{}\renewcommand{\rhtaWsPctLgbmWThreeHalf}{\ensuremath{-}7.2}  % CLOSE at W = 3,500 closes, 16:00, QLIKE % change full vs HAR, LightGBM
\providecommand{\rhtaWsPLgbmWThreeHalf}{}\renewcommand{\rhtaWsPLgbmWThreeHalf}{\ensuremath{p < 0.001}}  % ... its p statement
\providecommand{\rhtaWsPctXgbWThreeHalf}{}\renewcommand{\rhtaWsPctXgbWThreeHalf}{\ensuremath{-}7.3}  % CLOSE at W = 3,500 closes, 16:00, QLIKE % change full vs HAR, XGBoost
\providecommand{\rhtaWsPXgbWThreeHalf}{}\renewcommand{\rhtaWsPXgbWThreeHalf}{\ensuremath{p < 0.001}}  % ... its p statement
\providecommand{\rhtaWsPctRfWThreeHalf}{}\renewcommand{\rhtaWsPctRfWThreeHalf}{\ensuremath{-}3.4}  % CLOSE at W = 3,500 closes, 16:00, QLIKE % change full vs HAR, Random forest
\providecommand{\rhtaWsPRfWThreeHalf}{}\renewcommand{\rhtaWsPRfWThreeHalf}{\ensuremath{p = 0.001}}  % ... its p statement
\providecommand{\rhtaLinPctAllCloseMin}{}\renewcommand{\rhtaLinPctAllCloseMin}{\ensuremath{-}2.6}  % 16:00, full vs HAR, ALL48, penalized linear models: lowest of the range, %
\providecommand{\rhtaLinPctAllCloseMax}{}\renewcommand{\rhtaLinPctAllCloseMax}{\ensuremath{-}1.8}  % 16:00, full vs HAR, ALL48, penalized linear models: highest of the range, %
\providecommand{\rhtaTreePctAllCloseMin}{}\renewcommand{\rhtaTreePctAllCloseMin}{\ensuremath{-}6.5}  % 16:00, full vs HAR, ALL48, tree ensembles: lowest of the range, %
\providecommand{\rhtaTreePctAllCloseMax}{}\renewcommand{\rhtaTreePctAllCloseMax}{\ensuremath{-}5.6}  % 16:00, full vs HAR, ALL48, tree ensembles: highest of the range, %
\providecommand{\rhtaPctAllCloseMin}{}\renewcommand{\rhtaPctAllCloseMin}{\ensuremath{-}6.5}  % 16:00, full vs HAR, ALL48, all six classes: lowest of the range, %
\providecommand{\rhtaPctAllCloseMax}{}\renewcommand{\rhtaPctAllCloseMax}{\ensuremath{-}1.8}  % 16:00, full vs HAR, ALL48, all six classes: highest of the range, %
\providecommand{\rhtaLinPctClCloseMin}{}\renewcommand{\rhtaLinPctClCloseMin}{\ensuremath{-}6.5}  % 16:00, full vs HAR, CLOSE, penalized linear models: lowest of the range, %
\providecommand{\rhtaLinPctClCloseMax}{}\renewcommand{\rhtaLinPctClCloseMax}{\ensuremath{-}5.4}  % 16:00, full vs HAR, CLOSE, penalized linear models: highest of the range, %
\providecommand{\rhtaTreePctClCloseMin}{}\renewcommand{\rhtaTreePctClCloseMin}{\ensuremath{-}8.5}  % 16:00, full vs HAR, CLOSE, tree ensembles: lowest of the range, %
\providecommand{\rhtaTreePctClCloseMax}{}\renewcommand{\rhtaTreePctClCloseMax}{\ensuremath{-}3.7}  % 16:00, full vs HAR, CLOSE, tree ensembles: highest of the range, %
\providecommand{\rhtaPctClCloseMin}{}\renewcommand{\rhtaPctClCloseMin}{\ensuremath{-}8.5}  % 16:00, full vs HAR, CLOSE, all six classes: lowest of the range, %
\providecommand{\rhtaPctClCloseMax}{}\renewcommand{\rhtaPctClCloseMax}{\ensuremath{-}3.7}  % 16:00, full vs HAR, CLOSE, all six classes: highest of the range, %
\providecommand{\rhtaLinPctRthCloseMin}{}\renewcommand{\rhtaLinPctRthCloseMin}{\ensuremath{-}0.5}  % 16:00, full vs HAR, RTH13, penalized linear models: lowest of the range, %
\providecommand{\rhtaLinPctRthCloseMax}{}\renewcommand{\rhtaLinPctRthCloseMax}{3.3}  % 16:00, full vs HAR, RTH13, penalized linear models: highest of the range, %
\providecommand{\rhtaTreePctRthCloseMin}{}\renewcommand{\rhtaTreePctRthCloseMin}{\ensuremath{-}3.4}  % 16:00, full vs HAR, RTH13, tree ensembles: lowest of the range, %
\providecommand{\rhtaTreePctRthCloseMax}{}\renewcommand{\rhtaTreePctRthCloseMax}{\ensuremath{-}1.4}  % 16:00, full vs HAR, RTH13, tree ensembles: highest of the range, %
\providecommand{\rhtaPctRthCloseMin}{}\renewcommand{\rhtaPctRthCloseMin}{\ensuremath{-}3.4}  % 16:00, full vs HAR, RTH13, all six classes: lowest of the range, %
\providecommand{\rhtaPctRthCloseMax}{}\renewcommand{\rhtaPctRthCloseMax}{3.3}  % 16:00, full vs HAR, RTH13, all six classes: highest of the range, %
\providecommand{\rhtaLinPctClNatCloseMin}{}\renewcommand{\rhtaLinPctClNatCloseMin}{\ensuremath{-}6.6}  % 16:00, full vs HAR, CLOSE_NAT, penalized linear models: lowest of the range, %
\providecommand{\rhtaLinPctClNatCloseMax}{}\renewcommand{\rhtaLinPctClNatCloseMax}{\ensuremath{-}5.4}  % 16:00, full vs HAR, CLOSE_NAT, penalized linear models: highest of the range, %
\providecommand{\rhtaLinAllClosePMin}{}\renewcommand{\rhtaLinAllClosePMin}{\ensuremath{p = 0.15}}  % 16:00, ALL48, linear: smallest p
\providecommand{\rhtaLinAllCtClosePMin}{}\renewcommand{\rhtaLinAllCtClosePMin}{\ensuremath{p = 0.33}}  % 16:00, ALL48@16, linear: smallest p
\providecommand{\rhtaLinRthCtClosePMin}{}\renewcommand{\rhtaLinRthCtClosePMin}{\ensuremath{p = 0.38}}  % 16:00, RTH13@16, linear: smallest p
\providecommand{\rhtaLinPctRthBarMin}{}\renewcommand{\rhtaLinPctRthBarMin}{\ensuremath{-}6.2}  % every bar, full vs HAR, RTH13, penalized linear models: lowest of the range, %
\providecommand{\rhtaLinPctRthBarMax}{}\renewcommand{\rhtaLinPctRthBarMax}{\ensuremath{-}5.4}  % every bar, full vs HAR, RTH13, penalized linear models: highest of the range, %
\providecommand{\rhtaTreePctRthBarMin}{}\renewcommand{\rhtaTreePctRthBarMin}{\ensuremath{-}14.0}  % every bar, full vs HAR, RTH13, tree ensembles: lowest of the range, %
\providecommand{\rhtaTreePctRthBarMax}{}\renewcommand{\rhtaTreePctRthBarMax}{\ensuremath{-}11.1}  % every bar, full vs HAR, RTH13, tree ensembles: highest of the range, %
\providecommand{\rhtaLinPctAllBarMin}{}\renewcommand{\rhtaLinPctAllBarMin}{\ensuremath{-}5.2}  % every bar, full vs HAR, ALL48, penalized linear models: lowest of the range, %
\providecommand{\rhtaLinPctAllBarMax}{}\renewcommand{\rhtaLinPctAllBarMax}{\ensuremath{-}4.9}  % every bar, full vs HAR, ALL48, penalized linear models: highest of the range, %
\providecommand{\rhtaTreePctAllBarMin}{}\renewcommand{\rhtaTreePctAllBarMin}{\ensuremath{-}12.0}  % every bar, full vs HAR, ALL48, tree ensembles: lowest of the range, %
\providecommand{\rhtaTreePctAllBarMax}{}\renewcommand{\rhtaTreePctAllBarMax}{\ensuremath{-}7.0}  % every bar, full vs HAR, ALL48, tree ensembles: highest of the range, %
\providecommand{\rhtaAllVsRthClosePctMin}{}\renewcommand{\rhtaAllVsRthClosePctMin}{\ensuremath{-}6.0}  % 16:00, ALL48 vs RTH13, full set, all six classes: lowest of the range, %
\providecommand{\rhtaAllVsRthClosePctMax}{}\renewcommand{\rhtaAllVsRthClosePctMax}{\ensuremath{-}1.9}  % 16:00, ALL48 vs RTH13, full set, all six classes: highest of the range, %
\providecommand{\rhtaNAllVsRthCloseSig}{}\renewcommand{\rhtaNAllVsRthCloseSig}{five}  % 16:00, ALL48 vs RTH13: classes with p < 0.05 (lower)
\providecommand{\rhtaAllVsRthBarPctMin}{}\renewcommand{\rhtaAllVsRthBarPctMin}{1.9}  % every bar, ALL48 vs RTH13, full set, all six classes (RTH13 more accurate): lowest of the range, %
\providecommand{\rhtaAllVsRthBarPctMax}{}\renewcommand{\rhtaAllVsRthBarPctMax}{5.2}  % every bar, ALL48 vs RTH13, full set, all six classes (RTH13 more accurate): highest of the range, %
\providecommand{\rhtaAllVsRthBarPMax}{}\renewcommand{\rhtaAllVsRthBarPMax}{\ensuremath{p = 0.008}}  % every bar, ALL48 vs RTH13: largest p of the six
\providecommand{\rhtaLinClVsRthClosePctMin}{}\renewcommand{\rhtaLinClVsRthClosePctMin}{\ensuremath{-}11.8}  % 16:00, CLOSE vs RTH13, full set, penalized linear models: lowest of the range, %
\providecommand{\rhtaLinClVsRthClosePctMax}{}\renewcommand{\rhtaLinClVsRthClosePctMax}{\ensuremath{-}9.4}  % 16:00, CLOSE vs RTH13, full set, penalized linear models: highest of the range, %
\providecommand{\rhtaTreeClVsRthClosePctMin}{}\renewcommand{\rhtaTreeClVsRthClosePctMin}{\ensuremath{-}2.2}  % 16:00, CLOSE vs RTH13, full set, tree ensembles: lowest of the range, %
\providecommand{\rhtaTreeClVsRthClosePctMax}{}\renewcommand{\rhtaTreeClVsRthClosePctMax}{1.6}  % 16:00, CLOSE vs RTH13, full set, tree ensembles: highest of the range, %
\providecommand{\rhtaTreeClVsRthClosePMin}{}\renewcommand{\rhtaTreeClVsRthClosePMin}{\ensuremath{p = 0.42}}  % ... tree ensembles: smallest p
\providecommand{\rhtaTreeClVsRthClosePMax}{}\renewcommand{\rhtaTreeClVsRthClosePMax}{\ensuremath{p = 0.69}}  % ... tree ensembles: largest p
\providecommand{\rhtaAllThirtyNineKVsRthClosePctMin}{}\renewcommand{\rhtaAllThirtyNineKVsRthClosePctMin}{\ensuremath{-}3.6}  % 16:00, ALL48_N39k vs RTH13, full set, all six classes: lowest of the range, %
\providecommand{\rhtaAllThirtyNineKVsRthClosePctMax}{}\renewcommand{\rhtaAllThirtyNineKVsRthClosePctMax}{1.3}  % 16:00, ALL48_N39k vs RTH13, full set, all six classes: highest of the range, %
\providecommand{\rhtaNAllThirtyNineKVsRthCloseSig}{}\renewcommand{\rhtaNAllThirtyNineKVsRthCloseSig}{two}  % ... classes with p < 0.05
\providecommand{\rhtaAllThirtyNineKVsRthBarPctMin}{}\renewcommand{\rhtaAllThirtyNineKVsRthBarPctMin}{2.7}  % every bar, ALL48_N39k vs RTH13, full set, all six classes: lowest of the range, %
\providecommand{\rhtaAllThirtyNineKVsRthBarPctMax}{}\renewcommand{\rhtaAllThirtyNineKVsRthBarPctMax}{10.4}  % every bar, ALL48_N39k vs RTH13, full set, all six classes: highest of the range, %
\providecommand{\rhtaAllThirtyNineKVsAllClosePctMin}{}\renewcommand{\rhtaAllThirtyNineKVsAllClosePctMin}{1.7}  % 16:00, ALL48_N39k vs ALL48 (39,000 vs 144,000 rows), full set: lowest of the range, %
\providecommand{\rhtaAllThirtyNineKVsAllClosePctMax}{}\renewcommand{\rhtaAllThirtyNineKVsAllClosePctMax}{4.1}  % 16:00, ALL48_N39k vs ALL48 (39,000 vs 144,000 rows), full set: highest of the range, %
\providecommand{\rhtaNAllThirtyNineKVsAllCloseSig}{}\renewcommand{\rhtaNAllThirtyNineKVsAllCloseSig}{five}  % ... classes with p < 0.05 (higher loss)
\providecommand{\rhtaAllThreeKVsRthThreeKClosePctMin}{}\renewcommand{\rhtaAllThreeKVsRthThreeKClosePctMin}{\ensuremath{-}8.1}  % 16:00, ALL48_N3k vs RTH13_N3k, full set, all six classes: lowest of the range, %
\providecommand{\rhtaAllThreeKVsRthThreeKClosePctMax}{}\renewcommand{\rhtaAllThreeKVsRthThreeKClosePctMax}{\ensuremath{-}3.3}  % 16:00, ALL48_N3k vs RTH13_N3k, full set, all six classes: highest of the range, %
\providecommand{\rhtaNAllThreeKVsRthThreeKCloseSig}{}\renewcommand{\rhtaNAllThreeKVsRthThreeKCloseSig}{five}  % ... classes with p < 0.05 (lower)
\providecommand{\rhtaAllThreeKVsRthThreeKBarPctMin}{}\renewcommand{\rhtaAllThreeKVsRthThreeKBarPctMin}{2.7}  % every bar, ALL48_N3k vs RTH13_N3k, full set, all six classes: lowest of the range, %
\providecommand{\rhtaAllThreeKVsRthThreeKBarPctMax}{}\renewcommand{\rhtaAllThreeKVsRthThreeKBarPctMax}{5.8}  % every bar, ALL48_N3k vs RTH13_N3k, full set, all six classes: highest of the range, %
\providecommand{\rhtaClVsRthThreeKClosePctMin}{}\renewcommand{\rhtaClVsRthThreeKClosePctMin}{\ensuremath{-}18.0}  % 16:00, CLOSE vs RTH13_N3k (3,000 rows each), full set: lowest of the range, %
\providecommand{\rhtaClVsRthThreeKClosePctMax}{}\renewcommand{\rhtaClVsRthThreeKClosePctMax}{\ensuremath{-}8.7}  % 16:00, CLOSE vs RTH13_N3k (3,000 rows each), full set: highest of the range, %
\providecommand{\rhtaClVsRthThreeKClosePMax}{}\renewcommand{\rhtaClVsRthThreeKClosePMax}{\ensuremath{p = 0.004}}  % ... largest p of the six
\providecommand{\rhtaRthThreeKVsRthClosePctMin}{}\renewcommand{\rhtaRthThreeKVsRthClosePctMin}{6.2}  % 16:00, RTH13_N3k vs RTH13 (3,000 vs 39,000 rows), full set: lowest of the range, %
\providecommand{\rhtaRthThreeKVsRthClosePctMax}{}\renewcommand{\rhtaRthThreeKVsRthClosePctMax}{12.3}  % 16:00, RTH13_N3k vs RTH13 (3,000 vs 39,000 rows), full set: highest of the range, %
\providecommand{\rhtaRthThreeKVsRthBarPctMin}{}\renewcommand{\rhtaRthThreeKVsRthBarPctMin}{2.7}  % every bar, RTH13_N3k vs RTH13 (3,000 vs 39,000 rows), full set: lowest of the range, %
\providecommand{\rhtaRthThreeKVsRthBarPctMax}{}\renewcommand{\rhtaRthThreeKVsRthBarPctMax}{14.9}  % every bar, RTH13_N3k vs RTH13 (3,000 vs 39,000 rows), full set: highest of the range, %
\providecommand{\rhtaShFullRthMin}{}\renewcommand{\rhtaShFullRthMin}{1.47}  % Sharpe ratio (D19, mid), full set, RTH13: lowest of the six
\providecommand{\rhtaShFullRthMax}{}\renewcommand{\rhtaShFullRthMax}{1.88}  % ... highest of the six
\providecommand{\rhtaNTreeShSigRth}{}\renewcommand{\rhtaNTreeShSigRth}{two}  % tree classes whose full-vs-HAR Sharpe ratio difference has Ledoit-Wolf p < 0.05, RTH13
\providecommand{\rhtaShFullAllMin}{}\renewcommand{\rhtaShFullAllMin}{1.25}  % Sharpe ratio (D19, mid), full set, ALL48: lowest of the six
\providecommand{\rhtaShFullAllMax}{}\renewcommand{\rhtaShFullAllMax}{1.82}  % ... highest of the six
\providecommand{\rhtaNTreeShSigAll}{}\renewcommand{\rhtaNTreeShSigAll}{none}  % tree classes whose full-vs-HAR Sharpe ratio difference has Ledoit-Wolf p < 0.05, ALL48
\providecommand{\rhtaShFullClMin}{}\renewcommand{\rhtaShFullClMin}{1.26}  % Sharpe ratio (D19, mid), full set, CLOSE: lowest of the six
\providecommand{\rhtaShFullClMax}{}\renewcommand{\rhtaShFullClMax}{1.94}  % ... highest of the six
\providecommand{\rhtaNTreeShSigCl}{}\renewcommand{\rhtaNTreeShSigCl}{one}  % tree classes whose full-vs-HAR Sharpe ratio difference has Ledoit-Wolf p < 0.05, CLOSE
\providecommand{\rhtaShTreeHarClMin}{}\renewcommand{\rhtaShTreeHarClMin}{\ensuremath{-}0.03}  % Sharpe ratio (D19, mid), HAR benchmark, CLOSE, tree ensembles: lowest
\providecommand{\rhtaShTreeHarClMax}{}\renewcommand{\rhtaShTreeHarClMax}{0.40}  % ... highest
\providecommand{\rhtaNDShTests}{}\renewcommand{\rhtaNDShTests}{12}  % Sharpe ratio tests ALL48 - RTH13 and CLOSE - RTH13, full set (six classes each)
\providecommand{\rhtaNDShSig}{}\renewcommand{\rhtaNDShSig}{none}  % ... of them with Ledoit-Wolf p < 0.05
\providecommand{\rhtaDShPMin}{}\renewcommand{\rhtaDShPMin}{\ensuremath{p = 0.17}}  % ... the smallest Ledoit-Wolf p
\providecommand{\rhtaWsNCells}{}\renewcommand{\rhtaWsNCells}{36}  % window x class cells of the close-only sweep (16:00, full vs HAR)
\providecommand{\rhtaWsNLower}{}\renewcommand{\rhtaWsNLower}{36}  % ... with the lower QLIKE on the full set
\providecommand{\rhtaWsNSig}{}\renewcommand{\rhtaWsNSig}{32}  % ... with p < 0.05
\providecommand{\rhtaWsSigFrom}{}\renewcommand{\rhtaWsSigFrom}{2{,}000}  % smallest W from which every cell has p < 0.05
\providecommand{\rhtaWsLinPctMin}{}\renewcommand{\rhtaWsLinPctMin}{\ensuremath{-}6.7}  % close-only sweep, 16:00, full vs HAR, linear, all W: lowest of the range, %
\providecommand{\rhtaWsLinPctMax}{}\renewcommand{\rhtaWsLinPctMax}{\ensuremath{-}1.5}  % close-only sweep, 16:00, full vs HAR, linear, all W: highest of the range, %
\providecommand{\rhtaWsTreePctMin}{}\renewcommand{\rhtaWsTreePctMin}{\ensuremath{-}8.5}  % close-only sweep, 16:00, full vs HAR, trees, all W: lowest of the range, %
\providecommand{\rhtaWsTreePctMax}{}\renewcommand{\rhtaWsTreePctMax}{\ensuremath{-}2.1}  % close-only sweep, 16:00, full vs HAR, trees, all W: highest of the range, %
\fi
\begin{table}[p]
\centering
\footnotesize
\setlength{\tabcolsep}{1.5pt}
\begin{threeparttable}
\caption{Robustness: the training sample}
\label{tab:rh_train_arms}
\begin{tabular*}{\linewidth}{@{\extracolsep{\fill}}lrrrrrr@{}}
\toprule
 & \multicolumn{3}{c}{Penalized linear} & \multicolumn{3}{c}{Tree ensembles} \\
\cmidrule(lr){2-4}\cmidrule(l){5-7}
 & Ridge & Lasso & Elastic net & LightGBM & XGBoost & Random forest \\
\midrule
\multicolumn{7}{@{}l}{\textit{Panel A: full predictor set vs HAR benchmark, 16:00 bar}} \\
\quad RTH13 (main) & $0.8$ {\scriptsize (0.759)} & $3.3$ {\scriptsize (0.300)} & $-0.5$ {\scriptsize (0.832)} & $-2.7$ {\scriptsize (0.071)} & $-3.4$ {\scriptsize (0.030)} & $-1.4$ {\scriptsize (0.357)} \\
\quad ALL48 & $-2.2$ {\scriptsize (0.238)} & $-2.6$ {\scriptsize (0.147)} & $-1.8$ {\scriptsize (0.337)} & $-6.1$ {\scriptsize (\ensuremath{<}0.001)} & $-6.5$ {\scriptsize (\ensuremath{<}0.001)} & $-5.6$ {\scriptsize (\ensuremath{<}0.001)} \\
\quad CLOSE & $-6.1$ {\scriptsize (\ensuremath{<}0.001)} & $-5.4$ {\scriptsize (\ensuremath{<}0.001)} & $-6.5$ {\scriptsize (\ensuremath{<}0.001)} & $-7.0$ {\scriptsize (\ensuremath{<}0.001)} & $-8.5$ {\scriptsize (\ensuremath{<}0.001)} & $-3.7$ {\scriptsize (\ensuremath{<}0.001)} \\
\addlinespace[2pt]
\quad CLOSE-nat & $-6.5$ {\scriptsize (\ensuremath{<}0.001)} & $-5.4$ {\scriptsize (\ensuremath{<}0.001)} & $-6.6$ {\scriptsize (\ensuremath{<}0.001)} & -- & -- & -- \\
\quad ALL48@16 & $-1.9$ {\scriptsize (0.333)} & $-1.1$ {\scriptsize (0.594)} & $-0.8$ {\scriptsize (0.691)} & -- & -- & -- \\
\quad RTH13@16 & $0.8$ {\scriptsize (0.759)} & $-1.8$ {\scriptsize (0.385)} & $1.5$ {\scriptsize (0.556)} & -- & -- & -- \\
\addlinespace[2pt]
\quad ALL48-3k & $-4.7$ {\scriptsize (\ensuremath{<}0.001)} & $-1.6$ {\scriptsize (0.212)} & $0.5$ {\scriptsize (0.784)} & $-8.0$ {\scriptsize (\ensuremath{<}0.001)} & $-6.6$ {\scriptsize (\ensuremath{<}0.001)} & $-6.9$ {\scriptsize (\ensuremath{<}0.001)} \\
\quad RTH13-3k & $-1.9$ {\scriptsize (0.121)} & $-1.9$ {\scriptsize (0.196)} & $-1.6$ {\scriptsize (0.393)} & $-3.7$ {\scriptsize (0.031)} & $-2.3$ {\scriptsize (0.206)} & $-4.5$ {\scriptsize (\ensuremath{<}0.001)} \\
\quad ALL48-39k & $-0.9$ {\scriptsize (0.679)} & $-3.6$ {\scriptsize (0.028)} & $-3.7$ {\scriptsize (0.001)} & $-5.5$ {\scriptsize (\ensuremath{<}0.001)} & $-4.7$ {\scriptsize (\ensuremath{<}0.001)} & $-4.9$ {\scriptsize (\ensuremath{<}0.001)} \\
\addlinespace[4pt]
\multicolumn{7}{@{}l}{\textit{Panel B: full predictor set vs HAR benchmark, every bar}} \\
\quad RTH13 (main) & $-5.4$ {\scriptsize (\ensuremath{<}0.001)} & $-6.2$ {\scriptsize (\ensuremath{<}0.001)} & $-6.0$ {\scriptsize (\ensuremath{<}0.001)} & $-14.0$ {\scriptsize (\ensuremath{<}0.001)} & $-13.8$ {\scriptsize (\ensuremath{<}0.001)} & $-11.1$ {\scriptsize (\ensuremath{<}0.001)} \\
\quad ALL48 & $-4.9$ {\scriptsize (\ensuremath{<}0.001)} & $-5.1$ {\scriptsize (\ensuremath{<}0.001)} & $-5.2$ {\scriptsize (\ensuremath{<}0.001)} & $-11.2$ {\scriptsize (\ensuremath{<}0.001)} & $-12.0$ {\scriptsize (\ensuremath{<}0.001)} & $-7.0$ {\scriptsize (\ensuremath{<}0.001)} \\
\addlinespace[2pt]
\quad ALL48@16 & $-5.2$ {\scriptsize (\ensuremath{<}0.001)} & $-5.3$ {\scriptsize (\ensuremath{<}0.001)} & $-5.5$ {\scriptsize (\ensuremath{<}0.001)} & -- & -- & -- \\
\quad RTH13@16 & $-5.4$ {\scriptsize (\ensuremath{<}0.001)} & $-6.2$ {\scriptsize (\ensuremath{<}0.001)} & $-6.9$ {\scriptsize (\ensuremath{<}0.001)} & -- & -- & -- \\
\addlinespace[2pt]
\quad ALL48-3k & $-2.7$ {\scriptsize (\ensuremath{<}0.001)} & $-2.1$ {\scriptsize (\ensuremath{<}0.001)} & $-1.8$ {\scriptsize (0.010)} & $-5.9$ {\scriptsize (\ensuremath{<}0.001)} & $-5.9$ {\scriptsize (\ensuremath{<}0.001)} & $-5.0$ {\scriptsize (\ensuremath{<}0.001)} \\
\quad RTH13-3k & $-3.6$ {\scriptsize (\ensuremath{<}0.001)} & $-3.8$ {\scriptsize (\ensuremath{<}0.001)} & $-4.2$ {\scriptsize (\ensuremath{<}0.001)} & $-6.0$ {\scriptsize (\ensuremath{<}0.001)} & $-6.8$ {\scriptsize (\ensuremath{<}0.001)} & $-4.7$ {\scriptsize (\ensuremath{<}0.001)} \\
\quad ALL48-39k & $-4.4$ {\scriptsize (\ensuremath{<}0.001)} & $-3.8$ {\scriptsize (\ensuremath{<}0.001)} & $-3.0$ {\scriptsize (\ensuremath{<}0.001)} & $-7.2$ {\scriptsize (\ensuremath{<}0.001)} & $-10.1$ {\scriptsize (\ensuremath{<}0.001)} & $-5.6$ {\scriptsize (\ensuremath{<}0.001)} \\
\addlinespace[4pt]
\multicolumn{7}{@{}l}{\textit{Panel C: across training samples, full predictor set}} \\
\multicolumn{7}{@{}l}{16:00 bar} \\
\quad ALL48 vs RTH13 & $-3.5$ {\scriptsize (0.006)} & $-6.0$ {\scriptsize (\ensuremath{<}0.001)} & $-2.1$ {\scriptsize (0.101)} & $-2.3$ {\scriptsize (0.030)} & $-1.9$ {\scriptsize (0.027)} & $-4.3$ {\scriptsize (\ensuremath{<}0.001)} \\
\quad CLOSE vs RTH13 & $-10.2$ {\scriptsize (\ensuremath{<}0.001)} & $-11.8$ {\scriptsize (\ensuremath{<}0.001)} & $-9.4$ {\scriptsize (\ensuremath{<}0.001)} & $1.6$ {\scriptsize (0.650)} & $-1.1$ {\scriptsize (0.688)} & $-2.2$ {\scriptsize (0.419)} \\
\quad CLOSE vs ALL48 & $-7.0$ {\scriptsize (\ensuremath{<}0.001)} & $-6.2$ {\scriptsize (0.002)} & $-7.5$ {\scriptsize (\ensuremath{<}0.001)} & $4.0$ {\scriptsize (0.349)} & $0.8$ {\scriptsize (0.778)} & $2.2$ {\scriptsize (0.458)} \\
\quad CLOSE-nat vs RTH13 & $-10.4$ {\scriptsize (\ensuremath{<}0.001)} & $-11.5$ {\scriptsize (\ensuremath{<}0.001)} & $-9.3$ {\scriptsize (\ensuremath{<}0.001)} & -- & -- & -- \\
\quad ALL48@16 vs RTH13 & $-3.2$ {\scriptsize (0.019)} & $-5.4$ {\scriptsize (0.002)} & $-2.1$ {\scriptsize (0.101)} & -- & -- & -- \\
\quad RTH13@16 vs RTH13 & $0.0$ {\scriptsize (1.000)} & $-5.4$ {\scriptsize (\ensuremath{<}0.001)} & $0.0$ {\scriptsize (1.000)} & -- & -- & -- \\
\multicolumn{7}{@{}l}{Every bar} \\
\quad ALL48 vs RTH13 & $1.9$ {\scriptsize (\ensuremath{<}0.001)} & $2.5$ {\scriptsize (0.008)} & $2.2$ {\scriptsize (\ensuremath{<}0.001)} & $4.5$ {\scriptsize (\ensuremath{<}0.001)} & $3.2$ {\scriptsize (\ensuremath{<}0.001)} & $5.2$ {\scriptsize (\ensuremath{<}0.001)} \\
\quad ALL48@16 vs RTH13 & $1.6$ {\scriptsize (0.001)} & $2.2$ {\scriptsize (0.008)} & $2.2$ {\scriptsize (\ensuremath{<}0.001)} & -- & -- & -- \\
\quad RTH13@16 vs RTH13 & $0.0$ {\scriptsize (1.000)} & $0.2$ {\scriptsize (0.851)} & $0.0$ {\scriptsize (1.000)} & -- & -- & -- \\
\bottomrule
\end{tabular*}
\begin{tablenotes}[flushleft]
\footnotesize
\item \textit{Notes:} The pooled models of Table~\ref{tab:rh_oos} trained on other samples, everything else as in the main design: predictors on the full grid of \rhtaNIntervals{} half-hour intervals, overnight included; a refit at every training row; the tree ensembles' frozen configurations; the linear penalties chosen in each sample by the main design's rule on that sample's pre-sample rows (validation year \rhtaValYear). RTH13, the main design: the \rhtaNBarsDay{} regular-hours bars of each day (bar ends 10:00 to 16:00~ET), a rolling window of \rhtaWinDays{} trading days (\rhtaRowsRth{} rows). ALL48: all \rhtaNIntervals{} intervals of the futures' day, overnight, weekend and holiday intervals and early-close sessions included (with the overnight indicator), \rhtaWinDays{} days (\rhtaRowsAll{} rows). CLOSE: the 16:00 bar only (the 15:30 forecast of the closing half hour), \rhtaRowsCl{} closes. Penalized linear models only: CLOSE-nat builds the closing-bar design on the closing rows (its own scaling); ALL48@16 and RTH13@16 choose the penalties on the 16:00 validation rows (the primary sample's forecast where the choice is unchanged). Equal training rows: ALL48-3k and RTH13-3k train on the last \rhtaRowsSmall{} rows of their sample (about \rhtaSpanAllSmall{} and \rhtaSpanRthSmall{} days; CLOSE's \rhtaRowsCl{} closes span about \rhtaSpanClSmall{} years), ALL48-39k on the last \rhtaRowsLarge{} (about \rhtaSpanAllLarge{} days; RTH13 has as many rows). Every sample is scored on the same \rhtaNdays{} trading days, \rhtaStart{} to \rhtaEnd: the 16:00 bar and every regular-hours bar (\rhtaNbars{} bars). Entries: the percentage change in mean QLIKE of the full predictor set relative to the HAR benchmark of the same model and sample (Panels A and B), or of the first-named sample relative to the second, full predictor set (Panel C); negative: the full set, or the first sample, has the lower loss. In parentheses, the $p$-value of the Diebold--Mariano test (Harvey--Leybourne--Newbold correction, Newey--West lag \rhtaHacLag) on the daily 16:00 loss or the daily mean over the \rhtaNBarsDay{} bars. The RTH13 rows are the forecasts of Tables~\ref{tab:rh_oos} (every bar) and~\ref{tab:rh_robust_w3500} (16:00, $W = \rhtaWinDays$). --: no forecast; the tree ensembles have no CLOSE-nat, ALL48@16 or RTH13@16 arm (their configurations are not re-tuned).
\end{tablenotes}
\end{threeparttable}
\end{table}

\begin{table}[p]
\centering
\footnotesize
\setlength{\tabcolsep}{1.5pt}
\begin{threeparttable}
\ContinuedFloat  % caption package: the same table number (inside threeparttable, before \caption)
\caption{Robustness: the training sample (continued)}
\label{tab:rh_train_arms_cont}
\begin{tabular*}{\linewidth}{@{\extracolsep{\fill}}llrrrrrr@{}}
\toprule
 &  & \multicolumn{3}{c}{Penalized linear} & \multicolumn{3}{c}{Tree ensembles} \\
\cmidrule(lr){3-5}\cmidrule(l){6-8}
 &  & Ridge & Lasso & Elastic net & LightGBM & XGBoost & Random forest \\
\midrule
\multicolumn{8}{@{}l}{\textit{Panel D: equal training rows, first sample vs second, full predictor set}} \\
\multicolumn{8}{@{}l}{16:00 bar} \\
ALL48-3k & RTH13-3k & $-8.1$ {\scriptsize (\ensuremath{<}0.001)} & $-6.5$ {\scriptsize (\ensuremath{<}0.001)} & $-4.2$ {\scriptsize (0.033)} & $-3.9$ {\scriptsize (0.038)} & $-3.3$ {\scriptsize (0.072)} & $-5.7$ {\scriptsize (\ensuremath{<}0.001)} \\
CLOSE & RTH13-3k & $-17.4$ {\scriptsize (\ensuremath{<}0.001)} & $-16.9$ {\scriptsize (\ensuremath{<}0.001)} & $-18.0$ {\scriptsize (\ensuremath{<}0.001)} & $-9.5$ {\scriptsize (0.004)} & $-9.5$ {\scriptsize (\ensuremath{<}0.001)} & $-8.7$ {\scriptsize (0.002)} \\
CLOSE & ALL48-3k & $-10.2$ {\scriptsize (\ensuremath{<}0.001)} & $-11.1$ {\scriptsize (\ensuremath{<}0.001)} & $-14.5$ {\scriptsize (\ensuremath{<}0.001)} & $-5.8$ {\scriptsize (0.093)} & $-6.4$ {\scriptsize (0.026)} & $-3.2$ {\scriptsize (0.272)} \\
ALL48-39k & RTH13 & $0.5$ {\scriptsize (0.731)} & $-3.6$ {\scriptsize (0.038)} & $0.1$ {\scriptsize (0.946)} & $0.4$ {\scriptsize (0.721)} & $1.3$ {\scriptsize (0.148)} & $-2.7$ {\scriptsize (0.007)} \\
ALL48-3k & ALL48 & $3.5$ {\scriptsize (0.022)} & $5.6$ {\scriptsize (\ensuremath{<}0.001)} & $8.2$ {\scriptsize (\ensuremath{<}0.001)} & $10.4$ {\scriptsize (\ensuremath{<}0.001)} & $7.7$ {\scriptsize (\ensuremath{<}0.001)} & $5.6$ {\scriptsize (\ensuremath{<}0.001)} \\
RTH13-3k & RTH13 & $8.7$ {\scriptsize (\ensuremath{<}0.001)} & $6.2$ {\scriptsize (0.009)} & $10.6$ {\scriptsize (\ensuremath{<}0.001)} & $12.3$ {\scriptsize (\ensuremath{<}0.001)} & $9.3$ {\scriptsize (\ensuremath{<}0.001)} & $7.1$ {\scriptsize (\ensuremath{<}0.001)} \\
ALL48-39k & ALL48 & $4.1$ {\scriptsize (\ensuremath{<}0.001)} & $2.6$ {\scriptsize (0.008)} & $2.2$ {\scriptsize (0.059)} & $2.7$ {\scriptsize (0.004)} & $3.3$ {\scriptsize (\ensuremath{<}0.001)} & $1.7$ {\scriptsize (0.017)} \\
\multicolumn{8}{@{}l}{Every bar} \\
ALL48-3k & RTH13-3k & $4.5$ {\scriptsize (\ensuremath{<}0.001)} & $4.9$ {\scriptsize (\ensuremath{<}0.001)} & $5.8$ {\scriptsize (\ensuremath{<}0.001)} & $4.7$ {\scriptsize (\ensuremath{<}0.001)} & $5.1$ {\scriptsize (\ensuremath{<}0.001)} & $2.7$ {\scriptsize (\ensuremath{<}0.001)} \\
ALL48-39k & RTH13 & $2.7$ {\scriptsize (\ensuremath{<}0.001)} & $4.2$ {\scriptsize (\ensuremath{<}0.001)} & $4.9$ {\scriptsize (\ensuremath{<}0.001)} & $10.4$ {\scriptsize (\ensuremath{<}0.001)} & $6.4$ {\scriptsize (\ensuremath{<}0.001)} & $7.9$ {\scriptsize (\ensuremath{<}0.001)} \\
ALL48-3k & ALL48 & $5.4$ {\scriptsize (\ensuremath{<}0.001)} & $5.9$ {\scriptsize (\ensuremath{<}0.001)} & $6.6$ {\scriptsize (\ensuremath{<}0.001)} & $15.2$ {\scriptsize (\ensuremath{<}0.001)} & $12.7$ {\scriptsize (\ensuremath{<}0.001)} & $6.8$ {\scriptsize (\ensuremath{<}0.001)} \\
RTH13-3k & RTH13 & $2.7$ {\scriptsize (0.003)} & $3.5$ {\scriptsize (0.007)} & $2.9$ {\scriptsize (\ensuremath{<}0.001)} & $14.9$ {\scriptsize (\ensuremath{<}0.001)} & $10.7$ {\scriptsize (\ensuremath{<}0.001)} & $9.4$ {\scriptsize (\ensuremath{<}0.001)} \\
ALL48-39k & ALL48 & $0.8$ {\scriptsize (0.118)} & $1.6$ {\scriptsize (0.001)} & $2.7$ {\scriptsize (\ensuremath{<}0.001)} & $5.7$ {\scriptsize (\ensuremath{<}0.001)} & $3.2$ {\scriptsize (\ensuremath{<}0.001)} & $2.6$ {\scriptsize (\ensuremath{<}0.001)} \\
\addlinespace[4pt]
\multicolumn{8}{@{}l}{\textit{Panel E: Sharpe ratio of the sign strategy, mid-quotes}} \\
RTH13 & Full & $1.85$ {\scriptsize (0.161)} & $1.72$ {\scriptsize (0.227)} & $1.94$ {\scriptsize (0.135)} & $1.75$ {\scriptsize (0.051)} & $1.67$ {\scriptsize (0.141)} & $1.96$ {\scriptsize (0.007)} \\
 & HAR & $0.98$ & $0.98$ & $0.98$ & $0.61$ & $0.91$ & $0.01$ \\
ALL48 & Full & $1.69$ {\scriptsize (--)} & $1.58$ {\scriptsize (--)} & $1.61$ {\scriptsize (--)} & $1.42$ {\scriptsize (0.013)} & $1.71$ {\scriptsize (0.087)} & $1.99$ {\scriptsize (0.002)} \\
 & HAR & $0.97$ & $0.97$ & $0.90$ & $0.04$ & $0.72$ & $0.27$ \\
CLOSE & Full & $1.34$ {\scriptsize (--)} & $1.50$ {\scriptsize (--)} & $1.78$ {\scriptsize (--)} & $1.87$ {\scriptsize (\ensuremath{<}0.001)} & $1.56$ {\scriptsize (0.014)} & $1.24$ {\scriptsize (0.063)} \\
 & HAR & $0.87$ & $0.87$ & $0.87$ & $-0.47$ & $0.02$ & $-0.07$ \\
CLOSE-nat & Full & $1.37$ {\scriptsize (--)} & $1.22$ {\scriptsize (--)} & $1.84$ {\scriptsize (--)} & -- & -- & -- \\
 & HAR & $0.80$ & $0.80$ & $0.80$ & -- & -- & -- \\
ALL48@16 & Full & $1.43$ {\scriptsize (--)} & $1.37$ {\scriptsize (--)} & $1.61$ {\scriptsize (--)} & -- & -- & -- \\
 & HAR & $0.97$ & $1.00$ & $0.77$ & -- & -- & -- \\
RTH13@16 & Full & $1.85$ {\scriptsize (--)} & $1.91$ {\scriptsize (--)} & $1.94$ {\scriptsize (--)} & -- & -- & -- \\
 & HAR & $0.98$ & $1.00$ & $0.94$ & -- & -- & -- \\
ALL48-3k & Full & $1.33$ {\scriptsize (--)} & $1.54$ {\scriptsize (--)} & $1.57$ {\scriptsize (--)} & $1.29$ {\scriptsize (0.398)} & $1.19$ {\scriptsize (0.526)} & $1.49$ {\scriptsize (0.149)} \\
 & HAR & $1.58$ & $1.63$ & $1.54$ & $0.94$ & $0.91$ & $0.83$ \\
RTH13-3k & Full & $1.13$ {\scriptsize (--)} & $1.68$ {\scriptsize (--)} & $1.37$ {\scriptsize (--)} & $0.89$ {\scriptsize (0.242)} & $0.86$ {\scriptsize (0.785)} & $0.96$ {\scriptsize (0.201)} \\
 & HAR & $1.11$ & $1.11$ & $1.11$ & $0.16$ & $0.70$ & $0.33$ \\
ALL48-39k & Full & $0.58$ {\scriptsize (--)} & $1.31$ {\scriptsize (--)} & $1.75$ {\scriptsize (--)} & $1.86$ {\scriptsize (0.001)} & $1.81$ {\scriptsize (0.005)} & $0.86$ {\scriptsize (0.680)} \\
 & HAR & $1.13$ & $1.15$ & $1.19$ & $0.26$ & $0.40$ & $0.66$ \\
\addlinespace[2pt]
ALL48 & $-$ RTH13 & $-0.16$ {\scriptsize (0.718)} & $-0.14$ {\scriptsize (0.757)} & $-0.33$ {\scriptsize (0.427)} & $-0.33$ {\scriptsize (0.231)} & $0.04$ {\scriptsize (0.883)} & $0.03$ {\scriptsize (0.927)} \\
CLOSE & $-$ RTH13 & $-0.50$ {\scriptsize (0.219)} & $-0.22$ {\scriptsize (0.623)} & $-0.16$ {\scriptsize (0.711)} & $0.12$ {\scriptsize (0.749)} & $-0.11$ {\scriptsize (0.669)} & $-0.72$ {\scriptsize (0.278)} \\
\addlinespace[4pt]
\multicolumn{8}{@{}l}{\textit{Panel F: CLOSE on windows of $W$ closes, full predictor set vs HAR benchmark, 16:00 bar}} \\
$W = 1{,}000$ & & $-4.8$ {\scriptsize (\ensuremath{<}0.001)} & $-2.6$ {\scriptsize (0.122)} & $-1.5$ {\scriptsize (0.510)} & $-7.8$ {\scriptsize (\ensuremath{<}0.001)} & $-6.8$ {\scriptsize (\ensuremath{<}0.001)} & $-2.1$ {\scriptsize (0.069)} \\
$W = 1{,}500$ & & $-4.1$ {\scriptsize (0.036)} & $-5.0$ {\scriptsize (\ensuremath{<}0.001)} & $-4.3$ {\scriptsize (0.035)} & $-5.8$ {\scriptsize (0.001)} & $-7.2$ {\scriptsize (\ensuremath{<}0.001)} & $-2.2$ {\scriptsize (0.052)} \\
$W = 2{,}000$ & & $-3.9$ {\scriptsize (0.049)} & $-4.6$ {\scriptsize (\ensuremath{<}0.001)} & $-4.6$ {\scriptsize (0.007)} & $-7.4$ {\scriptsize (\ensuremath{<}0.001)} & $-7.8$ {\scriptsize (\ensuremath{<}0.001)} & $-3.2$ {\scriptsize (0.002)} \\
$W = 2{,}500$ & & $-4.8$ {\scriptsize (0.008)} & $-5.3$ {\scriptsize (\ensuremath{<}0.001)} & $-5.8$ {\scriptsize (\ensuremath{<}0.001)} & $-7.0$ {\scriptsize (\ensuremath{<}0.001)} & $-7.6$ {\scriptsize (\ensuremath{<}0.001)} & $-3.7$ {\scriptsize (\ensuremath{<}0.001)} \\
$W = 3{,}000$ & & $-6.1$ {\scriptsize (\ensuremath{<}0.001)} & $-5.4$ {\scriptsize (\ensuremath{<}0.001)} & $-6.5$ {\scriptsize (\ensuremath{<}0.001)} & $-7.0$ {\scriptsize (\ensuremath{<}0.001)} & $-8.5$ {\scriptsize (\ensuremath{<}0.001)} & $-3.7$ {\scriptsize (\ensuremath{<}0.001)} \\
$W = 3{,}500$ & & $-6.5$ {\scriptsize (\ensuremath{<}0.001)} & $-5.1$ {\scriptsize (0.011)} & $-6.7$ {\scriptsize (\ensuremath{<}0.001)} & $-7.2$ {\scriptsize (\ensuremath{<}0.001)} & $-7.3$ {\scriptsize (\ensuremath{<}0.001)} & $-3.4$ {\scriptsize (0.001)} \\
\bottomrule
\end{tabular*}
\begin{tablenotes}[flushleft]
\footnotesize
\item \textit{Notes:} Samples, entries and $p$-values as in Panels A--C. Panel D: the equal-row samples against one another and against their equal-calendar siblings. Panel E: the Sharpe ratio $\SR$ of the sign strategy of Section~\ref{sec:rh_econ_measures} on the \rhtaOptDays{} post-warmup option days, at mid-quotes, for the forecast of the full predictor set (Full) and of the HAR benchmark (HAR); in parentheses after Full, the bootstrap $p$-value of Full minus HAR (paired block bootstrap, Appendix~\ref{sec:rh_methods_tests}), computed for the linear models in RTH13 only (--). Rows ALL48 $-$ RTH13 and CLOSE $-$ RTH13: the difference in $\SR$ between the full-set strategies of two samples, with its bootstrap $p$-value. Panel F: CLOSE on rolling windows of $W$ closes, everything else unchanged (the linear penalties re-chosen at each $W$ by the same rule; at $W = 1{,}000$ the penalized models use 1{,}001 closes; the tree ensembles keep the configurations chosen on the pooled regular-hours bars); $W = \rhtaWinDays$ is the CLOSE row of Panel~A.
\end{tablenotes}
\end{threeparttable}
\end{table}


% Decision D15 (author, 2026-10-26): keep the NYSE-holiday sessions in the design and disclose them, with the measured
% linear 2 x 2 (FOMC fixes x holiday sessions; results/fomcfix_2026-10-26, RULE.txt addendum A, cells A-D) as
% robustness. NOTE FOR OFFICIAL PUBLICATION (plan D15): rerun every model with DROP_HOLIDAY_SESSIONS on before the
% final version.
\subsection{NYSE-holiday sessions}
\label{app:rh_robust_holiday}
\begin{itemize}
% 2026-10-26 (coordinator): unwrapped \IfFileExists{generated/rehaul_holiday_table.tex} (the file exists; the flat Overleaf build has no rehaul/ directory, so the wrapper skipped its body there). Its false branch was:
% | %
% | \item \textbf{The sessions.} On full NYSE holidays the E-mini trades a shortened session that ends at about
% |   13:00~ET. Its regular-hours bars are rows of the estimation sample and enter the target's diurnal factor, but they
% |   are not scored, since Table~\ref{tab:rh_oos} scores the bars of NYSE trading days
% |   (Appendix~\ref{app:rh_dp_sessions}). The main design keeps them.
% | \item \textbf{The sessions dropped.} The pooled linear models re-estimated without the holiday sessions are in the
% |   replication package.
%
\item \textbf{The sessions.} On full NYSE holidays the E-mini trades a shortened session whose last regular-hours
  bar ends at about \rhholLastBar~ET. The \rhholNSessions{} such sessions of the evaluation span have \rhholNBars{}
  regular-hours bars. They are rows of the estimation sample and enter the target's diurnal factor, but they are not
  scored, since Table~\ref{tab:rh_oos} scores the bars of NYSE trading days: each pooled model has \rhholNForecast{}
  forecasts and \rhholNScored{} scored bars (Appendix~\ref{app:rh_dp_sessions}). The main design keeps them.
\item \textbf{The sessions dropped.} Table~\ref{tab:rh_holiday} re-estimates the penalized linear models of
  Table~\ref{tab:rh_oos} with the holiday sessions dropped from the estimation rows and from the diurnal factor,
  every trailing average still computed on the full grid and every penalty frozen, with and without the FOMC fixes
  of Appendix~\ref{app:rh_dp_fomc}; every variant is scored on the same bars.
  \begin{itemize}
  \item With the FOMC fixes, dropping the sessions moves the full predictor set's change in every-bar QLIKE against
    the HAR benchmark by \rhholEffDBBarPpRng{} percentage points, and its change at 16:00 by
    \rhholEffDBClosePpRng{} (without the fixes, \rhholEffCABarPpRng{} and \rhholEffCAClosePpRng).
  \item The Sharpe ratio of the full-set sign strategy moves by \rhholEffDBShFullRng.
  \end{itemize}
% AUTO-GENERATED by writeup/make_rehaul_holiday.py -- do not edit; re-run the script.
% Table tab:rh_holiday (appendix, D15): the linear 2 x 2 of results/fomcfix_2026-10-26 (RULE.txt ADDENDUM A).
\begin{table}[htbp]
\centering
\footnotesize
\setlength{\tabcolsep}{3pt}
\begin{threeparttable}
\caption{NYSE-holiday sessions: kept or dropped}
\label{tab:rh_holiday}
\begin{tabular*}{\linewidth}{@{\extracolsep{\fill}}llllrrrrr@{}}
\toprule
 & & & & \multicolumn{2}{c}{Every bar} & \multicolumn{2}{c}{16:00} & \\
\cmidrule(lr){5-6}\cmidrule(lr){7-8}
Model & Cell & Holidays & FOMC fix & \% vs HAR & DM ($p$) & \% vs HAR & DM ($p$) & $\SR$ full / HAR \\
\midrule
Ridge & A & kept & before & \ensuremath{-}5.4 & \ensuremath{-}5.63 (\ensuremath{<}0.001) & +0.8 & 0.31 (0.759) & 1.85 / 0.98 \\
 & B & kept & after & \ensuremath{-}6.0 & \ensuremath{-}6.28 (\ensuremath{<}0.001) & +1.8 & 0.67 (0.500) & 1.75 / 0.98 \\
 & C & dropped & before & \ensuremath{-}5.2 & \ensuremath{-}5.39 (\ensuremath{<}0.001) & \ensuremath{-}0.2 & \ensuremath{-}0.07 (0.943) & 1.93 / 1.01 \\
 & D & dropped & after & \ensuremath{-}5.9 & \ensuremath{-}6.03 (\ensuremath{<}0.001) & +0.9 & 0.37 (0.712) & 1.79 / 1.01 \\
\addlinespace[3pt]
Lasso & A & kept & before & \ensuremath{-}6.2 & \ensuremath{-}4.72 (\ensuremath{<}0.001) & +3.3 & 1.04 (0.300) & 1.72 / 0.98 \\
 & B & kept & after & \ensuremath{-}9.0 & \ensuremath{-}6.34 (\ensuremath{<}0.001) & +4.2 & 1.35 (0.178) & 1.75 / 0.98 \\
 & C & dropped & before & \ensuremath{-}6.1 & \ensuremath{-}4.76 (\ensuremath{<}0.001) & +2.8 & 0.93 (0.353) & 2.22 / 1.01 \\
 & D & dropped & after & \ensuremath{-}8.7 & \ensuremath{-}6.08 (\ensuremath{<}0.001) & +4.0 & 1.33 (0.185) & 2.14 / 1.01 \\
\addlinespace[3pt]
Elastic net & A & kept & before & \ensuremath{-}6.0 & \ensuremath{-}6.98 (\ensuremath{<}0.001) & \ensuremath{-}0.5 & \ensuremath{-}0.21 (0.832) & 1.94 / 0.98 \\
 & B & kept & after & \ensuremath{-}6.7 & \ensuremath{-}7.36 (\ensuremath{<}0.001) & +0.6 & 0.24 (0.807) & 1.95 / 0.98 \\
 & C & dropped & before & \ensuremath{-}5.9 & \ensuremath{-}6.71 (\ensuremath{<}0.001) & \ensuremath{-}1.3 & \ensuremath{-}0.56 (0.579) & 1.67 / 1.01 \\
 & D & dropped & after & \ensuremath{-}6.6 & \ensuremath{-}7.04 (\ensuremath{<}0.001) & 0.0 & \ensuremath{-}0.02 (0.984) & 1.61 / 1.01 \\
\bottomrule
\end{tabular*}
\begin{tablenotes}[flushleft]
\footnotesize
\item \textit{Notes:} The penalized linear models of Table~\ref{tab:rh_oos} on four designs: the \rhholNSessions{} NYSE-holiday sessions of the evaluation span (\rhholNBars{} regular-hours bars, a shortened futures session ending about \rhholLastBar~ET) kept in or dropped from the estimation rows and the diurnal factor $B$, and the FOMC calendar corrections of Section~\ref{sec:rh_data} before or after; cell B is the design of the main text. Every cell keeps its frozen penalty and is scored on the same \rhholNScored{} bars of Table~\ref{tab:rh_oos} (the holiday bars are never scored) and on their 16:00 bar. \% vs HAR: the full set's QLIKE against the HAR set of the same class (negative: the full set has the lower loss); DM: the Diebold--Mariano statistic of the daily average (every bar) or 16:00 QLIKE, HLN correction, Newey--West lag \rhorHacLag, $p$-value in parentheses. $\SR$: the annualized Sharpe ratio of the sign strategy at mid-quotes, full set / HAR set, on the \rheSzDays{} post-warmup option days.
\end{tablenotes}
\end{threeparttable}
\end{table}

\end{itemize}
